\documentclass[12pt,a4paper,oneside]{report}

% Packages
\usepackage[T1]{fontenc}
\usepackage{ebgaramond}  % Garamond font
\usepackage{amsmath,amsfonts,amssymb}
\usepackage{graphicx}
\usepackage{xcolor}
\usepackage{hyperref}
\usepackage[numbers,compress]{natbib}
\usepackage{geometry}
\usepackage{setspace}
\usepackage{fancyhdr}
\usepackage{enumitem}
\usepackage{ragged2e}
\usepackage{lettrine}
\usepackage{tocloft}  % For customizing table of contents and lists
\usepackage{titlesec}  % For customizing chapter and section spacing

% Define subsubsubsection command
\titleclass{\subsubsubsection}{straight}[\subsection]
\newcounter{subsubsubsection}[subsubsection]
\renewcommand{\thesubsubsubsection}{\thesubsubsection.\arabic{subsubsubsection}}
\titleformat{\subsubsubsection}
  {\normalfont\normalsize\bfseries}{\thesubsubsubsection}{1em}{}
\titlespacing*{\subsubsubsection}{0pt}{3.25ex plus 1ex minus .2ex}{1.5ex plus .2ex}
\usepackage{atbegshi}  % For conditional page content
\usepackage{eso-pic}  % For adding content to pages (better browser PDF compatibility)
\usepackage{tikz}  % For drawing the footer line
\usepackage{caption}  % For customizing captions
\usepackage{booktabs}  % For professional table formatting
\usepackage{tabularx}  % For flexible table widths
\usepackage{multirow}  % For multirow cells in tables
\usepackage{refcount}  % For getting page numbers from labels
\usepackage{float}  % For better figure placement control
\usepackage{pdflscape}  % For landscape pages
\usepackage{rotating}  % For rotating tables and figures
\usepackage{placeins}  % For FloatBarrier to control figure placement
\usepackage{adjustbox}  % For automatically scaling tables to fit

% Setup for list of equations
\newlistof{equations}{loe}{}
\newcommand{\listequationsname}{Equations}

% Command to add equations to the list
% Usage: \myequation[optional label]{equation content}
\newcommand{\myequation}[2][]{%
    \refstepcounter{equations}%
    \addcontentsline{loe}{equations}{\protect\numberline{\theequations}$#2$}%
    \ifx\relax#1\relax\else\label{#1}\fi%
    #2%
}

% Page setup
\geometry{margin=1in}
\setstretch{1.8}  % 20% increase from 1.5
\justifying
\setlength{\parindent}{0pt}  % No indentation at start of paragraphs
\setlength{\headheight}{20.4pt}  % Fix fancyhdr warning

% Float placement parameters - prefer bottom placement (~1cm from footer) unless figure occupies whole page
\setlength{\topfraction}{0.85}  % Allow figures to take up to 85% of top of page
\setlength{\bottomfraction}{0.85}  % Allow figures to take up to 85% of bottom of page
\setlength{\textfraction}{0.15}  % Require at least 15% text on page
\setlength{\floatpagefraction}{0.7}  % Figures on float pages must take at least 70%
\renewcommand{\topfraction}{0.85}
\renewcommand{\bottomfraction}{0.85}
\renewcommand{\textfraction}{0.15}
\renewcommand{\floatpagefraction}{0.7}
\setlength{\dbltopfraction}{0.85}
\renewcommand{\dbltopfraction}{0.85}
\setlength{\dblfloatpagefraction}{0.7}
\renewcommand{\dblfloatpagefraction}{0.7}
\setlength{\abovecaptionskip}{0.5cm}  % Space above caption
\setlength{\belowcaptionskip}{0.5cm}  % Space below caption
% Table-specific float settings to ensure tables stay between header and footer
\setlength{\intextsep}{1.5cm}  % Space above and below floats in text (reduced to allow more space for tables)
\setlength{\floatsep}{1.5cm}  % Space between floats
\setlength{\textfloatsep}{1.5cm}  % Space between top/bottom floats and text
% Ensure tables don't exceed text area
\setlength{\dbltextfloatsep}{1.5cm}  % Space for double-column floats
\makeatletter
\setlength{\@fptop}{0pt plus 1fil}  % Allow tables at top
\setlength{\@fpsep}{1.5cm plus 2fil}  % Space between float pages
\setlength{\@fpbot}{0pt plus 1fil}  % Allow tables at bottom
\makeatother

% Set default text color to grey (darker for visibility)
\definecolor{textgrey}{gray}{0.0}

% Ensure body text is consistently 12pt throughout (normalsize = 12pt when documentclass is 12pt)
\normalsize

% Hyperref setup
\hypersetup{
    colorlinks=true,
    linkcolor=blue,
    filecolor=magenta,      
    urlcolor=cyan,
    citecolor=blue,
    pdftitle={Material detection with SPADs},
    pdfauthor={AMA},
    pdfstartpage=1,
    pageanchor=false
}

% Adjust chapter and section spacing
\titlespacing*{\chapter}{0pt}{0pt}{10pt}  % Reduce space after chapter
\titlespacing*{\section}{0pt}{5pt}{5pt}  % Reduce space before and after section

% Ensure every section starts on a fresh page
\let\oldsection\section
\renewcommand{\section}{\clearpage\oldsection}

% Set figure caption font size to 10pt (smaller than body text) and position at bottom
\captionsetup[figure]{font={footnotesize}, labelfont=bf, textfont=normalfont, position=bottom}
\captionsetup[table]{position=bottom}
\renewcommand{\figurename}{Figure}
\setlength{\captionmargin}{0pt}

% Header and footer
% Redefine chapter and section marks to not use uppercase
\renewcommand{\chaptermark}[1]{\markboth{#1}{}}
\renewcommand{\sectionmark}[1]{\markright{#1}}

\pagestyle{fancy}
\fancyhf{}
\fancyhead[L]{\fontsize{8}{9.6}\selectfont\itshape\leftmark}
\fancyhead[R]{\thepage}
\renewcommand{\headrulewidth}{0pt}  % Remove header line
\renewcommand{\headrule}{}  % Remove header rule
\renewcommand{\footrulewidth}{0pt}  % Remove footer line

% Flag to track if we're using arabic (integer) page numbering
\newif\ifarabicpagenum
\arabicpagenumfalse

% Flag to track if we're in bibliography
\newif\ifinbibliography
\inbibliographyfalse

% Ensure chapter pages use fancy style
\let\oldchapter\chapter
\renewcommand{\chapter}{\thispagestyle{fancy}\oldchapter}

% Include equations database
% Equations Database
% ==================
% This file contains all equations to be inserted in the thesis document.
% Each equation has a unique ID and can be inserted anywhere in the document.
%
% INSERTION INSTRUCTIONS:
% -----------------------
% To insert an equation in the document, use:
% \insertequation{EQUATION_ID}
%
% To insert a numbered equation (for list of equations), use:
% \insertequationnumbered{EQUATION_ID}{LABEL}
%
% Example: \insertequation{1}
% Example: \insertequationnumbered{1}{eq:example}
%
% The numbered version will appear in the List of Equations.

% ============================================
% EQUATION DEFINITIONS
% ============================================

% Equation ID: 1
% Status: Ready for insertion
% Description: Speed of light equation (wavelength-frequency relationship)
\newcommand{\equationone}{%
    c = \lambda v
}

% Equation ID: 2
% Status: Ready for insertion
% Description: Target distance calculation for indirect time-of-flight measurement
\newcommand{\equationtwo}{%
    d_{\text{target}} = \frac{c \cdot \varphi}{4 \pi \nu}
}

% Equation ID: 3
% Status: Ready for insertion
% Description: Softmax function (probability distribution for classification)
\newcommand{\equationthree}{%
    p_{\text{class}} = \frac{e^{\text{logit}_{i}}}{\sum_{j=1}^{K_c} e^{\text{logit}_{j}}}
}

% Equation ID: 4
% Status: Ready for insertion
% Description: Linear model with residual error
\newcommand{\equationfour}{%
    U = r_{m} \cdot p_{m} + e
}

% Equation ID: 5
% Status: Ready for insertion
% Description: Summation-based linear model with residual error
\newcommand{\equationfive}{%
    U = \sum_{m=1}^{M} r_{m} \cdot p_{m} + e
}

% Equation ID: 6
% Status: Ready for insertion
% Description: Gaussian mixture model with integration
\newcommand{\equationsix}{%
    \lambda_{i} = \sum_{j=1}^{M} \left\{ \int_{t_{i}}^{t_{i+1}} \frac{S_{j}}{\sqrt[\sigma_{j}]{2\pi}} e^{\left(\frac{(t-\mu_{j})^{2}}{2\sigma^{2}_{j}}\right)} dt\right\} + b
}

% Equation ID: 7
% Status: Ready for insertion
% Description: Single Gaussian component with integration
\newcommand{\equationseven}{%
    \lambda_{i} = \int_{t_{i}}^{t_{i+1}} \frac{S_{j}}{\sqrt[\sigma_{j}]{2\pi}} e^{\left(\frac{(t-\mu_{j})^{2}}{2\sigma^{2}_{j}}\right)} dt + b
}

% Equation ID: 8
% Status: Ready for insertion
% Description: Poisson probability distribution
\newcommand{\equationeight}{%
    p(h_{i} = k) = \frac{\lambda_{i}^{k}e^{-\lambda_{i}}}{k!}
}

% Equation ID: 9
% Status: Ready for insertion
% Description: Return signal as product of system components
\newcommand{\equationnine}{%
    R_{\text{return}} = P_{\text{laser}}(t) \cdot M_{\text{target}}(t) \cdot J_{\text{TDC}}(t) \cdot J_{\text{SPAD}}(t)
}

% Equation ID: 10
% Status: Ready for insertion
% Description: Round-trip time for time-of-flight measurement
\newcommand{\equationten}{%
    T_{\text{rt}} = 2 \frac{d_{\text{target}}}{c}
}

% Equation ID: 11
% Status: Ready for insertion
% Description: Loss function for height and width prediction (with constraint: α ∈ {0,1})
\newcommand{\equationeleven}{%
    \text{Loss}_{h,w} = \sum_{i=0}^{a} \sum_{j=0}^{g^{2}} \alpha_{ji}^{\text{trg}} \left[ (h_{\text{gt}} - h_{\text{pred}})^{2} + (w_{\text{gt}} - w_{\text{pred}})^{2} \right], \quad \alpha \in \{0,1\}
}

% Equation ID: 12
% Status: Ready for insertion
% Description: Classification loss function with double summation
\newcommand{\equationtwelve}{%
    \text{Loss}_{C} = \sum_{i=0}^{a} \sum_{j=0}^{g^{2}} \alpha_{ji}^{\text{trg}} \left[ (C_{\text{gt}} - C_{\text{pred}})^{2} \right], \quad C_{\text{gt}} \in \{1\}
}

% Equation ID: 13
% Status: Ready for insertion
% Description: Classification probability loss function with constraint
\newcommand{\equationthirteen}{%
    \text{Loss}_{p_{\text{gt}}^{\text{cls}}} = \sum_{i=0}^{a} \sum_{j=0}^{g^{2}} \alpha_{ji}^{\text{trg}} \left[ (p_{\text{gt}}^{\text{cls}} - p_{\text{pred}}^{\text{cls}})^{2} \right], \quad 0 \leq p_{\text{pred}}^{\text{cls}} \leq 1
}

% Equation ID: 14
% Status: Ready for insertion
% Description: Material classification probability loss function with constraint
\newcommand{\equationfourteen}{%
    \text{Loss}_{p_{\text{gt}}^{\text{mat\_cls}}} = \sum_{i=0}^{a} \sum_{j=0}^{g^{2}} \alpha_{ji}^{\text{trg}} \left[ (p_{\text{gt}}^{\text{mat\_cls}} - p_{\text{pred}}^{\text{mat\_cls}})^{2} \right], \quad 0 \leq p_{\text{gt}}^{\text{mat\_cls}} \leq 1
}

% Equation ID: 15
% Status: Ready for insertion
% Description: Coordinate loss function for x and y positions with constraint
\newcommand{\equationfifteen}{%
    \text{Loss}_{x,y} = \sum_{i=0}^{a} \sum_{j=0}^{g^{2}} \alpha_{ji}^{\text{trg}} \left[ (x_{\text{gt}} - x_{\text{pred}})^{2} + (y_{\text{gt}} - y_{\text{pred}})^{2} \right], \quad \alpha \in \{0,1\}
}

% ============================================
% INSERTION COMMANDS
% ============================================

% Command to insert an equation (unnumbered, inline or display)
% Usage: \insertequation{EQUATION_ID}
% Example: \insertequation{1}
%
% This will insert the equation in display mode (centered, on its own line)
\newcommand{\insertequation}[1]{%
    \begin{equation*}
        \csname equation#1\endcsname
    \end{equation*}%
}

% Command to insert a numbered equation (appears in List of Equations)
% Usage: \insertequationnumbered{EQUATION_ID}{LABEL}
% Example: \insertequationnumbered{1}{eq:example}
%
% This will:
% - Insert the equation with a number
% - Add it to the List of Equations
% - Create a label for cross-referencing with \ref{LABEL}
\newcommand{\insertequationnumbered}[2]{%
    \begin{equation}
        \csname equation#1\endcsname
        \label{#2}
    \end{equation}%
    \addcontentsline{loe}{equations}{\protect\numberline{\theequation}$\csname equation#1\endcsname$}%
}

% Command to insert an equation inline (within text)
% Usage: \insertequationinline{EQUATION_ID}
% Example: \insertequationinline{1}
%
% This will insert the equation inline with the text
\newcommand{\insertequationinline}[1]{%
    $\csname equation#1\endcsname$%
}

% Command to insert an equation in a specific environment
% Usage: \insertequationenv{EQUATION_ID}{ENVIRONMENT}
% Example: \insertequationenv{1}{align}
%
% Supported environments: equation, equation*, align, align*, gather, gather*, etc.
\newcommand{\insertequationenv}[2]{%
    \begin{#2}%
        \csname equation#1\endcsname
    \end{#2}%
}

% ============================================
% EQUATION REGISTRY (for reference)
% ============================================
% 
% ID | Location | Status | Description
% ---|----------|--------|------------
%  1 |  -       | Ready  | Speed of light equation (c = λv)
%  2 |  -       | Ready  | Target distance calculation for ToF measurement
%  3 |  -       | Ready  | Softmax function (probability distribution)
%  4 |  -       | Ready  | Linear model with residual error
%  5 |  -       | Ready  | Summation-based linear model with residual error
%  6 |  -       | Ready  | Gaussian mixture model with integration
%  7 |  -       | Ready  | Single Gaussian component with integration
%  8 |  -       | Ready  | Poisson probability distribution
%  9 |  -       | Ready  | Return signal as product of system components
% 10 |  -       | Ready  | Round-trip time for time-of-flight measurement
% 11 |  -       | Ready  | Loss function for height and width prediction
% 12 |  -       | Ready  | Classification loss function with double summation
% 13 |  -       | Ready  | Classification probability loss function with constraint
% 14 |  -       | Ready  | Material classification probability loss function with constraint
% 15 |  -       | Ready  | Coordinate loss function for x and y positions with constraint
%



% Include quotes database
% Quote command definitions
% These commands output the quote text with proper formatting

\newcommand{\quoteone}{%
\vspace{1.0\baselineskip}%
\begin{quote}\ttfamily
Spatial features encompass the two-dimensional geometric properties and spatial relationships of target objects, including their shapes and relative positions within the scene. While RGB imaging systems excel at capturing such spatial information, they exhibit limitations in acquiring structural features, which include three-dimensional surface profiles and the geometric composition of objects.%
\end{quote}%
\vspace{1.0\baselineskip}%
}

\newcommand{\quotetwo}{%
\vspace{1.0\baselineskip}%
\begin{quote}\ttfamily
A cat can remember, can understand the physical world, can plan complex actions, can do some level of reasoning—actually much better than the biggest LLMs. —Yann LeCunn (2024)%
\end{quote}%
\vspace{1.0\baselineskip}%
}

\newcommand{\quotethree}{%
\vspace{1.0\baselineskip}%
\begin{quote}\ttfamily
vision is the production of relevant uncluttered visual information representations of the external world. —Marr 1976, Marr et al 1978%
\end{quote}%
\vspace{1.0\baselineskip}%
}

\newcommand{\quotefour}{%
\vspace{1.0\baselineskip}%
\begin{quote}\ttfamily
AI is the science of making machines capable of performing tasks that would require intelligence if done by humans. —The Society of mind, Minsky, Marvin Lee, 1927 [13]%
\end{quote}%
\vspace{1.0\baselineskip}%
}

\newcommand{\quotefive}{%
\vspace{1.0\baselineskip}%
\begin{quote}\ttfamily
In the new setting of ideas, the distinction (between particles and waves) has vanished, because it was discovered that all particles have also wave properties, and vice-versa. Neither of the concepts must be discarded, they must be amalgamated. Which aspect obtrudes itself depends not on the physical object but on the experimental device set up to examine it. —Erwin Schrödinger (1867–1961). Science, Theory and Man (New York, Dover, 1957)%
\end{quote}%
\vspace{1.0\baselineskip}%
}

\newcommand{\quotesix}{%
\vspace{1.0\baselineskip}%
\begin{quote}\ttfamily
homogenization efficiency test is to determine if the instrument homogenizer is breaking fat globules down into small enough sizes to allow the instrument to give accurate test results without interference due to light scattering by large fat globules. —D. M. BARBANO et al, 1989%
\end{quote}%
\vspace{1.0\baselineskip}%
}

\newcommand{\quoteseven}{%
\vspace{1.0\baselineskip}%
\begin{quote}\ttfamily
… the Milko-tester (Foss Electric, Hillerød, Denmark) measured the Vis scattering after dispersion of the CN micelles and homogenization to produce a more uniform fat globule size distribution. —Aernouts et al., 2015%
\end{quote}%
\vspace{1.0\baselineskip}%
}


% Title page information
\title{Material detection with SPADs}
\author{Deborah Akuoko}
\date{\today}

\begin{document}
\color{textgrey}  % Set text color after document begins
%\let\cleardoublepage\clearpage
% Suppress any blank pages from document initialization
\makeatletter
\@twosidefalse
\@openrightfalse
\makeatother
% Blank page at the start of the project
\thispagestyle{empty}
\newpage
% Page 1 with UoE logo and text - start immediately to prevent blank page
\thispagestyle{empty}\mbox{}
\vspace*{\fill}
\begin{center}
    \includegraphics[width=0.6\textwidth]{figures/uoe_logo_1.png}
    
    \vspace{1cm}
    
    \begin{minipage}{0.8\textwidth}
        \justifying\sloppy
        \setlength{\parindent}{0pt}
        This thesis has been submitted in fulfilment of the requirements for a postgraduate degree (e.~g. PhD, MPhil, DClinPsychol) at the University of Edinburgh. Please note the following terms and conditions of use:
        
        \vspace{0.5cm}
        
        \begin{itemize}[leftmargin=*]
            \item This work is protected by copyright and other intellectual property rights, which are retained by the thesis author, unless otherwise stated.
            \item A copy can be downloaded for personal non-commercial research or study, without prior permission or charge.
            \item This thesis cannot be reproduced or quoted extensively from without first obtaining permission in writing from the author.
            \item The content must not be changed in any way or sold commercially in any format or medium without the formal permission of the author.
            \item When referring to this work, full bibliographic details including the author, title, awarding institution and date of the thesis must be given.
        \end{itemize}
    \end{minipage}
\end{center}
\vspace*{\fill}
\newpage

% Page 2 with University and degree text
\thispagestyle{empty}
\begin{center}
    {\huge\href{http://www.ed.ac.uk/}{The University of Edinburgh}}
    
    \vspace{0.5cm}
    
    {\huge Ph.D Thesis}
    
    \vspace{1.5cm}
    
    \rule{\textwidth}{0.5pt}
    
    \vspace{0.5cm}
    
    {\huge\textbf{Material detection with SPADs}}
    
    \vspace{0.5cm}
    
    \rule{\textwidth}{0.5pt}
    
    \vspace{1cm}
    
    \begin{minipage}{\textwidth}
        \textit{Author: Deborah Akuoko}
        \hfill
        \begin{minipage}[t]{0.5\textwidth}
            \raggedleft
            \textit{Advisor(s): Dr Istvan Gyongy (Principal)\\
            Dr Daniel Chitnis (Supporting)}
        \end{minipage}
    \end{minipage}
    
    \vspace{1.5cm}
    
    A thesis submitted in fulfilment of the requirements for the degree of Doctor of Philosophy
    
    \vspace{0.5cm}
    
    with
    
    \vspace{0.5cm}
    
    \textit{Institute for Integrated Micro and Nano Systems (IMNS)}
    
    \vspace{0.5cm}
    
    in
    
    \vspace{0.5cm}
    
    \href{https://eng.ed.ac.uk/}{School of Engineering}
    
    \vspace{0.5cm}
    
    \includegraphics[width=0.32\textwidth]{figures/uoe_logo_2.png}
    
    \vspace{0.5cm}
    
    December 2025
\end{center}
\newpage

% Page 3 with header (thin grey line and page number 'i')
\pagenumbering{roman}
\setcounter{page}{1}
\newpage
\thispagestyle{empty}
\vspace*{-1.5cm}
\hfill\textcolor{gray}{\textit{\roman{page}}}
\vspace{0.2cm}
\noindent\rule{\textwidth}{0.5pt}
\vspace{1cm}
{\fontsize{12pt}{18pt}\selectfont
{\fontsize{14pt}{21pt}\selectfont\textbf{Declaration}}

I, Deborah Akuoko, hereby declare that this thesis is the result of my own work, except where otherwise acknowledged. It contains no material previously published or written by another person, nor material that has been accepted for the award of any other degree or diploma at this or any other institution. All sources of information have been properly cited.

\vspace{1cm}
\begin{tabular}{@{}l@{}}
Signed: \rule{8cm}{0.4pt} \\
Date: \rule{8cm}{0.4pt} \\
\end{tabular}
}
\vspace*{\fill}
\newpage
\thispagestyle{empty}
\vspace*{-1.5cm}
\hfill\textcolor{gray}{\textit{\roman{page}}}
\vspace{0.2cm}
\noindent\rule{\textwidth}{0.5pt}
\vspace{0.5cm}
{\fontsize{12pt}{18pt}\selectfont
{\fontsize{14pt}{21pt}\selectfont\textbf{Abstract}}

\vspace{0.2cm}
Automatic pattern recognition (PR) systems have proven to be indispensable in machine perception, especially deep neural networks (DNN) for vision. The last few years have seen the development of sophisticated machine vision architectures as well as powerful process infrastructures. Nevertheless, reliability remains a crucial factor in the adoption such systems in safety critical applications such as self driving or advanced driving or advanced driving assistance systems. With the single photon avalanche diode (SPAD) detector, this work tackles the limitations of some of the best-performing vision systems and presents experimentally verified solutions founded on how structural features are sourced for visual scene inference. A combined instrument of SPAD and APS (active pixel sensor) imagers reveals significant improvements in how machines perceive and process real-world scenes and make reliable decisions in comparison to traditional \textnormal{2D} spatially resolved image inputs. While state-of-the-art vision systems are not perfect, certain criticisms of their performance could be ambiguous or inaccurate. To ensure clarity and assert the relevance of exact machine vision problems addressed by SPADs in this work, parts of the reviewed literature are dedicated to a thorough assessment of machine perception. Key deliverables of this work include a custom framework (Particle Wave Detection) that guides how the instrument perceives and processes the signal. Objectives one (material detection) and three (material purity detection) focus on solid and fluid material classification using machine learning models and \textnormal{1D} transient image inputs. Based on successful performance, objective two is used to validate the novel idea of combining both \textnormal{1D} transient and \textnormal{2D} spatially resolved images (of RGB encoding) for improved object detection. The dual detection model combines the material classifier with either of the selected best performing state of the art visual object detection models. Overall, machines demonstrated improved scene perception and reliable inferences with SPADs, especially in the dual detection model for vision applications. Potential applications include low-cost, compact and robust systems for robotic vision and food safety monitoring.
}
\vspace*{\fill}
\newpage

% Page 3 duplicate (Lay Summary)
\newpage
\thispagestyle{empty}
\vspace*{-1.5cm}
\hfill\textcolor{gray}{\textit{\roman{page}}}
\vspace{0.2cm}
\noindent\rule{\textwidth}{0.5pt}
\vspace{1cm}
{\fontsize{12pt}{18pt}\selectfont
{\fontsize{14pt}{21pt}\selectfont\textbf{Lay Summary}}

This project is focused on demonstrating how optical sensing with single-photon avalanche diodes (SPADs) can address machine perception problems. In the absence of a universal perception theory, PR systems rely on learning and decision models to try and make sense of the world. Thus, a fundamental framework is first developed to guide how signals (or percepts) are measured before feature extraction, learning, or inference takes place. In all three objectives, a 4x4 zone SPAD detector of an active invisible narrow band (~940nm) source is used to acquire a time-resolved signal and processed with no regard to its spectral composition. This processing approach is considered achromatic or non-spectral. In objective one, the signal is sourced from flat surface materials of varying homogeneous composition. The variations in absorption and scattering are extracted and processed by a machine learning algorithm for multi-label classification. High prediction accuracies are observed in validation tests (~98\%). In objective two, the time-resolved signal is combined with spatially resolved multispectral images (RGB encoded) to perform a combined task of target localisation and classification. While most visual models perform this task with RGB encoded images, the combination of high-performing state-of-the-art (SOTA) visual perception models with time-resolved classifier models significantly improved the quality of inferences made, making it a relevant breakthrough for machine vision. Finally, the device is used to evaluate the purity of a homogenised fluid material (milk), which also yields high assessment results (~97\% detection accuracy). Overall, material detection is essential in machine perception. Remotely detecting materials with SPADs proves to be not only reliable but also a key player in addressing some of the unsolved problems of vision, such as how structural features may be sourced for visual scene interpretation. The presented solution serves as a first call to the research community to realise the impact of time-resolved features on machine vision quality, which consequently determines how it counts towards the ultimate goal of building truly intelligent machines.
}
\vspace*{\fill}
\newpage

% Acknowledgement
\newpage
\thispagestyle{empty}
\vspace*{-1.5cm}
\hfill\textcolor{gray}{\textit{\roman{page}}}
\vspace{0.2cm}
\noindent\rule{\textwidth}{0.5pt}
\vspace{1cm}
{\fontsize{12pt}{18pt}\selectfont
{\fontsize{14pt}{21pt}\selectfont\textbf{Acknowledgement}}

% Acknowledgement text goes here
}
\vspace*{\fill}
\newpage

% Blank page (page 6)
\newpage
\thispagestyle{empty}
\vspace*{-1.5cm}
\hfill\textcolor{gray}{\textit{\roman{page}}}
\vspace{0.2cm}
\noindent\rule{\textwidth}{0.5pt}
\vspace*{\fill}
\newpage

% Contents (page 7)
\newpage
\thispagestyle{empty}
\vspace*{-1.5cm}
\hfill\textcolor{gray}{\textit{\roman{page}}}
\vspace{0.2cm}
\noindent\rule{\textwidth}{0.5pt}
\vspace{0.5cm}
{\fontsize{12pt}{18pt}\selectfont
{\fontsize{14pt}{21pt}\selectfont\textbf{Contents}}

\vspace{0.5cm}

\makeatletter
\@starttoc{toc}
\makeatother
}
\vspace*{\fill}
\newpage

% Figures (page 8)
\newpage
\thispagestyle{empty}
\vspace*{-1.5cm}
\hfill\textcolor{gray}{\textit{\roman{page}}}
\vspace{0.2cm}
\noindent\rule{\textwidth}{0.5pt}
\vspace{0.5cm}
{\fontsize{12pt}{18pt}\selectfont
{\fontsize{14pt}{21pt}\selectfont\textbf{Figures}}

\vspace{0.5cm}

\makeatletter
\@starttoc{lof}
\makeatother
}
\vspace*{\fill}
\newpage

% Tables (page 9)
\newpage
\thispagestyle{empty}
\vspace*{-1.5cm}
\hfill\textcolor{gray}{\textit{\roman{page}}}
\vspace{0.2cm}
\noindent\rule{\textwidth}{0.5pt}
\vspace{0.5cm}
{\fontsize{12pt}{18pt}\selectfont
{\fontsize{14pt}{21pt}\selectfont\textbf{Tables}}

\vspace{0.5cm}

\makeatletter
\@starttoc{lot}
\makeatother
}
\vspace*{\fill}
\newpage

% Equations (page 10)
\newpage
\thispagestyle{empty}
\vspace*{-1.5cm}
\hfill\textcolor{gray}{\textit{\roman{page}}}
\vspace{0.2cm}
\noindent\rule{\textwidth}{0.5pt}
\vspace{0.5cm}
{\fontsize{12pt}{18pt}\selectfont
{\fontsize{14pt}{21pt}\selectfont\textbf{Equations}}

\vspace{0.5cm}

\makeatletter
\@starttoc{loe}
\makeatother
}
\vspace*{\fill}
\newpage

% Mathematical notation (page 11)
\newpage
\thispagestyle{empty}
\vspace*{-1.5cm}
\hfill\textcolor{gray}{\textit{\roman{page}}}
\vspace{0.2cm}
\noindent\rule{\textwidth}{0.5pt}
\vspace{0.5cm}
{\fontsize{12pt}{18pt}\selectfont
{\fontsize{14pt}{21pt}\selectfont\textbf{Mathematical notation}}

\vspace{0.5cm}

{\setlength{\parindent}{0pt}%
\begin{description}[leftmargin=1.8cm, style=nextline, itemsep=4pt]
    \item[$c$] Speed of light ($\approx 3\times10^{8}$ m/s).
    \item[$\lambda$] Wavelength of optical signal (meters).
    \item[$\nu$] Optical frequency (Hz).
    \item[$\varphi$] Measured phase shift in indirect ToF (radians).
    \item[$d_{\text{target}}$] One-way target distance (meters).
    \item[$p_{\text{class}}$] Softmax probability for class prediction.
    \item[$\text{logit}_{i}$] Raw (unnormalised) score for class $i$.
    \item[$K_c$] Number of classes in classification tasks.
    \item[$h$] Planck's constant.
    \item[$E$] Photon energy.
    \item[$\mu_i$] Expected photon count in bin $i$.
    \item[$\mu_{i,\lambda}$] Expected photon count in bin $i$ at wavelength $\lambda$.
    \item[$S$] Total signal photons.
    \item[$S_m$] Signal strength for material $m$.
    \item[$S_\lambda$] Total signal photons at wavelength $\lambda$.
    \item[$S_{m,\lambda}$] Signal strength for material $m$ at wavelength $\lambda$.
    \item[$\sigma$] Standard deviation of Gaussian profile.
    \item[$\sigma_m$] Standard deviation for material $m$.
    \item[$\sigma_\lambda$] Wavelength-dependent standard deviation.
    \item[$\sigma_{m,\lambda}$] Standard deviation for material $m$ at wavelength $\lambda$.
    \item[$t_0$] Peak arrival time.
    \item[$t_{0,m}$] Peak time for material $m$.
    \item[$t_{0,\lambda}$] Peak arrival time at wavelength $\lambda$.
    \item[$t_{0,m,\lambda}$] Peak time for material $m$ at wavelength $\lambda$.
    \item[$t_i$] Start time of bin $i$.
    \item[$\Delta t$] Histogram bin width (seconds).
    \item[$b$] Average ambient noise photons per bin.
    \item[$b_\lambda$] Wavelength-dependent ambient noise level.
    \item[$h_i$] Photon count in bin $i$.
    \item[$k$] Photon count value.
    \item[$M$] Number of composite materials or surfaces.
    \item[$N$] Number of pixels in detector array.
    \item[$B$] Number of time bins per pixel.
    \item[$K$] Spectral dimensionality (number of wavelength channels).
    \item[$W$] Number of spectral bands or image width (spatial dimension).
    \item[$H$] Image height (spatial dimension).
    \item[$C$] Number of channels (e.g., 3 for RGB).
    \item[$U$] Measured pixel intensity.
    \item[$U_\lambda$] Measured intensity at wavelength $\lambda$.
    \item[$p_m$] Fractional proportion of material $m$.
    \item[$r_m$] Reflectance of material $m$.
    \item[$r_{m,\lambda}$] Spectral reflectance of material $m$ at wavelength $\lambda$.
    \item[$\epsilon$] Error term capturing measurement uncertainties.
    \item[$\epsilon_\lambda$] Wavelength-dependent error term.
    \item[$T_B$] Temporal spread over time span.
    \item[$I_{\text{pix},i}$] Initial state of pixel $i$.
    \item[$I_{\text{sig},i}$] Signal intensity in pixel $i$.
    \item[$\mathbf{D}$] Array of pixel elements.
    \item[$\mathcal{O}(\cdot)$] Big O notation for computational complexity.
\end{description}%
}
}
\vspace*{\fill}
\newpage

% Symbols (page 12)
\newpage
\thispagestyle{empty}
\vspace*{-1.5cm}
\hfill\textcolor{gray}{\textit{\roman{page}}}
\vspace{0.2cm}
\noindent\rule{\textwidth}{0.5pt}
\vspace{0.5cm}
{\fontsize{12pt}{18pt}\selectfont
{\fontsize{14pt}{21pt}\selectfont\textbf{Symbols}}

\vspace{0.5cm}

{\setlength{\parindent}{0pt}%
\begin{description}[leftmargin=1.8cm, style=nextline, itemsep=4pt]
    \item[$\pi$] Ratio of a circle's circumference to its diameter.
    \item[$e$] Base of natural logarithms ($\approx 2.71828$).
    \item[$\sum$] Summation operator over indexed terms.
    \item[$\int$] Integral operator over a continuous variable.
    \item[$\exp(\cdot)$] Exponential function.
    \item[$\exp\left(-\frac{(t_i - t_0)^2}{2\sigma^2}\right)$] Gaussian exponential term.
    \item[$\alpha_{ji}^{\text{trg}}$] Binary indicator for target presence in loss terms.
    \item[$*$] Convolution operator.
    \item[$\approx$] Approximately equal to.
    \item[$\times$] Multiplication operator.
    \item[$/$] Division operator.
    \item[$\sqrt{\cdot}$] Square root operator.
    \item[$\cdot$] Multiplication dot operator.
\end{description}%
}
}
\vspace*{\fill}
\newpage

% Abbreviation (page 13)
\newpage
\thispagestyle{empty}
\vspace*{-1.5cm}
\hfill\textcolor{gray}{\textit{\roman{page}}}
\vspace{0.2cm}
\noindent\rule{\textwidth}{0.5pt}
\vspace{0.5cm}
{\fontsize{12pt}{18pt}\selectfont
{\fontsize{14pt}{21pt}\selectfont\textbf{Abbreviations}}

\vspace{0.5cm}

\renewcommand{\descriptionlabel}[1]{\hspace\labelsep\normalfont #1}
\setlength{\itemsep}{0pt}
\setlength{\parsep}{0pt}
\setlength{\topsep}{0pt}
\setlength{\partopsep}{0pt}
\setlength{\baselineskip}{18pt}
\begin{description}[leftmargin=2.5cm, style=nextline, itemsep=0pt, parsep=0pt, topsep=0pt, partopsep=0pt]
    \item[ADAS] Advanced Driving Assistance Systems
    \item[AM-CW] Amplitude-Modulated Continuous Wave
    \item[APD] Avalanche Photodiode
    \item[APS] Active Pixel Sensor
    \item[CCD] Charge-Coupled Device
    \item[CMOS] Complementary Metal-Oxide Semiconductor
    \item[CNN] Convolutional Neural Network
    \item[CW iToF] Continuous Wave Indirect Time-of-Flight
    \item[DLL] Delay-Locked Loop
    \item[DNN] Deep Neural Networks
    \item[D-SiPM] Digital Silicon Photomultiplier
    \item[dToF] Direct Time-of-Flight
    \item[EBCCD] Electron-Bombarded Charge-Coupled Device
    \item[EBCMOS] Electron-Bombarded Complementary Metal-Oxide Semiconductor
    \item[EB-MCP-CMOS] Electron-Bombarded Multichannel Plate CMOS
    \item[EMCCD] Electron-Multiplying Charge-Coupled Device
    \item[FMCW] Frequency-Modulated Continuous Wave
    \item[FM-CW] Frequency-Modulated Continuous Wave
    \item[GAN] Generative Adversarial Network
    \item[GS/s] Gigasamples per second
    \item[HAPD] Hybrid Avalanche Photodiode
    \item[HPD] Hybrid Photodetector
    \item[ICCD] Imaging Intensifying Charge-Coupled Device
    \item[ICMOS] Imaging Intensifying Complementary Metal-Oxide Semiconductor
    \item[IMNS] Institute for Integrated Micro and Nano Systems
    \item[iToF] Indirect Time-of-Flight
    \item[LiDAR] Light Detection and Ranging
    \item[LSTM] Long Short-Term Memory
    \item[MCP] Microchannel Plate
    \item[MaPMT] Multi-Anode Photomultiplier Tube
    \item[PAWD] Particle Wave Detection
    \item[p-iToF] Pulsed Indirect Time-of-Flight
    \item[PMF] Probability Mass Function
    \item[PMT] Photomultiplier Tube
    \item[PLL] Phase-Locked Loop
    \item[pnCCD] p-n type Charge-Coupled Device
    \item[PR] Pattern Recognition
    \item[RGB] Red Green Blue
    \item[RNN] Recurrent Neural Network
    \item[RO TDC] Ring Oscillator Time-to-Digital Converter
    \item[ROI] Region of Interest
    \item[SEM] Secondary Electron Multiplication
    \item[SiPM] Silicon Photomultiplier
    \item[SNR] Signal-to-Noise Ratio
    \item[SOTA] State-of-the-Art
    \item[SPAD] Single Photon Avalanche Diode
    \item[SRTM] Space Resolved Temporal Measurements
    \item[SSD] Single-Shot Detector
    \item[TAC] Time-to-Amplitude Converter
    \item[TCSPC] Time-Correlated Single-Photon Counting
    \item[TDC] Time-to-Digital Converter
    \item[TGSPC] Time-Gated Single-Photon Counting
    \item[ToF] Time-of-Flight
    \item[YOLO] You Only Look Once
    \item[YOLOv3] You Only Look Once version 3
    \item[YOLOv8] You Only Look Once version 8
    \item[DINOv3] DINO version 3
\end{description}
}
\vspace*{\fill}
\newpage

% Glossary of terms (page 14)
\newpage
\thispagestyle{empty}
\vspace*{-1.5cm}
\hfill\textcolor{gray}{\textit{\roman{page}}}
\vspace{0.2cm}
\noindent\rule{\textwidth}{0.5pt}
\vspace{0.5cm}
{\fontsize{12pt}{18pt}\selectfont
{\fontsize{14pt}{21pt}\selectfont\textbf{Glossary of terms}}

\vspace{0.5cm}

\renewcommand{\descriptionlabel}[1]{\hspace\labelsep\normalfont #1}
\setlength{\itemsep}{0pt}
\setlength{\parsep}{0pt}
\setlength{\topsep}{0pt}
\setlength{\partopsep}{0pt}
\setlength{\baselineskip}{18pt}
\begin{description}[leftmargin=3cm, style=nextline, itemsep=0pt, parsep=0pt, topsep=0pt, partopsep=0pt]
    \item[Time-gated single-photon counting (TGSPC)] A pulsed iToF implementation where photon arrivals are routed into discrete time windows for histogramming.
    
    \item[Transient histogram] A one-dimensional time-series of photon arrivals (counts versus time) capturing depth and material-dependent scattering.
    
    \item[Time-of-flight (ToF)] Measurement of distance by timing the propagation of light between a source and a target; implemented as direct (dToF) or indirect (iToF) variants.
    
    \item[Direct time-of-flight (dToF)] ToF variant that measures photon arrival times directly with fine temporal resolution (e.g., SPAD-based timing histograms).
    
    \item[Indirect time-of-flight (iToF)] ToF variant that estimates distance from phase shifts of modulated optical signals (e.g., AM-CW or FM-CW demodulation).
    
    \item[Light detection and ranging (LiDAR)] Active optical sensing that estimates range by measuring light returns; implemented with dToF or iToF schemes.
    
    \item[Convolutional neural network (CNN)] A deep learning architecture using convolutional layers to extract hierarchical features from images or signals.
    
    \item[Achromatic processing] A signal processing approach that does not consider the spectral composition of the signal, focusing instead on temporal or structural features rather than wavelength-dependent properties.
    
    \item[Advanced driving assistance systems] Electronic systems in vehicles that assist drivers in driving and parking functions, designed to increase vehicle safety and reduce human error.
    
    \item[Dual detection model] A combined detection framework that integrates time-resolved material classification with spatially resolved visual object detection models to improve inference quality and reliability.
    
    \item[Feature extraction] The process of identifying and selecting relevant characteristics or patterns from raw sensor data that can be used for machine learning and decision-making.
    
    \item[Logistic regression] A statistical machine learning method used for binary or multi-class classification problems, modelling the probability of a categorical outcome based on input features.
    
    \item[Machine perception] The ability of artificial systems to interpret and understand sensory information from the environment, enabling machines to make sense of their surroundings.
    
    \item[Machine vision] A technology and method used to provide imaging-based automatic inspection and analysis for applications such as automatic inspection, process control, and robot guidance.
    
    \item[Material detection] The process of identifying and classifying different types of materials based on their physical or optical properties using sensor data.
    
    \item[Material purity detection] The assessment and classification of material composition quality, particularly for fluid materials, to determine contamination levels or compositional variations.
    
    \item[Multi-label classification] A machine learning task where each instance can be assigned multiple labels simultaneously, as opposed to single-label classification where each instance belongs to only one class.
    
    \item[Particle Wave Detection] A custom framework developed in this work that guides how the SPAD instrument perceives and processes time-resolved signals for material detection and classification.
    
    \item[Safety critical applications] Systems or applications where failure could result in loss of life, significant property damage, or environmental harm, requiring high reliability and fault tolerance.
    
    \item[Single Photon Avalanche Diode (SPAD)] A photodetector capable of detecting individual photons with high sensitivity and temporal resolution, operating in Geiger mode where a single photon can trigger an avalanche current.
    
    \item[Spatially resolved images] Image data that captures two-dimensional spatial information, representing the distribution of light or other signals across a plane, typically encoded in formats such as RGB.
    
    \item[Time-resolved signal] A signal that contains temporal information, capturing how a measured quantity (such as photon arrival times) changes over time, enabling analysis of transient phenomena.
    
    \item[Transient image] A one-dimensional representation of time-resolved data, where the signal is captured as a function of time rather than spatial coordinates, used for temporal feature analysis.
    
    \item[Visual scene inference] The process of interpreting and understanding the content and structure of a visual scene from sensor data, enabling machines to make informed decisions about their environment.
    
    \item[Gaussian mixture model] A probabilistic model that represents a probability distribution as a weighted sum of multiple Gaussian distributions, used to model multi-material temporal signals in time-of-flight systems.
    
    \item[Photon pileup] A phenomenon where multiple photon arrivals result in only one detectable event, degrading performance under high illumination conditions by reducing the probability of detecting target-reflected signal photons.
    
    \item[Geiger mode] An operating regime of SPADs where a single photon can trigger a self-sustaining avalanche process, enabling single-photon sensitivity with high temporal resolution.
    
    \item[Secondary electron multiplication (SEM)] Techniques by which electrons or photo-electrons are multiplied and guided towards readout circuits following the photoelectric effect.
    
    \item[Spectral unmixing] The process of separating material contributions from mixed pixels in spectral imaging systems by solving systems of linear equations across multiple spectral bands.
    
    \item[Subsurface scattering] The phenomenon where photons penetrate into a material and scatter internally before re-emerging, causing temporal spread in photon arrival times.
    
    \item[Computational complexity] A measure of the computational resources (time and memory) required to execute an algorithm, typically expressed using big O notation.
    
    \item[Big O notation] Mathematical notation describing the limiting behaviour of a function when the argument tends towards a particular value or infinity, used to characterise algorithmic complexity.
    
    \item[Shot noise] Statistical fluctuations in photon counts due to the quantum nature of light, following a Poisson distribution.
    
    \item[Poisson distribution] A discrete probability distribution that expresses the probability of a given number of events occurring in a fixed interval of time or space.
    
    \item[Gaussian distribution] A continuous probability distribution characterised by a bell-shaped curve, used to model temporal signal profiles in time-of-flight systems.
    
    \item[Histogram binning] The process of discretising continuous temporal signals into discrete time periods (bins) for processing and analysis.
    
    \item[Time-to-digital converter (TDC)] Electronic circuitry that converts time intervals between events into digital values, enabling precise temporal measurement in SPAD systems.
    
    \item[Phase-locked loop (PLL)] A control system that generates an output signal whose phase is related to the phase of an input signal, used in multiphase clock generation for enhanced timing resolution.
    
    \item[Delay-locked loop (DLL)] A circuit similar to a phase-locked loop but used specifically for clock synchronization and delay compensation in timing circuits.
    
    \item[Ring oscillator] A device composed of an odd number of NOT gates whose output oscillates between two voltage levels, used in time-to-digital converters for timing measurements.
    
    \item[Multiphase clock] A clocking scheme where multiple clock phases are generated and distributed, providing enhanced timing resolution by recording the state of each phase during detection events.
    
    \item[On-chip histogramming] The process of generating temporal histograms directly within the detector pixel array, reducing readout bandwidth requirements compared to timestamp-based approaches.
    
    \item[Spectral processing] Signal processing approaches that explicitly resolve or model wavelength-dependent structure in signals, requiring processing of multi-dimensional spectral data.
    
    \item[Non-spectral processing] Signal processing approaches that treat signals as achromatic, operating only on aggregate quantities such as total intensity or photon-arrival statistics without explicit wavelength decomposition.
    
    \item[Intensity-based processing] Signal processing methods that capture the spatial distribution of light intensity, encoding information about material appearance, texture, and surface properties through spatial patterns.
    
    \item[Temporal signal processing] Signal processing methods that capture the time-resolved behaviour of light--matter interactions, encoding information about material composition, optical properties, and internal structure through temporal patterns.
    
    \item[Wave--particle duality] A fundamental principle of quantum mechanics describing how light and matter exhibit both wave-like and particle-like properties depending on the measurement context.
    
    \item[Photoelectric effect] The generation of photoelectric current resulting from incident radiation on a detector, which can be internal (electron-hole pair generation) or external (electron emission).
    
    \item[Impact ionisation] A process in semiconductor devices where high-energy charge carriers create additional electron-hole pairs through collisions, enabling signal amplification in avalanche photodiodes.
    
    \item[Quantum efficiency] The ratio of the number of charge carriers collected to the number of incident photons, characterising the sensitivity of a photodetector.
    
    \item[Signal-to-noise ratio (SNR)] A measure of signal strength relative to background noise, quantifying the quality of a signal for reliable detection and processing.
    
    \item[Ambient noise] Background photon counts from sources other than the target signal, including dark counts, stray light, and environmental illumination.
    
    \item[Dark counts] Spontaneous photon detections that occur in the absence of incident light, arising from thermal effects and other noise sources in photodetectors.
    
    \item[Reflectivity profile] The temporal distribution of reflected photons as a function of arrival time, encoding material-specific optical properties including surface reflection and subsurface scattering.
    
    \item[Time-resolved measurement] A measurement technique that captures temporal information about signal evolution over time, enabling analysis of transient phenomena and material-dependent optical responses.
    
    \item[Axial resolution] Resolution along the sensor's line of sight (depth dimension), enabled by time-of-flight measurements that directly measure photon propagation times.
    
    \item[Lateral resolution] Resolution perpendicular to the sensor's line of sight (spatial dimensions), typically provided by two-dimensional pixel arrays in imaging systems.
    
    \item[Peak clustering] The process of grouping temporal signal peaks based on similarity criteria, enabling material classification and scene segmentation in time-of-flight systems.
    
    \item[Region of interest (ROI)] A subset of image pixels or signal bins that are identified as containing relevant information for a specific vision task, such as object boundaries or material regions.
    
    \item[Supervised learning] A machine learning paradigm where models are trained on labelled data, learning to map inputs to known outputs through optimisation of a loss function.
    
    \item[Self-supervised learning] A machine learning paradigm where models learn representations from unlabelled data by creating supervisory signals from the data itself, such as predicting masked portions or different views of the same input, enabling effective feature learning without manual annotations.
    
    \item[Unsupervised learning] A machine learning paradigm where models identify patterns and structure in data without explicit labels or ground truth annotations.
    
    \item[Spatiotemporal fusion] The integration of spatial (geometric) and temporal (structural) information from different sensor modalities to enhance scene understanding, combining two-dimensional spatially resolved images with one-dimensional time-resolved signals for improved detection and classification performance.
    
    \item[Feature extraction] The process of identifying and selecting relevant characteristics or patterns from raw sensor data that can be used for machine learning and decision-making.
    
    \item[Softmax activation] A mathematical function that converts raw scores (logits) into probability distributions over multiple classes, ensuring probabilities sum to unity.
    
    \item[Dropout regularisation] A technique in neural network training where a fraction of neurons are randomly set to zero during training to prevent overfitting and improve generalization.
    
    \item[Convolutional layer] A neural network layer that applies learnable filters to detect local patterns such as edges, textures, and material-specific features in input data.
    
    \item[Pooling layer] A neural network layer that reduces spatial dimensions by aggregating values within local regions, providing translation invariance and reducing computational complexity.
    
    \item[Fully connected layer] A neural network layer where each neuron is connected to all neurons in the previous layer, enabling global decision-making and high-level feature combination.
    
    \item[Space Resolved Temporal Measurements (SRTM)] Time-resolved signals that are also spatially resolved across two-dimensional pixel arrays, enabling temporal analysis at multiple spatial locations simultaneously.
    
    \item[Lock-in amplification] A signal processing technique that extracts phase information from modulated signals by performing coherent demodulation, selectively extracting the phase-shifted signal component at the modulation frequency while suppressing noise and interference.
    
    \item[Coherent demodulation] A signal processing method that extracts phase information from modulated signals by synchronously mixing the received signal with a reference signal at the modulation frequency, enabling precise phase measurement.
    
    \item[Multi-path reflections] Optical phenomena where light signals travel along multiple paths between source and detector, causing signal corruption and ambiguity in time-of-flight measurements.
    
    \item[Unambiguous range] The maximum distance that can be measured without ambiguity in time-of-flight systems, determined by the modulation frequency or temporal sampling rate.
    
    \item[Pulsed indirect time-of-flight (p-iToF)] An indirect time-of-flight technique that uses pulsed illumination and measures phase shifts from the returned pulses, combining aspects of both direct and indirect ToF methods.
    
    \item[Continuous wave indirect time-of-flight (CW iToF)] An indirect time-of-flight technique that employs continuously modulated optical signals and lock-in amplification to extract phase information for distance estimation.
    
    \item[Amplitude-modulated continuous wave (AM-CW)] A continuous wave indirect time-of-flight technique where the radiation source is continuously modulated using waveforms such as pseudo-noise sequences, sinusoidal, or rectangular signals, with phase difference measured through demodulation circuitry.
    
    \item[Frequency-modulated continuous wave (FM-CW)] An indirect time-of-flight technique that estimates target depth by measuring the frequency difference between continuously emitted and returned signals, utilising continuous wave frequency modulation rather than pulsed signal transmission.
    
    \item[Probability mass function (PMF)] A function that gives the probability that a discrete random variable is exactly equal to some value, used to model photon count distributions in time-of-flight systems.
    
    \item[p-n type charge-coupled device (pnCCD)] A specialised charge-coupled device with p-n junction architecture, originally developed for spectroimaging applications in X-ray astronomy, exhibiting excellent sensitivity characteristics.
    
    \item[Multi-window integration] A phase measurement technique in amplitude-modulated continuous wave systems where the integration period is partitioned into multiple time windows to enable phase extraction and distance calculation.
    
    \item[Coherence length] The distance over which a light source maintains its phase relationship, required for frequency-modulated continuous wave systems to enable accurate frequency difference measurements.
    
    \item[Phase difference] The angular difference between transmitted and received waveforms in indirect time-of-flight systems, directly related to the propagation time and target distance.
    
    \item[Demodulation circuitry] Electronic circuits that extract phase or frequency information from modulated optical signals, enabling distance estimation in indirect time-of-flight systems.
    
    \item[Modulation frequency] The frequency at which the optical signal is modulated in indirect time-of-flight systems, determining the unambiguous range and measurement precision.
    
    \item[Phase offset] The angular displacement between different sampling windows in multi-window integration techniques, enabling reconstruction of the phase relationship between transmitted and received waveforms.
    
    \item[Range resolution] The minimum distance separation that can be distinguished in time-of-flight measurements, determined by temporal resolution, modulation frequency, and signal processing capabilities.
    
    \item[Multi-frequency modulation] A technique for extending the unambiguous range in indirect time-of-flight systems by using multiple modulation frequencies, enabling resolution of phase ambiguities.
    
    \item[Signal integration period] The time duration over which optical signals are accumulated in time-of-flight detectors, affecting signal-to-noise ratio and measurement accuracy.
\end{description}
}
\vspace*{\fill}
\newpage

% Main content starts here - switch to arabic numbering
\pagenumbering{arabic}
\setcounter{page}{1}
\arabicpagenumtrue  % Enable footer line for pages with integer page numbers

% Chapter 1: Introduction
\chapter{Introduction}
\label{chap:introduction}
\newpage

\section{Optical imaging}
\label{sec:optical_imaging}

Optical sensing for vision begins with the interaction of electromagnetic energy and matter; the resulting frequency, intensity, and polarization content is captured by an optical chain that couples sources, detectors, and signal-processing blocks to form a task-ready representation \cite{elachi2021,bass2010,schott2007}. In practice, most deployed systems still hinge on compact RGB imagers: they sample three narrow spectral bands, and subsequent digital pipelines extract features for detection or classification. This architecture is cost-effective and well understood, yet its sensing bandwidth and spatial-only sampling leave persistent blind spots—material ambiguity, illumination sensitivity, occlusion and clutter, and weak textural cues—especially when inference must be robust at low signal-to-noise ratios or under complex scene geometry \cite{sharan2013,schwartz2013}.

Spectroscopic instruments widen the spectral aperture but at the price of bulk and calibration overhead, while scanning lidars add geometric coverage with substantial complexity and power draw \cite{inagaki2020,becker2023}. To address these limits, this work pivots to direct time-of-flight (dToF) single-photon avalanche diode (SPAD) sensing, which measures photon arrival statistics to recover scene structure and material cues with picosecond timing fidelity. Figure~\ref{fig:1_1} illustrates the proposed processing flow.

\newpage
\begin{figure}[p]
\centering
\begin{tikzpicture}[
    node distance=2.5cm, 
    auto,
    every node/.style={font=\small},
    main node/.style={rectangle, draw=black!60, fill=gray!10, text width=3cm, align=center, minimum height=1cm, rounded corners=3pt, line width=0.8pt},
    item node/.style={rectangle, draw=black!60, fill=gray!10, text width=2.2cm, align=center, minimum height=1cm, rounded corners=3pt, line width=0.8pt},
    process node/.style={rectangle, draw=black!60, fill=gray!10, text width=2.5cm, align=center, minimum height=0.8cm, rounded corners=3pt, line width=0.8pt}
]
    % Define nodes vertically
    \node[main node] (scene) {Observed scene};
    \node[main node, below of=scene] (signal) {Measured signal};
    \node[main node, below of=signal] (pawd) {PAWD framework};
    
    % Draw arrows vertically with gaps (not touching boxes)
    \draw[->, line width=1.2pt, color=black!70] ([yshift=-0.3cm]scene.south) -- ([yshift=0.3cm]signal.north);
    \draw[->, line width=1.2pt, color=black!70] ([yshift=-0.3cm]signal.south) -- ([yshift=0.3cm]pawd.north);
    
    % Short vertical line from PAWD framework
    \draw[line width=0.8pt, color=black!60] (pawd.south) -- ++(0,-0.5cm) coordinate (pawd_bottom);
    
    % Horizontal line spanning 4 items
    \draw[line width=0.8pt, color=black!60] ([xshift=-4cm, yshift=0cm]pawd_bottom) -- ([xshift=4cm, yshift=0cm]pawd_bottom);
    
    % Four items arranged horizontally
    \node[item node, below of=pawd_bottom, xshift=-4cm, yshift=1.2cm] (item1) {Intensity\\spectral};
    \node[item node, below of=pawd_bottom, xshift=-1.3cm, yshift=1.2cm] (item2) {Intensity\\non-spectral};
    \node[item node, below of=pawd_bottom, xshift=1.3cm, yshift=1.2cm] (item3) {Temporal\\spectral};
    \node[item node, below of=pawd_bottom, xshift=4cm, yshift=1.2cm] (item4) {Temporal\\non-spectral};
    
    % Short vertical lines from horizontal line to each item
    \draw[line width=0.8pt, color=black!60] ([xshift=-4cm, yshift=0cm]pawd_bottom) -- ++(0,-0.5cm);
    \draw[line width=0.8pt, color=black!60] ([xshift=-1.3cm, yshift=0cm]pawd_bottom) -- ++(0,-0.5cm);
    \draw[line width=0.8pt, color=black!60] ([xshift=1.3cm, yshift=0cm]pawd_bottom) -- ++(0,-0.5cm);
    \draw[line width=0.8pt, color=black!60] ([xshift=4cm, yshift=0cm]pawd_bottom) -- ++(0,-0.5cm);
    
    % Second horizontal line below the 4 items
    \coordinate (items_bottom) at ([yshift=-1.2cm]item2.south);
    
    % Four vertical lines pointing upward
    \draw[line width=0.8pt, color=black!60] ([xshift=-2.8cm, yshift=0.95cm]items_bottom) -- ++(0,-0.8cm);
    \draw[line width=0.8pt, color=black!60] ([xshift=-0.1cm, yshift=0.95cm]items_bottom) -- ++(0,-0.8cm);
    \draw[line width=0.8pt, color=black!60] ([xshift=2.5cm, yshift=0.95cm]items_bottom) -- ++(0,-0.8cm);
    \draw[line width=0.8pt, color=black!60] ([xshift=5.2cm, yshift=0.95cm]items_bottom) -- ++(0,-0.8cm);
    
    % Horizontal lines below items (just below the bottom vertical lines)
    % First horizontal line connecting the first 2 bottom vertical lines
    \draw[line width=0.8pt, color=black!60] ([xshift=-2.8cm, yshift=0.1cm]items_bottom) -- ([xshift=-0.1cm, yshift=0.1cm]items_bottom);
    % Second horizontal line connecting the 3rd and 4th bottom vertical lines
    \draw[line width=0.8pt, color=black!60] ([xshift=2.5cm, yshift=0.1cm]items_bottom) -- ([xshift=5.2cm, yshift=0.1cm]items_bottom);
    
    % Items under first bottom horizontal line (spatial)
    \node[process node, below of=items_bottom, xshift=-1.45cm, yshift=-0.8cm] (spatial_features) {Spatial features};
    \node[process node, below of=spatial_features] (spatial_pattern) {Spatial pattern\\processing};
    \node[process node, below of=spatial_pattern] (spatial_decisions) {Spatial decisions};
    
    % Arrow from horizontal line to spatial features (matching gap between items)
    \draw[->, line width=1.2pt, color=black!70] ([xshift=-1.45cm, yshift=-0.2cm]items_bottom) -- ([yshift=0.3cm]spatial_features.north);
    % Arrows for spatial flow
    \draw[->, line width=1.2pt, color=black!70] ([yshift=-0.3cm]spatial_features.south) -- ([yshift=0.3cm]spatial_pattern.north);
    \draw[->, line width=1.2pt, color=black!70] ([yshift=-0.3cm]spatial_pattern.south) -- ([yshift=0.3cm]spatial_decisions.north);
    
    % Items under second bottom horizontal line (structural)
    \node[process node, below of=items_bottom, xshift=3.85cm, yshift=-0.8cm] (structural_features) {Structural features};
    \node[process node, below of=structural_features] (structural_pattern) {Structural pattern\\processing};
    \node[process node, below of=structural_pattern] (structural_decisions) {Structural decisions};
    
    % Arrow from horizontal line to structural features (matching gap between items)
    \draw[->, line width=1.2pt, color=black!70] ([xshift=3.85cm, yshift=-0.2cm]items_bottom) -- ([yshift=0.3cm]structural_features.north);
    % Arrows for structural flow
    \draw[->, line width=1.2pt, color=black!70] ([yshift=-0.3cm]structural_features.south) -- ([yshift=0.3cm]structural_pattern.north);
    \draw[->, line width=1.2pt, color=black!70] ([yshift=-0.3cm]structural_pattern.south) -- ([yshift=0.3cm]structural_decisions.north);
    
    % Connect the last two items (spatial decisions and structural decisions)
    % Two vertical lines going up from each decision box
    \draw[line width=0.8pt, color=black!60] ([yshift=-0.3cm]spatial_decisions.south) -- ++(0,-0.5cm) coordinate (spatial_decisions_bottom);
    \draw[line width=0.8pt, color=black!60] ([yshift=-0.3cm]structural_decisions.south) -- ++(0,-0.5cm) coordinate (structural_decisions_bottom);
    
    % Horizontal line connecting the two vertical lines
    \draw[line width=0.8pt, color=black!60] (spatial_decisions_bottom) -- (structural_decisions_bottom);
    
    % Vertical line from centre of horizontal line (single continuous line, 0.5cm like first vertical line)
    % Create coordinate at midpoint: spatial at x=-1.45cm, structural at x=3.85cm, midpoint at x=1.2cm
    \path (spatial_decisions_bottom) -- (structural_decisions_bottom) coordinate[midway] (horizontal_mid);
    \draw[line width=0.8pt, color=black!60] (horizontal_mid) -- ++(0,-0.5cm) coordinate (spatiotemporal_top);
    
    % Spatiotemporal decision item (positioned below the vertical line)
    \node[process node, below of=spatiotemporal_top, yshift=0.2cm] (spatiotemporal_decision) {Spatiotemporal decision};
    
    % Arrow from vertical line to spatiotemporal decision (single continuous arrow)
    \draw[->, line width=1.2pt, color=black!70] (spatiotemporal_top) -- (spatiotemporal_decision.north);
\end{tikzpicture}
\caption[Flow chart showing the processing pipeline from observed scene through measured signal to the PAWD framework.]{\textbf{Flow chart showing the processing pipeline from observed scene through measured signal to the PAWD framework.} Intensity measured signal is often spatially resolved to represent scenes. Nevertheless, time resolved signal may also be spatially resolved. This is termed space resolved temporal measurements (SRTM). SRTM are not useful for structural decisions. However, they may play key roles in spatial decisions like pose estimation, scale estimation and others. A true spatiotemporal solution is one that combines the crude time resolved signal with intensity resolved signal for decision making.}
\label{fig:1_1}
\end{figure}
\FloatBarrier

The PAWD framework partitions sensing and processing into four regimes: intensity-spectral, intensity non-spectral, temporal-spectral, and temporal non-spectral. Intensity-spectral systems (e.g., RGB and hyperspectral imagers) operate on spatially resolved, wavelength-dependent measurements. Intensity non-spectral systems use spatial intensity without explicit spectral separation (e.g., greyscale or infrared cameras). Temporal-spectral systems resolve both photon arrival time and wavelength, while temporal non-spectral systems, such as the SPAD dToF instruments in this work, operate on one-dimensional time-resolved histograms from a monochromatic source. The framework demonstrates how spatial (geometric) and structural (material and depth) cues can be combined to form spatiotemporal decision pipelines.

% Quote at end of section 1.1
\vspace{1cm}
\noindent\begin{minipage}{\textwidth}
\setlength{\parindent}{0pt}%
\setlength{\leftskip}{0pt}%
\setlength{\rightskip}{0pt}%
\setlength{\emergencystretch}{3em}%
\setstretch{2.0}%
\footnotesize%
\ttfamily%
\color{textgrey}%
\justifying%
\sloppy%
\noindent\quoteone
\end{minipage}

\newpage
\section{SPAD Sensors}
\label{sec:spad_sensors}

SPAD direct time-of-flight imagers are now widely deployed because compact, low-cost devices deliver picosecond timing and single-photon sensitivity. Current ToF systems span consumer cameras, medical imaging, robotics, and industrial automation. The return signal could be used for true depth estimation, material discrimination, pose estimation, and several others. Distance, surface finish, subsurface scattering, detector angle to target surface normal and other factors may affect optical properties, leading to distinct photon-arrival signatures (Figure~\ref{fig:1_2}).

\begin{figure}[H]
\centering
\includegraphics[width=0.9\textwidth,height=0.7\textheight,keepaspectratio]{figures/1.2.png}

    \caption[Comparison of signal reflectivity profiles for translucent and opaque materials.]{\textbf{Comparison of signal reflectivity profiles for translucent and opaque materials.} Comparison of signal reflectivity profiles (photon time of arrival histograms obtained with SPAD) for translucent material (top) and opaque material (bottom). The observed curves indicate that positioned at the same distance, the translucent material presents a longer tail compared to the opaque material. This is resulting from the delayed arrival of sub-surface scattered photons from the translucent material.}
\label{fig:1_2}
\end{figure}

The distinct photon-arrival signatures shown in Figure~\ref{fig:1_2} demonstrate how translucent materials exhibit longer tails due to delayed arrival of subsurface-scattered photons compared to opaque materials at the same distance. This temporal signature difference enables material discrimination through time-resolved measurements.

Optical material sensing estimates composition from light–matter interactions captured by instruments such as SPADs \cite{elachi2021}. Under passive or active illumination, incident radiation induces surface reflection, absorption, and subsurface scattering; the relative mix depends on atomic structure and finish, yielding opaque, translucent, transparent, or fully absorbing responses \cite{elachi2021}. Active illumination sources deliver controlled, directional light that enhances signal-to-noise ratios and enables improved material discrimination (Figure~\ref{fig:1_3}).

The illumination source classification illustrated in Figure~\ref{fig:1_3} shows the distinction between passive and active illumination, as well as external versus internal source placement for remote material sensing. Active illumination provides controlled, directional light that enhances signal quality and material separability compared to passive illumination systems.

\begin{figure}[H]
\centering
\includegraphics[width=\textwidth]{figures/1.3.png}

    \caption[Visual illustration of remote material detection setup showing passive, active, external and internal sources.]{\textbf{Visual illustration of remote material detection setup showing passive, active, external and internal sources.}}
\label{fig:1_3}
\end{figure}

Surface topography further shapes angular reflection: rough interfaces yield diffuse lobes, while polished surfaces produce specular peaks; most materials are mixed and wavelength dependent \cite{elachi2021}. These interactions generate characteristic reflectivity profiles that SPAD receivers capture as time-resolved histograms or transient images. Detectors typically demodulate or digitise signals via ADC or TDC stages, then forward intensity or photon time-of-arrival data for downstream processing and readout.

Complete sensing stacks combine illumination (laser or LED; infrared or multispectral; passive or active) with depth-estimation optics (time-of-flight or stereo), detectors (photodiode, SPAD, APD), and task-specific processing. Spatial imagers (e.g., CMOS photodiodes) recover geometry, while impact-ionisation devices such as SPADs capture temporal signatures that reveal material-dependent optical paths. Chapters 4–6 apply these principles: transient-only material classification, hybrid spatio-temporal object detection, and purity assessment of consumable fluids. Related studies report similar gains when combining spatial and time-resolved cues for robust material and scene understanding \cite{becker2023,elachi2021}.

\newpage

\section{Computer vision}
\label{sec:computer_vision}

Computer vision extracts actionable structure from images by combining spatial features with scene semantics. Modern RGB pipelines deliver fast detection and classification, yet purely spatial cues struggle with material ambiguity, lighting variation, occlusion, and adversarial perturbations \cite{sharan2013,schwartz2013,rani2025}. Applications such as food safety or autonomous navigation require knowledge of what objects are made of—not only where they are—so decisions depend on composition and authenticity rather than appearance alone \cite{sharan2013,chen2024}.

Time-resolved sensing adds structural evidence. Direct ToF SPAD measurements capture photon-arrival statistics that encode depth, subsurface scattering, and other material-dependent effects, providing complementary features to RGB streams without the cost of full spectroscopy. Integrating these transients with spatial detectors improves robustness to spoofing and low-light conditions and reduces ambiguity in perspective or texture (Figure~\ref{fig:1_1}) \cite{chen2024}. Real-time detectors such as YOLO remain strong for spatial tasks \cite{redmon2015,redmon2018}, but fusing SPAD-derived structural cues addresses their susceptibility to purely appearance-based adversarial inputs \cite{rani2025}. Chapter 5 adopts that fusion approach: spatial networks supply localisation and semantics, while SPAD transients supply material and depth cues for resilient perception.

\subsection{Machine perception}
\label{subsec:machine_perception}

Machine perception has advanced from early electronic computing engines \cite{hartree1946} through modern large language models, yet it still lacks the general, materially grounded intelligence seen in biological systems \cite{chollet2019}. Scaling compute and model depth has not resolved the core limitation: current pipelines often operate on appearance-only cues, leaving them brittle to distribution shift and adversarial artefacts \cite{rani2025}. As LeCun argues, even a household animal exhibits richer world understanding than today’s largest models, underscoring the gap between statistical pattern fitting and embodied perception \cite{mims2024}.

Closing that gap requires better signal acquisition, not just larger architectures. When sensing captures only spatial texture, systems cannot distinguish genuine objects from convincing replicas or infer physical properties; adding structural cues such as direct ToF SPAD measurements provides the temporal and material evidence needed for robust decisions. Reliable machine perception therefore depends on sourcing informative measurements before scaling network complexity, so that learning operates on signals that encode the physical world rather than its superficial projections \cite{rani2025}.

% Quote at end of subsection 1.3.1
\vspace{1cm}
\noindent\begin{minipage}{\textwidth}
\setlength{\parindent}{0pt}%
\setlength{\leftskip}{0pt}%
\setlength{\rightskip}{0pt}%
\setlength{\emergencystretch}{3em}%
\setstretch{2.0}%
\footnotesize%
\ttfamily%
\color{textgrey}%
\justifying%
\sloppy%
\noindent\quotetwo
\end{minipage}

\subsection{Classical models of machine vision}
\label{subsec:classical_models}

Early theoretical frameworks in visual perception established foundational principles but lacked mechanisms for validating signal integrity. Gestalt psychology, emerging in the 1920s, proposed holistic perception where the whole exceeds the sum of parts through integrated top-down processing \cite{palmer1999}. However, this approach was criticised for its absence of neural processing foundations and mathematical rigour. Half a century later, Marr's computational theory of vision introduced a bottom-up modular framework, proposing specialised components for distinct visual tasks including edge detection, stereo vision, motion detection, object recognition, and colour processing \cite{marr2010}. Marr's theory, while influential, faced criticism for its rigid sequential structure and lack of clear neural mechanisms.

Without strong mathematical foundations, neither framework provided a complete explanation for vision, motivating the development of biologically inspired artificial neural networks. Machine learning approaches introduced testable models for both learning (neural networks) and classification (probabilistic methods) \cite{lecun1998,fukushima1980}. While neural networks can extract features from inputs, they lack guiding principles for validating input data quality as accurate representations of target scenes. In contrast, biological vision systems can perceive spatial, texture, and depth information whose quality depends on the perceiver's experience and cognitive capabilities, enabling effortless discrimination between authentic materials (e.g., stone versus glass) and their representations (e.g., cardboard images).

For decades, vision models have lacked such validation mechanisms, compromising inference quality and system robustness in visual machine learning models and pattern recognition applications \cite{sharan2013,schwartz2013}. This fundamental limitation means that even sophisticated neural architectures cannot distinguish between authentic physical objects and their representations, leading to potential failures in safety-critical applications. These challenges manifest as persistent unsolved problems in state-of-the-art vision literature: low-level detection limitations, perspective ambiguities, occlusion handling, and other measurable model constraints \cite{inagaki2020,edozie2025,depthpro2024}. The absence of mechanisms to validate signal integrity and distinguish real materials from representations represents a critical gap that must be addressed before vision systems can achieve the reliability required for autonomous operation in complex real-world environments \cite{rani2025}.


% Quote at end of section 1.3.2
\vspace{1cm}
\noindent\begin{minipage}{\textwidth}
\setlength{\parindent}{0pt}%
\setlength{\leftskip}{0pt}%
\setlength{\rightskip}{0pt}%
\setlength{\emergencystretch}{3em}%
\setstretch{2.0}%
\footnotesize%
\ttfamily%
\color{textgrey}%
\justifying%
\sloppy%
\quotethree
\end{minipage}

\subsection{Ambiguous challenges of machine vision}
\label{subsec:ambiguous_challenges}

Vision models face criticism for relying on pattern matching rather than true scene understanding, yet this critique may be misplaced. Pattern matching is not a limitation but a fundamental mechanism: biological vision systems, which inspire these models, perform sophisticated pattern matching operations \cite{marr2010,lecun1998,fukushima1980}. Sensory inputs are processed into learned patterns for storage and retrieval, enabling recognition, classification, and decision-making. The distinction between effective and ineffective pattern matching lies not in the mechanism itself, but in the richness and integrity of the input signals from which patterns are extracted. When vision systems rely solely on spatial appearance features from two-dimensional images, they learn patterns that may not capture fundamental structural properties of objects and materials, leading to limitations in scene understanding. The reliability of decisions derived from extracted patterns depends on whether those patterns capture meaningful, material-discriminative information or merely superficial appearance characteristics that can be easily deceived by visual illusions or representations. Spatial-only imagery forces models to infer depth, material, and structure from appearance proxies, making decisions fragile under illumination changes, occlusion, or synthetic spoofs \cite{sharan2013,schwartz2013}. The decisive factor is signal quality: discriminative patterns emerge when the input captures the underlying physics of the scene, not just its 2D appearance. Time-resolved transients from SPAD-based direct time-of-flight sensing add such physically grounded cues, exposing subsurface scattering and true range while remaining compatible with efficient pattern-recognition pipelines. By fusing spatial appearance with temporal structure (see chapter 5), vision systems resolve ambiguities that purely spatial features cannot, enabling more reliable material and scene level understanding.

\subsection{Low-level detection challenges}
\label{subsec:low_level_detection}

In biological systems, sensory loss serves as a primary indicator of cognitive decline, often correlated with ageing but also resulting from disease or disability \cite{bernabei2014,dammeyer2010}. This manifests as diminished perceptual capacity: hearing less information than is communicated, perceiving less representation than the actual scene, or experiencing reduced tactile sensitivity. In machine vision systems, an analogous sensory loss occurs when detectors receive only spatially resolved input and are tasked with constructing complete scene representations. Under such constraints, the system must synthesise all relevant features—including structural features such as true depth, material composition, and subsurface scattering—from incomplete measurements. This fundamental limitation creates significant sensory loss, as actual time-resolved structural features are absent from the input, leading to elevated processing costs and representation quality overhead.

A machine vision system's inability to correctly and fully receive all relevant features of the target leads to signal losses that manifest in the quality of decisions made, resulting in less intelligent inferences. The fundamental issue lies in the system's reliance on incomplete information, where missing structural features must be inferred or approximated rather than directly measured, creating a fundamental gap between perceived and actual scene properties. This gap becomes particularly problematic in safety-critical applications where accurate scene understanding is essential, as the system operates on assumptions rather than direct measurements of material properties and structural characteristics. The consequences of this sensory loss extend beyond individual detection errors to systemic failures in autonomous systems, where incorrect material identification or depth estimation can cascade into catastrophic decision-making failures. Unlike biological systems that can adapt and compensate for sensory degradation through experience and contextual reasoning, machine vision systems lack this adaptive capability when operating on fundamentally incomplete sensory data.

Stereoscopic depth estimation exemplifies this limitation: it produces synthetic values from incomplete three-dimensional information, triangulating two two-dimensional intensity images rather than measuring true depth directly \cite{boyer1989}. Though evaluation of vision algorithms is primarily based on statistical metrics, these metrics can be misleading under practical circumstances \cite{vanwinckelen2012}. The susceptibility of modern vision systems to limited low-level understanding and adversarial examples underscores such arguments \cite{rani2025}. When a system has no mechanism to self-validate the quality of data received as input, no level of testing can correct or ensure reliability, as the system cannot distinguish between accurate and inaccurate representations of the scene. Accurate signal and feature sourcing represents a major challenge in vision models, as it implicates higher-level limitations such as general intelligence, where the foundation of understanding must be built upon reliable, directly measured features like true depth rather than synthesised approximations. Direct time-of-flight sensing addresses this root cause: single-photon histograms capture absolute range and material-scattering signatures, supplying the structural features that downstream vision models cannot safely infer.

\subsection{Role of SPADs in addressing machine vision challenges}
\label{subsec:spads_machine_vision}

Machine vision presents a complex challenge, yet human cognition provides key indicators for evaluating intelligent behaviour. This work positions SPAD sensing as an instrumental approach that enables visual systems to satisfy one such indicator: the capacity to infer material composition beyond mere spatial localisation. This capability is fundamental to scene understanding, enabling systems to distinguish authentic materials from their representations and to reason about spatial relationships, including foreground and background context. A complementary indicator is the ability to measure true distance to targets and their surrounding objects, providing geometrically grounded scene structure.

The fundamental role of SPAD detectors in machine vision addresses visual intelligence from an instrumental perspective. While establishing a universal definition of intelligence remains contentious \cite{minsky1986,chollet2019}, a contextual framework provides clarity: if humans would not consider a task to require intelligence, then a machine's execution of that task may similarly fail to demonstrate intelligence. For instance, when humans cannot differentiate between a drawing of food and actual food, this indicates learning difficulties or perceptual disorders. When machines exhibit the same limitation—failing to distinguish objects at this fundamental level despite high processing speeds—they may not be considered intelligent. This highlights the gap between computational performance and genuine perceptual understanding, where speed and pattern recognition accuracy do not equate to intelligent behaviour.

Driven by the goal of more intelligent machine perception, this work employs both spatial and time-resolved optical signals to enhance scene understanding capabilities. The integrated instrument combines a single-photon avalanche diode (SPAD) sensor with near-infrared laser illumination and an RGB-encoded CMOS active pixel sensor (APS). The SPAD sensor provides structural features through temporal information that reveals material properties and true depth relationships inaccessible to spatial intensity images alone. This combined modality enables acquisition of high-quality structural features that support improved three-dimensional vision, bridging the gap between computational processing capabilities and genuine understanding of the physical world.

% Quote at end of subsection 1.3.4
\vspace{1cm}
\noindent\begin{minipage}{\textwidth}
\setlength{\parindent}{0pt}%
\setlength{\leftskip}{0pt}%
\setlength{\rightskip}{0pt}%
\setlength{\emergencystretch}{3em}%
\setstretch{2.0}%
\footnotesize%
\ttfamily%
\color{textgrey}%
\justifying%
\sloppy%
\quotefour
\end{minipage}

\newpage
\section{Aims and Objectives}
\label{sec:aims_objectives}

This work aims to develop compact sensing and signal-processing systems that enhance machine vision capabilities without requiring bulky or expensive instrumentation. Guided by a custom framework, the project integrates single-photon avalanche diode (SPAD) direct time-of-flight (dToF) sensor modules with conventional RGB sensors to enable material detection—the ability to infer material composition of target objects from their optical properties. The solutions address fundamental challenges outlined in Chapter~\ref{chap:introduction} by improving low-level visual analysis, enabling robots and autonomous systems to achieve more accurate environmental understanding and make more reliable decisions for safety-critical applications including autonomous surgical and mobility assistance, trustworthy perception in autonomous vehicles, and security surveillance.

The approach leverages complementary strengths of temporal and spatial information: SPAD sensors capture time-resolved signals that encode material-specific optical properties, while RGB sensors provide spatial context and appearance information. This spatiotemporal fusion enables comprehensive scene understanding that extends beyond traditional vision systems limited to spatial appearance features alone. Three key research areas are explored: long-range material detection from planar surfaces (Chapter~\ref{chap:flat_homogeneous_material_detection}), liquid composition assessment for food safety applications (Chapter~\ref{chap:material_purity_detection}), and improved object detection through combined transient signals from SPAD sensors and two-dimensional intensity-resolved signals from active pixel sensors (Chapter~\ref{chap:spatiotemporal_object_detection}). Each solution demonstrates the practical viability of compact, cost-effective SPAD-based systems for addressing fundamental limitations in current machine vision approaches \cite{sharan2013,schwartz2013}.

Objectives (signal-processing focus):
\begin{enumerate}[leftmargin=*]
    \item \textbf{Long-range material detection:} Classify planar solids from transient profiles across longer arbitrary continuous target distances, various orientations, and view angles exceeding fixed-range dToF classifiers and maintaining robustness outside laboratory geometries 
    \item \textbf{Spatiotemporal object detection:} Combine one-dimensional SPAD transients with two-dimensional APS images to supply structural cues that RGB-alone detectors miss, improving discrimination under spoofing, low light, or visual ambiguity
    \item \textbf{Liquid purity assessment:} Detect contaminants in homogenised milk via temporal signatures without spectral decomposition, offering a compact, non-invasive alternative to conventional chemical or spectroscopic assays 
\end{enumerate}

Application domains include autonomous navigation and manipulation, surgical or mobility assistance, secure surveillance, and food-quality monitoring, where materially informed decisions are required under constrained size, weight, and power budgets.

\subsection{Overview of SPAD Applications}
\label{subsec:spad_applications_overview}
Three experimental projects demonstrate the practical impact of SPAD time-of-flight sensors as cost-effective, compact instruments that deliver efficient solutions to fundamental machine vision challenges. These investigations establish the role of time-resolved signals in constructing intelligent machine vision systems, validating temporal information as a viable alternative to expensive spectral instrumentation while achieving comparable or superior performance in material characterisation tasks. The solutions address core limitations in current vision systems: the capacity to distinguish authentic materials from their representations, and the extraction of structural features that remain inaccessible through spatial intensity images alone. By integrating spatial and temporal information, this work establishes a foundation for more reliable and comprehensive scene understanding in safety-critical applications such as autonomous systems and robotic vision. A brief overview of each project follows.

\subsection{SPAD for flat solid surface detection}
\label{subsec:spad_flat_solid}
This work presents one-dimensional time-resolved signals from SPAD sensors for detecting material composition of flat homogeneous solid surfaces. Controlled target samples enable material feature extraction from geometrically simple objects, facilitating robust classification. Previous research has addressed similar tasks \cite{tachella2019,becker2023}, yet significant improvements are demonstrated through the experimental validation reported here. Most dToF material detection systems require fixed target positions, which limits classification quality when target distance varies continuously. The presented instrument operates at extended ranges (up to 4\,m), enabling continuous distance variation that exceeds the short-range detection capabilities of prior work. Additionally, many existing material detection systems were developed under highly controlled laboratory conditions, limiting their applicability to realistic everyday scenarios. This approach addresses these limitations by enabling material detection across varying distances without requiring fixed target positions, thereby improving practical applicability. The system's ability to maintain detection accuracy at extended ranges (up to 4\,m) demonstrates robustness for real-world deployment where target distances cannot be precisely controlled. Experimental validation under less controlled conditions further distinguishes this work from previous approaches that relied heavily on laboratory settings with fixed geometries and distances. Details of the experiments are presented in Chapter~\ref{chap:flat_homogeneous_material_detection}.

\subsection{SPAD for improved machine vision}
\label{subsec:spad_improved_vision}
SPAD direct time-of-flight (dToF) sensors, combined with active pixel sensor (APS) imagers, address fundamental limitations in machine vision systems. Vision systems are frequently challenged by low-level object detection failures, occlusion, perspective ambiguity, and other shortcomings that compromise decision-making reliability. Accurate three-dimensional perception requires correctly sourced and processed spatial and structural features. However, for several decades, most vision systems have relied exclusively on two-dimensional spatially resolved intensity images (e.g., RGB-encoded) to infer structural, depth, and spatial features \cite{marr2010}. This protocol creates a fundamental sensory deficit, leading to impoverished representations at both pixel and latent space levels.

This work addresses this limitation by extracting surface and subsurface scattering properties (optical signatures) from targets using single-photon avalanche diode (SPAD) sensors. The instrument generates distinct source features: spatial patterns (shape, geometric properties), depth (distance), and structural patterns (texture, surface roughness). The temporal information captured by SPAD sensors provides material-specific signatures that are fundamentally distinct from and complementary to spatial intensity patterns, enabling machines to distinguish between real materials and their representations, detect material composition, and understand structural properties inaccessible through spatial images alone \cite{schwartz2013,sharan2013,tafone2023}.

A combination of transient and RGB images was evaluated using state-of-the-art vision models. Experimental results demonstrate that even without sophisticated architectures or high processing power, machines exhibited enhanced perceptual understanding compared to some of the largest contemporary visual interactive models. The integration of temporal signals with spatial images enables more reliable decision-making, particularly in scenarios where spatial information alone is ambiguous or insufficient, such as distinguishing real objects from photographs or detecting material properties that are not visually apparent \cite{yang2024}. These findings suggest that SPAD time-of-flight detectors may be indispensable for building truly intelligent machines, as they provide the structural and depth information necessary for comprehensive scene understanding that current spatial-only approaches cannot achieve.

\subsection{SPAD for material purity detection}
\label{subsec:spad_purity_detection}
The final application employs transient signals from homogenised milk to assess liquid purity based on manufacturer-defined safety metrics for consumable products. Conventional contamination detection methods rely primarily on chemical assays, which are invasive, or on optical techniques that require complex instrumentation; spectrophotometric approaches, in particular, exhibit limited reliability. This work demonstrates that time-resolved SPAD transients capture microscopic scattering changes caused by milk adulterants, enabling non-invasive purity assessment without chemical or spectroscopic assays \cite{nielsen2013,gregory1961}. A supervised neural network classifier learns material-discriminative temporal patterns from transient signals, trained on ground-truth purity states collected using the instrument. The classifier enables accurate distinction between pure and contaminated milk samples without requiring spectral decomposition, generalising across different manufacturers, fat compositions, and temporal variations. Eight comprehensive experimental tests investigate purity measurements using matter-based contaminants categorised through prior internal studies, including invisible adulterants such as hydrogen peroxide that cannot be detected through visual inspection. These tests validate the system's ability to generalise under diverse conditions, demonstrating robust performance across manufacturer variations, fat content differences, and temporal drift. The solution provides a reliable, cost-effective, and accessible alternative to conventional purity measurement methods. This accessibility empowers consumers to perform their own milk adulteration screening, promoting food safety and enabling proactive quality control at the consumer level. Chapter 6 describes the details of each experimental test and the resulting model performance metrics.

\newpage
\section{Thesis Contributions}
\label{sec:thesis_contributions}

This thesis presents material detection research using single-photon avalanche diode (SPAD) sensors, delivering three primary objectives alongside supporting achievements including a custom theoretical framework, web-based deployment applications, a specialised data format (\texttt{.sto}), and complementary analytical studies. The following sections outline the key contributions.

\subsection{Objective 0 (Chapter 3): Particle Wave Detection framework}
The Particle Wave Detection (PAWD) framework categorises signal processing into four distinct regimes: spatial-spectral, spatial non-spectral, temporal-spectral, and temporal non-spectral. This classification distinguishes temporal non-spectral processing (direct time-of-flight) from spectrally decomposed techniques such as RGB or multispectral transient methods, positioning temporal non-spectral approaches as an independent research domain. The framework guides all experimental work, including subsurface scattering modelling (Figure~\ref{fig:1_1}) and detection of objects with non-planar multimaterial surfaces in Chapter~5, providing a structured foundation for understanding how structural cues in transients complement spatial cues.

\subsection{Objective 1 (Chapter 4): Flat homogeneous material detection}
This work presents material detection for flat homogeneous surfaces through modelling of subsurface scattering and reflectivity profiles of returned signals from planar objects positioned at arbitrary distances from the detector. Given a target of interest, the model extracts target material identity, target distance, foreground object material and distance, and background object material and distance. Training across varied standoff distances, pose orientations, and view angles enables robust classification beyond fixed laboratory configurations \cite{becker2023}. A Python Flask web application provides a user-friendly interface for testing and inference, enabling practical deployment and evaluation of the material detection system.

\subsection{Objective 2 (Chapter 5): Spatiotemporal object detection}
This project presents the first reported combination of two-dimensional spatially resolved images with one-dimensional direct time-of-flight (dToF) time-resolved signals of the same target for object detection. Unlike RGB-D systems that estimate depth via stereoscopy or indirect time-of-flight, this solution employs dToF signals (photon timing histograms) that provide more than depth: they capture optical properties useful for inferring material composition. A hybrid neural network combines time- and space-resolved signals, with object localisation and label features derived from spatially resolved images and structural features extracted from time-resolved signals. Experimental results demonstrate that even with narrow-band or single-wavelength achromatic signals, machines achieve improved structural and spatial understanding when structural or depth features originate from transient signals and topological features from spatially resolved images. Since extracting structural features from spatially resolved images is complex and has reported limitations, this work uses SPAD one-dimensional time-resolved signals (transients) to extract features from surface reflectivity and subsurface scattering to characterise material composition \cite{sharan2013,schwartz2013,yang2024}. Due to the model's novelty, a custom Python pickle object with \texttt{.sto} extension was developed to store spatial and structural features in a serialised, immutable format for efficient loading; details are provided in Section~1.4.3.1. Additionally, a Python Flask web application provides a user-friendly interface for testing and inference of the spatiotemporal object detection system.

\subsection{Objective 3 (Chapter 6): Milk purity detection}
Material purity detection of homogenised fluid (consumable milk) using SPAD time-of-flight sensors represents a practical application of single-photon detection technology for food safety assessment. The homogenisation process ensures uniform distribution of fat globules and particles throughout the milk, eliminating temporal instability issues that would otherwise make purity measurement unreliable, as the liquid nature of inhomogeneous milk causes particle configurations to change over time, making ground truth labels inconsistent when the liquid is shaken or moved. The SPAD sensor captures time-resolved transient signals from homogenised milk, where each photon's arrival time encodes information about the material's optical properties, scattering characteristics, and internal composition. When contaminants are present, they alter the temporal response signature of the milk, creating detectable differences in transient image profiles that can be distinguished from pure milk samples through machine learning classification. The system's ability to detect invisible contamination, such as hydrogen peroxide, demonstrates its effectiveness for real-world food safety applications where visual inspection fails. The approach offers significant advantages over traditional purity assessment methods: it is non-invasive, requiring no sample extraction or container opening; it is compact and low-cost compared to bulky optical spectroscopy setups; it provides rapid detection with processing times ranging from 1--5\,ms; and it generalises across different manufacturers and fat compositions, as validated through experiments sourcing milk products from various manufacturers, different fat contents, and across different days. The temporal signals capture material-specific information that cannot be easily replicated in images or synthetic samples, enabling the system to distinguish between pure consumable milk and contaminated or adulterated products, qualifying as a simple, cost-effective means to empower average non-technical consumers to test food safety while encouraging food safety practices and building consumer trust.

Prior to modelling time-resolved signals for milk impurity detection, classifying different contaminants was challenging at both instrumental and material levels. Chemical purity detection methods are invasive, and existing optical approaches are either less efficient (visible light) or complex and bulky to set up. At the material/contaminant level, the challenge stems from the infinite ways matter can combine. To address these issues, impurities were classified based on matter states, which proved useful in achieving the objective (see Chapter~6 for details). A web-based application for milk purity detection facilitates user-friendly testing and inference through a Python Flask web application developed for material detection tasks.

\subsection{Implications for machine intelligence}
Prominent scientists have described the difficulty of establishing a single definition of machine intelligence, agreeing that intelligent machines perform tasks that would require intelligence if done by humans \cite{minsky1986,chollet2019}. Differentiating between material composition of objects is considered an intelligent task by humans, requiring sophisticated perceptual and cognitive abilities to distinguish materials based on their physical properties, optical characteristics, and structural signatures. For the first time in literature, this work posits and demonstrates that without complex processing infrastructure or learning architecture, structural features extracted from time-resolved signals improve machines' ability to differentiate materials, which contributes significantly to their reasoning, cognition quality, and consequently their intelligence. This finding suggests that fundamental improvements in machine intelligence may be achieved not only through increasingly sophisticated learning architectures and computational resources, but also through more effective extraction and utilisation of structural information from appropriate signal modalities, establishing a pathway toward more intelligent vision systems that can reason about material properties and make decisions based on structural understanding rather than appearance alone.

\newpage
\section{Thesis Outline}
\label{sec:thesis_outline}

This thesis investigates how direct time-of-flight (dToF) single-photon sensing augments material understanding and vision reliability. The chapters are organised as follows.

\textbf{Chapter 1} provides an introduction to the entire thesis, offering a comprehensive overview of all major sections and establishing the research context. This chapter presents key highlights including thesis deliverables, aims and objectives, and the project value proposition. The introduction serves as a roadmap for readers, guiding them through the structure and significance of the work while emphasising contributions and practical applications of the research.

\textbf{Chapter 2} presents an in-depth review of related and relevant literature, examining similar work to the presented solutions to justify the theoretical and empirical foundations upon which the new solutions have been developed. The review positions the proposed methods within prior signal processing and computer vision research, establishing connections between existing approaches and the innovations presented in subsequent chapters.

\textbf{Chapter 3} describes the Particle Wave Detection (PAWD) framework, a formal framework that characterises signal processing systems with respect to electromagnetic energy--matter interactions and measurement. The framework formalises two major categories of signal measurement to guide objectives: intensity spectral or non-spectral versus temporal spectral or non-spectral. The framework combines spatially resolved signals from which geometric features may be suitably sourced with time-resolved signals, which encode structural features including composition and depth. The framework also describes major theoretical foundations backing the projects: the behaviour of signals from translucent materials (Figure~\ref{fig:1_1}) and multi-material object perception, grounding subsequent experiments in electromagnetic energy--matter interactions \cite{jungerman2022}.

\textbf{Chapter 4} explores a material-detection solution for less geometrically complex (planar) objects using transient images, where embedded within the returned signal profiles are signature structural patterns that uniquely categorise the material composition of the surface from which they emerged. These temporal signatures encode material-specific optical properties, scattering characteristics, and internal structure information that are not readily accessible through spatial intensity images alone, enabling accurate material discrimination even when objects share similar visual appearances. A lightweight classifier is employed to classify objects based on their material composition, leveraging learned temporal patterns to distinguish between different materials with high accuracy while maintaining computational efficiency suitable for real-time applications. The solution demonstrates that time-resolved signals provide a robust foundation for material detection, offering advantages over traditional spatial-only approaches by capturing material properties that are invariant to lighting conditions, surface textures, and spatial patterns \cite{jungerman2022,becker2023}.

\textbf{Chapter 5} describes a dual material and object detection solution that processes multi-sensor intensity and transient images to address computer vision problems. The approach uses a custom spatiotemporal neural network that differs from traditional object detection models in architecture and input format, requiring paired space- and time-resolved images. The two-dimensional spatially resolved image is processed to extract spatial features, while the transient or time-resolved image is processed for structural features. The chapter also describes the role of time-resolved signals in machine intelligence, explaining how multi-sensor features are learned through this integrated framework \cite{tachella2019,jungerman2022}.

\textbf{Chapter 6} presents a food safety solution that measures material purity by comparing the target material's transient profile (time-resolved measurements from the SPAD sensor) to a thresholded ground truth profile derived from previously learned material transient signatures. The solution is demonstrated using a complex homogenised liquid material (milk), where the ground truth signal is extracted from the consumable milk condition as advised by the milk producer, enabling reliable purity assessment through temporal pattern comparison.

\textbf{Chapter 7} presents a summary of contributions, solution generalisability, future directions and conclusions of the work. The conclusion summarises all solutions presented, while the future directions outline next steps for each solution and how they can be further developed into more efficient approaches. The novelty of the solutions and their generalisability are also detailed, emphasising the potential for broader application across material and vision tasks.

\newpage
\chapter{Literature Review}
\label{chap:literature_review}

\newpage
\section{Overview}
\label{sec:lit_review_overview}

This chapter presents a comprehensive review of related research and background literature contextualising the three solutions presented in this work. The review compares current optical sensing approaches against each solution, assessing both strengths and limitations of existing methods to demonstrate how the presented solutions address fundamental gaps in material detection and scene understanding. The analysis reveals that existing approaches predominantly rely on either spatial features or spectral information in isolation, overlooking the complementary benefits of temporal signal analysis for material characterisation \cite{schwartz2013,sharan2013,su2016,becker2023}.

Spectroscopy is reviewed as a signal processing paradigm based on multi-wavelength analysis (spectral signal processing), serving to contrast the primary signal processing approach employed in this work, which operates without considering variations in composite wavelengths of the illumination source. This non-spectral, achromatic approach, detailed in Chapter 3, demonstrates that temporal information can provide comparable material discrimination capability without the computational overhead associated with spectral analysis \cite{holt2016,elachi2021}. The comparison reveals that while spectral methods offer rich wavelength-dependent information, temporal methods capture material-specific optical properties through time-resolved interactions, enabling efficient material detection with reduced complexity \cite{su2016,becker2023,tafone2023}.

Material detection instruments are reviewed to establish various detector categories and justify why SPADs represent the most efficient sensor technology for this work. The review covers different detector technologies and their characteristics, establishing SPADs as optimal for time-resolved material detection due to their single-photon sensitivity, high temporal resolution, and compact form factor \cite{piron2020}. Further details on SPAD limitations are reviewed to inform readers of potential vulnerabilities in this work, including dark count rates, afterpulsing effects, and array size constraints that may affect detection performance in certain scenarios \cite{piron2020}.

Depth estimation techniques are reviewed, including stereoscopy, time-of-flight (ToF) methods, and interferometry, in comparison to direct time-of-flight (dToF) depth sensing employed in this project. Stereoscopy is reviewed as the most commonly used depth estimation approach and historically the most reliable until ToF methods gained prominence \cite{laga2022,hu2020,godard2019}. It remains fundamental to numerous vision systems, including machine learning applications, despite inherent limitations \cite{laga2022}. Interferometry, though less prominent, is reviewed to provide comprehensive coverage of depth sensing approaches \cite{huang1991,aumann2019}. Echolocation methods are reviewed to compare dToF with indirect time-of-flight (iToF), as both approaches employ different signal processing strategies \cite{bamji2022}. This comparison highlights why dToF is more efficient for material detection than iToF, with dToF providing superior temporal resolution and single-photon sensitivity that enable more accurate material characterisation \cite{su2016,becker2023}.

Under the role of SPADs in vision, details of the solution and its uniqueness relative to existing methods are presented, particularly the innovation of combining time-resolved information with spatially resolved images (Chapter 5) for enhanced scene understanding \cite{yang2023,yang2024,gadzicki2020}. This hybrid approach represents a significant advancement beyond traditional vision systems that rely solely on spatial or temporal information. Relevant literature employing SPAD technology for both depth estimation and material classification is reviewed, demonstrating growing recognition of SPADs as versatile sensors for comprehensive scene understanding \cite{su2016,becker2023,tafone2023,yang2024}. The review shows that while previous work has explored SPADs for depth sensing or material detection separately, the integration of both capabilities in a unified framework is novel and addresses fundamental limitations in current vision systems \cite{yang2023,yang2024,perez2024}.

The material detection with SPADs section, synonymous with the main project title, describes and reviews literature relating to experiments and methods employed in the design and implementation of the solution \cite{su2016,becker2023,tafone2023,yang2024}. This section provides a foundation for the material detection approaches presented in subsequent chapters, establishing connections between proposed methods and existing research while highlighting unique contributions of this work. Finally, the chapter concludes with summaries of unique value propositions for each of the three presented solutions, providing readers with a clear understanding of how each solution addresses specific challenges and contributes to advancing the field of material detection and scene understanding. 

\newpage
\section{Optical sensing for object and material detection}
\label{sec:optical_sensing_approaches}

Optical sensing techniques exist for both material detection and object detection, each addressing distinct aspects of scene understanding. Optical sensing techniques for material detection focus on determining the material composition of targets, distinguishing materials based on their intrinsic optical properties. In contrast, optical sensing techniques for object detection focus on identifying and localising objects as targets within a scene, enabling spatial recognition and classification of distinct entities.

Several signal processing techniques relate to the solutions presented in this work. Material detection methods based on RGB intensity images \cite{schwartz2013,sharan2013,depthpro2024}, non-RGB (infrared or achromatic) intensity images combined with depth maps \cite{sharma2023}, and LiDAR-based material detection using raster-scanned reflectivity measurements from target surfaces \cite{yang2024} have been investigated. Complementary object detection approaches include RGB intensity images with depth information \cite{chen2024}, and time-of-flight-based object detection systems \cite{tafone2023,fasolino2024,yang2023}. Significant work on imaging hidden objects using time-of-flight detectors has demonstrated promising applications in specialised domains such as medical imaging \cite{lyons2019}. Notable advances include the use of a single time-of-flight pixel to infer three-dimensional scene information \cite{turpin2019}, a problem that would be considered ill-posed if the single pixel measured signal as intensity. While these approaches demonstrate various capabilities in material detection and scene understanding, they often operate in isolation, focusing on either spatial or temporal information, but rarely integrating both modalities effectively. The existing literature reveals a fundamental gap: most methods rely primarily on spatial features extracted from intensity images, which capture appearance and texture but may not adequately represent intrinsic material properties that manifest in temporal signal characteristics. Furthermore, current benchmarks and evaluation frameworks tend to assess material detection and object detection separately, failing to capture the synergistic benefits of combined spatiotemporal analysis. This limitation becomes particularly evident when considering real-world applications where both object identity and material composition are essential for comprehensive scene understanding, such as in autonomous systems that must distinguish between real objects and their representations, or in quality control applications requiring both object detection and material purity assessment. This work addresses some of the most challenging questions about material detection and scene understanding in a manner not previously addressed, by demonstrating how temporal signals can provide structural features that complement spatial information, enabling more robust and reliable material detection that extends beyond appearance-based classification.

Spatial intensity image-based material recognition approaches, such as the material trait recognition method presented in \cite{schwartz2013}, detect pixel-level material traits including ``shiny'' or ``woven'' that encode characteristic material appearance properties. However, these traits represent physical properties secondary to optical or mechanical properties, creating fundamental limitations in material discrimination. Consider a scenario involving multiple samples of the same object category, each composed of different materials (as examined in Chapter 5). Shiny surfaces may be observed across multiple objects, yet this visual similarity does not imply material equivalence and may not require identical treatment from robotic systems. Models dependent on such material traits may misclassify materials or guide vision systems toward less reliable predictions. The fundamental limitation lies in the confusion between material composition and surface appearance, where visual traits can be shared across fundamentally different materials or artificially created through surface treatments, coatings, or lighting conditions. This approach fails to capture intrinsic optical properties that truly distinguish materials, such as their interaction with light at molecular or structural levels, which are more stable and material-specific than surface appearance characteristics \cite{elachi2021,bass2010,su2016}. Consequently, systems relying on such methods may struggle in real-world scenarios where materials share similar visual traits but exhibit different compositions, properties, or safety requirements. A similar approach based on monocular material detection is presented in \cite{depthpro2024}, and the same challenges from \cite{schwartz2013} may apply.

Previous material detection research based on intensity RGB images \cite{sharan2013} has been reviewed. Similar to \cite{schwartz2013} and \cite{depthpro2024}, this approach relies solely on spatial features rather than reflectivity data encoding valuable information on surface and subsurface optical properties and material composition \cite{su2016,becker2023}. As a result, the method is likely to be confounded by cases of the same material exhibiting different spatial patterns. For instance, consider a ceramic bowl of plain white colour and a ceramic cup with white and cream colour patterns. In this work, the aim is to detect that both objects are composed of the same ceramic material. However, based on the ground truth information used in \cite{sharan2013}, each object would be defined as a different material, primarily because the basis for selection derives from spatial patterns rather than actual structure or composition. Conversely, if a purple natural fruit and a corresponding 3D printed fruit with identical colour patterns were presented to their model, it would likely identify both as being composed of the same material, despite fundamental compositional differences.

The fundamental limitation of appearance-based material detection approaches lies in their conflation of material identity with visual appearance, which represent distinct physical concepts \cite{schwartz2013,sharan2013,chen2024}. Material composition is determined by intrinsic optical properties including absorption coefficients, scattering characteristics, and refractive indices, which manifest through fundamental light-matter interactions \cite{elachi2021,bass2010,su2016}. These properties are encoded in temporal response characteristics and reflectivity patterns rather than spatial colour distributions, which can be easily manipulated or replicated through surface treatments, coatings, or lighting conditions \cite{su2016,becker2023,tafone2023}. While appearance-based methods demonstrate robustness to lighting variations and generalise well to real photographs despite synthetic training data, this capability introduces a critical vulnerability in safety-critical applications: the inability to distinguish between authentic physical materials and their two-dimensional representations or synthetic replicas \cite{chen2024}. This limitation becomes particularly problematic in applications such as mobile device face unlock systems, where vision systems must differentiate between real objects and images of those objects. This is a common misconception in other material detection work \cite{chen2024}, but as demonstrated in this project, a material detection system should ideally be able to recognise authentic materials separately from other representations such as images. The distinction between appearance-based and property-based material detection becomes crucial when considering that real-world applications require systems capable of making decisions based on actual material properties rather than superficial visual similarities.

Alternative approaches employ non-RGB intensity images combined with depth information for material detection. The time-of-flight intensity and depth information-based material detection method \cite{sharma2023} combines infrared intensity images, depth images, and aligned point cloud data (often referred to as D-AI images) to perform material classification. This approach demonstrates the effectiveness of not relying on spectral composition for scene understanding, an idea promoted in this work. However, the D-AI image generation process involves complex alignment and registration operations that fuse multi-modal data, potentially introducing artefacts or losing temporal information valuable for material discrimination. Similar to appearance-based methods \cite{schwartz2013,sharan2013,chen2024}, intensity-based approaches face challenges in differentiating between authentic materials and their synthetic representations. For instance, such systems may struggle to distinguish between a 3D-printed fruit and its natural counterpart, despite fundamental compositional differences, when both exhibit similar spatial intensity patterns \cite{su2016,becker2023}. This inability to distinguish between real materials and their synthetic representations represents a fundamental limitation that could lead to misclassification in practical applications where material authenticity matters. Another challenge is related to the use of indirect time-of-flight (iToF) rather than single-photon direct time-of-flight (dToF) techniques. iToF reliance on non-single-photon techniques may result in less informative signals that barely capture unique material optical properties. The phase-based measurement approach in iToF systems provides lower temporal resolution and less detailed temporal signatures compared to dToF, which can directly measure individual photon arrival times with picosecond precision \cite{piron2020}. This temporal resolution limitation means that subtle material-specific optical properties manifesting in fine temporal structures may be lost or averaged out in iToF measurements \cite{su2016,becker2023}. Consequently, iToF-based systems may struggle to capture nuanced temporal patterns that distinguish materials with similar spatial appearance but different optical characteristics, limiting their effectiveness for applications requiring precise material discrimination.

The approach described in \cite{yang2024} estimates the optical action of irregular and internally inhomogeneous materials using a semi-infinite sample optical model based on Kubelka-Munk (K-M) theory. The K-M diffuse modification equation enables material characterisation by modelling how light interacts with scattering and absorbing media, converting reflectivity measurements into optical property estimates that are independent of surface geometry or viewing angle \cite{elachi2021,bass2010}. This theoretical foundation allows the system to extract material-specific information even when objects are partially hidden or viewed from non-optimal angles, addressing a fundamental limitation of appearance-based methods that rely on direct line-of-sight observations \cite{schwartz2013,sharan2013}. The ability to work with obscured or partially visible objects makes this approach particularly valuable for applications in cluttered environments, surveillance systems, or scenarios where objects may be intentionally concealed. This perspective-invariant capability shares common principles with the solutions described in Chapters 4 and 6 of this work, where material classification relies on temporal signal characteristics rather than spatial appearance.

Based on the above review, current object detection tasks, both ToF-based \cite{tafone2023,fasolino2024,yang2023} and RGB-based \cite{chen2024}, would greatly benefit from incorporating time-of-flight signals as sources of depth and structural features. The inclusion of depth information into salient object detection systems \cite{chen2024} can result in three-dimensional detection and improved scene understanding. The temporal information provided by ToF signals captures material-specific optical properties and structural characteristics that complement spatial appearance features, enabling more robust detection that is less susceptible to lighting variations, surface textures, or visual camouflage \cite{su2016,becker2023,tafone2023}. This integration of temporal and spatial information addresses the fundamental challenge of distinguishing real objects from their representations, as temporal signals encode physical material properties that cannot be easily replicated in images or synthetic models \cite{su2016}.

Other systems that combine space- and time-resolved signals for detection of hidden objects in highly scattering media reveal promising breakthroughs with time-of-flight detectors. In \cite{lyons2019}, computational algorithms spatially resolve transient signals into spatial intensity representations of targets within scattering media, with depth information conveying target distance from the detector. While this system provides true depth and hidden object position, limitations include coarse resolution for small array time-of-flight sensors and increased setup costs when scaling to larger arrays. Chapter 4 presents a solution addressing this limitation: using a 32×32 SPAD array targeted at specialised objects like internal body tissue, material detection can identify the material composition of hidden or obscured objects even with a single pixel and without knowing their shape \cite{su2016}. This capability is particularly valuable in medical imaging and security applications where objects may be obscured but material identification remains critical for decision-making \cite{lyons2019,tye2025}.

Inferring spatial scene configuration from a single intensity pixel is considered an ill-posed problem \cite{pizlo2001}. However, time-resolved signals in combination with machine learning systems demonstrate the use of transient images to infer spatial representations \cite{turpin2019,jungerman2022}. The temporal information provides additional constraints that transform the ill-posed problem into a solvable one, enabling spatial reconstruction from minimal sensor data \cite{turpin2019}. While this potentially improves scene analysis speed, image size, and computational efficiency, it could be further enhanced with material detection. The integration of material detection would enable systems to not only reconstruct spatial configurations but also identify material properties, providing richer semantic information for scene understanding \cite{su2016,becker2023,yang2024}. Specifically, either \cite{turpin2019} or \cite{lyons2019} would fit into the spatial detection side of the presented solution in Chapter 5, such that distance information from the material detection side is not estimated, as these systems already perform this function \cite{yang2023,yang2024}. This complementary integration leverages the strengths of each approach: spatial detection systems provide distance and position information, while material detection systems provide composition and property information, together enabling comprehensive scene understanding that neither approach could achieve independently.


\subsection{Summary of Techniques}
\label{subsec:signal_image_processing_summary}

Optical sensing approaches for material and object detection can be categorised into two primary modalities: spatial intensity-based methods and temporal time-of-flight-based methods (Table~\ref{tab:2_1}).

Material detection techniques include: (1) RGB intensity image-based methods that detect pixel-level appearance traits (e.g., shiny, woven) but suffer from conflating material composition with surface appearance \cite{schwartz2013,sharan2013,depthpro2024}; (2) non-RGB intensity images combined with depth maps (D-AI images) that fuse infrared intensity, depth, and point cloud data but face challenges in distinguishing authentic materials from synthetic representations \cite{sharma2023}; and (3) LiDAR-based material detection using Kubelka-Munk theory for perspective-invariant material characterisation from raster-scanned reflectivity measurements \cite{yang2024}. Object detection approaches encompass RGB images with depth information \cite{chen2024} and time-of-flight-based systems \cite{tafone2023,fasolino2024,yang2023}. Specialized techniques include hidden object imaging in scattering media \cite{lyons2019} and single-pixel 3D inference from transient signals \cite{turpin2019,jungerman2022}.

A fundamental limitation across most existing methods is their operation in isolation, focusing on either spatial or temporal information but rarely integrating both modalities effectively. Appearance-based approaches rely on spatial features that capture texture and appearance but inadequately represent intrinsic material properties manifesting in temporal signal characteristics. Indirect time-of-flight (iToF) systems provide lower temporal resolution compared to direct time-of-flight (dToF), limiting their ability to capture nuanced material-specific optical properties. The integration of temporal and spatial information addresses the critical challenge of distinguishing real objects from their representations, as temporal signals encode physical material properties that cannot be easily replicated in images or synthetic models \cite{su2016,becker2023,tafone2023}.

\newpage
\begin{sidewaystable}[p]
\centering
\footnotesize
\begin{tabularx}{\textwidth}{p{2.2cm}Xp{2.5cm}p{2cm}p{2cm}Xp{0.3cm}}
\toprule
\textbf{Technique Type} & \textbf{Input Data} & \textbf{Processing Category} & \textbf{Active Illumination} & \textbf{Depth Information} & \textbf{Primary Advantage} & \textbf{Refs} \\
\midrule
Material classification & RGB intensity images & Spatial processing & No (passive) & None & Object-invariant detection; Simple implementation; Limited by appearance-based traits; Cannot distinguish authentic materials from representations & \cite{schwartz2013,sharan2013,depthpro2024} \\
\midrule
iToF-based material detection & IR intensity + depth/point cloud (D-AI) & Multi-modal spatial processing & Yes (iToF) & True depth (iToF) & Combines multiple modalities; Doesn't rely on spectral composition; Lower temporal resolution than dToF; Challenges distinguishing synthetic representations & \cite{sharma2023} \\
\midrule
LiDAR-based material detection (K-M theory) & dToF data (histogram), raster-scanned reflectivity & Temporal signal processing & Yes (dToF) & True depth & Perspective-invariant material characterisation; Independent of surface geometry; Works with obscured objects & \cite{yang2024} \\
\midrule
Low-cost compact dToF material detection & dToF histogram (VL53L1X) & Temporal signal processing & Yes (dToF) & True depth & Compact and low cost; Material classification and depth imaging; Non-line-of-sight tracking & \cite{callenberg2021} \\
\midrule
RGB-D object detection & RGB + depth (stereoscopy) & Multi-modal spatial processing & No (passive) & Relative depth & 3D object detection; Improved scene understanding & \cite{chen2024} \\
\midrule
ToF object detection & ToF data & ToF signal processing & Yes (ToF) & True depth & True depth information; Material-specific optical properties; Robust to lighting variations & \cite{tafone2023,fasolino2024,yang2023} \\
\midrule
Hidden object imaging & ToF transient signal & Temporal signal processing & Yes (ToF) & True depth & Detects hidden objects in scattering media; True depth; Coarse resolution for small arrays; Medical imaging applications & \cite{lyons2019} \\
\midrule
Single-pixel 3D inference & Single ToF pixel & ML-based temporal processing & Yes (ToF) & True depth & Solves ill-posed problem; Fast processing; Reduced image size; Spatial reconstruction from minimal sensor data & \cite{turpin2019,jungerman2022} \\
\bottomrule
\end{tabularx}
\caption{Table of comparison for different optical sensing processing techniques for scene inference}
\label{tab:2_1}
\end{sidewaystable}



\newpage
\section{Spectroscopy}
\label{sec:spectroscopy}

All colors emerge from the radiation spectrum. The origins of spectroscopy trace to Sir Isaac Newton's discovery in 1666 \cite{su2016,holt2016}, where he observed wavelength-dependent refractive indices as light passed through a prism. Following this discovery, Newton demonstrated that radiation comprises multiple wavelengths, with red exhibiting the least refraction and blue the most as they pass through a prism.

Spectroscopy is an analytical method that examines the spectral composition of signals resulting from energy-matter interactions. For material detection applications, spectroscopic analysis focuses on energy measurements of the spectral composition of remotely detected signals. On the energy spectrum, waves are grouped into spectral bands exhibiting different frequencies, which influence the behavioural and ionising properties of the signal. The fundamental rationale behind spectroscopic methods is to leverage the wavelength-dependent optical responses of materials to incident radiation, enabling identification of unique or signature material properties through collection and analysis of spectral measurements from energy-matter interactions. Spectral measurements may be acquired as intensity values or photon arrival time distributions. Spectroscopy aims to understand energy-matter interactions per wavelength within the incident signal; thus, for both two-dimensional intensity or time-resolved images, wavelength variations are required for spectroscopic analysis to occur. Imaging spectroscopy, as applied in digital photography, may be used to analyse and determine colour spaces used by pixels through optical filtering.

Spectral measurements are categorised into three modalities: monospectral, multispectral, and hyperspectral \cite{heinz2001,vasile2023}. Monospectral measurements filter a single spectral band for processing, activating only one colour channel and assigning a single colour to all pixels within the image (e.g., red or green only). Multispectral imaging filters multiple wavelengths; RGB images represent a common type of multispectral images where multiple spectral bands (RGB or CYM) are activated. Hyperspectral imaging captures significantly more wavelengths from the wave composition \cite{heinz2001,vasile2023}. The distinct spectral bands in an image constitute its channels: an RGB image has three active channels, a monospectral image has one, and a hyperspectral image contains as many channels as wavelengths composed in the incident wave.

The choice between monospectral, multispectral, and hyperspectral approaches depends on specific material detection requirements, with hyperspectral imaging providing the richest spectral information at the cost of increased data complexity and computational overhead \cite{heinz2001,vasile2023}. Each material exhibits characteristic absorption and reflection patterns across different wavelengths, creating unique spectral signatures that enable material identification and classification \cite{holt2016,elachi2021}. Spectral resolution, defined as the ability to distinguish between closely spaced wavelengths, becomes increasingly important as the number of spectral bands increases, with hyperspectral systems capable of detecting subtle spectral variations that may be imperceptible in multispectral or monospectral measurements \cite{heinz2001,vasile2023}. However, the high dimensionality of hyperspectral data introduces challenges in data storage, processing speed, and feature extraction, often requiring sophisticated dimensionality reduction techniques such as principal component analysis or band selection algorithms to manage computational complexity while preserving material-discriminative information \cite{heinz2001,wold1987,bishop2006}. The effectiveness of spectroscopic material detection relies on the assumption that different materials exhibit measurably distinct spectral responses, which holds true for many materials but may be limited when materials have similar chemical compositions or when environmental factors such as illumination conditions or surface geometry significantly influence the measured spectral signatures \cite{schwartz2013,sharan2013,elachi2021}.

From Section 2.1.1, current research, especially those using colour images, appears disconnected from how resulting colour images relate to spectral analyses. However, all colour images result from spectral colour combinations. Chapter 5 of this work combines spectral (RGB) images with one-dimensional time-resolved data as input for dual object and material detection. The infrared spectral range (especially 780 nm to 1 mm according to ISO CIE) is known to be absorbed by most materials with little attenuation, which may have contributed as one of the major reasons infrared remains the most common illumination source for optical remote sensing \cite{elachi2021,bass2010,schott2007}.

Given a material illuminated by a polychromatic source, light particles interact differently through absorbance at different frequencies or wavelengths. Some frequencies may pass through the material unabsorbed; others may be partly absorbed, reflected, or scattered. These characteristics are measurable using quantities such as reflectance, absorbance, or transmittance and can cumulatively generate signature profiles with which the material can be reliably identified \cite{jungerman2022,becker2023}. The profiles are often analysed with empirical methods. When a certain spectral band is analysed by composition for material detection, the variable wavelengths form multiple dimensions. This differs from the solutions presented in this work, where returned photons from a pulsed, monochromatic source are processed with no regard to wavelengths (after optical band-pass filtering at the wavelength of the source). The returned photons are counted with respect to their arrival times using a time-correlated single-photon counting (TCSPC) instrument within the SPAD.

The achromatic approach eliminates the need for multi-wavelength analysis, reducing computational complexity from processing high-dimensional spectral data cubes to analysing one-dimensional temporal histograms \cite{su2016,becker2023,tafone2023}. This temporal information captures material-specific scattering and absorption characteristics through time-resolved photon arrival patterns, providing structural information that complements the spatial and colour information from RGB images \cite{jungerman2022,su2016}. While spectral methods require sophisticated unmixing algorithms and extensive calibration across multiple wavelengths \cite{heinz2001,vasile2023}, the time-resolved approach leverages the temporal signatures inherent in photon-material interactions, enabling material discrimination without spectral dimensionality \cite{su2016,becker2023,tafone2023}. The combination of RGB spatial information with temporal material signatures creates a complementary information space where spatial features identify what objects are present, while temporal features reveal what materials they are composed of \cite{yang2023,yang2024,gadzicki2020}. Objectives one through three in this project may be achieved by spectroscopy, but the achromatic approach presented in this work is considered less complex and comparably efficient, offering a practical alternative that maintains detection accuracy while reducing system complexity and computational requirements \cite{su2016,becker2023,tafone2023}.

\subsection{Non-spectral methods}
\label{subsec:non_spectral_methods}

The literature on material detection methods that operate without spectral composition analysis remains comparatively sparse. As described in Chapter~\ref{chap:pawd} and (Figure~\ref{fig:2_1}), signals may be measured continuously over time (time-resolved) or digitised as integrated intensity values resolved in two-dimensional space, with either modality operating with or without spectral decomposition. This work employs time-resolved signal measurements without regard to spectral composition, termed non-spectral processing. By excluding multi-wavelength analysis, computational complexity is significantly reduced, as the system processes single-channel temporal data rather than multi-dimensional spectral data cubes that require extensive computational resources for storage, analysis, and feature extraction \cite{su2016,becker2023,tafone2023}. The reduced dimensionality enables faster processing times, lower memory requirements, and more efficient algorithms suitable for real-time or resource-constrained environments \cite{su2016,becker2023,jungerman2022}. Despite this simplification, the quality of results and analysis outputs remain impressive (see Chapter~\ref{chap:pawd}), demonstrating that temporal information can effectively compensate for the absence of spectral information by capturing material-specific signatures through time-resolved photon-material interactions \cite{su2016,becker2023,tafone2023,yang2024}. The non-spectral approach proves particularly effective for material detection tasks where temporal patterns encode sufficient discriminative information, such as material composition, optical properties, and structural characteristics that manifest in time-of-flight measurements \cite{su2016,becker2023,tafone2023}. This finding challenges the conventional assumption that spectral information is necessary for accurate material characterisation, opening new possibilities for efficient, cost-effective material detection systems that maintain high performance while reducing computational overhead \cite{su2016,becker2023,tafone2023}. The success of non-spectral methods in this work validates that temporal signal processing can achieve comparable or superior results to spectral approaches in specific material detection applications, establishing a foundation for future research in efficient vision systems \cite{su2016,becker2023,tafone2023,yang2024}.

\newpage
\section{Material detection instruments}
\label{sec:material_detection_instruments}

Material detection instruments exploit unique optical properties arising from energy-matter interactions to characterise materials \cite{elachi2021} (Figure~\ref{fig:2_1}). These systems comprise radiation detectors and signal processing circuitry that transform raw sensor outputs into interpretable data formats \cite{bass2010}. Radiation detectors, typically analog input channels or receivers, capture photons or electromagnetic radiation that has interacted with target materials, converting these interactions into measurable electrical signals encoding material-specific information \cite{kolanoski2020,piron2020}. These detectors may incorporate secondary multiplication stages such as signal amplification circuitry for processing and buffering. The signal processing chain includes analog-to-digital conversion, signal conditioning, and feature extraction stages, transforming detector outputs into data formats suitable for analysis and classification \cite{bass2010}. Signal amplification stages enhance weak signals from low-intensity interactions or distant targets, preserving material-discriminative information above noise levels \cite{piron2020}. Buffering stages provide temporary storage and signal stabilisation, enabling reliable processing of transient or time-varying signals characteristic of time-of-flight and temporal material detection systems. The overall architecture maximises sensitivity to material-specific optical properties while minimising noise and maintaining signal integrity throughout the processing chain \cite{bass2010,piron2020}. Figure~\ref{fig:2_1} illustrates the fundamental classification of material detection instruments based on their detection mechanisms, categorising systems according to whether they process radiation as wave-like or particle-like phenomena. This classification framework provides a foundation for understanding the diverse approaches to material characterisation employed in optical sensing systems.

\newpage
\begin{figure}[H]
    \centering
    \includegraphics[width=0.9\textwidth,height=0.7\textheight,keepaspectratio]{figures/2.1.png}

    \caption[Classification of material detection instruments by wave-like or particle-like radiation processing.]{\textbf{Classification of material detection instruments by wave-like or particle-like radiation processing.} Visual illustration of different material detection instruments and their corresponding categories as processing wavelike or particle-like radiation. Material detection instruments are classified according to their fundamental detection mechanisms, which operate on either wave-like or particle-like interpretations of electromagnetic radiation. This categorization reflects the dual nature of light-matter interactions, where detectors may exploit wave properties (interference, diffraction, phase relationships) or particle properties (individual photon detection, quantum effects, discrete energy transfer) to extract material-specific information from optical signals.}
    \label{fig:2_1}
\end{figure}



\subsection{Light detectors}
\label{subsec:light_detectors}

Light detectors are categorised into two primary classes: vacuum detectors (including gas-filled types) and solid-state detectors (see Figure~\ref{fig:2_1}). Vacuum detectors operate in Geiger-Mueller, proportional, and ionisation regions, serving as important instruments for radiation measurement \cite{kolanoski2020}. However, the applications considered in this project have requirements for compactness, which is largely realised in solid-state technology. Thus, this review focuses on solid-state particle detectors, though vacuum detectors are briefly reviewed.

Particle detectors are categorised by the detector's state of matter (solid-state or gas-filled detectors) \cite{kolanoski2020}. For each category, detectors vary by radiosensitive material (phosphor, argon, silicon, etc.), photoelectric effects (internal or external), and detection mechanism (e.g., active mechanisms such as impact ionisation, or passive mechanisms such as PIN photodiode arrays) \cite{bass2010}. Solid-state detectors offer significant advantages for compact applications, including smaller physical dimensions, lower power consumption, and integration with semiconductor electronics, making them particularly suitable for portable and embedded vision systems \cite{bass2010,piron2020}. The solid-state category encompasses various detector types, including silicon-based photodiodes, avalanche photodiodes (APDs), and single-photon avalanche diodes (SPADs), each offering different sensitivity, temporal resolution, and noise characteristics \cite{piron2020}. The choice of detector type depends on specific application requirements such as temporal resolution, sensitivity, dynamic range, and operating conditions. In contrast, vacuum detectors, while offering excellent sensitivity and well-characterised performance, require larger physical volumes, higher operating voltages, and more complex support electronics, limiting their suitability for compact applications \cite{kolanoski2020}. The detection mechanisms fundamentally differ between active detectors, which generate internal gain through processes like impact ionisation in avalanche photodiodes, and passive detectors, which rely on direct photoelectric conversion without internal amplification, each mechanism offering distinct trade-offs between sensitivity, speed, and noise characteristics \cite{bass2010}.

The photoelectric effect refers to the generation of photoelectric current as a result of incident radiation on a detector \cite{bass2010}. The effect can be internal or external depending on the mechanism of photo detection, where external photoelectric effect involves electron emission from the detector surface into vacuum, while internal photoelectric effect involves electron-hole pair generation within the detector material \cite{bass2010}. The review of photoelectric effects of detectors contributes toward distinguishing solid-state and gas-filled detectors, each employing fundamentally different physical processes for converting incident photons into measurable electrical signals \cite{kolanoski2020,bass2010}.

Detection mechanisms of both solid-state and gas-filled detectors involve at least one process called secondary electron multiplication (SEM) \cite{bass2010}. This refers to techniques by which electrons or photo-electrons are multiplied and guided towards readout circuits following the photoelectric effect. In solid-state detectors, this can occur via impact ionisation in avalanche photodiodes (APDs) or through photodiode pixels, though standard photodiode pixels do not use electron multiplication and instead rely on direct charge collection \cite{bass2010}. In vacuum detectors, SEM occurs through dynodes in photomultiplier tubes (PMTs) or through scintillators that convert photons to electrons via secondary processes \cite{kolanoski2020}.

It should be noted that the impact ionisation process in SPADs (Single-Photon Avalanche Diodes) and EMCCDs (Electron-Multiplying Charge-Coupled Devices), as used for signal amplification, is not usually classified as SEM in the traditional sense, despite involving electron multiplication \cite{piron2020}. This distinction arises because SPADs operate in Geiger mode, where a single photon can trigger a self-sustaining avalanche, while EMCCDs use impact ionisation in a controlled, linear multiplication regime, both differing from conventional SEM mechanisms that involve discrete multiplication stages. The Geiger mode operation of SPADs represents a unique detection mechanism where the avalanche process is self-quenching, enabling single-photon sensitivity with high temporal resolution \cite{piron2020}.

SEM mechanisms for photodetectors of both matter types are compared and assessed with respect to the mechanism of Geiger mode impact ionisation, which is the detection mechanism used by SPADs. Geiger mode impact ionisation differs from conventional SEM in that it operates in a bistable regime where a single carrier can trigger a full breakdown, resulting in a digital output signal rather than an analog multiplication process \cite{piron2020}. For each SEM type, corresponding charge collection and readout circuitry of eligible detectors are compared to SPADs, highlighting the differences in signal processing, timing characteristics, and system integration requirements \cite{bass2010,piron2020}. The comparison reveals that SPADs offer advantages in single-photon detection, temporal resolution, and compactness, while traditional SEM-based detectors may provide benefits in gain linearity and dynamic range for certain applications \cite{piron2020}.


\subsection{Vacuum Detectors}
\label{subsec:vacuum_detectors}

Vacuum detectors, including gas-filled variants \cite{kolanoski2020}, comprise devices containing a vacuum chamber within which photosensitive materials (solid or fluid) receive incident radiation. Unlike solid-state detectors where detection occurs within the semiconductor material itself, vacuum detectors exploit external photoelectric effects, where incident photons liberate electrons from photocathode surfaces into the vacuum environment \cite{bass2010}. The vacuum medium enables electron multiplication through secondary emission processes, achieving extremely high gain and sensitivity for low-light and single-photon detection applications \cite{bass2010,kolanoski2020}, which are particularly useful for this project's low-light and single-photon detection requirements.

Dynode electron multipliers represent a primary class of vacuum detectors, consisting of electrode structures within vacuum tubes. Photomultiplier tubes (PMTs), multi-anode photomultiplier tubes (MaPMTs), and video camera tubes operate by accelerating photoelectrons through cascaded dynode stages, where each collision produces multiple secondary electrons through secondary emission \cite{bass2010}. This cascaded multiplication process yields signal amplification factors exceeding one million, enabling detection of signals below the noise floor of conventional solid-state detectors \cite{bass2010}.

Microchannel plate (MCP) detectors provide fast, high-gain amplification for single-photon or electron detection within vacuum chambers \cite{bass2010}. MCPs consist of arrays of microscopic channels functioning as independent electron multipliers, providing spatial resolution while maintaining high temporal resolution \cite{bass2010}. Examples include imaging intensifying charge-coupled devices and complementary metal-oxide semiconductor devices (ICCD and ICMOS), as well as electron-bombarded multichannel plate CMOS (EB-MCP-CMOS) detectors \cite{bass2010}.

Silicon-based detectors operating within vacuum chambers, termed silicon vacuum detectors, combine vacuum-based electron multiplication with silicon's quantum efficiency and spatial resolution capabilities \cite{bass2010}. Examples include electron-bombarded CCD (EBCCD), electron-bombarded CMOS (EBCMOS), hybrid photodetectors (HPD), and hybrid avalanche photodiodes (HAPD) \cite{bass2010}.

Scintillators are frequently paired with silicon detectors to extend spectral response to radiation types for which silicon exhibits limited sensitivity, such as X-rays and gamma rays \cite{bass2010}. The scintillator converts high-energy photons into visible light, which silicon detectors can then capture and measure, effectively extending the detector's spectral response range \cite{bass2010}. This combination enables detection of radiation types that would otherwise be undetectable with silicon-based detectors alone, making it valuable for applications in medical imaging, nuclear physics, and high-energy particle detection \cite{bass2010}.

Vacuum detectors offer significant advantages including high sensitivity and extremely high gain through electron multiplication, making them capable of detecting signals below the noise floor of conventional solid-state detectors \cite{bass2010,kolanoski2020}, which is particularly useful for this project's low-light and single-photon detection requirements. However, compared to solid-state detectors, vacuum detectors exhibit limitations including susceptibility to performance degradation from delicate anode wire contamination, deposition, or ageing effects that can degrade efficiency over time \cite{kolanoski2020}. Additionally, vacuum detectors require high-voltage power supplies, exhibit greater sensitivity to environmental factors such as temperature variations and mechanical shock, and typically have larger physical dimensions due to vacuum chamber requirements \cite{bass2010,kolanoski2020}. The applications considered in this project require miniaturized detectors, which have been largely realised through advanced silicon semiconductor technology, where modern fabrication techniques enable compact, integrated detector arrays with excellent performance characteristics \cite{piron2020}. Solid-state detectors also offer advantages in terms of lower power consumption, faster response times, better stability, and reduced maintenance requirements compared to vacuum-based systems \cite{piron2020}. Consequently, vacuum detectors are not considered for this project, as solid-state technology better aligns with the project's requirements for compact, portable, and reliable detection systems \cite{piron2020}.

\subsection{Solid State Detectors}
\label{subsec:solid_state_detectors}

Solid-state detectors (SSDs) are fabricated from solid materials, typically silicon semiconductors. A key advantage of SSDs over vacuum-based detectors lies in their higher material density and the overall advancement in semiconductor technology, enabling inexpensive and compact sensor implementations \cite{bass2010,kolanoski2020}. This section describes solid-state detectors by category: thermal detectors and quantum detectors such as SPADs \cite{piron2020}. Thermal detectors (such as thermocouples) sense radiation by measuring variations in temperature as the detector absorbs radiation. Thermal detectors are passive sensors, which means they do not perform any secondary amplification of photoelectrons after incidence \cite{bass2010}. CCD and CMOS detectors are generally passive in this sense. This project considers only active detectors that have a higher electron amplification mechanism, which is why SPADs are used rather than photodiodes \cite{piron2020}. High electron amplification leads to a higher probability of detecting even a single electron, producing an output signal that rises above the noise floor (high sensitivity). For this reason, thermal and other passive detectors are not considered for this work, and the review is brief on such devices.

SPADs achieve single-photon sensitivity through avalanche multiplication, where a single photoelectron triggers a self-sustaining avalanche process that amplifies the signal by several orders of magnitude (this is not the typical SEM mechanism observed in other SSDs) \cite{piron2020}. This internal gain mechanism enables SPADs to detect individual photons with high temporal resolution, making them ideal for time-of-flight applications where precise timing of photon arrival is critical. The active amplification process in SPADs provides significant advantages over passive detectors, including improved signal-to-noise ratio, reduced susceptibility to electronic noise, and the ability to operate at lower light levels \cite{piron2020}. Specific detector performance parameters are also highlighted, including impact ionisation, quantum efficiency, noise, uniformity, spectral response, responsivity, stability, and speed \cite{bass2010}. For the specific quantum detector used (SPAD) in this project, manufacturer specifications are shared.

\subsection{Quantum Detectors}
\label{subsec:quantum_detectors}

This section categorises photodetectors into two fundamental types: photodiode and impact ionisation pixels. The categorization is based on the underlying physical mechanisms by which these detectors convert incident photons into measurable electrical signals, with each type exhibiting distinct operational characteristics, sensitivity profiles, and temporal response capabilities \cite{bass2010}. The following sections review each type of photodetector and their mode of operation in comparison to SPADs (Single-Photon Avalanche Diodes), which represent a specialised class of impact ionisation detectors optimised for single-photon sensitivity and precise time-of-flight measurements \cite{piron2020,gyongy2022}.

Photodiode-based detectors operate through the photoelectric effect, where incident photons generate electron-hole pairs that are collected as photocurrent, providing continuous analog signals suitable for intensity measurements but typically limited in temporal resolution and single-photon sensitivity \cite{bass2010}. Impact ionisation pixels, including SPADs, leverage avalanche multiplication processes where a single photon can trigger a self-sustaining avalanche of charge carriers, enabling detection of individual photons with exceptional temporal precision and sensitivity \cite{piron2020}. The reason for this comparative approach is to highlight why the chosen detector (SPAD) for this project is the most suitable for the objectives, particularly for time-resolved material detection applications that require single-photon sensitivity, nanosecond-scale temporal resolution, and the ability to extract structural features from transient signals that are inaccessible through conventional photodiode-based intensity measurements \cite{piron2020,gyongy2022}.

\subsubsection{Photodiode pixel}
\label{subsubsec:photodiode_pixel}

The photodiode pixel, also referred to as an active pixel sensor (APS), represents the fundamental imaging technology ubiquitous in mobile devices and digital cameras \cite{bass2010}. These detectors operate in reverse-bias mode, establishing relatively small depletion regions compared to impact ionisation detectors such as SPADs. In contrast, SPADs utilise substantially higher reverse-bias voltages that extend well beyond the breakdown voltage, enabling single-photon detection through avalanche multiplication \cite{piron2020}.

When photons possessing sufficient energy interact with the depletion region of an APS, electron excitation to the conduction band occurs, generating electron-hole pairs \cite{bass2010}. The electric field at the p-n junction facilitates charge carrier separation, with electrons and holes migrating to their respective electrodes. The total number of electron-hole pairs generated exhibits proportionality to the absorbed photon energy divided by the bandgap energy required to elevate an electron to the conduction band. Consequently, the photocurrent generated by incident photons demonstrates linear proportionality to the irradiance (intensity) of the incident light \cite{bass2010}. This photocurrent typically charges a capacitor and is subsequently converted into a voltage signal for further processing. The linear response of APS detectors to light intensity renders them well-suited for capturing spatial intensity distributions in conventional imaging applications, where the objective is to reproduce the visual appearance of a scene \cite{bass2010}.



\subsubsection{Impact ionisation pixel}
\label{subsubsec:impact_ionisation_pixel}

SPADs belong to the class of impact ionisation pixels (IIP) (see Figure~\ref{fig:2_1}). This class of detectors exhibits some of the highest electron amplification properties and sensitivity due to operation in reverse bias mode, beyond the breakdown voltage of the p-n junction \cite{bass2010}. IIPs are categorised according to the level of internal gain during the carrier multiplication process. IIP gain depends largely on the rate of impact ionisation, which is further influenced by the electric field across the depletion region \cite{bass2010}.

Examples of Gain>1 detectors are avalanche photodiodes (APDs) and electron-multiplying CCDs (EMCCDs) \cite{bass2010}. Compared to photodiode pixels (where Gain = 1), IIPs have much higher sensitivity. However, they do not outperform Gain>>1 IIPs, which operate at significantly higher reverse bias voltages \cite{bass2010}.

When an IIP is reverse-biased, incident radiation generates photoelectrons in the depletion region. Due to the high electric field across the depletion region, electrons are pulled toward the negative terminal (cathode) while holes are pulled toward the positive terminal (anode) \cite{bass2010}. Higher electric field strength increases carrier drift velocity and expands the depletion region. At a certain point, carriers begin to collide with the crystal lattice. With a much stronger electric field, these photoelectrons undergo further collisions and generate additional electron-hole pairs through impact ionisation \cite{bass2010}.

The rate at which this process occurs (impact ionisation rate) depends on the magnitude of the reverse bias voltage, which is the distinguishing characteristic between Gain>1 IIPs and Gain>>1 IIPs \cite{bass2010}. Gain>>1 IIPs operate at reverse bias voltages significantly higher than the breakdown voltage. SPAD detectors belong to the Gain>>1 IIP category \cite{piron2020}.

Other examples of Gain>>1 IIPs include p-n type CCDs (pnCCDs), which match SPAD performance in terms of sensitivity \cite{bass2010}. The original purpose of pnCCDs was for spectroimaging in X-ray astronomy, and they exhibit excellent sensitivity \cite{bass2010}. However, pnCCDs are large and tend to be spectrometric. Given that this project is based on miniaturized detectors and follows an achromatic signal processing approach, pnCCDs are not considered for this work.

SPAD detectors have evolved and advanced beyond traditional single-photon counting features. The following paragraph provides details of SPAD sensors, as this device shapes the performance of remote material detection systems in this project.

The impact ionisation process in SPADs arises from biasing the p-n junction beyond the breakdown voltage, creating an avalanche multiplication effect that provides significantly higher sensitivity than conventional photodiodes \cite{piron2020}. When a single photon is absorbed, it triggers an avalanche of charge carriers, producing a measurable current pulse. This avalanche mechanism enables SPADs to detect individual photons with exceptional temporal resolution (picosecond scale), rendering them ideal for time-of-flight measurements where precise photon arrival times must be recorded \cite{piron2020,gyongy2022}. This work employs SPADs because the application requires a detector with sufficient sensitivity to detect and time individual photons, enabling the capture of rich datasets relating to the optical properties of targets \cite{su2016,becker2023}. The single-photon sensitivity and picosecond temporal resolution of SPADs enable the extraction of material-specific temporal signatures that are inaccessible through conventional intensity-based imaging \cite{su2016,becker2023,tafone2023}. The photodiode pixel (APS) is not considered the most suitable detector technology for the material detection aspect of this project, though it is considered for object detection tasks where spatial information is paramount \cite{bass2010}. Other classes of higher-sensitivity, impact ionisation detectors are reviewed in the following section.

\subsubsection{Single photon avalanche diode (SPAD)}
\label{subsubsec:spad}

Single-photon avalanche diodes (SPADs), also known as Geiger-mode avalanche photodiodes (APDs), are photodetector devices consisting of a p-n junction reverse-biased above its breakdown voltage, operating with gain significantly greater than unity \cite{piron2020}. When reverse-biased above breakdown, the device operates in Geiger mode, creating an unstable state that enables detection of individual photons. The reverse-biased configuration increases the depletion region of the junction, placing the device in this unstable Geiger-mode state. The probability of triggering a self-sustaining avalanche, and consequently detecting a single photon arrival, increases with the reverse-bias voltage relative to the breakdown voltage \cite{piron2020}.

Each photon arrival triggers a series of impact ionisation events, leading to an electron avalanche that generates a detectable current flow. These single-photon events, translated into detectable current flows, can be timed (as in direct time-of-flight, dToF) or counted (by the demodulation circuit of indirect ToF pixels) \cite{gyongy2022}. Since electromagnetic radiation travels at the speed of light, the round-trip distance of photons can be calculated from their timing, providing accurate depth measurements from the detector \cite{tachella2019}.

The advantage of SPAD dToF for this project lies in single-photon detection sensitivity and time-of-flight depth sensing circuitry, which are essential for material classification, true depth estimation, and structural feature extraction. The high sensitivity enables detection of weak signals from remote targets, while the temporal resolution provides precise depth information critical for material characterisation.

After each photon detection, the SPAD must be reset by reducing the bias voltage below breakdown to quench the avalanche and restore the device to a state capable of detecting the next photon. This reset period, known as dead time, typically lasts a few nanoseconds and limits the maximum photon detection rate \cite{piron2020}. Details of time-of-flight depth sensing are found in the optical echolocation section (Section~\ref{subsec:time_of_flight}).

\subsubsection{SPAD challenges}
\label{par:spad_challenges}

SPAD imagers have long been products of low-volume military and scientific research markets. However, their rapid evolution \cite{tachella2019,jungerman2022} has enabled high-availability, compact, and cost-effective versions. Although this evolution has introduced some limitations, these can be mitigated through software-based post-processing. The following sections examine the principal challenges of SPAD-based imaging systems and present mitigation strategies employed in this work.

\subsubsection{SPAD sensor jitter}
\label{par:spad_sensor_jitter}

The total timing performance of SPAD-based systems is determined by two primary sources of temporal uncertainty: time-to-digital converter (TDC) jitter and SPAD jitter \cite{itzler2011}. These components collectively define the achievable timing resolution in pulsed LiDAR applications.

\paragraph{TDC Jitter}
\label{par:tdc_jitter}

Photon arrival times in pulsed LiDAR systems are measured using Time-to-Digital Converters (TDCs) \cite{sesta2023}. Despite consistent photon arrival patterns, electronic noise and inherent circuit limitations in the time measurement circuitry introduce temporal imprecision that degrades the sensor's timing resolution. This temporal uncertainty, termed TDC jitter, manifests as variability in the time measurement response, where the intervals between successive photon arrivals exhibit temporal delays. Unlike SPAD jitter, TDC jitter can be quantified through clock cycle differential analysis, providing a direct assessment of the precision with which analog time intervals are converted to discrete digital representations.

\paragraph{SPAD Jitter}
\label{par:spad_jitter}

As the signal returns from the target, photon detections exhibit random temporal uncertainty with respect to the true photon arrival time, arising from the stochastic nature of the avalanche multiplication process \cite{itzler2011,gyongy2022}. This jitter is quantified as the Full Width at Half Maximum (FWHM) of the transient signal histogram, which aggregates photon arrival times collected over multiple laser cycles, once the temporal profile of the laser pulse and any TDC jitter contributions have been compensated for through deconvolution \cite{sesta2023,piron2020}. SPAD jitter imposes a fundamental limit on the precision with which the arrival time of individual photons can be ascertained and is typically on the order of 100~ps for silicon SPAD devices \cite{itzler2011}.

\subsubsection{Coarse TDC resolution}
\label{subsubsec:coarse_tdc_resolution}

A fundamental limitation in cost-effective SPAD-based imaging systems is the coarse temporal resolution of time-to-digital converters (TDCs), which results in wider time bins that constrain the system's ability to resolve closely spaced surfaces and limit the temporal resolution of histogram profiles \cite{sesta2023}. These constraints restrict the extraction of fine-grained optical characteristics from target materials.

In this work, the TDC resolution was approximately 250~ps, which is finer than the laser pulse width. The measured SPAD signal histogram $h(t)$ represents the temporal convolution of three components: the laser pulse temporal profile $p(t)$, the combined SPAD and TDC jitter $j(t)$, and the target material's temporal response $m(t)$:
\begin{equation}
h(t) = p(t) \ast j(t) \ast m(t)
\label{eq:spad_histogram_convolution}
\end{equation}
where $\ast$ denotes temporal convolution. Since this convolved signal extends across multiple TDC bins, surface depth can be estimated with sub-bin resolution through interpolation or centre-of-mass (CoM) techniques, partially compensating for the coarse bin width \cite{erdogan2017,zhang2019}.

\subsubsection{Noise}
\label{subsubsec:noise}

Noise in direct time-of-flight (dToF) SPAD systems manifests through multiple mechanisms, with two primary contributions: background noise from ambient illumination sources and Poisson-distributed shot noise arising from the quantum nature of photon detection \cite{piron2020}. In active monochromatic illumination configurations, such as the system employed in this work, the SPAD detector incorporates an optical bandpass filter centred at the transmitted signal wavelength, which substantially attenuates ambient light contributions and reduces the background noise floor. The subsequent sections examine these noise sources and their implications for signal processing and material classification performance.

\paragraph{Ambient noise}
\label{par:ambient_noise}

SPAD detectors operating without spectral filtering are exposed to broadband ambient illumination spanning wavelengths beyond the emitted laser band, resulting in elevated incident photon rates that combine both ambient and signal photons \cite{piron2020}. When the total photon flux exceeds the detection and timing circuitry bandwidth, photon pile-up distortion manifests in the temporal histogram, degrading depth estimation accuracy and material classification performance \cite{tontini2024}. The SPAD sensor employed in this work incorporates a near-infrared (NIR) bandpass filter centred at the pulsed laser wavelength, which substantially attenuates out-of-band ambient contributions. However, in outdoor ranging scenarios, solar illumination contains significant NIR spectral components that can pass through the filter, contributing to the ambient noise floor. To mitigate ambient contributions, the average ambient photon counts across the histogram bins are estimated and subtracted, effectively removing the baseline offset while preserving the underlying signal structure \cite{tontini2024}. This baseline subtraction eliminates the DC component of ambient illumination but does not reduce the associated Poisson-distributed shot noise, which remains a fundamental limitation in low-light and high-ambient scenarios \cite{piron2020}.

\paragraph{Photon shot noise}
\label{par:shot_noise}

Shot noise, also termed Poisson noise, represents a fundamental noise source in SPAD-based time-of-flight systems arising from the quantum nature of photon detection \cite{piron2020}. Unlike deterministic signal components, photon arrival events exhibit inherent randomness with no predictable temporal ordering, a characteristic observed across all radiation detectors \cite{bass2010}. In material classification applications, shot noise can obscure relevant temporal features in histogram data, degrading the discriminative information required for accurate material identification.

Under static illumination conditions, the Poisson-distributed shot noise in a given histogram bin is characterised by a standard deviation equal to the square root of the mean photon count. For instance, a bin reporting a mean count of 100 photons over repeated measurements exhibits an estimated shot noise standard deviation of $\sqrt{100} = 10$ photons \cite{bass2010}. This fundamental relationship constrains the achievable signal-to-noise ratio in photon-limited scenarios, where the temporal resolution and material classification performance are ultimately limited by the statistical fluctuations inherent in the photon counting process.

\paragraph{Read noise}
\label{par:read_noise}

Conventional photodiode-based detectors exhibit read noise during signal readout, arising from multiple sources including kTC noise from charge storage capacitor reset operations, amplifier electronic noise, and quantization noise introduced during analog-to-digital conversion \cite{bass2010}. In contrast, SPADs demonstrate inherently low read noise characteristics due to their avalanche multiplication mechanism, which produces discrete output pulses with amplitudes significantly above the electronic noise floor \cite{piron2020}. This fundamental difference in signal generation enables SPAD-based systems to achieve superior signal-to-noise ratios in photon-limited detection scenarios, making them particularly advantageous for time-resolved material classification applications where precise temporal feature extraction is critical.

\paragraph{Dark counts and afterpulsing}
\label{par:dark_counts_and_afterpulsing}

SPADs operate with extreme sensitivity, enabling single-photon detection through avalanche multiplication. However, this high sensitivity introduces two fundamental noise mechanisms that can degrade time-of-flight (ToF) measurement accuracy. Dark counts arise from thermally generated charge carriers that spontaneously trigger avalanche events in the absence of incident photons. In ToF SPAD sensors, these spurious events are recorded as photon detections and are typically distributed uniformly across the histogram time bins, introducing background noise that is uncorrelated with the target signal \cite{piron2020}.

During an avalanche event, non-first-photon charges can become trapped in defect states within the semiconductor material. Following the SPAD reset period (dead time), these trapped charges may be randomly released, triggering subsequent false avalanche events known as afterpulsing \cite{piron2020}. In ToF SPAD systems, afterpulsing manifests as spurious secondary peaks in the photon time-of-arrival histogram that follow the primary signal peak from the target object, potentially causing range estimation errors.

Both dark counts and afterpulsing can distort histogram data quality and interfere with accurate temporal feature extraction. However, for the commercial SPAD sensor employed in this work, both effects are maintained at low levels (approximately 100 counts per second dark count rate, and less than 1\% afterpulsing probability), resulting in minimal impact on material classification performance \cite{piron2020}.

\paragraph{Optical cross-talk}
\label{par:optical_cross_talk}

Optical cross-talk represents an additional noise mechanism in SPAD arrays, where avalanche events in one pixel can emit secondary photons through electroluminescence that are subsequently detected by neighbouring pixels, resulting in spurious photon counts and potential range estimation errors \cite{piron2020}. This interference mechanism can degrade temporal histogram quality and introduce false detections that are spatially correlated with genuine signal events. However, the SPAD sensor module employed in this work incorporates optimised pixel isolation structures designed to minimise optical cross-talk, ensuring that this noise source does not significantly impact the temporal feature extraction and material classification performance .

\newpage
\section{Depth estimation techniques}
\label{sec:depth_estimation_techniques}

Depth estimation methods can be broadly categorised into stereoscopic, time-of-flight (ToF), and interferometric approaches. Stereoscopic methods reconstruct three-dimensional scene geometry from two-dimensional spatially resolved images by exploiting geometric parallax between multiple viewpoints \cite{laga2022,hu2020,godard2019}. These techniques can operate on monocular, binocular, or multi-view image sequences, with contrast-based feature matching enabling depth inference from equivalent 2D image pairs \cite{rock2015}. Despite their historical reliability and widespread adoption in computer vision systems, stereoscopic methods suffer from inherent limitations including computational complexity and significant post-processing overhead during depth map generation \cite{bass2010,laga2022}.

Recent research has explored hybrid approaches that combine multiple depth estimation modalities to improve accuracy and robustness. However, the evolution of optical time-of-flight technology, particularly SPAD-based direct time-of-flight (dToF) sensors, has demonstrated superior efficiency and effectiveness for material depth detection compared to traditional stereoscopic methods \cite{mcmanamon2019}. The following sections provide a brief review of stereoscopic depth estimation, followed by a more detailed examination of echolocation-based ToF depth sensing, and a discussion of interferometric methods including optical coherence tomography (OCT).

\subsection{Stereoscopy}
\label{subsec:stereoscopy}

The human visual system achieves depth perception through stereopsis, enabling qualitative representations of relative positions, distances, scale, and texture from scene features \cite{pizlo2001,wheatstone1838}. Stereoscopic depth estimation operates through three primary mechanisms: binocular vision, which exploits parallax from dual viewpoints \cite{hu2020,wheatstone1838}; monocular vision, which infers depth from single-view geometric and photometric cues \cite{godard2019}; and structure-from-motion, which estimates depth from temporal sequences of images \cite{schonberger2016}.

Stereoscopic depth estimation encompasses two fundamental paradigms: live scene depth estimation for real-time applications, and three-dimensional scene reconstruction for offline processing \cite{laga2022}. Live stereoscopy finds extensive application in autonomous vehicles, robotics, food inspection systems, and other real-time three-dimensional vision applications \cite{laga2022,hu2020}. Despite promising progress and widespread adoption, stereoscopic methods suffer from significant computational limitations, most notably the high post-processing overhead required for three-dimensional information extraction, which diminishes their attractiveness for many real-time systems \cite{laga2022,bass2010}.

While stereoscopy demonstrates effectiveness for biometric feature extraction, the extracted features rely predominantly on spatial cues such as edges, lines, and geometric patterns \cite{hsu2024}. This spatial dependency limits the ability to extract organic structural features encoding material composition, rendering stereoscopic systems less reliable for security clearance applications requiring material-based discrimination \cite{sharan2013,schwartz2013}. In practical implementations, this limitation has necessitated complementary validation mechanisms; for instance, facial recognition systems have been augmented with pulse metering to validate target humanness, compensating for the absence of material-based discrimination capabilities \cite{hsu2024}.

The fundamental principle underlying optical stereoscopic systems relies on extracting contrast-related features through geometric parallax analysis \cite{bass2010,laga2022}. This approach performs effectively for high-contrast scenes or images, whether achromatic or colour-encoded \cite{bass2010}. However, in low-contrast scenarios, the quality of extracted features may be insufficient for accurate depth estimation \cite{laga2022}. Achieving adequate contrast for long-distance depth estimation typically requires wider triangulation baselines, necessitating larger optical systems \cite{laga2022}. This fundamental trade-off between system size and performance limits the practical applicability of stereoscopy in many deployment scenarios \cite{bass2010,laga2022}. Despite these recognised limitations, stereoscopy has significantly influenced the design of optical depth sensing systems, establishing foundational principles that continue to inform depth sensing research and development \cite{laga2022,wheatstone1838}.

\subsection{Time of flight (ToF)}
\label{subsec:time_of_flight}

Time-of-flight (ToF) depth estimation, also known as echolocation, represents a more reliable approach to three-dimensional sensing compared to methods that infer depth from two-dimensional images \cite{tachella2019,jungerman2022}. The fundamental challenge of extracting depth information from spatially resolved intensity images has remained a persistent limitation in machine vision systems \cite{chen2024,laga2022}. In contrast, ToF methods utilise radiation propagation—whether mechanical, vibrational, or electromagnetic waves—to directly measure target distance, demonstrating superior effectiveness for three-dimensional sensing across numerous applications 

Beyond efficient depth estimation, ToF signals inherently encode reflectivity information that enables characterisation of material composition from the surfaces from which they emerge \cite{su2016,becker2023,yang2024,tafone2023}. This material discrimination capability, which forms the central theme of this work across objectives one through three, represents a significant advantage over purely spatial imaging approaches.

In the literature, optical ToF depth estimation is frequently associated with high-gain sensors, commonly referred to as ToF sensors \cite{piron2020}. However, within the broader context of optical sensing (see Chapter 3), any high-gain sensor with appropriate architecture—such as silicon photomultiplier (SiPM) arrays—can be configured for either stereoscopic or ToF depth sensing. SiPM architectures provide high gain and single-photon sensitivity through arrays of avalanche photodiodes operating in Geiger mode. (Figure~\ref{fig:2_2}).

\begin{figure}[H]
\centering
\includegraphics[width=0.95\textwidth]{figures/2.2.png}

    \caption[Direct time-of-flight operation illustrating signal transmission, energy-matter interaction, and detection.]{\textbf{Direct time-of-flight operation illustrating signal transmission, energy-matter interaction, and detection.} Transmitted signals (red arrows) represent strong illumination, with pulsed signals shown as discrete wave packets and continuous signals as uninterrupted waveforms. Following energy-matter interaction, scattered and partially absorbed return signals exhibit reduced intensity and propagate to the detector. The demodulation circuit within each pixel, employing either direct or indirect approaches, times photon arrivals for pulsed signals or counts phases for continuous signals to determine signal travel time, which is subsequently used to derive target material depth or distance.}
\label{fig:2_2}
\end{figure}


ToF sensors can be implemented using custom CMOS demodulation pixels, where timing precision is fundamentally limited by the demodulation bandwidth, typically on the order of 100 MHz \cite{bamji2022}. High-gain sensors overcome this bandwidth limitation and offer superior photon efficiency in ToF ranging applications \cite{piron2020}. Furthermore, while CMOS pixels are predominantly employed in indirect time-of-flight (iToF) systems, where precision and unambiguous range represent competing design constraints, high-gain sensors are well-suited for direct time-of-flight (dToF) implementations without such trade-offs \cite{bamji2022}. Consequently, high-gain detectors with single-photon sensitivity, such as single-photon avalanche diodes (SPADs), have become closely associated with ToF depth sensing applications \cite{piron2020}. The direct time-of-flight operation illustrated in Figure~\ref{fig:2_2} demonstrates how pulsed or continuous signals interact with target materials and return to the detector, enabling precise distance measurement through temporal analysis of photon arrival times.

Since optical ToF methods depend fundamentally on the constant speed of light, the following sections first review the speed of light before discussing the operating mechanisms and classifications of ToF depth estimation.

\subsubsection{Speed of Light}
\label{subsubsec:speed_of_light}

The determination of the speed of light has evolved through a series of experimental milestones that progressively refined measurement precision \cite{wikipedia2019}. Galileo Galilei conducted the first documented attempt in the 17th century, employing two observers positioned one kilometer apart who signaled each other by covering and uncovering lanterns. While this methodology lacks the rigor of modern experimental standards, Galileo's observations established a fundamental conclusion: light either propagates instantaneously or at a velocity far exceeding human temporal perception. The latter interpretation remains valid, as light travels at approximately $3 \times 10^8$ m/s.

The first quantitative evidence that light does not propagate instantaneously was provided by Rømer in 1676, who estimated the speed of light through astronomical observations of Jupiter's moon eclipses. This time-of-flight measurement confirmed that light propagation exhibits finite velocity, though the precision remained limited compared to contemporary standards. A more rigorous terrestrial measurement was achieved by Fizeau in 1849, who employed a rotating cogwheel as a mechanical shutter to modulate a pulsed light beam projected to a mirror approximately 8 km distant. By synchronizing the shutter rotation rate with the round-trip travel time, such that the returning signal was blocked by a cogwheel tooth, Fizeau estimated the speed of light to be approximately $3.153 \times 10^8$ m/s.

The speed of light is established as a fundamental physical constant with significance extending beyond classical electromagnetism, as formalized in Einstein's theory of special relativity. In modern signal processing and time-of-flight applications, this constant governs electromagnetic wave propagation and enables precise temporal-to-spatial conversion. The relationship between the speed of light $c$, wavelength $\lambda$, and frequency $\nu$ is expressed as:

\begin{equation}
c = \lambda \nu
\label{eq:speed_of_light}
\end{equation}
\addcontentsline{loe}{equations}{\protect\numberline{\theequation}$c = \lambda \nu$}

This fundamental relationship enables precise depth estimation in time-of-flight systems, where the round-trip propagation time of optical signals is directly converted to distance measurements through the constant speed of light.

\subsubsection{Time of flight (ToF) operation}
\label{subsubsec:tof_operation}

Optical time-of-flight measurement relies on the fundamental relationship between photon propagation distance and the constant speed of light to determine the temporal duration required for signals to travel from the illumination source to the target material and return to the detector. Distance calculation follows the relationship: $d = c \cdot t/2$, where $d$ represents distance, $c$ is the speed of light, and $t$ is the round-trip time, with the factor of one-half accounting for the return path.

For a distance resolution of 1~mm, the fundamental timing precision requirement is approximately 6.7~picoseconds, derived directly from the speed of light constant. However, through repeated measurements and statistical signal processing techniques, such resolution can be achieved with lower individual timing precision by employing temporal averaging or interpolation methods that reduce measurement uncertainty through ensemble statistics .

Conventional ToF and LIDAR systems historically employed bulky laser scanners with opto-mechanical scanning mechanisms to cover the field of view, resulting in large system footprints that required careful alignment and were susceptible to mechanical failures from shocks and vibrations \cite{piron2020}. Modern time-of-flight systems, including the imager used in this work, utilise laser sources that flood-illuminate the entire scene simultaneously, eliminating mechanical scanning requirements and enabling more compact, robust system architectures \cite{piron2020}.

Unlike stereoscopic systems, where wider baseline separation between cameras improves depth estimation accuracy, increasing the distance between the illumination source and receiver in ToF systems can introduce unwanted artefacts such as shadows due to parallax effects. Therefore, ToF systems typically employ co-located or closely positioned illumination and detection components to minimise such artefacts \cite{bamji2022}.

The choice of demodulation method directly influences the temporal signal characteristics and output data format. Different demodulation techniques—including direct time-of-flight (dToF), indirect time-of-flight (iToF), and frequency-modulated continuous wave (FMCW)—produce distinct temporal signal profiles, each with specific advantages and limitations \cite{bamji2022}. The following sections examine each mechanism in detail, along with their respective advantages and limitations relevant to this project.

\subsubsection{Direct Time of Flight}
\label{subsubsec:direct_tof}

Direct time-of-flight (dToF) depth estimation determines target distance by directly measuring the round-trip propagation time of optical signals. The system employs a pulsed laser source to illuminate the scene while registering photon arrival times, enabling precise temporal measurement of individual photon events in single-photon dToF implementations. By compressing laser energy into short pulses, the system reduces the impact of uncorrelated ambient photons, improving signal-to-noise ratio. The received temporal signal profile encodes optical properties of the target material, including surface reflectivity, absorption, and subsurface scattering characteristics, motivating the selection of dToF for this project \cite{su2016,becker2023}.

The operational range of dToF systems scales with laser pulse energy due to the inverse-square law governing backscattered signal intensity. However, maximum permissible laser power is constrained by eye safety regulations. The SPAD detector employed in this work utilises a Class 1 laser source operating at 905 nm wavelength, ensuring eye-safe operation with an unambiguous range of 5 m.

Signal timing and readout circuitry constitute critical components of dToF systems. In the detector used for this work, each pixel incorporates an electrical circuit with a high-voltage supply and a series resistor connected to the photodiode for avalanche quenching. The resulting signals are digitised through an inverter stage and subsequently shortened to prevent avalanche overlap when multiple SPAD outputs are combined. Digital events are temporally registered using a time-to-digital converter (TDC) synchronized with the pulsed laser source \cite{sesta2023,nissinen2009,richardson2009}. Photon timestamps are accumulated over multiple laser cycles during the exposure period to construct photon timing histograms, which are subsequently read out for analysis \cite{piron2020}.

Two primary exposure modes are employed: global shutter and rolling shutter. In global shutter operation, all pixels activate simultaneously for a fixed exposure duration, followed by row-by-row readout during which photon detection is disabled. Rolling shutter operation activates and reads out pixel rows sequentially, enabling back-to-back exposures with minimal readout dead time. In both configurations, the acquired data undergoes decoding and processing, typically on a host computer or microcontroller, to extract depth information or perform histogram-based analysis as implemented in this project \cite{piron2020}.

\paragraph{Direct time of flight timing circuitry}
\label{par:direct_tof_timing_circuitry}

Timing circuits are fundamental to depth measurement in direct time-of-flight (dToF) systems. A high-speed digital clock initiates upon laser pulse emission and terminates upon detection of the first returning photon, with the resulting timestamp providing the distance measurement. While effective for single-detector configurations, deploying fast clocks across large detector arrays incurs significant power consumption. Consequently, alternative designs employ simplified, lower-precision timing architectures to reduce power and memory requirements. Several timing methodologies have been developed, each offering distinct trade-offs between temporal resolution, system complexity, processing throughput, power efficiency, and silicon area. These include time-to-amplitude converters (TACs), ring oscillator time-to-digital converters (RO TDCs), multi-event time-to-digital converters (TDCs), and multiphase clock architectures.

Time-to-amplitude converters (TACs) represent an analog timing approach offering reduced power consumption relative to digital alternatives. Despite susceptibility to noise and component mismatch, TACs achieve high temporal precision. Studies by Kumagai \textit{et al.} \cite{kumagai2021} and Parmesan \textit{et al.} \cite{parmesan2014} employed this methodology, where each detected photon samples a shared ramp signal across the detector array. The implementation by Parmesan \textit{et al.} \cite{parmesan2014} achieved a timing resolution of 6.66 picoseconds (ps).

Ring oscillator time-to-digital converters (RO TDCs) operate by generating per-pixel clock signals upon timing initiation. A counter increments with each oscillator cycle completion, and upon timing termination, the oscillator's phase position within its cycle is converted to temporal information. The first RO TDC implementation utilised a 16-stage differential ring oscillator with a 10-bit loop counter \cite{nissinen2003}. Subsequent improvements achieved 250 ps precision using delay lines and expanded to a complete range-finding system \cite{nissinen2009}. The work by Nissinen \textit{et al.} \cite{nissinen2009} represents the first dToF SPAD implementation employing the RO TDC methodology.

Multi-event time-to-digital converters (multi-event TDCs) enable multiple timestamp recording per measurement frame, enhancing photon processing capability. This capability is critical because single-photon-per-pixel-per-laser-cycle detectors are susceptible to photon pileup, where multiple photon arrivals result in only one detectable event, degrading performance under high illumination conditions by reducing the probability of detecting target-reflected signal photons. The SPAD detector employed in this work is believed to utilise a multi-event TDC architecture. Richardson \textit{et al.} \cite{richardson2009} implemented a multi-event TDC for a digital silicon photomultiplier (D-SiPM) capable of timestamp recording at 10 gigasamples per second (GS/s), addressing the photon pileup limitation.

In certain dToF SPAD implementations, multiphase clocks are distributed to each time-to-digital converter. A phase-locked loop (PLL) or delay-locked loop (DLL) generates these multiple clock phases. During a SPAD detection event (Figure~\ref{fig:2_3}), the state of each clock phase is recorded, providing enhanced timing resolution.

The multiphase clock approach employs a simple counter to track complete clock cycles from start to stop. Although this method requires additional circuitry for multiphase clock distribution, it offers improved timing stability compared to ring oscillator implementations \cite{alabbas2018}. Figure~\ref{fig:2_3} shows the SPAD signal detection path and voltage progression during the avalanche and quenching process, illustrating the key operational stages of single-photon detection in time-of-flight applications.


\paragraph{Direct time of flight timing histograms}
\label{par:direct_tof_timing_histograms}

While photons returning from target surfaces arrive in continuous time, detected photons are allocated into discrete time periods, referred to as bins, for processing purposes. These bins typically match the time-to-digital converter (TDC) resolution and span a specific time range (approximately 5 meters for the SPAD detector used in this work). In the signal preprocessing pipeline, each bin is populated with the total number of reflected photons detected within that time period to create a histogram. This process constitutes time-of-flight histogramming in SPAD-based systems, including the detector employed in this project.

The primary advantage of histogramming for this work lies in its facilitation of material composition feature extraction. From the 144-bin time-resolved signal, it is possible to extract target reflectivity bins for the material detection processing pipeline. Traditionally, histograms were processed as a post-readout step;

\begin{figure}[H]
\centering
\includegraphics[width=0.95\textwidth]{figures/2.3.png}

    \caption[SPAD signal detection path and voltage progression during avalanche and quenching.]{\textbf{SPAD signal detection path and voltage progression during avalanche and quenching.} Left: SPAD signal detection path. Centre: voltage progression from high voltage till quenched. Right: Key for left and centre diagrams.}
\label{fig:2_3}
\end{figure}

however, the overhead of frame readout rate scales with detector array size, necessitating on-chip implementation of the histogramming process. Memory requirements for embedded histograms are often the limiting factor in array resolution, as they scale with desired precision and detectability \cite{gyongy2022}.

On-chip per-pixel histogramming was initially demonstrated by \cite{piron2020} using a 10-bit ripple counter for 32 bins. This approach was further advanced by \cite{erdogan2017}, who implemented a multiple-event flash TDC with a single digital silicon photomultiplier (D-SiPM), enabling per-phase histogram generation from a 33-phase ring oscillator. Rather than recording separate event timestamps, this architecture creates a histogram for each phase from the one-hot output of the folded flash time-to-digital converter.

In 2018, \cite{alabbas2018} implemented a D-SiPM array where histograms are generated for each illumination pulse, representing a significant advancement in reducing readout rates for direct time-of-flight SPAD systems. This work employs on-chip histograms for signal processing inputs. Readout data bandwidth has been further reduced by \cite{hutchings2019}, who reported an SPAD direct time-of-flight image sensor with in-pixel, partial histogramming. This approach enables reading only histogram bins corresponding to the peak and its immediate surroundings (cropped histogram), which is particularly relevant for signal processing applications.

Using a direct time-of-flight flash LiDAR system, \cite{zhang2019} achieved depth images at short distances with a frame rate of 30 frames per second. In the present work, however, the task of peak extraction is performed off-chip using custom-built algorithms of low computational cost. Retrieval of high-quality information (cropped histogram with centre of mass) off-chip is demonstrated at both lower computational cost and longer distance ranges compared to \cite{zhang2019}. Currently, several commercial sensors \cite{zhang2019,vl53l5cx2023,tmf8801} incorporate on-chip histogram features that utilise I2C interfaces for data readout.

\subsubsection{Indirect Time of Flight}
\label{subsubsec:indirect_tof}

Unlike direct time-of-flight (dToF) systems that employ high-precision timers to measure the direct temporal delay of photon pulses, indirect time-of-flight (iToF) systems typically estimate distance by measuring the phase difference between transmitted and received amplitude-modulated or frequency-modulated optical signals \cite{bamji2022}. At the demodulation circuit, received signals are integrated over different time windows defined by reference clocks synchronized with the transmitted signal. The integrated signals from these time gates are subsequently converted into a phase difference measurement \cite{bamji2022}. The measured phase difference encodes the target distance, which is directly proportional to the phase shift and can be calculated using the modulation frequency and the speed of light.

The distance measurement in iToF systems is derived from the modulation frequency $\nu$ and the measured phase difference $\varphi$ according to
\begin{equation}
d_{\text{target}} = \frac{c \cdot \varphi}{4 \pi \nu}
\label{eq:itof_distance}
\end{equation}
\addcontentsline{loe}{equations}{\protect\numberline{\theequation}$d_{\text{target}} = \frac{c \cdot \varphi}{4 \pi \nu}$}%
where the phase shift is converted to time delay, and $d_{\text{target}}$ represents the one-way distance to the target. Higher modulation frequency $\nu$ ensures more precise measurements but results in a shorter unambiguous range. However, frequency selection is constrained by bandwidth limitations in either the transmission or receiving channels \cite{bamji2022}. The accuracy of iToF depth estimation is affected by several factors: detector sensitivity (quantum efficiency for CMOS pixels), noise sources (photon shot noise, readout noise), readout rate, integration window length, integration time, and optical signal power. Additional architecture-dependent limitations also exist \cite{bamji2022}.

For most iToF sensors, the unambiguous range extends to a few meters, which is often extended through multi-frequency modulation techniques \cite{bamji2022}. However, return signals from target objects within the unambiguous range can still be corrupted by multi-path reflections, depending on the target surroundings. Multi-path reflections are typically less problematic in dToF sensors, as they appear as separate histogram peaks that can be distinguished and processed independently. The performance of iToF depth estimation depends on the modulation technique employed. These techniques are primarily categorised as pulsed (p-iToF) and continuous wave (AM-CW, FM-CW) methods, as summarised in the following sections.

\paragraph{Pulsed iToF}
\label{par:pulsed_itof}

Pulsed indirect time-of-flight (iToF) employs periodic radiation pulse bursts to estimate time-of-flight through temporal signal analysis. In high-gain detectors such as single-photon avalanche diodes (SPADs), signal integration utilises a digital counting methodology where individual photon arrivals are counted and routed to dedicated counters within the demodulation circuitry based on their detection time window \cite{piron2020}. This time-gated single-photon counting (TGSPC) process represents a specific implementation of pulsed iToF, enabling precise temporal resolution through windowed photon detection \cite{bamji2022}.

The implementation architecture determines the processing efficiency: single-counter configurations require sequential TGSPC operation, processing one time window per measurement cycle. In contrast, multi-counter architectures enable simultaneous signal integration through parallel reference phase-shift clocks, facilitating concurrent processing of multiple time windows \cite{bamji2022}. Alternatively, CMOS demodulation pixels employ multiple charge collection gates that activate during distinct time windows synchronized with the transmitted signal, providing an additional pathway for simultaneous multi-window signal integration \cite{bamji2022}.

\paragraph{Continuous wave (CW) iToF}
\label{par:continuous_wave_itof}

Continuous wave indirect time-of-flight (CW iToF) systems employ lock-in amplification to extract phase information from modulated optical signals, enabling distance estimation through phase difference measurements between transmitted and received waveforms \cite{bamji2022}. The lock-in amplifier performs coherent demodulation, selectively extracting the phase-shifted signal component at the modulation frequency while suppressing noise and interference. This approach enables precise phase measurement and subsequent distance calculation based on the phase delay. CW iToF implementations are primarily categorised into amplitude-modulated (AM-CW) and frequency-modulated (FM-CW) techniques, each offering distinct signal processing characteristics and performance trade-offs \cite{bamji2022}.

\subparagraph{Amplitude-modulated continuous wave (AM-CW)}
Amplitude-modulated continuous wave (AM-CW) represents a prominent approach for optical time-of-flight imaging, where the radiation source is continuously modulated using waveforms such as pseudo-noise sequences, sinusoidal, or rectangular signals at frequencies typically ranging from approximately 10 to 100 MHz \cite{bamji2022}. The phase difference between the continuously transmitted and received signals is measured through demodulation circuitry, with the integration period commonly partitioned into four time windows to enable phase measurement and distance calculation \cite{bamji2022}. This multi-window integration approach facilitates robust phase extraction by sampling the modulated signal at different phase offsets, allowing reconstruction of the phase relationship between transmitted and received waveforms.

\subparagraph{Frequency-modulated continuous wave (FM-CW)}
Frequency-modulated continuous wave (FM-CW) iToF estimates target depth by measuring the frequency difference between continuously emitted and returned signals, representing an alternative indirect time-of-flight technique that utilises continuous wave frequency modulation rather than pulsed signal transmission \cite{bamji2022}. The evolution of FM-CW techniques began approximately seven years after the invention of the laser, with early development at MIT Lincoln Laboratory in the 1960s \cite{gschwendtner2000,patel1968}. FM-CW iToF systems require laser sources with extended coherence lengths compared to other ToF methods, imposing constraints on system design and component selection \cite{bamji2022}. These requirements increase system complexity and cost, while known limitations in range resolution and detection capabilities for smaller objects may restrict applicability in certain scenarios, particularly for long-range measurements and large-scale data collection applications \cite{bamji2022}.

\subsection{Interferometry}
\label{subsec:interferometry}

Interferometry represents a fundamentally distinct yet conceptually analogous approach to depth sensing compared to echolocation. This technique exploits the interference properties of superimposed electromagnetic waves to extract depth information from target surfaces through phase difference analysis \cite{huang1991,aumann2019}. The underlying principle relies on the coherent superposition of two waves of identical frequency, generating an intensity modulation pattern that encodes the phase relationship between the emitted and reflected waveforms. By analysing these interference patterns, the phase difference can be precisely determined, enabling accurate depth estimation of the target surface.

Unlike direct time-of-flight (dToF) methods, which measure the round-trip propagation time of discrete optical pulses, interferometric depth sensing employs continuous wave (CW) illumination sources, typically characterised by low temporal coherence \cite{huang1991,aumann2019}. The low coherence length is a critical design parameter that constrains interference to occur only when the optical path lengths of the reference and sample arms are matched within the coherence length, thereby enabling precise depth resolution and localisation of reflecting interfaces. This coherence-gated detection mechanism provides superior depth resolution compared to conventional CW interferometry while maintaining the advantages of continuous illumination for signal integration and noise reduction.

\paragraph{Optical coherence tomography (OCT)}
\label{par:optical_coherence_tomography}

Optical coherence tomography (OCT) \cite{huang1991} employs coherent illumination sources to perform depth-resolved imaging through cross-sectional reconstruction (tomography). The technique utilises low-coherence interferometry to generate high-resolution volumetric representations of target structures, with primary applications in medical imaging for non-invasive visualisation of internal tissue architecture \cite{aumann2019}. OCT implementations are categorised into three principal modalities based on their signal acquisition and processing architectures: time-domain OCT (TD-OCT), spectral-domain OCT (SD-OCT), and swept-source OCT (SS-OCT) \cite{aumann2019}. Each modality employs distinct signal processing strategies that yield different performance characteristics in imaging speed, axial resolution, and signal-to-noise ratio.

TD-OCT represents the initial implementation of OCT technology, utilising sequential depth scanning through mechanical translation of a reference mirror to probe different optical path lengths. This scanning mechanism requires temporal sampling of interference signals at each depth position, resulting in inherently slow acquisition rates and elevated computational processing requirements. The sequential nature of TD-OCT signal acquisition limits its applicability in time-critical applications such as clinical diagnostics, where both imaging speed and signal quality are essential.

SD-OCT addresses the speed limitations of TD-OCT by implementing parallel spectral detection. The system employs a broadband light source and a stationary spectrometer to simultaneously capture interference spectra across all depth positions in a single acquisition. By detecting the full interference spectrum without mechanical scanning, SD-OCT achieves significantly higher acquisition speeds and improved signal-to-noise ratio compared to TD-OCT, enabling real-time cross-sectional imaging capabilities.

SS-OCT employs a wavelength-swept laser source that rapidly scans across a broad spectral range, with depth information extracted through Fourier transformation of the time-resolved interference signal detected by a single photodetector. This approach eliminates the spectrometer, enabling millimeter-range axial resolution with high-speed imaging capabilities. The swept-source architecture provides enhanced penetration depth compared to SD-OCT, making it particularly suitable for deep tissue imaging applications in ophthalmology and cardiovascular imaging, where extended imaging range is required.

\newpage
\section{Role of SPAD time of flight sensors in machine vision}
\label{sec:spad_tof_machine_vision}

SPAD sensors, encompassing both direct and indirect time-of-flight (ToF) modalities, have established significant roles in machine vision systems. The predominant implementations leverage depth estimation capabilities to enable diverse applications, including robotic navigation systems for industrial logistics, autonomous lawn mowers, and unmanned aerial vehicles, as well as consumer electronics such as security systems, coffee machines, and mobile devices featuring face recognition and people counting functionalities. This work explores an unconventional application of SPAD technology: the extraction of material composition information from time-resolved signals.

While the transient signal inherently encodes reliable depth information, the temporal distribution of photon arrival times—manifested as reflectivity profiles—contains richer material-dependent information. These profiles are influenced by both surface and subsurface optical properties of the observed target, rendering them sensitive to material composition \cite{su2016}. Three experimental demonstrations are presented to validate this application of SPAD time-resolved signals in machine vision.

The first experiment focuses on material feature extraction from geometrically simple planar objects. The returned signal profiles embed signature structural patterns that correlate with the material composition of the surface from which they originate. Previous work has compared physics-based models with neural network approaches for classifying these features \cite{su2016}. In \cite{su2016}, neural network models demonstrated superior performance over physics-based models in material classification tasks, motivating the adoption of machine learning classifiers for material classification objectives in this work.

The second demonstration integrates SPAD time-of-flight signals with conventional spatially resolved RGB images for object detection. The machine vision research community recognises that accurate three-dimensional environmental perception requires the extraction of both spatial and structural features \cite{chen2024}. Spatial features encompass lines, vectors, angles, two-dimensional Cartesian geometry, and higher-level topological properties of the target. Structural features, in contrast, consider depth and the internal and external composition of the target as it interacts with the illumination source. In this work, two-dimensional spatially resolved images are processed for spatial feature extraction, while SPAD transient or time-resolved signals are analysed for structural features.

Finally, transient signals from SPAD sensors are employed to investigate material purity. The test material is homogenised milk, with purity estimation achieved by comparing sample transient profiles against previously learned ground truth measurements from pure milk samples. Eight experiments are conducted to evaluate accurate purity detection across a range of different contaminants.

\newpage
\section{Material detection with SPADs}
\label{sec:material_detection_spads}

Prior to the availability of compact and cost-effective SPAD sensors, material detection relied on alternative optical approaches, with spectroscopy serving as a prominent foundation for numerous material detection innovations \cite{holt2016}. However, spectroscopy-based methods often exhibit significant complexity, requiring bulky equipment, expensive instrumentation, and sophisticated spectral analysis \cite{holt2016}. This complexity motivates the adoption of compact, cost-effective alternatives such as single-photon avalanche diode (SPAD) sensors \cite{piron2020}.

SPADs operate as range sensors employing an echolocation mechanism for target sensing, enabling efficient distance estimation through round-trip time-of-flight measurements of emitted light pulses and the speed of light. Beyond depth estimation, the temporal reflectivity profile of the returned signal encodes material-specific information that enables material discrimination \cite{su2016,becker2023,tafone2023}. Different materials exhibit distinct temporal reflectivity patterns based on their optical properties, scattering characteristics, and internal structure, enabling material classification through temporal signal analysis \cite{su2016,becker2023,tafone2023,yang2024}.

Across all three objectives in this work, SPAD sensors serve as the primary instrument for acquiring time-resolved signals in the form of photon time-of-arrival histograms. These histograms capture the temporal distribution of reflected photons, providing rich information about material properties. The temporal data is processed to extract two key features: the centre of mass (CoM), representing the average arrival time and relating to target distance, and a cropped array of target reflectivity (histogram peak) per SPAD pixel (zone), which captures material-specific reflection characteristics. These extracted features are subsequently fed to a classifier to learn material composition features, enabling automated material detection and classification.

Detailed mechanisms are presented in Objective One (Chapter 4), where reflectivity profiles of translucent and opaque materials were employed for flat surface material composition classification. The experiments demonstrate that temporal reflectivity patterns contain sufficient information to distinguish between different materials, even when spatial appearance features are similar \cite{su2016,becker2023,tafone2023}. Previous research has often focused on short-range and/or fixed distances under controlled environments, limiting practical applicability \cite{su2016}. The presented experiments utilised longer distances and less controlled environments, demonstrating the robustness and practical viability of the approach under real-world conditions.

Comprehensive comparisons of all experiments with corresponding state-of-the-art solutions are detailed in the respective chapters (see Chapters 4, 5, and 6). Each chapter presents specific applications of SPAD-based material detection, validates the approach through comprehensive experiments, and compares performance with existing methods \cite{yang2024,tafone2023,becker2023}.

\subsection{Models for material detection}
\label{subsec:models_material_detection}

As established in Chapter 1 (Figure~\ref{fig:1_1}), visual detection models may be optimised for spatial or structural feature extraction. The majority of state-of-the-art visual models employ neural network architectures due to their effectiveness in modelling nonlinear relationships inherent in complex real-world objects.

The selection of deep neural networks in this work is motivated by the requirement to overcome signal non-idealities, including noise and SPAD timing jitter, while extracting discriminative features from material temporal responses. Neural networks excel at learning complex, nonlinear relationships in data, making them particularly well-suited for extracting material-discriminative features from noisy temporal signals where traditional signal processing methods may struggle.

This capability is not unique to neural network classifiers, as physics-based optical models have demonstrated comparable performance in certain scenarios \cite{becker2023}. Physics-based models leverage known physical principles and material optical properties to predict temporal responses, providing interpretable results grounded in optical theory. However, in material detection experiments, neural networks have been shown to outperform physics-based models \cite{becker2023}, particularly when dealing with complex, heterogeneous materials or when signal quality is compromised by noise and measurement uncertainties.

Throughout this project, logistic regression classifiers are employed for modelling time-resolved features in material detection tasks. Logistic regression provides a computationally efficient baseline that effectively classifies materials based on extracted temporal features while maintaining interpretability and requiring minimal training data. For complex, data-intensive spatial models operating on RGB images, transfer learning with pre-trained networks \cite{bozinovski2020} enables faster training by leveraging features learned from large-scale image datasets while adapting to the specific material detection task. In this work, pre-trained YOLO and DINO models were fine-tuned on the custom natural objects dataset for the experiments, ensuring that while the models benefited from pre-training, they were specifically adapted to the selected natural objects and experimental conditions.

\subsection{Neural network approaches}
\label{subsec:neural_network_approaches}

At lower levels of abstraction, machine vision systems detect edges in images, which correspond to boundaries and outlines defining the structural composition of visual scenes. Edge detection algorithms, including gradient-based operators (Sobel, Canny) and learned filters in deep networks, identify discontinuities in pixel intensity that correspond to object boundaries, texture transitions, and geometric features. These low-level processing tasks are particularly valuable for image segmentation, where pixels are grouped into coherent regions based on similarity criteria such as colour, texture, or intensity. Segmentation techniques enable the isolation and identification of distinct objects or regions within complex scenes, with critical applications in medical imaging for delineating tumours or regions of interest within diagnostic images, where precise boundary detection is essential for accurate diagnosis and treatment planning. Similarly, in material detection and scene understanding, edge detection and segmentation serve as foundational preprocessing steps that enable the separation of different material regions and objects before further analysis.

Classification operates at a higher level of abstraction, assigning semantic labels to detected regions or objects after successful identification or segmentation. This enables machines to understand not only what is present in a scene, but also the category or class to which each element belongs. In non-linear datasets where traditional linear classifiers fail to capture complex feature relationships, convolutional neural networks (CNNs) \cite{lecun1998,lecun1989} have demonstrated superior efficiency in extracting hierarchical features for classification tasks. CNNs excel at learning spatial hierarchies of features through their convolutional layers, which apply learnable filters that detect local patterns such as edges, textures, and shapes. These local features are progressively combined through deeper layers to form more complex, abstract representations that capture object-level and scene-level semantics. The hierarchical feature learning capability of CNNs, combined with their translation invariance properties and parameter sharing mechanisms, makes them particularly well-suited for visual recognition tasks where spatial relationships and local patterns are crucial.

The adoption of CNNs in custom material detection models has become standard practice. Material detection presents unique challenges compared to object detection, as it requires the identification of intrinsic material properties (such as optical reflectivity, surface composition, or mechanical characteristics) rather than merely recognising object shapes or appearances. Each objective in this work employs a custom CNN architecture tailored for material-specific processing and classification, designed to leverage the unique characteristics of the input data modalities (such as time-of-flight signals, intensity images, or multi-modal fusion). These custom architectures are optimised to extract material-discriminative features that distinguish between different material compositions, even when objects share similar spatial patterns or appearances. The CNN architectures are carefully designed to handle the specific characteristics of the sensing modalities employed, whether processing transient time-of-flight signals, intensity images, or fused multi-modal data. Detailed descriptions and expansions on CNN architecture per experiment are found in the corresponding chapters (see Chapter 4, Chapter 5, and Chapter 6), where each chapter presents the specific network design, training procedures, and performance evaluation for the material detection task at hand.

Chapters 4, 5, and 6 each instantiate a custom CNN variant tuned to the ground-truth datasets of their respective studies, ensuring that architecture choices and training schemes are dataset-specific rather than one-size-fits-all. Figure~\ref{fig:2_4} shows a generic CNN template; layer depths, kernel sizes, normalization, and task heads are specialised per dataset.


\begin{figure}[H]
    \centering
    \includegraphics[width=0.8\textwidth]{figures/2.4.png}
    \caption[Core CNN architecture for N-dimensional pixel input.]{\textbf{Core CNN architecture for N-dimensional pixel input.} Illustration of core CNN architecture for N-dimensional pixel input; model input sequence, convolution of input sequence and subsampling of convolved input to fully connected layer and logits for prediction. Kernel size and number of filters are varied per layer.}
    \label{fig:2_4}
\end{figure}


Other neural network architectures that facilitate segmentation and classification tasks include encoder-decoder networks and transformer-based models with attention mechanisms \cite{vaswani2017}. These architectures are particularly effective for tasks requiring precise spatial localisation, such as semantic segmentation, where each pixel must be classified. Sequential models represent another specialised architecture most suitable for sequential data, commonly found in applications for natural language processing (NLP) and time series analysis. Common sequential models include recurrent neural networks (RNNs) \cite{schmidt2019} and Long Short-Term Memory networks (LSTMs) \cite{hochreiter1997}; neither of which have been used in this work.

Generative networks may be used to synthesise features from objects or whole scenes. Chapter 2 considers Generative Adversarial Networks (GANs) \cite{goodfellow2020} for synthesising material features. Details of the purpose for which they were chosen alongside the classifiers can be read under the experiment chapters (Chapters 4 and 6). Neural networks may be supervised (trained on previously labelled data) or unsupervised (with no prior ground truth training). Supervised learning has been used throughout experiments to demonstrate the learnability of material features in the presented solutions.

Beyond individual vision tasks, machines are sometimes trained to perform hybrid decisions that combine multiple objectives. The most prominent such decision is object detection: a combination of target localisation and classification tasks in a single framework. Chapter 5 of this project performs enhanced object detection by including material classification. The most common object detection systems are single-shot detectors (SSDs) \cite{liu2016}; specifically, YOLO (You Only Look Once) variants \cite{redmon2015,redmon2016,redmon2018}. Chapter 5 implements YOLOv3 \cite{redmon2018}, DINOv3 \cite{simeoni2025}, and YOLOv8 as demonstrations of multiple spatial object detection systems being enhanced by the presented material detection solution.

\subsubsection{Convolutional neural networks}
\label{subsubsec:convolutional_neural_networks}

This section outlines the CNN architecture, feature extraction, and classification pipeline for material detection. The CNN model accepts spatial RGB images or fused multimodal images with dimensions $H \times W \times C$, where $H$ and $W$ represent spatial dimensions (height and width) and $C$ denotes the number of channels (e.g., 3 for RGB, 1 for greyscale, or more for multimodal fusion). The overall architecture is illustrated in Figure~\ref{fig:2_4}.
Convolutional blocks progressively extract hierarchical features using learnable kernel filters \cite{lecun1998}. Each convolutional layer applies filters that detect local patterns such as edges, textures, and material-specific features. Lower layers capture simple patterns (edges, basic textures), while deeper layers combine these into complex, material-discriminative patterns. Each convolutional layer is followed by a ReLU (Rectified Linear Unit) activation function, which introduces non-linearity by setting negative values to zero, enabling the network to learn complex material-discriminative patterns that would be impossible with linear transformations alone.

Max pooling layers apply strides (typically $2 \times 2$) to reduce spatial dimensions by selecting the maximum value within each $2 \times 2$ window. This operation decreases computational load while maintaining important features and providing translation invariance—the ability to recognise features regardless of their exact spatial position. This reduces feature map dimensions by half at each pooling stage (e.g., $H \times W \rightarrow H/2 \times W/2 \rightarrow H/4 \times W/4$), progressively compressing spatial information while preserving the most salient features.

The flattening layer converts the 2D feature maps from the last convolutional layer into a 1D vector, preparing the multi-dimensional feature representations for the fully connected classification layers. This transformation reshapes the feature maps into a single vector that can be processed by dense layers. Fully connected (FC) layers are dense layers that process the flattened features for high-level material classification. These layers combine features from across the entire spatial extent of the input, enabling the network to make global decisions about material identity. They also use ReLU activation to introduce non-linearity in the classification path, allowing the network to learn complex decision boundaries. Dropout regularisation is applied between fully connected layers with a probability threshold (e.g., 0.2 or 0.5) to randomly set a fraction of neurons to zero during training. This helps prevent overfitting by forcing the network to learn redundant representations and improving generalisation to unseen material samples. During inference, all neurons are active, but their outputs are scaled to account for the dropout probability used during training. The output layer produces material class predictions using softmax activation, which converts raw scores (logits) into probability distributions over material classes. The number of output neurons equals the number of material classes to be detected.

\paragraph{Understanding Logits vs. Probabilities}
Logits and probabilities represent different stages of the prediction process. Logits are raw, unnormalised scores that can be positive, negative, or zero, are not bound by any range, and unlike probabilities, do not sum to 1 or 100\%. They represent the relative confidences before normalization. For instance, an 8-class output may generate logits such as $[2.3, -1.5, 0.8, 4.2, -0.3, 1.1, 0.5, 3.0]$ that do not sum to one. When softmax (Equation~\ref{eq:softmax}) is applied to the logits, values become normalised between the range of 0 and 1, and all probabilities sum to 1.0. The highest logit values are translated into the highest probabilities. For example, the logits above might become probabilities $[0.15, 0.02, 0.05, 0.73, 0.03, 0.08, 0.04, 0.20]$, where the class with logit 4.2 receives the highest probability (0.73). Probabilities are easier to interpret as they represent the model's confidence that the input belongs to each material class, with values directly comparable and summing to 100\%.

\begin{equation}
\label{eq:softmax}
p_{\text{class}} = \frac{e^{\text{logit}_{i}}}{\sum_{j=1}^{K_c} e^{\text{logit}_{j}}}
\end{equation}
\addcontentsline{loe}{equations}{\protect\numberline{\theequation}$p_{\text{class}} = \frac{e^{\text{logit}_{i}}}{\sum_{j=1}^{K_c} e^{\text{logit}_{j}}}$}

\newpage
\section{Conclusion}
\label{sec:conclusion_chapter2}

This chapter has reviewed the landscape of material detection and scene understanding approaches, revealing fundamental limitations in existing methods and establishing the foundation for the solutions presented in subsequent chapters. The analysis demonstrates that current optical sensing approaches predominantly operate in isolation, focusing on either spatial intensity features or spectral information, while overlooking the complementary benefits of temporal signal analysis for material characterisation \cite{schwartz2013,sharan2013,su2016,becker2023}.

Spatial intensity-based methods, including RGB and non-RGB approaches, capture appearance and texture features but exhibit critical limitations in distinguishing material composition from visual appearance \cite{schwartz2013,sharan2013,chen2024}. These methods conflate material identity with surface characteristics, making them vulnerable to misclassification when materials share similar visual traits or when synthetic representations mimic authentic materials—a fundamental limitation in safety-critical applications requiring material authenticity verification \cite{su2016,becker2023,yang2024}.

Spectroscopic methods offer rich wavelength-dependent information but require sophisticated instrumentation, extensive computational resources, and multi-dimensional spectral analysis \cite{holt2016,elachi2021,heinz2001,vasile2023}. The review establishes that temporal, non-spectral (achromatic) approaches can achieve comparable material discrimination capability without the computational overhead associated with spectral analysis \cite{su2016,becker2023,tafone2023}. This finding validates the non-spectral signal processing paradigm employed throughout this work, demonstrating that material-specific optical properties can be effectively captured through time-resolved interactions rather than wavelength decomposition.

Depth estimation techniques, including stereoscopy and time-of-flight methods, have been examined with particular emphasis on direct time-of-flight (dToF) approaches \cite{laga2022}. While stereoscopic methods provide relative depth information, they suffer from computational complexity and limited capability to extract material-specific structural features \cite{laga2022}. Time-of-flight methods, particularly dToF implementations using SPAD sensors, offer superior temporal resolution and single-photon sensitivity that enable both accurate depth estimation and material characterisation through temporal reflectivity profiles \cite{su2016,becker2023,tafone2023}.

The review of material detection instruments establishes SPADs as optimal sensors for time-resolved material detection due to their single-photon sensitivity, picosecond temporal resolution, and compact form factor \cite{piron2020}. Unlike conventional photodiode-based detectors that capture intensity distributions, SPADs enable extraction of material-specific temporal signatures that encode structural and compositional information inaccessible through spatial appearance alone \cite{su2016,becker2023}.

The fundamental gap identified in current literature lies in the limited exploitation of spatiotemporal fusion, where spatial detection systems provide distance and position information while material detection systems provide composition and property information \cite{gadzicki2020,yang2023,yang2024}. Most existing approaches operate in isolation, focusing on either spatial or temporal information rather than effectively integrating both modalities. This work addresses this limitation by demonstrating how temporal signals provide structural features that complement spatial information, enabling robust material detection that extends beyond appearance-based classification and addresses the critical challenge of distinguishing real objects from their representations in safety-critical applications \cite{su2016,becker2023,yang2024}.

The neural network approaches reviewed, particularly convolutional neural networks, provide the computational framework for extracting material-discriminative features from temporal signals \cite{lecun1998}. The combination of SPAD temporal sensing with machine learning classifiers enables automated material detection that overcomes signal non-idealities while maintaining computational efficiency \cite{becker2023,su2016}.

In summary, this literature review establishes that: (1) existing material detection methods exhibit fundamental limitations in distinguishing material composition from visual appearance; (2) temporal signal analysis provides a viable alternative to spectral methods with reduced complexity; (3) SPAD sensors offer optimal characteristics for time-resolved material detection; and (4) the integration of spatial and temporal modalities represents an underexplored opportunity for enhanced scene understanding. The solutions presented in subsequent chapters address these gaps by demonstrating practical implementations of SPAD-based material detection that leverage temporal signatures for robust, efficient material characterisation across diverse applications.


\newpage
\chapter{Particle Wave Detection Framework}
\label{chap:pawd}


\newpage
\section{Introduction}
\label{sec:pawd_introduction}

The behaviour of electromagnetic energy underpins all optical sensing and signal-processing systems, yet its dual wave–particle nature makes it difficult to impose a single, unified modelling paradigm. Classical theories evolved from Newton’s corpuscular view of light to wave-based formulations and, eventually, to modern quantum electrodynamics, where energy exhibits both particle-like and wave-like characteristics depending on the measurement context. In practice, signal-processing pipelines instantiated in imaging, ranging, and spectroscopy instruments implicitly choose one side of this duality: they either operate in a spectral (wave-based) representation or in a particle-like, count-based representation. This chapter formalises a framework—Particle Wave Detection (PAWD)—that makes this choice explicit and uses it to organise signal-acquisition and processing strategies for material-aware computer vision.

PAWD is motivated by the observation that many modern systems, especially those used in robotics and machine perception, are hybrids. A conventional digital camera, for example, measures photon counts at each pixel (a particle-like process) but interprets them through colour filter arrays and demosaicing in a spectral (wave-based) space. Time-of-flight (ToF) imagers, single-photon avalanche diode (SPAD) dToF sensors, and LiDAR units similarly convert incident photons into digital counts, then aggregate or filter these measurements in either spectral or non-spectral domains, depending on the task \cite{su2016,becker2023}. Existing literature shows that combining temporal and spatial measurements can dramatically improve depth and object detection performance compared with intensity-only pipelines, particularly in cluttered or safety-critical scenes \cite{moramartin2021,lyons2019}. However, the underlying energy-model assumptions of these systems are rarely articulated, which obscures how and why certain processing regimes succeed.


The PAWD framework treats the wave–particle duality as a design axis rather than a purely philosophical construct. It interprets every sensing and processing chain as choosing one of four regimes defined by two orthogonal decisions: (1) whether measurements are resolved in intensity (spatial amplitude) or time (temporal response), and (2) whether processing is spectral (wavelength-resolved) or non-spectral (achromatic). These four regimes—intensity-spectral, intensity non-spectral, temporal-spectral, and temporal non-spectral—map directly onto common optical instruments and signal-processing pipelines used in computer vision and remote sensing. For example, RGB cameras and hyperspectral imagers inhabit the intensity-spectral regime, while many SPAD-based dToF systems used in this work operate in the temporal non-spectral regime, measuring only photon arrival times from a monochromatic source without resolving the underlying wavelength composition \cite{su2016}.

\begin{quote}
    \footnotesize
    \setstretch{2.0}%
    \ttfamily
    \color{textgrey}%
    \quotefive
    \end{quote}
    

\section{Historical background}
\label{sec:historical_background}

From a signal-processing perspective, each PAWD regime defines a feature space and an associated computational cost. Spectral methods operate on high-dimensional channels—wavelength bands or colour components—that must be separated, calibrated, and often unmixed, incurring storage and processing overhead that grows with spectral resolution \cite{su2016,becker2023}. In contrast, non-spectral methods collapse wavelength dependence and instead exploit temporal or spatial structure; for example, temporal histograms measured by SPAD dToF sensors encode range and subsurface scattering without any spectral decomposition. Empirically, many material-classification and depth-estimation tasks can be solved effectively using temporal non-spectral features, with lower complexity than hyperspectral pipelines and with hardware that is more compact and energy efficient.

PAWD also clarifies how "hybrid" systems should be constructed. Modern machine-vision pipelines frequently fuse RGB imagery with depth or transient measurements, but the fusion is often heuristic—concatenating channels or stacking feature maps—without a principled account of which modality contributes which information. Within PAWD, spatial intensity encodes geometric layout, texture, and appearance, while temporal non-spectral signals encode structural properties such as subsurface scattering, optical thickness, and absolute range. Works on high-speed single-photon ToF imaging and SPAD-based material classification demonstrate that temporal non-spectral cues can support robust object detection and material discrimination even when spatial appearance is ambiguous or adversarially spoofed \cite{moramartin2021,lyons2019,becker2023}. In this thesis, the PAWD framework guides the design of three solutions: long-range material detection, spatiotemporal object detection, and liquid purity assessment, all of which rely on temporal non-spectral processing as the primary structural signal channel.

Although the original motivation for PAWD is optical material detection with SPAD dToF sensors, the framework is agnostic to carrier and extends naturally to other sensing domains where energy–matter interactions are measured over time, frequency, or space. Radar, sonar, and terahertz systems, for example, can be classified into the same four regimes depending on whether they exploit spectral content, temporal impulse responses, or both. By expressing these systems in a common language, PAWD makes explicit the trade-offs between spectral and non-spectral processing, between spatial and temporal representations, and between model complexity and information content. The remainder of this chapter formalises the PAWD models, states the assumptions under which they hold, and links each regime to concrete signal-processing and computer-vision pipelines used in the subsequent chapters.

The behavioural nature of light has challenged scientists for centuries. Early optical experiments with prisms revealed that white sunlight decomposes into a continuous spectrum of colors, motivating corpuscular models in which light consisted of small particles while also hinting at underlying frequency-like structure \cite{crew1900}. Newton's corpuscular theory successfully explained rectilinear propagation and sharp shadows, but it struggled to account for interference and diffraction phenomena observed in more refined experiments \cite{crew1900}.

Wave-based explanations emerged to address these limitations. Huygens' principle formalized the idea that every point on a wavefront acts as a secondary source of spherical wavelets, providing a powerful framework for understanding reflection, refraction, and diffraction \cite{crew1900}. Young's double-slit experiment then offered decisive evidence for the wave nature of light by demonstrating stable interference fringes that could only be explained if light propagated as a coherent wave field \cite{young1804}. Fresnel extended this program by developing a quantitative wave theory of diffraction and polarization, replacing corpuscular intuition with a mathematically precise wave optics.

The unification of electricity and magnetism by Maxwell produced the decisive theoretical synthesis. Maxwell's equations predict self-propagating electromagnetic waves travelling at a speed numerically equal to the measured speed of light, implying that visible light is one band of the electromagnetic spectrum \cite{griffiths2017,schott2007}. This insight connected optical experiments with a broader field theory in which electric and magnetic fields are coupled and evolve jointly, establishing wave propagation, polarization, and radiation as consequences of a single electromagnetic framework \cite{griffiths2017,sadiku}.

At the turn of the twentieth century, quantum theory reintroduced particle-like behaviour in a complementary way. Einstein's interpretation of the photoelectric effect proposed that electromagnetic radiation is quantized into discrete energy packets (photons) with energy $E$ proportional to frequency $\nu$, reconciling black-body radiation and threshold effects in photoemission \cite{bohr1923,griffiths2017}. Formally,
\begin{equation}
E = h \nu,
\label{eq:photon_energy}
\end{equation}
where $h$ denotes Planck's constant. This relation underpins modern photon-counting detectors and links spectral content directly to quantized energy transfer.

The historical progression from corpuscles to waves and then to wave--particle duality suggests that neither description is sufficient in isolation: electromagnetic energy exhibits wave-like propagation in free space while interacting with matter through discrete quanta. Contemporary signal-processing systems exploit whichever representation best matches the sensing regime: spectroscopic instruments emphasise wave-based, frequency-resolved models, whereas SPAD-based direct time-of-flight sensors operate in a particle-like, photon-counting regime while still obeying the same electromagnetic constraints \cite{griffiths2017,elachi2021}. The PAWD framework in this chapter formalises that complementarity, treating wave and particle formalisms as design choices along a continuum rather than mutually exclusive theories.

\section{Practical motivation for the framework}
\label{sec:pawd_practical_motivation}

\noindent
The solutions developed in this thesis address three practical challenges in optical sensing and computer vision: long-range material detection, spatiotemporal object detection, and liquid purity assessment. Each application requires design choices about how radiation is measured and processed: whether signals are treated as spectrally resolved waves, as achromatic photon counts, or as hybrids of both. Without an explicit framework, these choices tend to be implicit and ad hoc, making it difficult to reason about what information is actually being exploited and why certain architectures generalise better than others. The Particle Wave Detection (PAWD) framework makes these assumptions explicit by organising sensing pipelines along two axes: intensity- versus time-resolved measurement, and spectral versus non-spectral processing. As summarised by the core PAWD models in Figure~\ref{fig:3_1}, every system in this work can be placed in one of four regimes, which then serves as a blueprint for feature extraction, learning, and decision making.

In this context, \emph{spectral} processing refers to any pipeline that explicitly resolves or models wavelength-dependent structure in the signal. Classical spectroscopy provides the canonical examples: instruments separate incident radiation into spectral components and then analyse absorption, emission, or scattering as a function of wavelength or wavenumber to infer material composition \cite{charles1924,holt2016,duckett2000}. These measurements may be spatially resolved (e.g., hyperspectral imaging) or non-imaging (e.g., fiber probes), and may operate in either intensity or time domains, but in all cases the output depends on how energy is distributed over a spectrum of frequencies. By contrast, \emph{non-spectral} processing treats signals as achromatic, operating only on aggregate quantities such as total intensity or photon-arrival statistics, regardless of how that energy is distributed over wavelength. A system that receives RGB or infrared data but collapses all channels into a single-band representation before analysis is therefore non-spectral from the PAWD perspective, even if the physical source spans many wavelengths.

\begin{figure}[H]
\centering
\includegraphics[width=0.7\textwidth]{figures/3.1.png}

    \caption[Core models of the Particle Wave Detection (PAWD) framework.]{\textbf{Core models of the Particle Wave Detection (PAWD) framework.} The framework partitions sensing and processing into four regimes: intensity-spectral, intensity non-spectral, temporal-spectral, and temporal non-spectral. Intensity-spectral systems (e.g., RGB and hyperspectral imagers) operate on spatially resolved, wavelength-dependent measurements. Intensity non-spectral systems use spatial intensity without explicit spectral separation (e.g., greyscale or infrared cameras). Temporal-spectral systems resolve both photon arrival time and wavelength, while temporal non-spectral systems, such as the SPAD dToF instruments in this work, operate on one-dimensional time-resolved histograms from a monochromatic source. The right-hand side highlights how spatial (geometric) and structural (material and depth) cues can be combined to form spatiotemporal decision pipelines.}
\label{fig:3_1}
\end{figure}


In this context, \emph{spectral} processing refers to any pipeline that explicitly resolves or models wavelength-dependent structure in the signal. Classical spectroscopy provides the canonical examples: instruments separate incident radiation into spectral components and then analyse absorption, emission, or scattering as a function of wavelength or wavenumber to infer material composition \cite{charles1924,holt2016,duckett2000}. These measurements may be spatially resolved (e.g., hyperspectral imaging) or non-imaging (e.g., fiber probes), and may operate in either intensity or time domains, but in all cases the output depends on how energy is distributed over a spectrum of frequencies. By contrast, \emph{non-spectral} processing treats signals as achromatic, operating only on aggregate quantities such as total intensity or photon-arrival statistics, regardless of how that energy is distributed over wavelength. A system that receives RGB or infrared data but collapses all channels into a single-band representation before analysis is therefore non-spectral from the PAWD perspective, even if the physical source spans many wavelengths.

Non-spectral approaches are attractive when the goal is to reduce complexity while preserving the structural cues most relevant to the task. In optical material sensing, it is common to assume that accurate characterisation requires spectral variables such as per-wavelength absorption coefficients or band-specific photon counts, so the problem is framed as inherently spectroscopic \cite{holt2016,baly1912}. However, both prior work and the experiments in this thesis show that time-resolved, achromatic measurements can encode rich reflectivity information (including surface scattering, subsurface transport, and geometric path length) without explicit spectral decomposition \cite{su2016,becker2023}. In PAWD terms, this places the proposed detectors and algorithms predominantly in the \emph{temporal non-spectral} regime: they operate on one-dimensional photon-arrival histograms from monochromatic or narrowband sources, yet still recover material-discriminative signatures comparable to more expensive hyperspectral pipelines.

The distinction between spectral and non-spectral processing is conceptually separate from the wave--particle duality reviewed earlier. Fundamental properties of visible light extend to the broader electromagnetic spectrum and other radiation sources, so energy--matter interactions in solids and fluids share the same underlying principles \cite{elachi2021,griffiths2017,charles1924}. Detectors designed around a predominantly wave-based view of radiation typically contain optics and electronics for splitting, grating, or otherwise resolving incident energy into spectral components, and then process those components to recover quantities such as phase, coherence, or spectral line structure. Examples include interferometric sensors, optical coherence tomography systems, and many forms of multi- or hyperspectral imagers \cite{huang1991,zeitler2007,vasile2023}. In PAWD, such systems occupy the \emph{spectral} quadrants—either intensity-spectral or temporal-spectral—because the downstream signal processing is explicitly conditioned on spectral composition.

Conversely, particle-focused detectors represent radiation as discrete photon events and are often engineered to be agnostic to wavelength, beyond broad-band filtering for eye safety or hardware convenience. Silicon photodiode pixels measure accumulated charge as intensity images, while single-photon avalanche diodes (SPADs) measure photon arrivals in time, frequently with per-pixel time-to-digital converters and histogram builders \cite{piron2020,bass2010}. When these devices are operated with monochromatic or narrowband illumination and no subsequent spectral unmixing, their outputs (greyscale images or one-dimensional temporal histograms) are inherently achromatic. High-gain SPAD arrays excel in this regime, delivering picosecond-scale timing and single-photon sensitivity that make them well-suited for direct time-of-flight (dToF) sensing and for capturing the transient signatures exploited in this work \cite{alabbas2018}. There are, however, important exceptions where SPAD architectures are arranged into spatially resolved imagers or combined with additional modalities, enabling high-speed object detection and tracking while still operating on non-spectral temporal data \cite{moramartin2021,lyons2019}.

From a systems perspective, the practical motivation for PAWD is to provide a common language for these design choices. Intensity and temporal measurements can each be spatially organised to form laterally resolved scene representations; they can also be fused with spectral channels or with learned feature maps from deep networks. Hybrid systems that combine RGB cameras with dToF SPAD sensors illustrate this point: spatial imagers deliver appearance and geometry, while temporal histograms deliver depth and material cues, and both can be integrated within a unified detection or classification pipeline \cite{su2016,becker2023,yang2023}. PAWD clarifies which parts of such a pipeline are fundamentally spectral versus non-spectral, and whether they operate in intensity or temporal domains. This, in turn, guides the architecture of the proposed solutions in Chapters~4--6, where non-spectral temporal features are deliberately favoured to reduce computational complexity while retaining the information necessary for robust material-aware perception.

\section{Optical properties of light}
\label{sec:optical_properties_of_light}

\noindent
Most objects encountered in everyday scenes are non-luminous; they become visible only because incident radiation is reflected from their surfaces toward the sensor. The resulting optical appearance is governed by how energy is partitioned between reflection, transmission, and absorption at the interface, so surface reflectance profiles provide the primary cue for localising targets and inferring their material composition \cite{elachi2021,schott2007}. In practice, these profiles are strongly shaped by surface topography and curvature (Figure~\ref{fig:3_2}).

Smooth, polished, or engineered coatings tend to return light in narrow, mirror-like lobes aligned with the law of reflection (specular reflection), whereas rough, textured, or unpolished (matte) materials scatter light over a wide range of angles (diffuse reflection). Many real-world surfaces combine both behaviours, or incorporate retro-reflective structures that direct energy back toward the source.

Surface texture can be described in terms of peaks, valleys, and waviness that deviate from an ideal planar interface. Crucially, the perceived roughness is wavelength dependent: when the characteristic height of surface irregularities is small relative to the incident wavelength, the surface appears optically smooth and reflections become more directional; when those features are comparable to or larger than the wavelength, reflections become increasingly diffuse. Thus, the same physical surface can manifest specular behaviour at longer wavelengths and diffuse behaviour at shorter wavelengths, a property routinely exploited in remote sensing and material analysis \cite{elachi2021,schott2007}. In this thesis, material reflectance is treated as a signature that remains stable once wavelength, illumination geometry, and viewing configuration are controlled, enabling remote or direct electro-optical sensing of material class from measured reflectance.

\begin{figure}[H]
\centering
\includegraphics[width=0.8\textwidth]{figures/3.2.png}

    \caption[Effect of surface texture and curvature on specular and diffuse reflection.]{\textbf{Effect of surface texture and curvature on specular and diffuse reflection.} Smooth, planar regions concentrate reflected energy into narrow specular lobes, while rough or curved regions redistribute energy over broader angular ranges, increasing diffuse components and reducing view dependence.}
\label{fig:3_2}
\end{figure}



\subsection{Smooth surface reflection (specular)}
\label{subsec:smooth_specular_reflection}

The following subsections briefly distinguish smooth (specular), rough (diffuse), and subsurface scattering regimes, which together define the optical context for the material detection systems developed later in this work.

Specular reflection arises when radiation encounters a surface whose microstructure is smooth relative to the incident wavelength, so that local surface normals are effectively parallel across the illuminated region. In this regime, scattering is minimal and most energy is redirected into a narrow lobe that obeys the law of reflection, producing the familiar mirror-like highlights observed on polished metals, glass, or well-finished dielectrics (Figure~\ref{fig:3_2}) \cite{elachi2021,schott2007}. The directional and concentrated nature of this returned signal is visually apparent in both spatial intensity images and time-resolved measurements from single-photon time-of-flight (ToF) sensors \cite{su2016}.%

In a simple intensity-based 2D pixel array (for example, a $4\times4$ grid), specular reflection concentrates photon flux onto pixels whose lines of sight align most closely with the specular direction, yielding bright central pixels surrounded by progressively darker neighbours as viewing angle departs from the specular axis. The greyscale pattern therefore encodes the angular concentration of the reflected lobe, with the brightest (often saturated) pixels at the centre and weaker intensities distributed across adjacent pixels that capture off-axis or partially scattered photons (left of Figure~\ref{fig:3_3}). An equivalent behaviour appears in ToF histograms: pixels near the specular direction exhibit narrow, high-amplitude peaks corresponding to strong, early photon arrivals, whereas surrounding pixels record broader, lower-amplitude distributions with elongated tails due to reduced photon counts and increased contributions from multiply scattered paths (right of Figure~\ref{fig:3_3}) \cite{su2016,tafone2023}. These temporal signatures provide a compact descriptor of specular behaviour that can be exploited for material discrimination and for separating specular highlights from diffuse background in downstream processing \cite{su2016,becker2023}.%

\begin{figure}[H]
\centering
\includegraphics[width=\textwidth]{figures/3.3.png}

    \caption[Specular reflection in intensity and time-of-flight domains.]{\textbf{Specular reflection in intensity and time-of-flight domains.} Left: a 4 × 4 intensity image where a smooth, planar specular surface concentrates returned photons into central pixels. Right: corresponding time-of-flight histograms, with central pixels exhibiting narrow, high-amplitude peaks and surrounding pixels showing broader, lower-amplitude responses with elongated tails.}
\label{fig:3_3}
\end{figure}

Smooth curved surfaces follow the same physical principles but modify the mapping between incident and reflected directions. Convex geometries tend to spread the specular lobe over a wider range of viewing angles, while concave or free-form geometries can focus or redirect energy away from the sensor, sometimes producing multiple distinct specular regions across the field of view (Figure~\ref{fig:3_4}). Real objects typically comprise assemblies of such curved patches, so their observed specular structure depends jointly on local curvature, illumination geometry, and sensor pose. For the purposes of this thesis, smooth specular reflection is treated as the limiting case of surface behaviour: it defines the most view-dependent component of the bidirectional reflectance distribution function (BRDF), against which rough (diffuse) reflection and subsurface scattering are contrasted in the following subsections.%

\begin{figure}[H]
\centering
\includegraphics[width=\textwidth]{figures/3.4.png}

    \caption[Effect of smooth curvature on specular reflection.]{\textbf{Effect of smooth curvature on specular reflection.} Convex, concave, and free-form smooth surfaces redistribute the specular lobe over different viewing directions, altering which pixels receive the strongest returns while preserving the underlying mirror-like behaviour.}
\label{fig:3_4}
\end{figure}



\subsection{Rough surface reflection (diffuse)}
\label{subsec:rough_diffuse_reflection}

Most real-world objects present hybrid rough surfaces whose microstructure departs from an ideal plane through peaks, valleys, and mesoscale waviness. When these surface irregularities are comparable to or larger than the illumination wavelength, incident photons are scattered over a broad range of directions rather than being confined to a narrow specular lobe \cite{elachi2021,schott2007}. The resulting diffuse reflection reduces directional intensity and makes appearance less view dependent, but it preserves material-dependent structure in both spatial intensity images and time-resolved measurements.%

In an intensity image acquired with a small array (for example, a $4\times4$ pixel patch observing a rough surface), scattered returns are distributed more uniformly across central and surrounding pixels: each pixel receives a similar, typically low, photon count because the reflected energy is spread across the detector's field of view. The left panel of Figure~\ref{fig:3_5} illustrates this behaviour, where diffuse reflection produces a flattened intensity pattern with muted contrast between centre and periphery. An analogous effect appears in direct time-of-flight histograms (right of Figure~\ref{fig:3_5}): photon arrivals are temporally dispersed by multiple scattering events, yielding broader, lower-amplitude distributions with elongated tails compared to the sharp peaks associated with smooth specular surfaces. Across the pixel array, this produces a family of histograms with similar low peaks and wide spreads, encoding both surface roughness and subsurface structure in the temporal domain \cite{su2016}.%
\begin{figure}[H]
\centering
\includegraphics[width=\textwidth]{figures/3.5.png}

    \caption[Diffuse reflection in intensity and time-of-flight domains.]{\textbf{Diffuse reflection in intensity and time-of-flight domains.} Left: a 4 × 4 intensity image of a rough, diffusely reflecting surface, where scattered returns distribute photons relatively evenly across central and surrounding pixels. Right: corresponding time-of-flight histograms, with all pixels exhibiting low, broadened peaks and extended tails due to multiple scattering and delayed photon arrivals.}
\label{fig:3_5}
\end{figure}

\subsection{Subsurface scattering (SSS)}
\label{subsec:subsurface_scattering}

Reflectance and material composition are jointly governed by both surface texture and intrinsic optical properties such as gloss, translucency, and absorption coefficients. Two nominally smooth materials (for example, polished nylon and a metallic mirror) can therefore exhibit distinct combinations of specular and diffuse components, producing unique reflectance signatures once illumination wavelength, geometry, and viewing configuration are fixed \cite{elachi2021,schott2007}. In spectral systems, these signatures are further decomposed across wavelength bands, increasing descriptive power but also expanding data dimensionality and computational cost. By contrast, achromatic time-of-flight sensing operates on single-band temporal histograms, leveraging the physics of energy–matter interaction to encode material-dependent scattering and absorption into the shape of the return without explicit spectral analysis \cite{su2016}. Empirical studies show that these achromatic temporal signatures are often sufficient for robust material discrimination, making non-spectral processing an efficient alternative to high-dimensional spectral pipelines in many vision tasks.%

\noindent
Subsurface scattering (SSS) occurs when incident photons penetrate beneath a material's surface, undergo multiple scattering and absorption events within the volume, and re-emerge at displaced positions and delayed times. In translucent and semi-transparent media this process redistributes radiant energy over a broader angular range and temporal window than pure surface reflection, producing characteristic low-amplitude, long-tailed return profiles in direct time-of-flight (dToF) histograms \cite{su2016,lyons2019}. For SPAD-based dToF sensors, these temporal signatures are particularly informative: the relative balance between prompt, surface-reflected photons and delayed, subsurface-scattered photons encodes material-dependent optical properties such as scattering mean free path and absorption length, enabling discrimination between materials that share similar macroscopic appearance \cite{becker2023,tafone2023} (Figure~\ref{fig:3_6}).
\begin{figure}[H]
\centering
\includegraphics[width=\textwidth]{figures/3.6.png}

    \caption[Subsurface scattering and its effect on time-of-flight histograms.]{\textbf{Subsurface scattering and its effect on time-of-flight histograms.} Incident photons (red) penetrate beneath a translucent surface, undergo multiple internal scattering events, and re-emerge at displaced locations and later times than photons reflected at the surface. The resulting single-photon time-of-flight histogram exhibits a broadened peak with an elongated tail compared to that of an opaque, purely surface-scattering material.}
\label{fig:3_6}
\end{figure}

\section{Material detection}
\label{sec:material_detection_overview}

Choice of illumination wavelength strongly influences the strength and depth of SSS. Near-infrared (NIR) sources, as used in this work, generally penetrate deeper into many organic and polymeric materials than visible (VIS) wavelengths because scattering coefficients decrease and absorption bands shift with increasing wavelength \cite{inagaki2020,aernouts2015}. NIR illumination therefore increases the contribution of subsurface paths to the detected signal while also reducing sensitivity to surface colour patterns, making the measured temporal response more directly related to bulk optical properties, a property exploited in diffuse optical tomography and medical imaging systems \cite{lyons2019,huang1991}.

Within the proposed SPAD material-detection pipeline, SSS manifests as a systematic elongation of the photon-arrival histogram around the main depth peak: translucent materials exhibit broader peaks and more pronounced exponential tails than opaque references observed at the same geometric range (see also Figure~\ref{fig:1_2}). This elongation, though subtle, is highly repeatable across measurements and distances, and prior work has shown that simple features derived from peak width and tail decay constants are sufficient for robust plastic and surface material classification using low-cost dToF sensors \cite{becker2023,tafone2023}. In this thesis, those temporal features constitute a primary cue for distinguishing material classes under fixed geometry and illumination, with subsurface scattering acting as the mechanism that converts microstructural differences into measurable histogram shape variations \cite{su2016,becker2023}.

Material detection represents the primary application domain guided by the optical sensing framework presented in this work. This section elaborates on material detection as the fundamental application, focusing on the investigation and characterisation of signature optical properties exhibited by materials under radiation exposure. These properties are unique per material due to inherent variations in radiation-matter interactions, particularly in absorption characteristics that manifest as distinct temporal and spectral signatures \cite{elachi2021,su2016}.

Prior to the advent of impact ionisation detectors that enable photon counting and precise temporal resolution, most optical sensing applications relied predominantly on spectral analyses operating in frequency and amplitude domains. This spectral approach remains the dominant paradigm in contemporary systems, where material characterisation typically exploits wavelength-specific variables such as absorption coefficients per wavelength or frequency to analyse material absorbance and infer composition \cite{holt2016,elachi2021}. However, one-dimensional time-resolved signals encode reflectivity information that equally infers material absorbance without requiring separation of constituent wavelength components within the signal (see Section~\ref{sec:optical_properties_of_light} and Figure~\ref{fig:1_2}) \cite{su2016,becker2023}.

Material optical responses exhibit distinct characteristics based on their interaction with incident radiation. Opaque materials demonstrate high absorption or reflectivity depending on surface texture, with matte surfaces exhibiting diffuse reflection and glossy surfaces producing specular reflection \cite{elachi2021}. Transparent materials exhibit minimal absorption, allowing most incident radiation to transmit through the material with minimal interaction. Translucent materials display partial absorption combined with partial scattering and reflectivity, resulting in complex interaction patterns that encode material-specific information in the temporal response \cite{su2016,becker2023,tafone2023}.

The SPAD detector employed throughout this project utilises a monochromatic infrared source operating at ~940nm wavelength. Despite this monochromatic illumination, signal processing adheres to achromatic principles, meaning the constituent wavelengths present in the signal are irrelevant to the processing methodology \cite{su2016}. This achromatic approach eliminates the computational overhead associated with spectral decomposition while maintaining material discrimination capability through temporal signature analysis \cite{su2016,becker2023,tafone2023}.

The signal input requirements vary across objectives. For Objective 2, a combined image object fusing multispectral RGB spatial information with one-dimensional time-resolved temporal signals creates a spatiotemporal representation for enhanced scene understanding \cite{yang2024,moramartin2021}. The time-resolved signal component captures reflectivity data encoding material composition and subsurface scattering characteristics, while the two-dimensional intensity image provides target localisation in spatial coordinates \cite{su2016}. This fusion approach addresses material composition analysis objectives that are commonly approached using spectrometric equipment, yet achieves comparable performance through non-spectrometric means that reduce system complexity and computational requirements \cite{su2016,becker2023,tafone2023}.

\section{Signal measurement approaches}
\label{sec:signal_measurement_approaches}

Material detection and sensing systems measure electromagnetic energy through two primary modalities: intensity-resolved measurements, which assess the overall power or amplitude of signals, and time-resolved measurements, which record the arrival times of composite photons \cite{su2016}. These measurement approaches differ fundamentally in the information they encode and the applications they enable.

Intensity-resolved measurements are characteristic of spatial scene resolution, where point-wise signal intensities over two-dimensional space enable scene interpretation and object identification. Conventional photographic imaging exemplifies this approach: each pixel records the integrated intensity of incident radiation, producing a two-dimensional intensity map that captures spatial patterns, textures, and appearance features. These intensity distributions encode geometric layout and surface properties, making them well-suited for object localisation and appearance-based classification tasks \cite{schwartz2013,sharan2013}.

Time-resolved measurements, by contrast, capture temporal signatures that encode depth and structural features beyond what spatial intensity alone provides \cite{su2016}. Direct time-of-flight (dToF) systems, such as those employing single-photon avalanche diode (SPAD) sensors, record photon arrival times relative to pulsed illumination, producing temporal histograms that encode both target distance and material-dependent optical properties \cite{becker2023}. These temporal signatures capture subsurface scattering, optical thickness, and absorption characteristics that manifest as distinct histogram shapes, enabling material discrimination even when spatial appearance is ambiguous \cite{su2016,becker2023}.

Both intensity and time-resolved approaches may operate with or without explicit spectral decomposition. Spectral systems resolve wavelength-dependent structure, enabling material characterisation through absorption and emission signatures across multiple frequency bands \cite{elachi2021,holt2016}. Non-spectral (achromatic) systems, such as the SPAD dToF instruments employed in this work, operate on single-band temporal histograms from monochromatic sources, eliminating the computational overhead of spectral decomposition while maintaining material discrimination capability through temporal signature analysis \cite{su2016,becker2023}.

The complementary nature of spatial intensity and temporal signals enables hybrid spatiotemporal fusion, as illustrated in the right-hand side of Figure~\ref{fig:3_1}. Spatial intensity encodes geometric layout, texture, and appearance, while temporal non-spectral signals encode structural properties such as subsurface scattering, optical thickness, and absolute range \cite{moramartin2021}. This fusion approach addresses material composition analysis objectives that are commonly approached using spectrometric equipment, yet achieves comparable performance through non-spectrometric means that reduce system complexity and computational requirements \cite{su2016,becker2023}. Modern machine-vision pipelines that fuse RGB imagery with co-registered dToF transients demonstrate improved robustness where spatial appearance is ambiguous or adversarially spoofed \cite{moramartin2021,yang2024}. Detailed descriptions of the processing models and computational characteristics for each measurement regime are provided in subsequent sections of this chapter.

\subsection{Achromatic vs mono-spectral}
\label{subsec:achromatic_vs_monospectral}

While achromatic and mono-spectral signals both operate on single-channel data, they represent fundamentally distinct processing paradigms that differ in their dependence on wavelength information \cite{elachi2021,schott2007}. This distinction becomes evident when examining signal processing requirements from the perspective of each approach.

Consider a multi-wavelength scene captured by an imaging system. When processed through a mono-spectral pipeline, the signal must first be filtered or decomposed to isolate a specific wavelength of interest (e.g., red) before further processing can proceed \cite{elachi2021}. The successful extraction and processing of the red channel image fundamentally depends on correctly identifying and isolating the target wavelength, making wavelength estimation and spectral filtering essential components of the signal processing chain \cite{schott2007}. In contrast, when the same multi-wavelength signal is processed through an achromatic (non-spectral) pipeline, the output is inherently greyscale, regardless of the wavelength composition of the incident radiation \cite{su2016}. Even when wavelength estimation capabilities are available, this information plays no meaningful role in determining the final signal output, as achromatic processing operates independently of spectral composition \cite{su2016,becker2023}.

Both mono-spectral and achromatic systems produce single-channel outputs, yet the processing techniques required differ substantially based on the role wavelength plays in the signal processing task \cite{elachi2021}. A system is considered mono-spectral when successful signal processing requires identification and selection of a specific frequency or wavelength \cite{schott2007}. Conversely, a system is classified as achromatic or non-spectral when signal processing succeeds without requiring wavelength estimation or frequency awareness, operating instead on integrated temporal or intensity information that aggregates across the spectral content of the source \cite{su2016,becker2023}.

\section{Assumptions}
\label{sec:pawd_assumptions}

The Particle Wave Detection (PAWD) framework rests on fundamental assumptions about signal properties and detection mechanisms that govern how electromagnetic energy is measured and processed. These assumptions formalize the design choices underlying all sensing pipelines developed in this work.

\begin{itemize}
    \item Signals have particle and/or wave properties. Thus, signal detectors use either of these behaviours as detection mechanism.
    
    \item Detection without spectral awareness or relevance are particle based and termed non-spectral while detection with spectral awareness or relevance are wave based and termed spectroscopy.
    
    \item Detectors may resolve signals by arrival times as transient along axial axes. This is often done with high gain sensors (e.g., SPADs). However, a coarse lateral resolution is possible.
    
    \item Detectors may resolve signals by measuring signal intensity per pixel. Often laterally resolved with low gain sensors like the active pixel sensor (APS).
    
    \item During signal detection, sensor may receive signal from pure material surfaces (homogeneous surface texture often composing only one element) or hybrid materials (composing multiple materials).
\end{itemize}

The following sections demonstrate possible combinations of non-spectral and spectroscopic signal detections. For each category, spatial intensity as well as corresponding temporal or transient resolutions are demonstrated.

As illustrated in the right-hand side of Figure~\ref{fig:3_1}, spatial intensity measurements encode geometric layout, texture, and appearance features that enable object localisation and appearance-based classification \cite{schwartz2013,sharan2013}. Temporal non-spectral signals, by contrast, encode structural properties such as subsurface scattering, optical thickness, and absolute range that provide material-discriminative information beyond spatial appearance alone \cite{su2016,becker2023}. The fusion of these complementary modalities enables hybrid spatiotemporal decision pipelines where spatial cues identify what objects are present while temporal cues reveal what materials they are composed of, addressing fundamental limitations in appearance-only vision systems \cite{yang2024,moramartin2021}.

\section{Non-spectral (achromatic) signal measurements}
\label{sec:non_spectral_achromatic_measurements}

Non-spectral signal measurements employ detectors whose operation does not depend on resolving the spectral composition of incident radiation. Even when wavelength-dependent information is, in principle, accessible, it is deliberately discarded in favour of an achromatic processing pipeline. At the pixel level, non-spectral measurements are typically represented either as single-band intensity values (for example, greyscale or narrowband infrared images) or as time-resolved photon arrival distributions that encode temporal or transient responses.

These achromatic spatial and temporal channels can be used independently or fused to form richer sensing modalities. In practice, combining per-pixel intensity with time-of-flight histograms within SPAD-based imagers and related architectures enables high-speed object detection, diffuse optical tomography, and other three-dimensional sensing tasks without explicit spectral decomposition \cite{moramartin2021,lyons2019}. For both spatial and temporal resolutions, signal processing may operate on individual pixels or across full arrays, allowing algorithms to exploit local structure and global context as required by the application.


\subsection{Non-spectral intensity signal}
\label{subsec:non_spectral_intensity_signal}

When measurable signals from targets arrive at detectors as particles, pixels accumulate electrons via internal or external photoelectric effects \cite{bass2010}. Unlike spectral detectors, non-spectral intensity pixels do not filter or decipher the wave composition of the signal. Instead, collected electrons are decoded to digital values representing measured charge. Intensity signal detectors consist of circuitry that convert collected electrons in pixel wells to their equivalent voltage (scalar) values (Figure~\ref{fig:3_7}).
\begin{figure}[H]
\centering
\includegraphics[width=\textwidth]{figures/3.7.png}

    \caption[Intensity signal detector circuitry converting collected electrons to voltage values.]{\textbf{Intensity signal detector circuitry converting collected electrons to voltage values.} The pixel well accumulates electrons through photoelectric effects, and readout circuitry converts the collected charge to equivalent voltage values for digital representation.}
\label{fig:3_7}
\end{figure}

The voltage values produced by the detector circuitry are represented as pixel intensities, where each pixel's accumulated charge corresponds to a digital intensity value. These values are represented as pixel intensities (Figure~\ref{fig:3_8}). Pixels organised in two-dimensional architectures are spatially resolved to form two-dimensional images (Figure~\ref{fig:3_8}).

In specialised modalities such as tomography, greyscale images gain axial resolution from multiple slices captured along the depth dimension. Two-dimensional laterally resolved greyscale images are known to be relevant for key machine vision and decision tasks, including spatial semantic segmentation, object detection, and scene understanding \cite{schwartz2013,sharan2013,chen2024}. A single achromatic intensity pixel, however, is rarely employed for signal perception decision tasks. A distinguishing attribute of depth estimation in laterally resolved achromatic images is the singular dimensionality of the wavelength spectrum. Spatial and structural organisation of light points in space are captured and represented as intensity measurements. 

\begin{figure}[H]
    \centering
    \includegraphics[width=0.7\textwidth]{figures/3.8.png}

    \caption[Spatially resolved two-dimensional intensity image formation from pixel arrays.]{\textbf{Spatially resolved two-dimensional intensity image formation from pixel arrays.} Individual pixel intensity values are organised in a two-dimensional architecture to form greyscale images that capture spatial patterns and geometric layout.}
    \label{fig:3_8}
    \end{figure}

Various decision tasks can be performed with two-dimensional achromatic images: depth estimation (stereoscopy), image segmentation, scale estimation, classification, and object detection. The following subsections provide brief descriptions of each task.

\subsubsection{Depth estimation}
\label{subsubsec:depth_estimation_intensity}

Achromatic intensity pixels often rely on two-dimensional lateral resolution for monocular or binocular depth estimation. Stereoscopic systems estimate pixel-wise depth and spatial coordinates to generate point clouds \cite{laga2022,hu2020}. Single-pixel depth estimates may be more useful for target depth approximations rather than precise whole-scene representation. While it is possible to average multiple pixels for depth estimation, there is limited evidence in the literature demonstrating its role in decision making \cite{laga2022}. Although individual pixels may be insufficient for whole-scene representation, the combined network of pixels' intensity and spatial coordinates may offer near-accurate three-dimensional representations of entire scenes \cite{godard2019,laga2022}.

\subsubsection{Segmentation}
\label{subsubsec:segmentation_intensity}

Image segmentation is a fundamental task in signal processing \cite{haralick1984,blaschke2010}. In two-dimensional achromatic images, segmentation at the pixel level is rarely meaningful as individual pixels contain unit intensity values. However, pixels in different parts of the image can be annotated to represent the spatial organisation of objects within the scene. The purpose of segmentation is to localise relevant parts of the image for interpreting the target scene. Accurately annotated two-dimensional images are useful for higher-level scene analysis tasks such as classification or detection \cite{schwartz2013,sharan2013}. With a target region of interest, relative segments can be localised only towards horizontal or vertical planes (lateral resolution). Although pixel intensities of some regions may be observed as positioned in front of or behind the target, such estimations are relative and can rarely be specified in true distance.

\subsubsection{Scale estimation}
\label{subsubsec:scale_estimation_intensity}

Scale estimation is a challenging problem in computer vision \cite{lindeberg1998,lowe2004}. The definition of target scale is expressed in terms of ratios of signal spatial distribution to target depth. For two-dimensional intensity images, the task is achievable with stereoscopy (binocular or monocular) \cite{laga2022}. The limitation of scale estimation with two-dimensional images centres around the high complexity of stereoscopy setup and the quality of triangulation, which has a direct impact on the precision of estimated distance \cite{laga2022,hu2020}. This makes the problem even harder for distant and poorly illuminated objects. A regular active pixel sensor (APS) may capture an image in darkness, but the image data is rarely usable for scene visualisation, segmentation, depth estimation, or scale estimation. Thus, this approach is generally considered a less reliable means of estimating target scale.

\subsubsection{Classification}
\label{subsubsec:classification_intensity}

Classification is a higher-level task partly enabled by segmentation \cite{schwartz2013,sharan2013}. From single pixels through annotated regions of interest (ROIs) or even whole images, labels can be assigned for classification and detection (Figure~\ref{fig:3_9}).
\begin{figure}[H]
\centering
\includegraphics[width=\textwidth]{figures/3.9.png}

    \caption[Achromatic image segmentation with bounding box and deformable outline.]{\textbf{Achromatic image segmentation with bounding box and deformable outline.} left: achromatic image segmentation with bounding box(green) and assigned label(blue) Right: image segmentation with deformable outline and assigned label. Segmented and labelled images may be used as inputs for object detection models.}
\label{fig:3_9}
\end{figure}

Whole-image classes are useful for generic scene description. Depending on segmentation quality, ROI delineations may be better for lower-level scene localisations or other complex vision tasks.

\subsubsection{Object detection}
\label{subsubsec:object_detection_intensity}

The definition of an object in signal processing literature is subjective \cite{redmon2015,redmon2018}. However, objects may represent an organised collection of features that communicate meaningful or tangible information. The ability to reliably localise or annotate an object region associated with a label constitutes object detection. Detection in two-dimensional greyscale images relies on feature variables such as annotation shape and variations in intensity values of pixels within the region \cite{redmon2015,redmon2018,liu2016}. Similar to segmentation, object detection is rarely possible from a single pixel with two-dimensional intensity images.

\subsubsection{Material detection}
\label{subsubsec:material_detection_intensity}

Material detection tasks attempt to identify annotations or targets by a label that defines their material composition. This is based on the assumption that reflected signals encode unique material properties of their surface of emergence. The Hounsfield scale used in the medical domain is a classic example of methods for material detection with greyscale images \cite{huang1991}. However, this approach is quite limited to body tissues and difficult to generalise for everyday objects. While reflectivity information can be encoded in pixel intensity, it is even more evident in the timed sequence of radiation particles (as seen in time-of-flight measurements) \cite{su2016,becker2023}. Hence, the concept of material detection in achromatic intensity images is considered less feasible compared to the corresponding task in transient images \cite{su2016,becker2023,tafone2023}.


\subsection{Processing models of non-spectral intensity signal}
\label{subsec:processing_models_spectral_intensity}

Signal intensity in non-spectral intensity measurements is defined as the estimated digital value of target reflectance within each pixel. Under active illumination, unabsorbed photons incident on a target surface are either reflected back toward the detector or transmitted through the material. A detectable fraction of reflected photons travel back into the detector's field of view without interacting with other scene objects. However, the intensity of the returned signal is subject to the inverse square law \cite{griffiths2017}, which states that electromagnetic radiation intensity decreases proportionally to the square of the distance from the source. This means that for a round-trip path (source to target and back to detector), the received signal intensity scales inversely with the square of the target distance, creating a fundamental relationship between distance and measured intensity that must be accounted for in signal processing models. The total number of photons contained within each pixel per returned signal can be described using a linear model \cite{heinz2001,schott2007}, though this model implicitly incorporates the distance-dependent intensity scaling.

In achromatic signal processing, there is no spectral awareness; thus, the signal intensity of a pixel from a hybrid material surface is represented by the variable $U$. For a target material composed of pure elements, the number of composite materials $M = 1$. In this case, pixel intensity is estimated from the reflectance of a single material $m$. The reflectance of the signal from material $m$ is given by $r_m$, and the fractional proportion of $m$ out of all $M$ materials is represented by $p_m$. The error term capturing variation between the actual observed surface reflectance and the computed pixel reflectance value from the model is expressed as $\epsilon$. Factors influencing $\epsilon$ include ambient noise, electronic noise, photon shot noise, and other measurement uncertainties \cite{schott2007}.

For a pure material case where $M = 1$, the pixel intensity value is represented by a linear model:

\begin{equation}
U = p_m \cdot r_m + \epsilon
\label{eq:3_1}
\end{equation}
\addcontentsline{loe}{equations}{\protect\numberline{\theequation}$U = p_m \cdot r_m + \epsilon$}

where $p_m$ represents 100\% of $M$ since there is only one material within the illumination area. The resulting intensity value of the pixel after incident signal interaction is the incident signal plus possible preexisting noise in the pixel. This ideal case, however, is rarely true for most real-world target surfaces.

Practically, target surfaces are often composed of hybrid materials that vary in reflectivity. Given an illuminated hybrid surface containing $M$ detectable materials, the return signal is expected to contain photons from each exposed material in $M$. Thus, the resulting signal in each pixel is a linear proportional sum of reflectances from each detectable material $m$. If a material is completely occluded by another opaque material within the same object, such material is considered undetectable, as the reflectance will only emerge from the occluding opaque material to the detector \cite{heinz2001,schott2007}.

Considering an impure hybrid surface of $M$ detectable materials and noise factors (electronic or ambient), a mixture model is considered:

\begin{equation}
U = \sum_{m=1}^{M} p_m \cdot r_m + \epsilon
\label{eq:3_2}
\end{equation}
\addcontentsline{loe}{equations}{\protect\numberline{\theequation}$U = \sum_{m=1}^{M} p_m \cdot r_m + \epsilon$}

where $p_m$ represents the fractional proportion of each material $m$, constrained such that $\sum_{m=1}^{M} p_m = 1$. The pure fractional proportion per composite material can be solved as a system of linear equations \cite{heinz2001}. In a signal detection experiment, multiple observations of an impure surface containing two materials can be made. For each observation, two material reflectance values are recorded ($r_{m_1}$ and $r_{m_2}$) and assigned to the total reflectance $r_M$ of the signal. The fractional composition of $r_{m_1}$ and $r_{m_2}$ can be solved using a system of linear equations representing multiple observations of the hybrid surface \cite{heinz2001,schott2007}.

When fed to a regression model, the probability values of $r_{m_1}$ and $r_{m_2}$ should be output with high reliability. In the case of intensity images, probability values per material can be scaled between the lowest and highest pixel brightness values (0--255), resulting in two output images, each representing different surface fractions. This method is applicable to hybrid surface detection of exposed composite materials, enabling material discrimination through linear unmixing of reflectance contributions \cite{heinz2001,schott2007,elachi2021}.


\subsection{Processing complexity of non-spectral intensity signal}
\label{subsec:processing_complexity_spectral_intensity}

This section presents a machine-independent proof of computational efficiency for processing non-spectral intensity signals, which are commonly employed to represent laterally resolved scenes in computer vision applications \cite{schwartz2013,sharan2013}. Consider a detector array comprising $N$ pixels, where each pixel $i$ has an initial intensity state $I_{\text{pix},i}$ and receives a returned signal intensity $I_{\text{sig},i}$. In the absence of input analog signals, pixel intensities are represented by an empty array $\mathbf{D}$ of $N$ elements, where each element contains either zero or a negligible noise component $c$. Digital signal input per pixel is modelled as an additive operation, where each pixel's intensity is updated by the incident signal value.

In the optimal pixel update scenario, a single pixel $I_{\text{pix},0}$ with noise intensity $c$ receives a returned signal $I_{\text{sig},0}$, resulting in a straightforward accumulation operation. The computational complexity of signal input operations, whether processing hybrid or pure surface materials, grows linearly with the size of the pixel array \cite{knuth1976}. In the worst-case scenario, signal accumulation across the entire array occurs once per frame, resulting in a space-time complexity of $\mathcal{O}(N)$ \cite{knuth1976}. This linear scaling property demonstrates the fundamental efficiency advantage of non-spectral intensity signal processing compared to spectral methods, which typically exhibit higher-dimensional complexity due to wavelength-dependent processing requirements \cite{heinz2001,schott2007}.


\subsection{Non-spectral temporal signal}
\label{subsec:non-spectral_temporal_signal}

Non-spectral temporal signal processing represents an alternative measurement approach that captures photon arrival times rather than intensity distributions. Similar to achromatic intensity imaging, temporal processing operates without explicit wavelength or frequency decomposition, making variations in the incident signal's spectral composition irrelevant to the signal processing task. While all time-of-flight detectors can theoretically be used for transient measurements, SPAD-based direct time-of-flight (dToF) sensors are more commonly employed for temporal signal acquisition. Temporal data is stored as time-stamped sequences of photon arrivals, which can be binned to form histograms for analysis. Encoded within these photon arrival signals is signature information that can reliably identify materials by composition through their unique temporal reflectivity profiles (see Sections~\ref{subsec:smooth_specular_reflection},~\ref{subsec:rough_diffuse_reflection}, and~\ref{subsec:subsurface_scattering}).

Detection systems may correctly identify object categories, such as distinguishing a person from other objects. However, algorithms often struggle to determine lower-level material details, such as whether an identified person object is actually a human or a human-shaped bronze statue. These scene perception tasks represent inverse problems in optics, yet they can be addressed to a high degree using temporal signal analysis \cite{pizlo2001}. While transient data improves scene perception quality through depth estimation, the resolution is limited to the axial plane. Consequently, temporal signals should be considered an important complement to lateral scene resolution rather than a replacement for spatial imaging in machine vision systems \cite{chen2024}.

In related research by Su \textit{et al.} \cite{su2016}, a classifier validation scene featured targets positioned at different spatial locations but at approximately the same depths. Given the core operational principles of ToF detectors, unless targets exhibit significant variation in albedo, the returned signals are bound to occlude, resulting in signals that should not belong to either class. However, the results made no mention of how possible occlusions were handled, highlighting a limitation in the experimental methodology.

In multi-pixel ToF sensors, transient signals per pixel can be spatially organised to form a two-dimensional matrix (Figure~\ref{fig:3_10}).
\begin{figure}[H]
\centering
\includegraphics[width=\textwidth]{figures/3.10.png}

    \caption[Comparison of 2D laterally resolved pixel intensity with APS and temporal resolution ToF sensor.]{\textbf{Comparison of 2D laterally resolved pixel intensity with APS and temporal resolution ToF sensor.} Left: 2D laterally resolved pixel intensity with APS. Right: 2D laterally resolved pixels each of 1D temporal resolution ToF sensor. On the right side, red signal peaks represent the ToF sensor reference signal (see chapter 4) and the right peak in blue represent time of flight signal from target object.}
\label{fig:3_10}
\end{figure}

Multiple or all pixels can be averaged to improve signal-to-noise ratio for whole-scene depth estimation. Other common LiDAR applications of depth estimation with two-dimensional temporal data include non-line-of-sight imaging, pose estimation, scale estimation, orientation or perspective estimation of target materials, and related tasks. Decision tasks that can be performed with one-dimensional transient images include depth estimation (echolocation or time-of-flight), image segmentation (multi-signal peak detection), scale estimation (signal-to-depth ratio), classification (peak labeling), object material detection, and two-dimensional object recovery from obscured diffusive media \cite{lyons2019}. Inferring scale variability, perspective, or material composition of annotated signals is also possible with non-spectral transient processing.

\subsubsection{Depth estimation}
\label{subsubsec:temporal_depth_estimation}

Depth estimation with achromatic transient signals relies on echolocation techniques, commonly termed time-of-flight (ToF). ToF is considered a more reliable and precise depth estimation approach compared to stereoscopy \cite{tachella2019,laga2022}. ToF systems estimate depth from each pixel for different viewpoints of the scene. Each pixel with sufficient signal can reliably represent the whole scene along the axial plane, providing absolute distance measurements rather than relative depth information.

\subsubsection{Segmentation}
\label{subsubsec:temporal_segmentation}

Segmentation for temporal signals is typically represented as localised signal peak clusters. Contrary to laterally resolved achromatic images, where segmented regions of interest consist of multiple pixels, it is possible to segment temporal signals based on a single pixel (Figure~\ref{fig:3_11}). This capability arises from the temporal nature of the signal, where different peaks in the time-of-flight histogram correspond to different targets at varying distances along the axial range.

\begin{figure}[H]
    \centering
    \includegraphics[width=\textwidth]{figures/3.11.png}

    \caption[Annotated achromatic target detection in 1D temporal data from SPAD sensor.]{\textbf{Annotated achromatic target detection in 1D temporal data from SPAD sensor.}}
    \label{fig:3_11}
    \end{figure}


Localisation in temporal signals is comparable to segmentation in laterally resolved intensity images, with the key difference being the direction of the axis along which clusters are created (Figure~\ref{fig:3_12}). Transient signals can only be localised in terms of range distance; that is, identification of scene components is typically within the range of the signal propagation path. This resolution helps depict key targets of the scene with respect to relative foreground and background targets (Figure~\ref{fig:3_12}), each with accurately estimated distance from the sensor.

\begin{figure}[H]
\centering
\includegraphics[width=\textwidth]{figures/3.12.png}
    \caption[2D laterally resolved intensity image and temporal data from SPAD along axial plane.]{\textbf{2D laterally resolved intensity image and temporal data from SPAD along axial plane.} Dark yellow (APS signal): 2D laterally resolved intensity image from scene. Red (high gain sensor): Temporal data by SPAD along axial plane. First peak represents mailbox in foreground. Highest peak represents target (boy) and last peak indicative of background object (bus).}
\label{fig:3_12}
\end{figure}

Axial clustering represents a highly informative approach to scene stereopsis and thus serves as a major complementary component to two-dimensional laterally resolved images. Though individual ToF pixels are sufficient to describe scenes, different architectures may enable further lateral resolution of pixels into a two-dimensional matrix (see Figures~\ref{fig:3_10} and~\ref{fig:3_13}).

\begin{figure}[H]
\centering
\includegraphics[width=\textwidth]{figures/3.13.png}
    \caption[2D laterally resolved pixels with 1D temporal data and single pixel temporal signal resolution.]{\textbf{2D laterally resolved pixels with 1D temporal data and single pixel temporal signal resolution.} Left: 2D laterally resolved pixels each of 1D temporal data. Right: single pixel of 1D temporal signal resolution. The red curve represents reference signal (see chapter 4) and the blue curve represent return signal from target.}
\label{fig:3_13}
\end{figure}


This spatial organisation of temporal signals enables enhanced scene analysis by combining both axial depth information and lateral spatial resolution. The pattern of variations in signal distribution (see Sections~\ref{subsec:smooth_specular_reflection},~\ref{subsec:rough_diffuse_reflection}, and~\ref{subsec:subsurface_scattering}) across pixels encodes clues to the position of targets across different viewpoints of the range sensor. This is relevant for applications such as pose estimation and non-line-of-sight sensing \cite{lyons2019}.

\subsubsection{Scale estimation}
\label{subsubsec:temporal_scale_estimation}

Estimating target scale in achromatic transient signals is relatively less complex. A reasonable approach to this task is measuring the signal-to-depth ratio, either based on a single pixel or with the whole array. One key benefit of scale estimation with one-dimensional transient signals is the low computational complexity yet accurate and reliable results. Even in spectral scale estimates, a single dimension should be sufficient for estimating target scale, demonstrating the efficiency of temporal approaches.

\subsubsection{Classification}
\label{subsubsec:temporal_classification}

Unlike photometric classification, a single pixel can be localised and labelled to indicate different targets (identified as signal peaks) in the axial range. An unsupervised approach to cluster identification may be used to detect and label multiple centre-of-mass (CoM) instances. Single-pixel classification is essential for whole-scene representations. When clusters in range are localised, it becomes straightforward to reliably classify the main target, target background, and/or foreground for in-depth scene descriptions. Peak-wise classification can be highly informative as it helps distinguish different objects within the scene as perceived at different positions within the sensor's line of sight (Figure~\ref{fig:3_13}). Correctly detecting the target peak is subject to the quality of the clustering or segmentation algorithm. The distribution of signal peaks can be generally approximated as Gaussian \cite{koerner2021} (Figure~\ref{fig:3_14}).
\begin{figure}[H]
\centering
\includegraphics[width=\textwidth]{figures/3.14.png}

    \caption[Single ToF histogram showing effects of ambient and Poisson noise versus Gaussian approximation.]{\textbf{Single ToF histogram showing effects of ambient and Poisson noise versus Gaussian approximation.} Left: illustration of a single ToF histogram showing less smoothness in return signal curve as effects of ambient (noise from surrounding light sources) and Poisson noise (noise from the randomness in photon arrival). Start bin at time position ($t_s$), end bin at time position ($t_e$), corresponds to the time of flight of photons to the target surface. Ambient noise is assumed to be uniformly distributed forming a baseline (Base error). Poisson noise is not uniform, thus embedded within the return signal. Right: illustration of approximated as Gaussian (smoother) single ToF histogram showing return signal. Thus, less obvious effects of ambient (noise from surrounding light sources) and Poisson noise (noise from the randomness in photon arrival). Start bin at time position ($t_s$), end bin at time position ($t_e$), $\mu$ corresponds to the time of flight of photons to the target surface.}
\label{fig:3_14}
\end{figure}

Additionally, subsurface scattering effects may elongate the tail due to delayed photon arrivals (see Chapter~\ref{chap:introduction}). Most generally, the resulting signal can be described as a convolution:
\begin{equation}
h(t) = p(t) * r(t) * j_{\text{TDC}}(t) * j_{\text{SPAD}}(t)
\label{eq:temporal_convolution}
\end{equation}
where $p(t)$ represents the temporal profile of the laser pulse, $r(t)$ is the target material response, $j_{\text{TDC}}(t)$ is the time-to-digital converter jitter, and $j_{\text{SPAD}}(t)$ is the SPAD jitter. However, other potential factors such as atmospheric scattering or photon pileup (which occurs under high ambient levels or highly intense reflection) may contribute to overall histogram quality. The histograms measured in the experimental chapters of this work (see Chapter~\ref{chap:introduction}) are not impacted by photon pileup. In translucent or diffusive materials (most of which have been used in this work), the signal distribution is comparable to heat energy diffusion, as described by Lyons \textit{et al.} \cite{lyons2019}. Other optical models consider diffusive subsurface scattering and background illumination in modelling signal distribution, as presented by Becker and Koerner \cite{becker2023}.

\subsubsection{Material detection}
\label{subsubsec:temporal_material_detection}

Temporal data is considered a precise depth estimation approach. Sourced from target surface reflectivity information, it encodes reflectance properties unique to the material (see Sections~\ref{subsec:smooth_specular_reflection},~\ref{subsec:rough_diffuse_reflection}, and~\ref{subsec:subsurface_scattering}). Thus, temporal signals are more relevant for inferring material composition of target objects rather than localisation in lateral space. Detecting objects in lateral space requires multiple pixels. However, a unit pixel with sufficient transient signal could describe an entire scene along the axial axis. Figures~\ref{fig:3_11} and~\ref{fig:3_13} describe segmentation and classification of signal peaks with time-of-flight pixels. However, these peaks are better suited for inferring material composition or localisation in space, providing complementary information to spatial imaging systems \cite{su2016,becker2023,yang2024}.



\subsection{Processing models of non-spectral temporal signal}
\label{subsec:processing_models_non-spectral_temporal}

The most common detectors for photon timing are high-gain sensors, such as time-of-flight SPAD sensors, which often exhibit coarse lateral resolution. To model photon arrival times, active illumination is projected onto a target surface. Unabsorbed photons are reflected back or transmitted through the surface material. The signal peak in the histogram is assumed to have a Gaussian temporal profile, centred around time $t = t_0$ (corresponding to the time-of-flight of photons to the target surface), with an average of $S$ total signal photons in the peak. The signal $S$ scales linearly with target reflectivity and inversely with target distance (for diffuse reflection surfaces) \cite{tontini2020}. Equations (9) and (10) in \cite{tontini2020} can be used to estimate the levels of $S$ and $b$ from the measured histogram data. Equation~\ref{eq:pure_material_gaussian} provides the per-bin approximation of the expected photon count for a pure material surface, where the continuous Gaussian profile is discretised into individual time bins.

The approximated Gaussian peak is discretised into $B$ bins, where each bin $i$ has a start time $t_i$ and end time $t_{i+1}$. The start time is given by $t_i = t_0 + (i - B/2) \cdot \Delta t$, where $i$ ranges from 1 to $B$ (the maximum number of bins $B$ in histograms from Chapters 4, 5, and 6 is 144) and $\Delta t$ represents the histogram bin width (approximately 250~ps for the histograms in Chapters 4, 5, and 6). Ambient noise, including dark counts, is assumed to follow a uniform distribution with an average of $b$ photons per bin (Figure~\ref{fig:3_14}).

For a pure material surface, the temporal signal per bin is modelled as a Gaussian distribution:
\begin{equation}
\mu_i = \frac{S}{\sqrt{2\pi\sigma^2}} \exp\left(-\frac{(t_i - t_0)^2}{2\sigma^2}\right)
\label{eq:pure_material_gaussian}
\end{equation}
\addcontentsline{loe}{equations}{\protect\numberline{\theequation}$\mu_i = \frac{S}{\sqrt{2\pi\sigma^2}} \exp\left(-\frac{(t_i - t_0)^2}{2\sigma^2}\right)$}

where $\mu_i$ represents the expected photon count approximation in bin $i$ (not the whole peak), $S$ is the total signal photons across the entire peak, $\sigma$ is the standard deviation of the Gaussian profile, and $t_0$ is the peak arrival time. This equation provides a per-bin approximation of the continuous Gaussian distribution, where each bin $i$ receives a fraction of the total signal $S$ based on its temporal position relative to the peak.

Since all bins are subject to shot noise, the individual photon count $h_i$ for each bin $i$ is drawn from a Poisson distribution with mean $\mu_i + b$, where $\mu_i$ is the per-bin expected photon count from Equation~\ref{eq:pure_material_gaussian}. The probability mass function (PMF) for each bin is expressed as:
\begin{equation}
P(h_i = k) = \frac{(\mu_i + b)^k e^{-(\mu_i + b)}}{k!}
\label{eq:poisson_pmf}
\end{equation}
\addcontentsline{loe}{equations}{\protect\numberline{\theequation}$P(h_i = k) = \frac{(\mu_i + b)^k e^{-(\mu_i + b)}}{k!}$}

When the ToF detector observes a multi-material surface (comprising $M$ total surfaces), the resulting histogram peak is modelled as a superposition of multiple Gaussian peaks, each representing one of the $M$ different surfaces. The per-bin expected photon count approximation is derived from a Gaussian mixture model, where each material $m$ contributes a fraction of its total signal $S_m$ to each bin $i$ based on the Gaussian distribution. The expected photon count per bin is given by:
\begin{equation}
\mu_i = \sum_{m=1}^{M} \frac{S_m}{\sqrt{2\pi\sigma_m^2}} \exp\left(-\frac{(t_i - t_{0,m})^2}{2\sigma_m^2}\right) + b
\label{eq:gaussian_mixture}
\end{equation}
\addcontentsline{loe}{equations}{\protect\numberline{\theequation}$\mu_i = \sum_{m=1}^{M} \frac{S_m}{\sqrt{2\pi\sigma_m^2}} \exp\left(-\frac{(t_i - t_{0,m})^2}{2\sigma_m^2}\right) + b$}

where $S_m$ is the total signal photons for material $m$ across the entire peak, $\sigma_m$ is the standard deviation, and $t_{0,m}$ is the peak time for material $m$. This equation provides a per-bin approximation where the expected photon count in each bin $i$ is the sum of contributions from all $M$ materials, each following a Gaussian distribution (as in Equation~\ref{eq:pure_material_gaussian}), plus the ambient noise $b$. The assumption that each material peak is Gaussian also explains the elongation of subsurface scattering peaks due to delayed photon returns \cite{koerner2021,su2016}.

Under controlled circumstances, the detector may receive signals from a pure material. In such instances, the target is considered a non-multi-material surface such that $M = 1$. The per-bin expected photon count approximation is then given by the Gaussian model of a single pure material (Equation~\ref{eq:pure_material_gaussian}), where each bin receives a fraction of the total signal $S$ based on its temporal position relative to the peak.


\subsection{Processing complexity of non-spectral temporal signal}
\label{subsec:processing_complexity_non-spectral_temporal}

This section presents a machine-independent analysis of the computational complexity for processing non-spectral temporal signals in SPAD-based time-of-flight systems \cite{piron2020}. Consider a detector array comprising $N$ pixels, where each pixel $i$ maintains an initial state $I_{\text{pix},i}$ representing the absence of photon return time records. Each pixel contains $B$ time bins, representing the temporal spread $T_B$ over a given time span. In the absence of input signal, pixel intensities are represented by an empty array $\mathbf{D}$ of $N$ elements, where each element contains either zero or a negligible noise component $c$.

Time-to-digital converter (TDC) signal input per pixel is represented by a distribution of incident signal across $B$ time bins \cite{sesta2023,alabbas2018}. To analyse the computational complexity of the resulting addition operation, the pixel update function is projected onto a two-dimensional plane with axes representing time (number of bins) and array size (number of pixels). In the optimal pixel update scenario, a single pixel $I_{\text{pix},0}$ with noise intensity $c$ receives a returned signal distributed over $B$ time bins. The resulting temporal distribution values after incident signal accumulation represent the reflectivity per unit time plus any preexisting noise in the pixel.

The computational complexity of signal input operations, whether processing hybrid or pure surface materials, grows with both the size of the pixel array $N$ and the number of time bins $B$ per pixel \cite{knuth1976}. In the worst-case scenario, signal accumulation occurs once across the total array size $N$ and once across the total number of bins $B$ for the temporal signal spread per pixel, resulting in memory and time complexity of $\mathcal{O}(N \times B)$ \cite{knuth1976}. This quadratic scaling in the two-dimensional parameter space demonstrates the fundamental trade-off between temporal resolution and computational efficiency in non-spectral temporal signal processing \cite{su2016}.

Two boundary cases illustrate the relationship between temporal and spatial complexity. If all photon arrival times are compressed into a single bin ($B = 1$), the complexity reduces to $\mathcal{O}(N)$, matching that of achromatic intensity pixel updates, since the access operation scans the total number of pixels once with no additional operations required for temporal binning \cite{knuth1976}. Conversely, when processing a single pixel with a large number of bins ($N = 1$), the algorithm scans the total number of bins once, again resulting in linear complexity $\mathcal{O}(B)$ that matches multi-pixel achromatic intensity updates. The quadratic complexity $\mathcal{O}(N \times B)$ emerges from the combined processes needed to update all bins per pixel and all pixels of the sensor, representing the general case for temporal histogram processing in SPAD-based systems \cite{hutchings2019,zhang2019}.

\newpage
\section{Spectral signal measurements}
\label{sec:spectral_signal_measurements}

Spectral signal measurements exploit wavelength-dependent variations in electromagnetic radiation to characterise materials through their unique spectral signatures \cite{elachi2021,schott2007}. Unlike non-spectral (achromatic) approaches that discard wavelength information, spectral detectors incorporate specialised circuitry—including optical filters, dispersive elements, or wavelength-selective sensors—to resolve and process the spectral composition of incident signals \cite{heinz2001,vasile2023}. This wavelength-resolved capability enables material discrimination based on characteristic absorption, reflection, and transmission properties that manifest across different spectral bands \cite{elachi2021,inagaki2020}.

Similar to non-spectral processing, spectral signals can be measured as either intensity distributions or temporal (time-resolved) photon arrival patterns \cite{moramartin2021,inagaki2020}. For intensity-based measurements, spectral information is typically represented as multi-channel images where each channel corresponds to a specific wavelength band or colour space component (e.g., RGB images with three spectral channels) \cite{schott2007}. Temporal spectral measurements, by contrast, capture wavelength-resolved photon arrival time distributions, enabling joint spectral–temporal analysis of material-dependent optical responses \cite{inagaki2020}.

Spectral measurement systems are classified according to the number of resolved wavelength bands: monospectral systems filter a single spectral band, multispectral systems capture multiple discrete bands (typically 3–10 channels), and hyperspectral systems acquire continuous spectral information across hundreds of narrow wavelength bands \cite{heinz2001,vasile2023}. The computational complexity of spectral processing scales with the number of spectral channels, as each band requires separate calibration, feature extraction, and often spectral unmixing algorithms to separate material contributions from mixed pixels \cite{heinz2001}. This dimensionality overhead distinguishes spectral approaches from achromatic methods, which operate on single-channel intensity or temporal data without explicit wavelength decomposition \cite{su2016,becker2023}.

Despite this increased complexity, spectral signal measurements enable reliable performance across diverse machine vision tasks including depth estimation, scale estimation, segmentation, classification, and object and material detection \cite{vasile2023,elachi2021}. The distinguishing characteristic of spectral processing—the explicit exploitation of wavelength-dependent material properties—provides complementary information to spatial and temporal cues, enabling robust material characterisation that addresses limitations in appearance-only vision systems \cite{heinz2001,vasile2023}.

\subsection{Spectral intensity signal}
\label{subsec:spectral_intensity_signal}

Spectral intensity signals exploit wavelength-dependent variations in electromagnetic radiation for spectroscopy and material characterisation tasks \cite{elachi2021,schott2007}. During signal decoding, incident radiation is filtered by wavelength, with each channel corresponding to a specific spectral dimension or colour component present in the detected band. The collected photoelectrons are converted to equivalent voltage values, known as intensity measurements (Figure~\ref{fig:3_15}).
\begin{figure}[H]
\centering
\includegraphics[width=\textwidth]{figures/3.15.png}

    \caption[Multispectral illumination source and chromatic pixel-well conversion to voltage values.]{\textbf{Multispectral illumination source and chromatic pixel-well conversion to voltage values.} From left: multispectral illumination source towards chromatic pixel-well (abstract illustration of electrons by incoming photons in pixel). For spectral analyses, pixel filters or splitters and other circuitry convert electrons from each identified spectrum to equivalent voltages; numeric voltage values (red:208, green:177, blue:212) known as pixel intensity are output for each channel. Representation of intensity as visualised in resulting coloured multispectral pixel (square tile). The diagram uses red, green, blue channels. However, this is possible with single or full spectrum. Also other colour combinations beside RGB are possible.}
\label{fig:3_15}
\end{figure}

However, the measured intensity is governed by the inverse square law \cite{griffiths2017}, which dictates that electromagnetic radiation intensity decreases proportionally to the square of the distance from the source. This fundamental relationship means that for remote sensing applications, the received signal intensity scales inversely with the square of the target distance, requiring distance compensation in spectral analysis to accurately characterise material properties independent of range. In two-dimensional pixel architectures, the resulting image is laterally resolved across spatial dimensions, with each pixel encoding intensity values for one or more spectral channels \cite{schott2007}.

Monospectral images contain a single channel, representing one colour or frequency dimension (Figure~\ref{fig:3_16}).
\begin{figure}[H]
\centering
\includegraphics[width=\textwidth]{figures/3.16.png}

    \caption[Comparison of 2D laterally resolved multispectral (RGB) pixels from APS and temporal resolution from SPAD sensor.]{\textbf{Comparison of 2D laterally resolved multispectral (RGB) pixels from APS and temporal resolution from SPAD sensor.} Left: 2D laterally resolved multispectral (RGB) pixels each of scaler intensity vales from APS. Right: 2D laterally resolved pixels each of 1D temporal resolution from SPAD sensor.}
\label{fig:3_16}
\end{figure}

However, single-channel or monospectral signals should not be confused with achromatic processing, as discussed in Section~\ref{subsec:non_spectral_methods}. In non-spectral image processing, tasks do not require explicit awareness of wavelength composition to function effectively. Even when it is possible to calculate or estimate signal frequency, such information is not essential for task realisation. In contrast, awareness of the specific colour or spectral band is essential for processing monospectral tasks. For example, in two-dimensional monospectral imaging (e.g., red-channel images), without awareness of the channel frequency, it becomes difficult to label the resolved image by its colour—whether red, green, or another spectral band \cite{elachi2021}.

Multispectral images combine more than one channel yet fewer than the full spectrum represents. A common multispectral representation is the RGB image, which captures three discrete spectral bands corresponding to red, green, and blue wavelengths \cite{schott2007}. Other combinations of multiple channels are possible, depending on the specific application requirements. Hyperspectral images, as implied by the name, utilise all available channels in the signal to represent images, typically comprising hundreds of narrow, contiguous wavelength bands that provide continuous spectral coverage \cite{heinz2001,vasile2023}.

\subsubsection{Depth estimation}
\label{subsubsec:spectral_depth_estimation}

Depth estimation in spectral signals follows similar approaches as in achromatic signals: stereoscopy, interferometry, and time-of-flight methods are applicable \cite{laga2022}. Averaging pixels in intensity images (colour or greyscale) is not typically a relevant operation for depth estimation. However, channel averaging and other dimensionality reduction operations may be necessary to enhance signal processing efficiency and reduce computational complexity \cite{heinz2001}.

\subsubsection{Segmentation}
\label{subsubsec:spectral_segmentation}

Automatic annotation or localisation tasks in two-dimensional spectral images often occur across laterally resolved arrays of pixels \cite{blaschke2010}. The operation involves detection of abrupt changes in signal with respect to spatial cues such as edges, outlines, and shapes (Figure~\ref{fig:3_17}).
\begin{figure}[H]
\centering
\includegraphics[width=\textwidth]{figures/3.17.png}

    \caption[Annotated spectral object detection in 2D laterally resolved multispectral (RGB) image.]{\textbf{Annotated spectral object detection in 2D laterally resolved multispectral (RGB) image.} The single annotation of gold colour contains multiple pixels each of three RGB channels; red:253, green:208, blue:23. A single intensity pixel (unlike SPAD detector) could not represent the circle object illustrated in the image.}
\label{fig:3_17}
\end{figure}

In colour images, there is no evidence in the literature to support the necessity of segmentation at the individual pixel level \cite{haralick1984}. Regions of interest (ROI) spanning multiple pixels can be segmented to represent objects or material boundaries. Segmentation helps localise relevant aspects of the image for higher-level vision tasks such as classification and object detection \cite{blaschke2010}. Due to the multichannel colour variable, two-dimensional spectral images are often considered more suitable for semantic scene analysis compared to greyscale images \cite{schott2007}. Individual pixel intensities themselves exhibit no significant internal edge variations to be segmented. The ROI annotations are useful for whole or partial scene representations, enabling material discrimination and object identification through spectral feature extraction \cite{heinz2001,vasile2023}.

\subsubsection{Scale estimation}
\label{subsubsec:spectral_scale_estimation}

Similar to achromatic target scale estimation, most spectral images (e.g., RGB) rely on stereoscopy methods to estimate scale using target depth-to-intensity ratios, which have known limitations (see Chapter~\ref{chap:literature_review}) \cite{laga2022}. This approach is less preferable for long-distance or low-illumination scenes and targets, where depth estimation accuracy degrades significantly .

\subsubsection{Classification}
\label{subsubsec:spectral_classification}

Pixels in spectral images may be labelled and categorised as belonging to some class. Reliable classifiers depend on quality annotation tasks that enable differentiation of different class features \cite{bishop2006}. While most classification algorithms require labeling of pixels or regions of interest (ROI), there are instances where it is possible to classify whole scenes or images without the laborious task of annotations; such algorithms are unsupervised \cite{bishop2006}. Most spectral classification algorithms are supervised and employ automatic pattern recognition models, such as deep learning architectures (see Chapter~\ref{chap:literature_review}) \cite{lecun1998}.

\subsubsection{Object detection}
\label{subsubsec:spectral_object_detection}

The definition of digital objects in signal processing literature is highly subjective. However, annotated objects of some class may be learned by automatic pattern recognition algorithms for generalisation on other unseen data \cite{redmon2015,redmon2018}. Object detection identifies specific image components via segmentation and labelling of regions of interest. Most machine vision algorithms perform this task with spectral images of commonly seen objects \cite{redmon2018}. In some instances, the task is complemented with depth information for better scene interpretation \cite{moramartin2021,yang2023,yang2024}. A key variable for determining material composition is surface reflectivity (see reflectance properties). Though it is possible to infer material type from spectral images, it cannot be reliably achieved with laterally resolved intensity signals alone \cite{su2016,becker2023}. Thus, spectral temporal signals become essential for making such inferences, as they encode both wavelength-dependent and time-resolved material properties \cite{inagaki2020,su2016,becker2023}.

\subsection{Processing models of spectral intensity signal}
\label{subsubsec:processing_models_spectral_intensity_detailed}

Intensity estimation models for spectral intensity signals differ from achromatic intensity signals by explicitly accounting for the spectrum of wavelengths present in the measurement. The total number of photons contained within each pixel per returned signal can be described using a multichannel linear model that accounts for spectral composition \cite{heinz2001,schott2007}. 

For a target material composed of pure elements, where the number of composite materials $M = 1$, pixel intensity is estimated from the reflectance of a single material $m$. The signal reflectance from material $m$ is given by $r_m$, and the fractional proportion of $m$ is represented by $p_m$. The error term $\epsilon$ captures variation between the actual observed surface reflectance and the computed pixel reflectance value from the model. Factors influencing $\epsilon$ include ambient noise, electronic noise, photon shot noise, and other measurement uncertainties \cite{elachi2021,schott2007}. This ideal pure material case, however, is rarely true for most real-world target surfaces.

Practically, target surfaces belong to objects composed of hybrid materials that vary in reflectivity. When a spectral pixel receives a signal from such a surface, the return signal is expected to have emerged from each material within the set of $M$ detectable materials. Thus, the resulting signal in each pixel is a linear proportional sum of reflectances from each detectable surface material $m$ \cite{heinz2001,vasile2023}. In cases of total occlusion by an opaque material within the same object, the occluded material is considered undetectable, as the reflectance will only emerge from the occluding opaque material \cite{schott2007}.

For a hybrid surface containing $M$ detectable materials, the multichannel spectral intensity model is expressed as:

\begin{equation}
U_\lambda = \sum_{m=1}^{M} p_m \cdot r_{m,\lambda} + \epsilon_\lambda
\label{eq:spectral_mixture}
\end{equation}
\addcontentsline{loe}{equations}{\protect\numberline{\theequation}$U_\lambda = \sum_{m=1}^{M} p_m \cdot r_{m,\lambda} + \epsilon_\lambda$}

where $U_\lambda$ represents the measured intensity at wavelength $\lambda$, $r_{m,\lambda}$ denotes the spectral reflectance of material $m$ at wavelength $\lambda$, $p_m$ represents the fractional proportion of each material $m$ (constrained such that $\sum_{m=1}^{M} p_m = 1$), and $\epsilon_\lambda$ accounts for wavelength-dependent noise and measurement uncertainties \cite{heinz2001,schott2007}. The fractional proportion per composite material can be solved using a system of linear equations across multiple spectral bands, enabling spectral unmixing to separate material contributions from mixed pixels \cite{heinz2001,vasile2023}. This approach forms the foundation for material discrimination in hyperspectral imaging systems, where the high-dimensional spectral signature enables identification of material composition through characteristic reflectance patterns across wavelength bands \cite{elachi2021,schott2007}.

\subsection{Processing complexity of spectral intensity signal}
\label{subsubsec:processing_complexity_spectral_temporal_pre}

This section analyses the algorithmic efficiency of processing spectral intensity signals, demonstrating how the additional spectral dimension increases computational complexity compared to achromatic intensity processing (see Section~\ref{subsec:processing_complexity_spectral_intensity}). Unlike non-spectral intensity signals, which operate on a two-dimensional spatial representation, spectral intensity processing introduces a third dimension corresponding to wavelength composition, fundamentally altering the computational requirements for signal processing operations \cite{heinz2001,schott2007}.

The digital input per pixel in spectral systems is represented by the initial sum of incident signal intensity along the spectral axis, followed by the fractional component contributions from each material $m$ within the set of $M$ detectable materials, as expressed in Equation~\ref{eq:spectral_mixture}. To assess the computational complexity of the resulting spectral unmixing operation, the worst-case scenario for pixel intensity update must account for both spatial pixel indexing and spectral band processing \cite{heinz2001,vasile2023}. 

For a detector array comprising $N$ pixels, where each pixel requires processing across $W$ spectral bands, the computational complexity scales as $\mathcal{O}(N \times W)$ in the general case \cite{knuth1976}. When spectral unmixing is performed to separate material contributions from mixed pixels, the complexity further increases due to the system of linear equations that must be solved across multiple spectral bands, resulting in worst-case complexity of $\mathcal{O}(N \times W \times M)$ where $M$ represents the number of composite materials \cite{heinz2001,vasile2023}. This quadratic scaling with respect to the number of spectral bands contrasts sharply with the linear complexity $\mathcal{O}(N)$ achieved in achromatic intensity processing, where wavelength-dependent operations are absent \cite{knuth1976}.

Compared to the $\mathcal{O}(N)$ complexity of achromatic intensity estimation presented in Section~\ref{subsec:processing_complexity_spectral_intensity}, spectral processing can be considered computationally costly, particularly for hyperspectral systems where $W$ may exceed hundreds of narrow, contiguous wavelength bands \cite{heinz2001,vasile2023}. The increased dimensionality requires extensive computational resources for storage, analysis, and feature extraction, with processing overhead that grows proportionally with spectral resolution \cite{schott2007}. Despite this complexity penalty, spectral signal processing remains preferable for natural image analysis tasks where wavelength-dependent material discrimination is essential \cite{elachi2021,schott2007}. However, for applications such as depth estimation, where temporal information provides comparable performance quality at lower computational cost, achromatic approaches offer a practical alternative that maintains detection accuracy while reducing system complexity \cite{su2016,becker2023}.

\subsection{Spectral temporal signal}
\label{subsec:spectral_temporal_signal}

Spectral temporal measurements combine echolocation-based time-of-flight analysis with wavelength-resolved photon arrival time detection across multiple spectral channels \cite{inagaki2020}. This approach extends achromatic temporal signal processing by incorporating spectral decomposition, enabling material discrimination through wavelength-dependent temporal signatures \cite{su2016,becker2023}. Pulsed direct time-of-flight (dToF) implementations for spectral temporal sensing employ multi-frequency illumination sources with wavelength-separation accessories to demodulate signals by spectral band. Similar to spectral intensity imaging, the detector explicitly accounts for variations in wavelength or frequency of the incident signal, providing enhanced material characterisation capabilities compared to achromatic temporal methods \cite{inagaki2020,su2016}.


\subsubsection{Depth estimation}
\label{subsubsec:spectral_temporal_depth_estimation}

Depth estimation in spectral temporal signals follows principles analogous to achromatic temporal depth estimation, with the additional multichannel attribute enabling enhanced signal processing \cite{koerner2021}. Averaging pixels across temporal channels can improve signal-to-noise ratio (SNR) in depth estimation tasks \cite{koerner2021}. However, individual pixels can also be utilised for depth estimation from different viewpoints, enabling multi-perspective analysis \cite{moramartin2021}. LiDAR applications of spectral temporal depth estimation include pose estimation, scale variability analysis, orientation estimation, and perspective estimation—challenges that are difficult to address using conventional two-dimensional intensity images alone \cite{moramartin2021,yang2024}.

\subsubsection{Segmentation}
\label{subsubsec:spectral_temporal_segmentation}

Segmentation of spectral temporal data employs peak clustering methodologies similar to achromatic temporal signal processing, with peak clusters exhibiting homogeneity within material classes while varying significantly between different material types \cite{su2016,becker2023}. Peak clusters from a single target in one spectral channel may exhibit different temporal distributions compared to peaks from the same target in other channels, reflecting variations in material interactions with different radiation frequencies \cite{inagaki2020,su2016}. Since each pixel receives the full spectrum of incident signals, whole-scene segmentation or target clustering can be performed based on a single pixel's spectral temporal response .

Spectral temporal signals resolve along the sensor's line of sight (Figure~\ref{fig:3_16}, Figure~\ref{fig:3_20}), enabling segmentation or peak clustering at three distinct positions: scene target, background, and foreground while time-domain spectroscopy offers enhanced segmentation capabilities through spectral discrimination, it incurs greater computational complexity compared to non-spectral temporal signal processing equivalents \cite{inagaki2020,heinz2001}. Two-dimensional laterally resolved temporal pixels in spectral systems share similar applications with achromatic temporal processing, augmented by the multichannel variable, including non-line-of-sight detection and pose estimation \cite{moramartin2021,lyons2019}.

\subsubsection{Classification}
\label{subsubsec:spectral_temporal_classification}

When clustered peaks are labelled according to their surface of emergence, they can be assigned to predefined material classes, enabling learning by pattern recognition algorithms \cite{su2016,becker2023}. In spectral temporal clustering, target peaks in range can be classified based on the spectral composition of individual pixels, defined regions of interest across multiple pixels, or the full pixel array—particularly for cases where the target peak represents the only cluster in range \cite{inagaki2020,su2016}. Similar to all temporal signal processing approaches, single pixels with sufficient signal quality can describe entire scenes along the axial dimension .

Illumination travels along the axial plane of the detector's line of sight, enabling spectral temporal data classification into target, foreground, or background components (Figure~\ref{fig:3_18}). Figure~\ref{fig:3_20} illustrates variations between axial and lateral resolutions of the same scene.


\begin{figure}[H]
    \centering
    \includegraphics[width=\textwidth]{figures/3.18.png}
        \caption[Annotation of multiple signal peaks in range for single multichannel ToF pixel.]{\textbf{Annotation of multiple signal peaks in range for single multichannel ToF pixel.} Left annotated target indicative of main target with higher intensity (reflectivity) whereas second annotation indicated another target in background (behind main target).}
    \label{fig:3_18}
    \end{figure}



\begin{figure}[H]
\centering
\includegraphics[width=\textwidth]{figures/3.20.png}

    \caption[2D laterally resolved pixels with unique wavelength and temporal resolution from SPAD sensor.]{\textbf{2D laterally resolved pixels with unique wavelength and temporal resolution from SPAD sensor.} From the left: 2D laterally resolved pixels each of unique wavelength and corresponding temporal resolution from SPAD sensor. Bottom: single pixel of 1D temporal signal resolution. From the right: 2D multispectral intensity image from scene with active pixel sensor (located above) and recorded multifrequency temporal data by single SPAD pixel along axial plane. First peak representing mailbox in foreground. Highest peak representing target (boy) and last peak indicative of background object (bus).}
\label{fig:3_20}
\end{figure}

Demonstrating how spectral temporal processing provides complementary information to spatial imaging systems \cite{yang2024,moramartin2021}. Depending on annotation quality, spectral segments can serve as reliable inputs for understanding material composition of emerging surfaces, with wavelength-dependent temporal signatures encoding material-specific optical properties \cite{su2016,becker2023,inagaki2020}.

\subsubsection{Scale estimation}
\label{subsubsec:spectral_temporal_scale_estimation}

Estimating target scale in multichannel temporal signals is relatively less complex than stereoscopic methods, with signal-to-depth ratios estimable using all spectral channels from a single pixel or across the entire sensor architecture in a laterally resolved two-dimensional fashion \cite{koerner2021}. The primary difference between scale estimation with one-dimensional transient signals versus spectral temporal signals lies in computational complexity, though both approaches yield accurate and reliable results \cite{koerner2021}. A single dimension of the spectrum can be sufficient for estimating target scale, demonstrating the efficiency of temporal approaches even when spectral information is available .

\subsubsection{Material detection}
\label{subsubsec:spectral_temporal_material_detection}

In spectral temporal signals, peaks represent target reflectivity, with material composition derived from temporal reflectivity values per spectral channel \cite{su2016,becker2023,inagaki2020}. Such temporal data can reliably describe material composition of surfaces from which signals are reflected, with temporal annotations possible from a single pixel given sufficient observed targets \cite{su2016}. Coarse lateral resolution of temporal data is possible in time-of-flight sensors with two-dimensional architectures, though these configurations have specialised use cases such as pose estimation \cite{moramartin2021,yang2024}. A combination of spectral intensity and non-spectral temporal signals (Figure~\ref{fig:3_19}) can be employed for enhanced object detection with material identification.
\begin{figure}[H]
\centering
\includegraphics[width=\textwidth]{figures/3.19.png}

    \caption[Spectral signal resolution across different measurements (intensity and temporal).]{\textbf{Spectral signal resolution across different measurements (intensity and temporal).}}
\label{fig:3_19}
\end{figure}

Demonstrating the complementary benefits of integrating spatial and temporal modalities \cite{yang2024,moramartin2021}.



\subsection{Processing models of spectral temporal signal}
\label{subsubsec:processing_models_spectral_temporal}

Spectral temporal signal processing extends the non-spectral temporal models described in Section~\ref{subsec:processing_models_non-spectral_temporal} by incorporating wavelength-dependent variations in object reflectivity and ambient levels \cite{inagaki2020,su2016}. Under normal circumstances, a detector may receive signals from multi-material surfaces, where the temporal spread of photon return times is represented by per-bin approximations using Gaussian distributions and fractional proportions of materials, analogous to the non-spectral temporal signal formulation (Equations~\ref{eq:pure_material_gaussian},~\ref{eq:poisson_pmf}, and~\ref{eq:gaussian_mixture}) \cite{tontini2020,koerner2021}.

For spectral temporal measurements, object reflectivity and ambient photon levels vary per wavelength $\lambda$, requiring explicit wavelength-dependent modelling across multiple spectral channels \cite{inagaki2020,heinz2001}. The spectral temporal versions of the signal processing equations account for these wavelength-dependent variations as follows, providing per-bin approximations analogous to the non-spectral temporal models.

For a pure material surface at wavelength $\lambda$, the per-bin expected photon count approximation is modelled as a wavelength-dependent Gaussian distribution:
\begin{equation}
\mu_{i,\lambda} = \frac{S_\lambda}{\sqrt{2\pi\sigma_\lambda^2}} \exp\left(-\frac{(t_i - t_{0,\lambda})^2}{2\sigma_\lambda^2}\right)
\label{eq:spectral_temporal_pure_material}
\end{equation}
\addcontentsline{loe}{equations}{\protect\numberline{\theequation}$\mu_{i,\lambda} = \frac{S_\lambda}{\sqrt{2\pi\sigma_\lambda^2}} \exp\left(-\frac{(t_i - t_{0,\lambda})^2}{2\sigma_\lambda^2}\right)$}

where $\mu_{i,\lambda}$ represents the expected photon count approximation in bin $i$ at wavelength $\lambda$ (not the whole peak), $S_\lambda$ is the total signal photons at wavelength $\lambda$ across the entire peak (scaling with wavelength-dependent target reflectivity), $\sigma_\lambda$ is the wavelength-dependent standard deviation of the Gaussian profile, and $t_{0,\lambda}$ is the peak arrival time at wavelength $\lambda$ \cite{tontini2020,koerner2021}. This equation provides a per-bin approximation where each bin $i$ receives a fraction of the total signal $S_\lambda$ based on its temporal position relative to the peak. The wavelength-dependent signal strength $S_\lambda$ reflects material-specific optical properties that vary with incident radiation frequency, enabling enhanced material discrimination through spectral-temporal signatures \cite{inagaki2020,su2016,becker2023}.

When the spectral temporal detector observes a multi-material surface comprising $M$ total surfaces, the resulting histogram peak is modelled as a superposition of wavelength-dependent Gaussian peaks, each representing one of the $M$ different materials. The per-bin expected photon count approximation is derived from a wavelength-dependent Gaussian mixture model, where each material $m$ at wavelength $\lambda$ contributes a fraction of its total signal $S_{m,\lambda}$ to each bin $i$ based on the Gaussian distribution. The expected photon count per bin is given by:
\begin{equation}
\mu_{i,\lambda} = \sum_{m=1}^{M} \frac{S_{m,\lambda}}{\sqrt{2\pi\sigma_{m,\lambda}^2}} \exp\left(-\frac{(t_i - t_{0,m,\lambda})^2}{2\sigma_{m,\lambda}^2}\right) + b_\lambda
\label{eq:spectral_temporal_mixture}
\end{equation}
\addcontentsline{loe}{equations}{\protect\numberline{\theequation}$\mu_{i,\lambda} = \sum_{m=1}^{M} \frac{S_{m,\lambda}}{\sqrt{2\pi\sigma_{m,\lambda}^2}} \exp\left(-\frac{(t_i - t_{0,m,\lambda})^2}{2\sigma_{m,\lambda}^2}\right) + b_\lambda$}

where $S_{m,\lambda}$ is the total signal photons for material $m$ at wavelength $\lambda$ across the entire peak, $\sigma_{m,\lambda}$ is the wavelength-dependent standard deviation, $t_{0,m,\lambda}$ is the peak time for material $m$ at wavelength $\lambda$, and $b_\lambda$ denotes the wavelength-dependent ambient noise level per bin \cite{tontini2020,koerner2021,heinz2001}. This equation provides a per-bin approximation where the expected photon count in each bin $i$ at wavelength $\lambda$ is the sum of contributions from all $M$ materials, each following a wavelength-dependent Gaussian distribution (as in Equation~\ref{eq:spectral_temporal_pure_material}), plus the wavelength-dependent ambient noise $b_\lambda$. Each spectral channel provides independent temporal signatures that encode material-specific optical interactions \cite{inagaki2020,su2016,becker2023}. The wavelength-dependent variations in signal strength, temporal spread, and ambient levels enable spectral temporal systems to achieve superior material discrimination compared to achromatic temporal processing, at the cost of increased computational complexity due to multi-channel processing requirements \cite{heinz2001}.


\subsection{Processing complexity of spectral temporal signal}
\label{subsubsec:processing_complexity_spectral_temporal}

The computational complexity of spectral temporal signal processing scales with three primary dimensions: the detector array size $N$ (total number of pixels), the temporal resolution $B$ (number of time bins per pixel), and the spectral dimensionality $K$ (number of wavelength channels) \cite{knuth1976}. In the most general case, pixel update operations must accumulate signals across all three dimensions, resulting in complexity that grows as $O(N \times B \times K)$ \cite{knuth1976}.

For a single pixel receiving signals distributed over $B$ time bins across $K$ spectral channels, the algorithm must process each spectral-temporal combination independently. The worst-case complexity occurs when signal accumulation requires scanning the entire array size $N$ once, the full spectral dimension $K$ once, and all temporal bins $B$ once per pixel, yielding a total computational cost proportional to the product of these three factors \cite{tontini2020,koerner2021}.

Two limiting cases demonstrate how spectral temporal complexity reduces to simpler regimes. When all photon arrival times are compressed into a single temporal bin ($B = 1$), the temporal distribution collapses and complexity becomes comparable to achromatic intensity pixel updates, scaling as $O(N \times K)$. In this regime, the algorithm scans the total number of pixels once and processes each spectral channel independently, eliminating temporal binning overhead. Conversely, when only a single pixel is considered with large $B$, the algorithm scans all $B$ bins once for that pixel, yielding complexity $O(B)$ that matches multi-pixel achromatic intensity updates for a single spectral channel \cite{koerner2021}.

The fundamental trade-off in spectral temporal processing lies in the increased dimensionality: while spectral resolution enables enhanced material discrimination through wavelength-dependent signatures \cite{inagaki2020,heinz2001}, the computational overhead grows multiplicatively with spectral channels, temporal bins, and array size. This complexity scaling motivates the non-spectral temporal approach employed in this work, where material discrimination is achieved through temporal signatures alone without explicit wavelength decomposition, reducing complexity from $O(N \times B \times K)$ to $O(N \times B)$ while maintaining comparable material detection performance \cite{su2016,becker2023}.



\section{Future research and development}
\label{sec:future_research_directions}

Future research will investigate how different materials influence Gaussian and linear mixture models (Equations~\ref{eq:gaussian_mixture} and~\ref{eq:spectral_temporal_mixture}) within the described methods \cite{tontini2020,koerner2021}. This will involve systematic studies across material categories to characterise how material properties—including optical constants, scattering coefficients, and absorption characteristics—manifest in mixture model parameters \cite{su2016,becker2023}. The goal is to establish material-specific model formulations that improve detection accuracy and enable more precise material characterisation. Research will explore how material heterogeneity, surface roughness, and internal structure affect mixture model assumptions, potentially leading to adaptive models that adjust their parameters based on material type \cite{koerner2021}.

Significant effort will be dedicated to formalizing fundamental material categories that can be sensed using non-spectral approaches \cite{su2016,becker2023,tafone2023}. This will involve developing a comprehensive taxonomy of materials based on their temporal response characteristics, optical properties, and interaction mechanisms with time-of-flight signals. The categorization will consider factors such as material composition (organic versus inorganic), optical properties (transparent versus opaque, reflective versus absorptive), structural characteristics (homogeneous versus heterogeneous), and temporal response signatures \cite{elachi2021}. This formalization will provide a foundation for systematic material detection research and enable the development of category-specific processing strategies.

Corresponding quantum physics equivalences of the models will be investigated in future research. The interaction between photons and materials, as measured in time-of-flight systems, has fundamental quantum mechanical origins \cite{elachi2021,griffiths2017}. Future work will explore how classical signal processing models relate to quantum optical processes, including photon-material interactions, quantum coherence effects, and the quantum nature of light-matter interactions \cite{elachi2021}. This investigation may reveal deeper insights into the physical processes underlying material detection, potentially leading to more physically grounded models and improved understanding of the fundamental limits of material characterisation using temporal signals. Research will examine how quantum effects such as photon statistics, coherence, and entanglement might influence or be leveraged in material detection systems.

Further development of the framework will progress toward understanding how a combination of both transient and spatial images can elevate machines to a higher level of perception \cite{moramartin2021,yang2023,yang2024}. Current systems primarily focus on ``what'' is perceived—identifying objects and materials present in a scene. Future research will advance toward understanding ``why'' and ``how'' the perceived entities exhibit certain behaviours \cite{gadzicki2020,yang2024}. This involves developing hybrid features that capture not only the presence and identity of materials and objects but also their functional properties, interaction mechanisms, and behavioural characteristics. For instance, research will investigate how temporal signals reveal material degradation, structural integrity, or dynamic processes, while spatial information provides context for understanding these temporal behaviours \cite{moramartin2021,lyons2019}. The integration of these modalities will enable machines to reason about material properties, predict material behaviour under different conditions, and understand causal relationships between material composition and observed phenomena. This progression from ``what'' to ``why'' and ``how'' represents a fundamental shift toward more sophisticated scene understanding, where machines can not only detect and classify materials but also understand their underlying physics, predict their behaviour, and reason about their functional properties \cite{yang2024,gadzicki2020}. This will require the development of new architectures that can jointly process spatial and temporal information, extract causal relationships, and perform higher-level reasoning about material properties and behaviours.




\section{Conclusion}
\label{subsec:conclusion_chapter_3}

The fundamental mechanism underlying all signal processing systems conforms to the wave--particle duality, a principle of quantum mechanics describing how light and matter exhibit both wave-like and particle-like properties \cite{bohr1923,elachi2021,griffiths2017}. In optical sensing and material detection, this duality manifests in how signals are measured and processed: the wave nature enables spectral analysis and interference phenomena, while the particle (photon) nature enables discrete, time-resolved measurements \cite{elachi2021}. Understanding this dual nature is essential for developing comprehensive signal processing frameworks, as different measurement approaches leverage different aspects of this duality to extract material information.

The distinction between intensity-based and temporal signal processing is fundamentally nominal rather than categorical. Both approaches measure the same underlying physical phenomena (the interaction of light with matter) but extract different aspects of this interaction \cite{su2016,becker2023}. Intensity-based methods capture the spatial distribution of light intensity, encoding information about material appearance, texture, and surface properties through spatial patterns. Temporal signal processing captures the time-resolved behaviour of light--matter interactions, encoding information about material composition, optical properties, and internal structure through temporal patterns \cite{su2016,becker2023}. While these approaches appear distinct, they measure complementary aspects of the same physical process, with the choice between them determined by which aspect of material information is most relevant for the specific application.

Each category exhibits strengths that satisfy specific purposes. Intensity-based processing excels in capturing spatial patterns, texture information, and appearance-based material traits, making it well-suited for applications where visual appearance or spatial distribution is paramount \cite{schwartz2013,sharan2013}. Temporal signal processing excels in capturing material composition, optical properties, and structural characteristics that manifest in time-resolved interactions, making it well-suited for applications where material identity or optical properties are critical \cite{su2016,becker2023,tafone2023}. The complementary nature of these approaches means that neither is universally superior; rather, their effectiveness depends on the specific requirements of the application, the nature of the materials being analysed, and the type of information needed for decision-making.

However, the spectral and non-spectral (achromatic) versions of each approach differ by ordinal levels of worst-case processing complexity. Spectral approaches, which analyse signals across multiple wavelengths or frequency bands, inherently require processing higher-dimensional data \cite{heinz2001,vasile2023}. Each wavelength or frequency band adds an additional dimension to the data space, exponentially increasing the computational requirements for feature extraction, dimensionality reduction, and classification \cite{knuth1976}. Spectral intensity images, for example, require processing three-dimensional data (spatial dimensions $\times$ spectral bands), while spectral temporal signals require processing time-resolved data across multiple wavelengths, creating four-dimensional data structures (time $\times$ wavelength $\times$ spatial dimensions) \cite{tontini2020,koerner2021}.

Spectral signal processing approaches exhibit significantly higher computational resource consumption compared to their achromatic counterparts \cite{su2016,becker2023}. This increased complexity arises from multiple factors: the need to process and store multidimensional data, the requirement for more sophisticated feature extraction algorithms to handle spectral correlations, the increased memory requirements for storing spectral data cubes, and the computational overhead of spectral analysis operations such as spectral unmixing, wavelength selection, and spectral feature extraction \cite{heinz2001,vasile2023}. In contrast, achromatic (non-spectral) approaches process single-channel or greyscale data, resulting in lower-dimensional data structures that require less computational resources for storage, processing, and analysis \cite{su2016,becker2023}. This computational advantage makes achromatic approaches more suitable for real-time applications, resource-constrained environments, and systems where processing speed is critical \cite{moramartin2021,lyons2019}.

The core essence of these models is to establish a theoretical and practical basis upon which each approach can be developed and improved for rigorous signal processing systems. This foundation encompasses understanding the fundamental physical principles underlying each approach, characterising their strengths and limitations, quantifying their computational requirements, and establishing frameworks for selecting the most appropriate approach for specific applications \cite{koerner2021}. By formalising these foundations, researchers and practitioners can make informed decisions about approach selection, optimise system design for specific requirements, and develop hybrid systems that leverage the complementary strengths of different approaches \cite{moramartin2021,yang2023,yang2024}. This foundational understanding also enables the development of standardised evaluation metrics, benchmarking protocols, and performance comparison frameworks that facilitate fair and meaningful assessment of different signal processing approaches across diverse applications and material types.

\newpage
\chapter{Flat Homogeneous Material Detection}
\label{chap:flat_homogeneous_material_detection}

\newpage
\section{Introduction}
\label{sec:flat_material_introduction}
\noindent Remote material detection presents fundamental challenges comparable to understanding matter itself, yet remains essential across diverse applications. Materials exhibit unique optical properties that serve as key identifiers for classification \cite{cleveland2009}. Indirect signal detection for inferential analysis constitutes a remote sensing task, though the term is more commonly associated with planetary sensing in geography \cite{cleveland2009}. This work presents experiments demonstrating automated remote detection of flat, homogeneous material surfaces using low-cost hardware and non-spectral signals.

A fundamental question arises: what makes material detection challenging for machines, and how does this work address such difficulties? The majority of challenges relate to material properties such as homogeneity and isotropy levels. However, this work focuses on challenges posed by instrumentation and signal processing, hence the selection of flat-surface homogeneous targets. To address these problems, single-photon avalanche diodes (SPADs) and achromatic illumination sources are employed.

The natural complexity of material structures and surface curvatures influences how they interact with energy. In material detection contexts, variations in material structures can be narrowed to homogeneity and isotropy. A structurally homogeneous material exhibits identical physical properties at every infinitesimal unit area. Consider a moulded clay object shattered into pieces: each fragment contains the same information to identify the material. This structural consistency found in every piece characterises homogeneity. Homogeneity manifests through a material's response to mechanical energy (stress or strain) and electromagnetic interaction metrics such as reflectivity. This project relies on the latter to build signature temporal profiles of material surfaces for classification and remote detection. While some naturally occurring materials exhibit structural consistency, synthetic materials also demonstrate structural homogeneity. This project sampled only synthetic homogeneous materials for experimentation.

Unfortunately, most material objects encountered in everyday life are structurally inhomogeneous, implying possible variations in reflectivity profiles from different surface areas. In object contexts, combinations of inhomogeneous materials with surface curvatures may result in further complexity. As the objective is reliable material detection with low-cost hardware, only solid flat-surface materials with homogeneous properties are sampled for experiments. Future work will consider inhomogeneous multi-curvature material detection for objects.

Beyond homogeneity lies isotropy, an even finer material property that subtly influences reflectivity. Material isotropy refers to the energy--matter behaviour of an infinitesimal unit area of material surface along its axes ($x$, $y$, and $z$). In this project context, the referred energy--matter behaviour is optical. Materials with high isotropic levels tend to behave consistently along each of the three axes. Unlike homogeneity, isotropy is rarely considered a valid basis for high-level material category distinction in experimental contexts. For instance, different wood types tend to exhibit different isotropy levels, yet all remain wooden materials at higher levels. Selected target materials for this experiment exhibit identical homogeneity and higher isotropy levels. Thus, the project reports no known significant material complexity problems. Future projects will address material complexity.

Instrumentation challenges relate to illumination source, signal measurement approach, and detection hardware factors. Most detection applications rely on illumination sources with multiple wavelengths \cite{su2016}. While this approach is widely adopted and efficient for most applications, it tends to be computationally more expensive compared to non-spectral counterparts. Thus, this work employs an achromatic illumination source for detection. Regarding signal measurement, two main approaches exist: intensity and temporal. Temporal measurements generate transient images superior to intensity images for depth estimation and other machine perception tasks \cite{su2016}. Transient images encode unique material albedos \cite{jungerman2022} and signal profile patterns that may serve as signature properties of the material. Automatic pattern recognition algorithms learn these patterns for decision-making.

Single-photon avalanche diodes (SPADs) are currently the leading technology for temporal or time-of-flight measurements. They have long been products of low-volume military and scientific research markets. However, their rapid evolution has enabled compact, miniaturized, and cost-effective versions \cite{vl53l5cx2022} that have enabled several applications in research domains, especially machine vision. SPAD evolution comes with hardware limitations: noise, coarse lateral resolution, and others. Thus, this project resolves such challenges with software approaches. Ambient noise, for instance, is estimated algorithmically from recorded timed photon intensity prior to illumination. Shot noise is addressed by estimating Poisson error per time bin.

During data collection, target materials are varied in range distance, orientation, viewing perspective, scale, and view angles to assess and test transient image limitations and invariances. Target materials can generate ambiguous return signals at different distances or view angles to the detector \cite{tontini2024}, which significantly challenges reliability of decision-making with such data. Two in-house algorithms are developed to address these challenges: multi-pixel transient stacking and depth-integrated single-pixel transients. The former processes all detector arrays per range cycle and has no awareness of target depth. The latter first selects the most significant pixel based on defined threshold and signal metrics, followed by normalization of the sequence with the centre of mass of the target in range.

To make material detection decisions, transients are processed by a lightweight neural network classifier. Impressive predictive accuracies were recorded in both approaches for classification. Validation with untrained surfaces of unseen classes (e.g., $z$-axis) resulted in high prediction accuracy ($\sim$90\%), indicating high levels of material homogeneity in targets. To ensure the classifier learns relevant features specific to seen classes, zero-shot learning is performed by feeding unseen classes to the model. Finally, two autoencoders (restricted Boltzmann machines (RBM) \cite{ackley1985} and convolutional autoencoders (CAE)) were used to learn latent spaces of trained classes. Both successfully reported very low losses between generator and real input signals, further proving that non-spectral transient images from target surfaces encode signature features of the material. All other assumptions of this experiment are also proven, which further supports the originality of the hypothesis as well as the quality and rigor of experimental methods developed.

From robotics applications through safety-critical machine perception systems, material detection is indispensable. This work aims to perform reliable remote detection of homogeneous material surfaces with non-spectral illumination sources and low-cost direct time-of-flight sensors. The problem of material complexity is mitigated by sampling only solid homogeneous materials with smooth surfaces. Challenges of remote detection systems extend beyond detector hardware, spanning variations in material structure, illumination source, signal measurement approaches, and detection hardware. The presented experiments show that even with homogeneous surfaces, non-spectral illumination, and low-cost compact detectors (SPADs), significant hardware limitations exist that are resolvable with software approaches.

Of all material structure dynamics, flat smooth surfaces are the least complex. Flat homogeneous materials exhibit relatively consistent reflectivity profiles across all unit areas. The lack of curvatures or irregular peaks and troughs allows incident surface normals of all reflected photons to be parallel \cite{su2016}, resulting in high similarity in all signals emerging from each point on the surface. This work fundamentally focuses on the viability of low-budget time-of-flight detectors for material detection. Thus, the choice of materials for experiments is limited to flat-surface homogeneous samples that keep material structural complexity low. In effect, the experiments are not challenged by material complexity factors.

Active illumination sources tend to be spectral or non-spectral. The former, often in spectroscopy, can be computationally costly \cite{su2016} for remote detection tasks. Thus, non-spectral illumination is sourced from a ~940nm class 1 laser source. Transient images are obtained from measuring photon arrival times. Even with the possibility of single-pixel inference, the quality of temporal data measurements can impact signal profiles and consequently affect inference reliability. Factors affecting measurement quality per pixel include noise sources, digitization, signal decay over longer distances, coarse TDC resolution, and target ambiguities (e.g., albedo--depth ambiguity, albedo--normal ambiguity, view angle (see Figure~\ref{fig:4_1}), and orientation ambiguity \cite{jungerman2022}).

%\vspace*{\fill}
%\vspace{-1cm}  % Position figure ~1cm from footer (fill pushes to bottom, -1cm positions it)
\begin{figure}[H]
    \centering
    \includegraphics[width=0.7\textwidth]{figures/4.1.png}
    \caption[View angle ambiguity in transient measurements.]{\textbf{View angle ambiguity in transient measurements.} Different view angles can generate similar transient profiles, creating challenges for material classification.}
    \label{fig:4_1}
\end{figure}
%\vspace{0.5cm}

Relevant noise sources in this work include ambient noise from detector background. Algorithmically, ambient noise is estimated as the average of initial bins of the transient histogram where target signal does not accumulate. This value accounts for a reliable estimate of background signal per bin. By subtracting across bins, the measured signal-to-background ratio (SBR) is improved. The nondeterministic nature of light particles also introduces noise. Due to the Poisson-distributed nature of incident photons, measured return photons in a given time range ($\Phi$) may differ from actual returned photons. (The term Poisson noise may also be referred to as shot noise.) Thus, photon arrival timestamps or histograms represent approximated versions of actual quantized returned signals. This holds true even under zero pile-up effect assumptions where every incident photon is captured by the SPAD. Sufficient levels of Poisson noise (e.g., multiple unrecorded laser cycles) may affect reliability of depth estimation in remote material detection. Thus, approximate Poisson mean value is estimated and subtracted across all histogram bins. The resulting sequence is fed to feature extraction algorithms for latent space extraction and decision-making. Other noise sources such as dark current or internally generated circuit noises are not major challenges for this work. High detector sensitivity and sampling rates enable recording of most laser cycles.

Nonetheless, returned signals are not measurable as full continuous waveforms. Thus, resulting transient histograms are discretised approximate representations of actual signals per time. Time-correlated single-photon counting (TCSPC) is the technique used by time-of-flight detectors in making such approximations. This project uses a ~940nm class 1 laser source, up to 5\,m range distance. The laser light could travel further than 5\,m; however, the limit is imposed partly by the time range of the histogram (TDC time range).

Often approximated as Gaussian, signal strength per bin is highest at the centre of the peak, which distributes normally outward. Coarse TDC resolution relates to the width of histogram bins. For example, a 250\,ps width would be coarser resolution than a 35\,ps width. Shorter pulse widths contribute to higher numbers of histogram bin events and denser TDC resolution, which provides good signal quality for target classification. The SPAD used for this work has 144 histogram bins.

Scene ambiguities are factors that may, under different circumstances, impair or create ambiguities in measured signals or transient images, leading to less reliable decision-making. Transient histograms encode multiple features: surface albedo, target depth, and surface normal; thus, they may not retain uniqueness across different surface configurations or lighting. Hence, the possibility exists of different target materials generating identical or extremely similar transient images. In \cite{jungerman2022}, these ambiguities are grouped into albedo--depth, albedo--normal, and orientation. Signal peaks are influenced by albedos and geometric flow of incident radiation (radiometric properties) on the surface. An albedo ambiguity could be caused by a diffuse surface located close to the detector, generating strong signal peaks that are either very similar or identical to transients from a highly reflective material surface at a longer distance. In other scenarios, tilted reflective surfaces may reflect less light than the same or other diffuse surfaces where the detector is at zero degrees to the target surface normal. The resulting transients are also invariant to the orientation of the surface. Thus, rotational symmetry does not affect the signal profile. In this work, ambiguities are resolved by two approaches: multi-pixel transient stacking and depth-integrated single-pixel transients. In single-pixel detection, the estimated depth (or distance) is integrated across the transient image. The resulting signal is a transformed signature profile that cannot be similar to any other surface. In multi-pixel detection, depth plays no role in encoding. Instead, pixels with different shifts in temporal spread are stacked into one sequence to represent the returned signal. Impressive results from signal classifiers validate the quality of signal processing and pattern recognition methods used.

Remote material detection is a form of remote sensing (RS) that focuses on indirect characterisation and acquisition of material composition data from target surfaces. The broad subject of material detection is strongly founded on the complex nature of matter itself, yet it also involves signal sources and instrumentation for detection. Historically, major related works in RS are centred around planetary sensing \cite{cleveland2009}. Nonetheless, there has been relevant remote material detection research, especially for non-spectral illumination tasks \cite{becker2023,su2016}. The target materials used in this work were carefully selected such that they would pose little to no challenges to signal processes. However, the ability to correctly sample such materials is only possible with a concise understanding of structural material properties and how they influence optical metrics during electromagnetic interactions. In this regard, material reflectance and other relevant structural properties are reviewed in terms of how they affect remote detection.

Unique surface profiles are based on surface properties that influence reflectance or albedo of returned signals (reflectivity). Details of how surface reflectance properties impact signal profiles can be found in the framework section (Chapter~\ref{chap:pawd}).

\subsection{Material structure and reflectivity}
\label{subsec:material_structure_reflectivity}

The fundamental structure of matter represents one of the most complex concepts in science. In engineering literature, material structures are typically characterised by their response to mechanical forces (stress and strain). However, material structure also profoundly influences electromagnetic interactions, which manifest through measurable optical properties such as reflectivity, absorptivity, scattering, and transmissivity \cite{elachi2021,griffiths2017}. This behavioural variation arises from the massless nature of photons, which cannot mechanically deform material surfaces. Instead, reflected photons can be temporally resolved to construct transient images representing temporal reflectivity tensors \cite{su2016,jungerman2022}. Understanding how material structural properties influence resulting transients is essential for selecting appropriate target materials for remote detection applications.

The most fundamental material property is its physical state: solid or fluid. Solid materials possess rigid structures that resist applied forces, while fluid materials (liquids, gases, plasmas) exhibit less rigid structures that flow or conform to container shapes. For ground truth signal acquisition, target materials must be confined to fixed positions where particles do not flow or move during the measurement cycle. The stability and immobility of solid particles make them the preferred choice for studying material properties in detection systems.

At the physical level, material composition may be uniform or non-uniform. Homogeneous materials exhibit identical structure at every infinitesimal volume throughout the entire object. As demonstrated in the target materials used in this work, every infinitesimal unit area of the surface shares the same composition as all other unit areas. This uniformity enables consistent ground truth signal collection and analysis from non-spectral radiation without requiring complex mixture models (Chapter~\ref{chap:pawd}). Material homogeneity remains invariant under all energy--matter interactions and detection conditions, including distance, viewing angle, perspective, and scale.

In solid homogeneous materials, electromagnetic radiation penetration varies significantly (see Figure~\ref{fig:4_2}); this property is commonly termed absorptivity \cite{elachi2021}.

\begin{figure}[H]
    \centering
    \includegraphics[width=0.7\textwidth]{figures/4.2.png}
    \caption[Material translucency: illustration of penetration levels of photons through matter.]{\textbf{Material translucency: illustration of penetration levels of photons through matter.} Top: transparent material that full transmission of all incident radiation through with no reflection or returned signal. Middle: translucent material with some absorption, scattering, transmission and reflection. Bottom: opaque materials with zero incident signal transmission and full reflectance.}
    \label{fig:4_2}
\end{figure}

Transparent materials allow all incident radiation to pass through, rendering them impractical for remote sensing applications where no measurable returned signal would be available for processing. Translucent materials permit partial transmission, with some incident signal scattered or reflected back, while opaque materials resist any signal penetration and produce full reflectance of incident radiation \cite{su2016,becker2023}. The materials selected for remote detection in this work are either opaque or translucent to ensure measurable reflectance for further analysis.

\subsection{Time of Flight (ToF) sensing}
\label{subsec:tof_sensing}

In remote material detection, material composition is inferred from recorded albedo. While several reliable approaches exist, including stereoscopy, state-of-the-art research reveals time-of-flight (ToF) sensing to be a leading technology for this task \cite{su2016,becker2023}. In ToF sensing, a target object is actively illuminated, and the returned signal is timed according to photon arrival and binned to form a histogram. Histogram values represent albedos at different time points and encode both the estimated depth of the target object and its material composition.

While measured albedo from the target encodes surface reflectivity and depth information, the unique material composition of the target surface may be inferred from the temporal signal distribution \cite{jungerman2022,su2016}. The most efficient and inexpensive approach to learning unique non-linear patterns from non-spectral transient images is through the use of time-of-flight (ToF) detectors, as formalized in the Particle Wave Detection (PAWD) framework (see Chapter~\ref{chap:pawd}). ToF sensors were previously limited to low-volume military and research projects for capturing photons at picosecond levels. Their evolution into compact, miniaturized, and low-cost versions \cite{vl53l5cx2022} has made them more accessible for tackling unsolved depth-related problems in machine perception.

In pulsed direct time-of-flight (dToF), the operating principle involves flooding a target or whole scene with an illumination source timed by a highly precise clock at the picosecond level (see Section~\ref{subsubsec:direct_tof}). An ideal pulse has key properties: laser power and pulse width. The amount of energy contained within a single pulse is characterised by laser power. Pulse width determines the break interval between each pulse. Pulse energy ($\Phi$) at initial time $t_0$ is zero due to delays (laser rise and fall times) in getting laser intensity to rise to expected energy and falling to zero before a new pulse is sent. The minimum time the laser stays on before reaching the desired intensity for the pulse is the pulse width.

\subsection{ToF modelling}
\label{subsec:tof_modelling}

An ideal model for each ToF pulse should account for the dynamics of laser rise and fall times, timing jitter, and noise sources such as Poisson noise. In practice, this is approximated as a Poisson mean-modified Gaussian function \cite{tontini2020,koerner2021} (see Chapter~\ref{chap:pawd}). The Gaussian approximation employed in this work results from the convolution of the temporal profile of the laser pulse with the instrument response function (IRF), which encompasses the combined effects of laser pulse characteristics and timing jitter from the time-to-digital converter (TDC) and SPAD components \cite{tontini2020}. The IRF is typically modelled as a Gaussian profile with a full width at half maximum (FWHM) on the order of nanoseconds, though non-Gaussian shapes can occur depending on system characteristics \cite{tontini2020}. In addition to Poisson noise, temporal uncertainties such as significant jitter can broaden the signal width, resulting in less reliable target depth estimation \cite{koerner2021}. To simulate more practical experimental settings, lighting conditions are varied across low and bright passive sources. Bright sources, including natural light and LED or tungsten illumination, emit infrared wavelengths similar to those used in SPAD systems, introducing baseline errors in the returned signal (see Chapter~\ref{chap:pawd}).

\subsection{Transient images}
\label{subsec:transient_images}

Transient imaging originated in the late 1970s with ``light-in-flight recording'' techniques \cite{abramson1978,abramson1983}, where illuminated scenes were recorded and reconstructed at different points in time. Subsequent work by Nilsson and Carlsson \cite{nilsson1998} applied similar approaches to measure object shape. Following recent evolution of time-of-flight (ToF) detectors, research has expanded transient imaging applications for both short and long-range distances, utilising spectral and non-spectral sources. This work investigates the viability and limitations of non-spectral transient waveforms for remote material detection and classification, particularly the impact of known ambiguities on classifier reliability.

Direct time-of-flight (dToF) sensors have emerged as efficient depth measurement techniques. In dToF systems, transient image acquisition begins with pulsed illumination directed toward the target scene, encompassing the target material and other artefacts within the foreground, background, or occluding the target. At a constant speed of light, the transient profile is used to infer photons' round-trip travel time:
\begin{equation}
T_{\text{rt}} = 2 \frac{d_{\text{target}}}{c}
\label{eq:round_trip_time}
\end{equation}
\addcontentsline{loe}{equations}{\protect\numberline{\theequation}$T_{\text{rt}} = 2 \frac{d_{\text{target}}}{c}$}

The resulting waveform (a one-dimensional temporal distribution derived from multiple photon cycles) is referred to as a transient image and encodes albedo, depth, and material surface composition. The temporal signal is discretised into short interval bins, forming transient histograms that represent the entire scene in range (see axial resolution in Chapter~\ref{chap:pawd}). These histograms encode albedo and depth information of perceived targets. To infer material composition, subtle variations in the returned signal profile are learned as input features for the classification model.

Depth information is crucial for remote material detection. Unlike other ToF applications \cite{tontini2024}, it is necessary to localise distinct targets in range; thus, the ToF information cannot simply be averaged in this work. Knowing the approximate or exact position of a target is essential for localising other objects behind (background) or in front (foreground) of it (Figure~\ref{fig:4_3}). Signals from all objects in the detector field of view are necessary for accurate perception of the scene.

\begin{figure}[H]
    \centering
    \includegraphics[width=0.7\textwidth]{figures/4.3.png}
    \caption[Intensity based transient image across multiple pixels of ToF sensor.]{\textbf{Intensity based transient image across multiple pixels of ToF sensor.} Annotation of transient image indicating variations in intensity for different signal peaks in range. Red annotations indicate non targets in foreground (objects perceived before target) and background (objects perceived after target).}
    \label{fig:4_3}
\end{figure}

In classifier input preprocessing, the region of interest (ROI) is cropped as a unit input to represent the targeted object in range. In \cite{su2016}, a classifier was validated on the grounds that it can classify both intensity and transient-based textures equally. However, this approach appears to contradict the subject of study in this context (reliable material detection) as well as the value proposition of ToF sensors. When capturing a spatially resolved image of wax material printed on paper, the perceivable material by the detector is paper rather than wax.

Often when multiple signals are returned, there is a single ROI (target). Targets may be preceded by foreground signals or followed by background signals (Figure~\ref{fig:4_3}). In instances where multiple distinct objects co-occur around the same points in time, the returned signal is a sum of all waveforms from each of the objects. In spatially resolved intensity images, this is termed occlusion. Fundamentally, signal occlusion in transient imaging differs significantly from intensity imaging. This is because the resulting transient of occluding objects is a composite transient that truly represents neither object (Figure~\ref{fig:4_4}).

\begin{figure}[H]
    \centering
    \includegraphics[width=0.7\textwidth]{figures/4.4.png}
    \caption[Illustration of target-background overlap (occlusion) in transient image.]{\textbf{Illustration of target-background overlap (occlusion) in transient image.} The isolated target object is assumed to return signal off the surface with no interference or mix up with other objects. On the contrary (often seen in background signal), signal from closely positioned multiple objects tend to partially or fully transform into a smear like signal spanning more bins than a single target should.}
    \label{fig:4_4}
\end{figure}

The core aim of this work does not focus on signals from occluded objects. Thus, input data is only sourced from objects at controlled target positions that return distinct ground-truth signals from the material of interest (MOI) only.

\subsection{Noise}
\label{subsec:noise}

In addition to the fundamental Poisson or shot noise inherent in photon counting, ambient light from passive sources (natural lighting or white light illumination) within the measurement environment establishes a baseline offset in transient waveforms \cite{tontini2020}. Unlike Poisson noise, which exhibits statistical variation, ambient levels are approximately uniform across the axial plane. This uniformity enables subtraction of the average ambient level across all time bins to recover the clean returned signal from active illumination.

However, variations in material reflectivity can introduce systematic errors in ambient estimation. Averaging time bins may overestimate ambient levels when highly reflective objects are present, while the converse occurs with low-reflectivity materials. The SPAD dead time typically exhibits no detectable pulse from the active illumination source (Figure~\ref{fig:4_5}), due to the non-zero wait time required for the SPAD to emit a signal of sufficient power and the finite pulse width.

\begin{figure}[H]
    \centering
    \includegraphics[width=0.56\textwidth]{figures/4.5.png}
    \caption[Illustration of SPAD event timeline (pulse start and stop for both reference and returned signal).]{\textbf{Illustration of SPAD event timeline (pulse start and stop for both reference and returned signal).} The histogram is constructed by time-stamping the leading edge of SPAD pulses (ambient and signal), and accumulating these timestamps over multiple timing/laser cycles. There are certain periods of time over which no signal photons are received.}
    \label{fig:4_5}
\end{figure}

Nonetheless, passive ambient signals remain detectable during this period, resulting in non-zero bin values across the SPAD dead time interval. Reflectivity values from bins spanning this dead time are averaged to represent the overall ambient noise level, which is subsequently subtracted from all other bins to obtain a cleaner waveform \cite{tontini2024,tontini2020}.

\subsection{Ambiguities}
\label{subsec:ambiguities}

Target orientation relative to the detector remains invariant to measured albedo (Figure~\ref{fig:4_6}) when view angle and distance are maintained.

\begin{figure}[H]
    \centering
    \includegraphics[width=0.7\textwidth]{figures/4.6.png}
    \caption[Orientation invariance: illustration of lat target surface orientation invariance to returned temporal signal from detector.]{\textbf{Orientation invariance: illustration of lat target surface orientation invariance to returned temporal signal from detector.}}
    \label{fig:4_6}
\end{figure}

Although \cite{jungerman2022} characterises this property as an ambiguity, experimental results in this work demonstrate that it may represent a key solution to the unsolved machine vision problem of perspective invariance.

Conversely, other ambiguities such as distance--albedo variabilities (Figure~\ref{fig:4_7}) pose significant threats to signal profile uniqueness \cite{tontini2020} due to the inverse square law governing electromagnetic radiation propagation \cite{griffiths2017}. According to the inverse square law, the intensity of radiation decreases proportionally to the square of the distance from the source, meaning that measured reflectivity scales inversely with the square of the target distance. This relationship likely explains why range distances in other non-spectral studies \cite{becker2023,jungerman2022} are limited to short distances.

\begin{figure}[H]
    \centering
    \includegraphics[width=0.7\textwidth]{figures/4.7.png}
    \caption[Distance–albedo ambiguity in transient measurements.]{\textbf{Distance–albedo ambiguity in transient measurements.} The inverse square law \cite{griffiths2017} dictates that signal intensity decreases proportionally to the square of the distance, creating a relationship where measured reflectivity scales inversely with the square of target distance. This inverse square relationship between target distance and measured reflectivity creates variability that threatens signal profile uniqueness.}
    \label{fig:4_7}
\end{figure}

This work employs algorithmic processing techniques that incorporate depth information into transient image measurements to resolve both distance and view angle ambiguities (Figure~\ref{fig:4_1}). This approach enables retention of distinct temporal sequences in both short- and long-range detection with achromatic signals. While depth is not a direct variable of the transient profile, it serves other critical functions: target localisation, and foreground and background signal spread across pixel arrays. With a custom automatic target localiser, regions of interest (ROIs) are identified from the centre of mass (CoM). The signal delay experienced upon interaction with the target surface constitutes the scene response function (SRF), which accounts for additional signal propagation factors such as interference or other non-target influences on the signal. The scene response also enables automatic, unsupervised target depth extraction, making the application feasible for real-time safety-critical systems. This work relies on a custom clustering algorithm that extracts target information with reference to its foreground and background.

\subsection{Classifier models}
\label{subsec:classifier_models_chapter_4}

Preprocessed transient signals encode material albedos, depth information, and residual noise components. While pixel stacking and depth-integrated transient processing address signal ambiguities \cite{jungerman2022,tontini2024}, non-deterministic factors such as residual Poisson noise may hinder the linearity of learnable patterns, necessitating non-linear learning algorithms. Deep neural networks provide the required non-linear mapping capabilities for extracting material-discriminative features from these complex temporal signatures \cite{lecun1998,su2016}.

Material surface profiles exhibit remarkable consistency across homogeneous surfaces. Variations in the approximate impulse response function (IRF) curve per material tend to be significantly low as surface texture decreases, even for different colored surfaces \cite{becker2023}. This consistency enables automatic pattern recognition approaches to learn robust material signatures from temporal signals.

Two distinct classifier architectures are employed in this work, corresponding to the two preprocessing strategies. For stacked pixel inference, a two-dimensional convolutional neural network (2D CNN) is designed to extract key reflectivity features from spatially resolved transient arrays, enabling the network to learn material-specific patterns across multiple detector pixels simultaneously (Section~\ref{sec:multizone_signal_classifier}). Depth-integrated single transients are processed by a one-dimensional convolutional network (1D CNN) that operates on temporal sequences extracted from individual pixels (Section~\ref{sec:single_zone_signal_classifier}). Both architectures leverage the temporal structure of transient signals to learn material-discriminative features \cite{su2016,becker2023}.

To validate the uniqueness and generalisability of learned features, zero-shot learning is performed by evaluating model performance on untrained perspective transients (see Section~\ref{subsec:multizone_experiment_results} for details). The high prediction accuracy recorded on untrained perspective transients confirms the perspective invariance of homogeneous material surface composition, demonstrating that the classifier learns material-intrinsic features rather than view-dependent artefacts. Evaluation on untrained surfaces further demonstrates the robustness of the classifier in recognising features from learned material classes, indicating successful extraction of material-discriminative temporal signatures \cite{jungerman2022}.

\newpage
\section{Assumptions}
\label{sec:assumptions}

The following assumptions guide the experimental design and analysis presented in this work. These assumptions formalize expectations about how electromagnetic signals interact with material surfaces and how temporal signal processing can enable material classification.

\begin{itemize}
    \item Flat homogeneous surfaces of varying surface reflectivity may respond differently to specific electromagnetic spectral bands.
    
    \item Flat homogeneous surfaces of varying colors may respond differently to specific electromagnetic spectral bands.
    
    \item Non-spectral signal processing may be more computationally efficient than spectral signal processing.
    
    \item Temporal signal profiles may encode features that can be used to identify material composition irrespective of detector angle of view, orientation, distance to target, view axis, and/or perspective. These profiles also remain robust to variations in target surface scale, shape, and/or form.
    
    \item Classification between a given set of materials may be reliably performed with time-of-flight (ToF) detectors. This capability extends to single-pixel configurations, where temporal histograms from individual detector elements contain sufficient information to discriminate between material classes.
\end{itemize}

\newpage
\section{Experiments' overview}
\label{sec:experiments_overview}

\noindent Two experiments were conducted using a single-photon avalanche diode (SPAD) direct time-of-flight (dToF) detector and a monochromatic, non-spectral illumination source. The experiments differ by the dimensionality of the inference data: (i) single-pixel classification from one-dimensional transient histograms, and (ii) multi-pixel classification from transient images. In both cases, the objective is to learn robust, material-dependent temporal signatures of target surfaces that can be processed by neural network classifiers, extending prior work on time-resolved material sensing with ToF systems \cite{su2016,tafone2023,callenberg2021}.

Transient image data were acquired from four planar, homogeneous materials: a black nylon sheet (BNT, Figure~\ref{fig:4_9}), a white nylon sheet (WNT, Figure~\ref{fig:4_8}), a white foam board (WGF, Figure~\ref{fig:4_11}), and a black cardboard sheet (BCB, Figure~\ref{fig:4_10}).

\begin{figure}[H]
\centering
\includegraphics[width=0.7\textwidth]{figures/4.8.png}

    \caption[White nylon sheet (WNT) target used in the SPAD experiments.]{\textbf{White nylon sheet (WNT) target used in the SPAD experiments.} Example capture showing the planar, homogeneous surface of the white nylon target under controlled illumination.}
\label{fig:4_8}
\end{figure}

Experimental factors considered for material detection included target orientation, view angle relative to the surface normal, sensor perspective, illumination condition, and sensor–target distance. Four discrete orientations were used: \texttt{pu}, \texttt{pd}, \texttt{pr}, and \texttt{pl}, where \texttt{p} denotes portrait, and \texttt{u}, \texttt{d}, \texttt{r}, and \texttt{l} indicate upright, downward, rightward, and leftward rotations, respectively. Three view angles were selected with respect to the target surface normal: \texttt{0dtsn}, \texttt{+15dtsn}, and \texttt{-15dtsn} (\texttt{d}: degrees; \texttt{tsn}: target surface normal). Materials were observed from three perspectives corresponding to the sensor translation axes: $x$-, $y$-, and $z$-axis viewpoints. Illumination conditions comprised bright scenes (with both passive and active components, including natural daylight) and dark scenes, with an equal number of observations acquired in each regime to balance the dataset.

\begin{figure}[H]
\centering
\includegraphics[width=0.7\textwidth]{figures/4.9.png}

    \caption[Black nylon sheet (BNT) target used in the SPAD experiments.]{\textbf{Black nylon sheet (BNT) target used in the SPAD experiments.} Example capture of the black nylon target, highlighting its low-albedo, weak-return response relative to the white nylon sheet.}
\label{fig:4_9}
\end{figure}


Ranging was performed over a discrete set of sensor–target distances from 20\,cm to 400\,cm in 20\,cm increments: 20\,cm, 40\,cm, 60\,cm, 80\,cm, 100\,cm, 120\,cm, 140\,cm, 160\,cm, 180\,cm, 200\,cm, 220\,cm, 240\,cm, 260\,cm, 280\,cm, 300\,cm, 320\,cm, 340\,cm, 360\,cm, 380\,cm, and 400\,cm. At each distance, approximately 10 observations were acquired, and the corresponding transient histograms were persisted for subsequent analysis. To improve signal-to-noise ratio at longer ranges, frames were temporally averaged, with up to eight independent transients combined for distances approaching 4\,m. On average, one second of acquisition time was required to capture a single transient histogram at a given range. For each range cycle, a corresponding RGB intensity image was also recorded to verify that the measured transients were correctly associated with the intended material target and to support cross-modal analysis of spatial and temporal cues.

\begin{figure}[H]
\centering
\includegraphics[width=0.7\textwidth]{figures/4.10.png}

    \caption[Black cardboard (BCB) target used in the SPAD experiments.]{\textbf{Black cardboard (BCB) target used in the SPAD experiments.} Example capture of the black cardboard surface, which serves as a low-reflectivity, weak-scattering reference material.}
\label{fig:4_10}
\end{figure}


\begin{figure}[H]
\centering
\includegraphics[width=0.7\textwidth]{figures/4.11.png}

    \caption[White foam (WGF) target used in the SPAD experiments.]{\textbf{White foam (WGF) target used in the SPAD experiments.} Example capture of the white foam material, whose porous, highly scattering structure induces broadened transient responses relative to the nylon sheets.}
\label{fig:4_11}
\end{figure}

\subsection{Instruments and setup}
\label{subsec:instruments_setup}

The experimental setup employs a custom SPAD sensor (VL53L8, MZ-Plus) operating in 4×4 zone mode, providing coarse lateral resolution with a non-spectral illumination source \cite{vl53l5cx2022}. The system features active illumination using a ~940nm class 1 laser source with 2\,ns pulse width, suitable for opaque surfaces perpendicular to the detector with negligible SPAD jitter and noise \cite{tontini2020}. The SPAD detector exhibits a 25° field of view (Figure~\ref{fig:4_12}), while the illumination source has a 35° field of view, ensuring comprehensive coverage of the target region. A summary of the experimental parameters, including materials, orientations, view angles, illumination conditions, and distance ranges, is provided in Table~\ref{tab:4_1}.

\begin{table}[H]
\centering
\adjustbox{max width=\textwidth,center}{
\begin{tabular}{p{0.28\textwidth}p{0.6\textwidth}}
\hline
\textbf{Factor} & \textbf{Values} \\
\hline
Materials & Black nylon sheet (BNT), white nylon sheet (WNT), white foam (WGF), black cardboard (BCB) \\
Orientations & \texttt{pu}, \texttt{pd}, \texttt{pr}, \texttt{pl} (portrait upright, downward, rightward, leftward) \\
View angles & \texttt{0dtsn}, \texttt{+15dtsn}, \texttt{-15dtsn} relative to target surface normal \\
Illumination & Bright (passive + active, including natural daylight) and dark scenes, balanced observations \\
Distances & 20\,cm to 400\,cm in 20\,cm steps (20 measurement levels) \\
Observations per distance & $\approx 10$ transient captures per material and configuration \\
SNR enhancement & Temporal averaging of up to 8 frames at longer distances ($\approx 4$\,m) \\
Auxiliary data & Co-registered RGB images per range cycle for verification and cross-modal analysis \\
\hline
\end{tabular}
}
\caption[Summary of experimental parameters for SPAD-based material sensing.]{\textbf{Summary of experimental parameters for SPAD-based material sensing.} Factors controlled during data acquisition for single-pixel and multi-pixel classification experiments.}
\label{tab:4_1}
\end{table}

%\vspace*{\fill}
%\vspace{-1cm}  % Position figure ~1cm from footer (fill pushes to bottom, -1cm positions it)
\begin{figure}[H]
    \centering
    \includegraphics[width=0.7\textwidth]{figures/4.12.png}

    \caption[SPAD detector field of view configuration.]{\textbf{SPAD detector field of view configuration.} The 25° field of view of the SPAD detector, shown in relation to the 35° illumination source field of view, ensuring adequate spatial coverage for material detection experiments.}
    \label{fig:4_12}
\end{figure}

A Sony RGB camera is co-mounted with the SPAD detector on a Raspberry Pi platform, providing greater than 80\% field of view intersection with the SPAD detector. The RGB camera exhibits a 35° field of view, enabling synchronized multi-modal data acquisition. Remote access to the sensor serial interface is facilitated through an external Wi-Fi adapter, while a four-direction remote moving robot enables automated positioning and data collection across varying sensor--target configurations.

Both SPAD and RGB sensors operate at 30\,fps with 2\,ms exposure time, enabling synchronized temporal acquisition of transient histograms and intensity images \cite{tontini2020}. The measured signal-to-background ratio (SBR) is approximately 0.05, reflecting the challenging low-light conditions and the need for sophisticated signal processing to extract material-discriminative features from noisy temporal signatures \cite{tontini2024,su2016}. This low SBR necessitates algorithmic approaches for ambient noise estimation and signal enhancement, as detailed in subsequent preprocessing sections.

\subsection{Data collection}
\label{subsec:data_collection}

Accurate feature extraction for surface material signatures requires ground truth data from known materials. To acquire this, a manual ranging experiment was conducted in a realistic natural scene setting comprising foreground and background artefacts. The experimental setup employed a SPAD module \cite{vl53l5cx2022} with an integrated illumination source. The module consists of a 2D array with integrated Time-Correlated Single-Photon Counting (TCSPC) and Time-to-Digital Converter (TDC) capabilities for generating timestamped photon counts \cite{tontini2020}. The system outputs 144-bin photon arrival time histograms, providing temporal resolution suitable for material discrimination \cite{su2016,jungerman2022}.

Non-spectral signals were actively sourced from target surfaces under several controlled conditions: scale, orientation, lighting, distance, and perspective. This systematic variation enables assessment of material signature invariance across different acquisition parameters, which is critical for robust material classification \cite{becker2023,su2016}. Data acquisition was performed via the module's serial interface, reading hexadecimal-format data at approximately 19\,MHz through an I2C interface connected to a Raspberry Pi host. The collected raw data was subsequently processed through a preprocessing pipeline to generate clean feature representations of target surfaces, as detailed in Section~\ref{subsec:data_preprocessing}.

\subsection{Data preprocessing}
\label{subsec:data_preprocessing}

The preprocessing pipeline generates four categories of distribution outputs: full-range signals with and without reference illumination, and cropped-range signals with and without reference illumination (Figure~\ref{fig:4_13}).

%\vspace*{\fill}
%\vspace{-1cm}  % Position figure ~1cm from footer (fill pushes to bottom, -1cm positions it)
\begin{figure}[H]
    \centering
    \includegraphics[width=0.7\textwidth]{figures/4.13.png}

    \caption[Signal peak extraction.]{\textbf{Signal peak extraction.} Top: multi pixel stacking approach considers multiple pixels from sensor incorporating slightly different views of same target in a 2D tensor without reference to depth or distance information. Bottom: depth integrated single pixel approach transforms transient sequence with target depth information and normalised to retain only relevant frequency information.}
    \label{fig:4_13}
\end{figure}

The core processing stages comprise signal transformations, noise removal, depth extraction, and transient stacking \cite{tontini2024,su2016}. Custom in-house algorithms were developed to handle ambient photon contributions in raw range data \cite{tontini2024}. To remove the background or ambient baseline, the baseline is estimated from bin values and subsequently subtracted across all bins. Baseline correction values are estimated independently for each frame, as ambient levels may vary during data acquisition \cite{tontini2024,tontini2020}. This per-frame estimation ensures robust handling of dynamic illumination conditions that could otherwise introduce systematic errors in signal recovery.

Algorithms for automatic target depth extraction, including foreground and background separation, are implemented. An edge detector traverses each histogram to identify sharp rises above a defined threshold. Upon detection, the edge is evaluated to distinguish between noise (non-target) and target signal components. Once a target signal is detected, it is parsed to a centre-of-mass (CoM) extractor that returns the bin identifier corresponding to the signal's centre of mass \cite{erdogan2017,zhang2019}. This value is used to extract the signal peak at full-width at half-maximum (FWHM), typically spanning 8 bins. An additional four bins are padded to the left and right sides of the extracted peak to introduce low-level signal augmentation. This padding is essential for feature extraction and pattern recognition models to generalise effectively and avoid overfitting to narrow signal windows \cite{su2016,becker2023}.

\subsection{Learning algorithm}
\label{subsec:learning_algorithm}

To uncover underlying patterns in material composition data, automatic pattern recognition algorithms are explored. Similar research has demonstrated that classical machine learning approaches can achieve impressive results for transient image classification \cite{su2016,becker2023}. However, for non-spectral signals acquired over long distance ranges with multiple orientations and multi-angle views, complex ambiguities and non-deterministic noise often persist even after Poisson and ambient denoising. Consequently, handcrafted features for classical physics models and machine learning may be less efficient for achieving optimal performance.

Multilayer neural networks exploit the assumption that most signal compositions exhibit hierarchical structure, enabling inference of signature material profiles from transients \cite{jungerman2022}. In intensity imaging, edges combine to form features, which further combine to form objects. Since intensity image resolution is typically spatial, the compositional hierarchies of transient images are more analogous to sound processing. Just as sounds combine to form phones, syllables, words, and sentences, returned photons per time bin for a particular target form histogram bins with unique spectral response functions (SRF). Signal peaks are filtered through convolution with kernels, while pooling operations downsample the tensor before activation. These steps are repeated for multiple refinement stages. The final steps unite all refined features into a fully connected layer where inference is made of the output class \cite{lecun1998}.

\newpage
\section{Multi-zone signal classifier (2D CNN)}
\label{sec:multizone_signal_classifier}

The multizone input to the classifier is a stack of $4 \times 4$ transient histograms, where each of the 16 zones contributes a one-dimensional cropped transient. These 16 one-dimensional sequences are arranged into a $16 \times 16$ pseudo-image, forming a low-resolution spatial--temporal representation of the scene. Each pixel in this $16 \times 16$ array represents a normalised, cropped temporal peak from a detector zone, so the input is already highly feature-rich despite its coarse spatial resolution.

From a convolutional neural network (CNN) standpoint, this configuration is extremely lightweight compared with typical image-classification settings. State-of-the-art CNNs for natural images often operate on inputs of $224 \times 224 \times 3$ or larger, while even early architectures such as LeNet used $28 \times 28$ inputs \cite{lecun1998}. In contrast, the $16 \times 16 \times 3$ inputs used here are approximately an order of magnitude smaller in spatial extent. Recent work has also demonstrated that ultra-low-resolution time-of-flight (ToF) sensors can support meaningful multiclass recognition tasks with compact CNNs \cite{fasolino2024}, and that ToF histograms carry material-discriminative information even at coarse spatial resolution \cite{su2016,becker2023}. This motivates a deliberately compact architecture that matches the small input size and embedded implementation constraints, while retaining sufficient capacity to separate subtle material differences.

The proposed network is a feedforward 2D CNN designed for four-class material classification of planar surfaces. The input tensor has shape $(N_b,3,16,16)$, where $N_b$ is the batch size and the three channels correspond to three stacked transient-derived modalities per zone. The model outputs logits for four material classes: white glass fiber (WGF), black cardboard (BCB), black non-transparent nylon (BNT), and white non-transparent nylon (WNT).

\subsection{Guiding principles and models}
\label{subsec:guiding_principles_models}

The experimental configuration employs a ~940nm class 1 laser source with achromatic processing, guided by the processing models of non-spectral temporal signals established in Chapter~\ref{chap:pawd} \cite{tontini2020}. Given the homogeneous composition of the target materials, the temporal signal adheres to the per-bin Gaussian approximation model described by Equation~\eqref{eq:pure_material_gaussian} \cite{tontini2020,koerner2021,su2016}. This single-peak Gaussian formulation provides an accurate per-bin representation of photon arrival times for uniform material surfaces, where the temporal distribution of reflected photons follows a normal distribution centred at the time-of-flight $t_0$ with standard deviation $\sigma$ \cite{becker2023}.

\subsection{Multizone experiment results}
\label{subsec:multizone_experiment_results}

The multizone transient processing pipeline successfully removes baseline offsets and suppresses Poisson shot noise, producing denoised histograms that serve as inputs to the deep neural network material classifier \cite{tontini2020}. Using stacked transient images as input, the multi-zone model attains high accuracy across both short- and long-range standoffs for single-distance and multi-distance configurations. When all sample variations are combined (multiple distances, orientations, view angles, and perspectives), the classifier achieves an overall accuracy exceeding 99\% (Figure~\ref{fig:4_14}), with robust performance across all four target material classes \cite{su2016,becker2023}.

%\vspace*{\fill}
%\vspace{-1cm}  % Position figure ~1cm from footer (fill pushes to bottom, -1cm positions it)
\begin{figure}[H]
    \centering
    \includegraphics[width=0.8\textwidth]{figures/4.14.png}
    \caption[Confusion matrix for CNN of multi-stack transient image input for four material classes.]{\textbf{Confusion matrix for CNN of multi-stack transient image input for four material classes.} Confusion matrix for CNN of multi-stack transient image input for four material classes (WNT, BNT, WGF and BCB) with over 99\% accuracy.}
    \label{fig:4_14}
\end{figure}

Validation on held-out samples from the trained classes yields approximately 90\% accuracy. Inspection of the corresponding confusion matrices shows that residual errors are dominated by misclassifications between materials with similar appearance, such as black nylon sheet versus black cardboard, indicating that the model is most challenged when both colour and transient signatures are closely aligned. To probe how well the network has internalized the material-specific transient structure, additional experiments were conducted using samples drawn from untrained background classes. For these out-of-distribution inputs, the desired behaviour is uniformly low confidence across all trained classes. In around 80\% of such cases, the classifier outputs low logits for every class (Figure~\ref{fig:4_15}), whereas the remaining 20\% of samples exhibit at least one non-negligible class probability (greater than 0.05), suggesting that some background transients partially overlap the learned decision boundaries.

%\vspace*{\fill}
%\vspace{-1cm}  % Position figure ~1cm from footer (fill pushes to bottom, -1cm positions it)
\begin{figure}[H]
    \centering
    \includegraphics[width=0.8\textwidth]{figures/4.15.png}
    \caption[Confusion matrix for never seen transient images.]{\textbf{Confusion matrix for never seen transient images.} Included background samples from either of the four material classes (WNT, BNT, WGF and BCB) and untrained materials (UNKNOWN) yielding ~19\% accuracy out of expected 0\%. Indicative of low confidence levels in over 80\% of predictions.}
    \label{fig:4_15}
\end{figure}

To further assess how learnable the transient features are, two generative autoencoder models are trained to capture the underlying latent space of the material classes and to encode transients into a compact representation that can be resampled or decoded back to the original data domain. These models provide an independent check on the structure and separability of the observed transients, complementing the discriminative CNN analysis; full architectural and training details are presented in Section~\ref{sec:generative_modelling}.


\subsection{Multizone model performance justification}
\label{subsec:multizone_model_performance_justification}

\noindent The substantial size of the convolutional neural network (CNN) and the strength of the reported multizone classification results naturally raise concerns about potential overfitting. However, the experimental design and evaluation protocol were explicitly structured to promote generalisation rather than memorisation. Training data were acquired across a wide range of standoff distances, viewing angles and object orientations, ensuring that each material class was represented under multiple geometries and signal-to-noise regimes, consistent with best practice in dToF-based material classification \cite{becker2023,su2016}. In addition to this diversity in acquisition, model selection relied on held-out validation sets, and zero-shot tests were performed on newly collected measurements gathered in distinct environments with backgrounds and illumination conditions that differed from those used during training. As illustrated in Figure~\ref{fig:4_15}, the CNN maintains low confidence on most out-of-distribution samples while preserving stable class rankings for in-distribution inputs, indicating that its decision boundaries extend beyond the specific training examples rather than exploiting spurious correlations. Collectively, these elements justify the reliability of the reported multizone model performance and argue against the hypothesis that it is an artefact of overfitting.

The achieved performance (over 99\% accuracy for in-distribution samples and approximately 90\% for validation on held-out samples) is consistent with previous relevant work in material classification using time-of-flight measurements. The Princeton work by Su et al.~\cite{su2016} demonstrated that machine learning models can achieve high accuracy (80.9\% validation accuracy using RBF SVM) for material classification using raw time-of-flight measurements, validating the feasibility of learning material-discriminative features from temporal signal characteristics. Similarly, the MDPI Sensors paper by Becker and Koerner~\cite{becker2023} reported classification performance of 96\% accuracy for plastic material classification using optical parameter features measured with direct time-of-flight depth sensors. The consistency of high performance across these independent studies, employing different machine learning approaches and material classification tasks, provides strong evidence that the observed accuracy in this work stems from the inherent material-discriminative information encoded in transient time-of-flight signals rather than dataset-specific artefacts or overfitting.

\section{Extra details of model accuracy}
\label{sec:extra_details_of_model_accuracy}

The model architecture is designed for the unique characteristics of transient-based material classification. The input dimensionality (16×16×3) is substantially smaller than standard CNN inputs: state-of-the-art architectures (VGG, ResNet, ImageNet models) typically process 224×224×3 images (196× larger), while custom CNNs commonly use 128×128×3 inputs (64× larger). Even LeNet, the earliest CNN architecture, operated on 28×28 inputs (3.06× larger). This extreme input compactness, combined with the fact that transients represent pre-processed signal peaks with inherent feature structure, motivates a correspondingly lightweight architecture with reduced convolutional complexity relative to standard image classification networks.

The network implements a lightweight convolutional architecture for multi-class material classification of flat surfaces, processing 16×16×3 dToF transient arrays and outputting logits for four material classes: WGF (White Glass Fiber), BCB (Black Cardboard), BNT (Black Non-Transparent), and WNT (White Non-Transparent). The architecture contains 1,061,284 trainable parameters, representing a compact design compared to standard deep learning architectures.

The architecture employs a feedforward design with two stages: (a) feature extraction via convolutional blocks performing hierarchical spatial feature learning, dimension reduction through max pooling, and feature map vectorization; and (b) classification via fully connected layers performing high-dimensional feature integration, dimensionality reduction, and multi-class logit generation.

\subsection{Layer-by-Layer Architecture}

\textbf{Layer 1: Initial Feature Detection}
The first convolutional layer \texttt{Conv2d(3→32, kernel=3×3, stride=1, padding=1)} applies 32 learnable 3×3 filters to the RGB input, preserving spatial dimensions (16×16) and extracting low-level features (edges, colour gradients, local texture patterns). This layer contains 928 parameters (3×32×3×3 weights + 32 biases), producing output shape (batch, 32, 16, 16). ReLU activation $f(x) = \max(0, x)$ introduces non-linearity and sparsifies activations.

\textbf{Layer 2: Feature Refinement}
The second convolutional layer \texttt{Conv2d(32→32, kernel=3×3, stride=1, padding=1)} refines and combines features from the first layer, containing 9,248 parameters (32×32×3×3 weights + 32 biases) and producing output shape (batch, 32, 16, 16). Following ReLU activation, \texttt{MaxPool2d(kernel=2×2, stride=2)} performs spatial downsampling, reducing dimensions from 16×16 to 8×8 and providing translation invariance. A flatten operation reshapes the tensor (batch, 32, 8, 8) to (batch, 2048), converting spatial feature maps to a 1D vector for fully connected processing.

\textbf{FC Layer 1: High-Dimensional Feature Integration (2048 → 512)}
The first fully connected layer \texttt{Linear(2048 → 512)} contains 1,049,088 parameters (98.9\% of total), integrating spatial features into a global material representation with output shape (batch, 512). Following ReLU activation, \texttt{Dropout(p=0.5)} randomly zeros 50\% of neurons during training to prevent overfitting; dropout is disabled during inference.

\textbf{FC Layer 2: Classification Logits (512 → 4)}
The final layer \texttt{Linear(512 → 4)} contains 2,052 parameters and outputs unnormalised class logits. Softmax is applied during inference to produce class probabilities, yielding output shape (batch, 4).

\subsection{Architectural Design Principles}

\textbf{Hierarchical Feature Learning:} The architecture employs progressive abstraction from edges to textures to material signatures through two-stage convolution, enabling increasing receptive fields via stacked convolutions to capture both local and global patterns.

\textbf{Regularisation Strategy:} Dropout ($p=0.5$) in the fully connected layers balances model capacity and generalisation, preventing memorisation while promoting robust feature representations.

\textbf{Spatial Processing:} With 16-zone transients stacked into a 2D space-resolved matrix, the architecture uses 3×3 kernels for efficient local pattern capture. Padding=1 preserves spatial dimensions, while max pooling provides translation invariance and computational efficiency.

\textbf{Parameter Efficiency:} The architecture contains 1,061,284 parameters, with 98.9\% concentrated in the first fully connected layer where high-dimensional feature integration occurs. Weight sharing in convolutions ensures efficient parameter utilisation.

\subsection{Information Flow Summary}

The information flow through the network proceeds as follows:
\begin{itemize}
    \item Input: (batch, 3, 16, 16) normalised RGB transient image
    \item Conv2d: 3→32, 3×3, padding=1 → (batch, 32, 16, 16)
    \item ReLU activation
    \item Conv2d: 32→32, 3×3, padding=1 → (batch, 32, 16, 16)
    \item ReLU activation
    \item MaxPool2d: 2×2, stride=2 → (batch, 32, 8, 8)
    \item Flatten → (batch, 2048)
    \item Linear: 2048→512 → (batch, 512)
    \item ReLU activation
    \item Dropout: p=0.5 (training only)
    \item Linear: 512→4 → (batch, 4)
    \item Output: Logits (batch, 4)
\end{itemize}

\subsection{Training Configuration and Overfitting Prevention}

The model employs several mechanisms to prevent overfitting:

\textbf{Weight Initialization:} Kaiming uniform initialization (PyTorch default for ReLU) ensures proper gradient flow and prevents vanishing gradients, contributing to stable training.

\textbf{Dropout Regularization:} Dropout ($p=0.5$) in the fully connected layer forces robust, redundant feature representations, particularly effective given the large parameter count (1,049,088) in the first fully connected layer where overfitting risk is highest.

\textbf{Loss Function and Optimisation:} Cross-entropy loss with SGD (learning rate=0.01, momentum=0.6) and batch size 130 provides stable gradient estimates. The moderate learning rate prevents aggressive weight updates, while momentum aids convergence.

\textbf{Input Normalization:} Normalization with mean=(0.5, 0.5, 0.5) and std=(0.5, 0.5, 0.5) ensures consistent input distribution, reducing spurious correlations related to input scale.

\textbf{Architecture Complexity:} The lightweight design (2 convolutional + 2 fully connected layers) is matched to the small input size (16×16). Unlike architectures designed for 224×224 or 128×128 inputs, this compact design avoids excessive capacity that would encourage memorisation, ensuring the model learns generalisable patterns.

\subsection{Justification for Non-Overfitting}

The lightweight architecture is suitable for small inputs and beneficial for embedded systems. With fewer layers (2 conv + 2 FC versus typical 3--5 conv + 2--3 FC) and simpler structure (no batch normalisation, minimal regularisation layers), the model balances capacity and generalisation.

Several factors mitigate overfitting risk: (1) extremely small input size (16×16) relative to standard CNN inputs, naturally limiting learnable pattern complexity; (2) transient signals are pre-processed signal peaks with inherent feature structure, reducing complex feature extraction requirements; (3) architectural complexity matched to input size, preventing excessive model capacity; (4) explicit regularisation via dropout ($p=0.5$) in the fully connected layer containing most parameters; (5) moderate learning rate and batch size in training configuration; and (6) compact parameter count (1,061,284) relative to input dimensionality, ensuring insufficient capacity for memorization while maintaining discriminative learning capability.

The combination of appropriate architecture complexity, explicit regularisation, and training configuration ensures the model generalises to unseen material samples while maintaining high classification accuracy, demonstrating that lightweight design does not equate to overfitting when properly matched to task and input characteristics.

\newpage
\section{Single-zone signal classifier (1D CNN)}
\label{sec:single_zone_signal_classifier}

The one-dimensional convolutional neural network (1D CNN) was adapted to operate on depth-integrated, single-pixel transients of length $1 \times 16$. Recovering material composition from such collapsed temporal profiles is inherently challenging, as subtle differences in the transient impulse response must be preserved despite integration and noise \cite{su2016,becker2023}. For planar, homogeneous targets, a single depth-integrated transient per pixel is, in principle, sufficient to characterise the surface, provided that ambiguities induced by distance, orientation, and albedo are controlled.

Single-zone experiments were therefore designed using individual pixels whose transients exceeded a minimum energy threshold, ensuring reliable albedo estimates. Two regimes were considered. In the first regime, targets were observed at a fixed distance while varying only orientation and view angle within the range $-15^{\circ} < \theta < +15^{\circ}$. Under these conditions, distance-related ambiguity was effectively removed, and the 1D CNN achieved high material classification accuracy that was comparable to the multi-zone 2D CNN model (Section~\ref{sec:multizone_signal_classifier}).

In the second regime, targets were observed over multiple distances while maintaining the same orientation and perspective variations used for the 2D stacked-input experiments. Introducing distance variation reintroduced depth--albedo ambiguity and substantially degraded performance: classification accuracy dropped by approximately $30\,\%$ relative to the fixed-distance case (see Section~\ref{sec:multizone_signal_classifier} for comparative results). This sensitivity to distance indicates that, for single-zone inputs, residual ambiguity in the transient profile cannot be fully resolved by the current preprocessing and network design.

To mitigate these effects, the same preprocessing steps used for the multi-zone stacked inputs were applied to each single-pixel transient. The key additional operation was depth integration via a centre-of-mass (CoM) estimate of the transient histogram,
\begin{equation}
z_{\mathrm{CoM}} = \frac{\sum_{t} t\,h(t)}{\sum_{t} h(t)},
\label{eq:com_depth_alignment}
\end{equation}
where $h(t)$ is the photon-count histogram at time bin $t$. Each transient was temporally normalised by shifting it such that its CoM aligned to a reference depth, yielding a distance-normalised sequence per pixel. This procedure was repeated independently for all sequences, producing a set of depth-aligned, single-zone inputs for the 1D CNN.

Despite these measures, the fixed-distance experiments consistently outperformed the multi-distance configuration, underscoring the dominant role of distance-induced ambiguity in single-zone inference. As a result, the primary experiments in this work rely on the multi-zone 2D CNN classifier, which is more robust to distance variation. Nevertheless, the single-zone 1D CNN remains an important direction for future research, where more expressive architectures and richer normalisation strategies may recover the lost performance and fully exploit the efficiency advantages of single-pixel transients.


\subsection{Single-zone experiment results}
\label{subsec:single_experiment_results}

\noindent The single-pixel 1D CNN achieved high material classification accuracy when evaluated under fixed-distance conditions, with overall performance of approximately 91\%. In this regime, all samples were acquired at a single standoff distance while view angle and orientation were varied within the range $-15^{\circ} < \theta < +15^{\circ}$, ensuring that distance-induced ambiguities were effectively removed. Under these constraints, the network learned stable, material-discriminative features from depth-integrated transients, supporting the conclusion that the underlying feature extraction and classification strategy is sound when distance variation is controlled.

When the same classifier was evaluated over a range of distances spanning 0--100~cm, while maintaining the same angular variations, accuracy dropped substantially to approximately 68--70\%. This reduction confirms that distance is a dominant source of ambiguity in single-zone processing: changes in propagation distance alter both the amplitude and temporal structure of the transient response, and these distance-dependent effects are difficult to disentangle from material-specific signatures when only a single spatial zone is available. In particular, attenuation, temporal broadening and background contributions can cause different materials at similar distances to exhibit comparable temporal profiles, or the same material at different distances to appear dissimilar, thereby confusing the classifier.

Depth normalisation using the centre-of-mass alignment defined in (\ref{eq:com_depth_alignment}) partially mitigated these effects by re-registering transients to a common reference depth, but did not fully restore the performance observed in the fixed-distance case. Residual distance-dependent variations in signal-to-noise ratio, multipath contributions and scattering behaviour remain present after normalisation and continue to interfere with material discrimination, especially for weakly reflecting or highly scattering surfaces \cite{su2016,tafone2023}. These observations are consistent with prior work in time-of-flight material classification, which has reported increased confusion between classes when distance and viewing geometry are not tightly controlled \cite{callenberg2021}.

By contrast, the multi-zone stacked-input CNN (Section~\ref{sec:multizone_signal_classifier}) maintained high accuracy across various range of distances and view angles. Within the experimentally tested range of $-15^{\circ} < \theta < +15^{\circ}$, neither moderate view-angle variation nor distance changes posed a substantial challenge to the multi-zone classifier, underscoring its suitability for practical deployments in which sensor-to-target distance and orientation cannot be precisely controlled \cite{fasolino2024}.

Processing complexity analysis revealed a clear trade-off between robustness and computational cost. Single-pixel modelling requires only one transient per decision, avoiding spatial correlation operations and operating on low-dimensional feature vectors; this results in reduced memory usage, lower latency and modest power demands, which are attractive for embedded or resource-constrained platforms. In contrast, the multi-zone stacked approach processes multiple spatial channels simultaneously and operates in a higher-dimensional feature space, increasing both memory footprint and computation time. These findings are consistent with the complexity analysis in Section~\ref{subsec:processing_complexity_non-spectral_temporal}, which predicted lower asymptotic cost for single-zone processing relative to multi-zone architectures \cite{knuth1976,bishop2006}.

Overall, the single-zone 1D CNN offers strong performance when distance is fixed, but its sensitivity to distance and angle variation limits its robustness in unconstrained settings. The multi-zone CNN delivers markedly higher resilience to these factors at the cost of greater computational complexity, highlighting a fundamental design trade-off between efficiency and robustness. Future work will explore hybrid schemes that combine depth normalisation with distance-aware feature transformations, adaptive filtering and selective zone processing, aiming to recover near multi-zone performance while retaining as much of the computational efficiency of single-pixel methods as possible.

\newpage
\section{Generative modelling}
\label{sec:generative_modelling}

Generative networks, also referred to as autoencoders in certain contexts, are unsupervised neural networks that learn two complementary functions: an encoder network that maps transient images into a compact latent space representation encoded as vectors, and a decoder network that synthesises new transient images similar to the original data from which the latent information was extracted. The objective of this experiment is to make reliable inferences about material composition by learning nonlinear patterns in the form of probability distributions encoded within transient image profiles.

While generative networks are commonly used for data augmentation in most domains, in remote material detection they serve a dual purpose: they help assert the quality and reproducibility of transient image features, and more critically, they enable the encoding of extracted material features such that similar materials can be synthesised with minimised reconstruction loss. Given that the total combinations of materials that can exist in any space is infinite, training detection or classifier models on all possible materials is infeasible in terms of time and processing cost. Therefore, there is a need for reliable inherent features that can assess never-before-seen materials as belonging to trained classes or not.

The essence of autoencoding in material detection is to learn a latent space of class members and use generative units to synthesise new transient images with the latent features. The learning algorithm optimises the generative task by minimising the loss between the synthesised and real transient signals \cite{ackley1985}. Two generative models were evaluated: restricted Boltzmann machines (RBM) and convolutional autoencoders (CAE). As implied by the name, RBMs have units restricted to direct (one-to-many) connections with each unit in the hidden layer. The stochasticity of RBM networks also suits the non-deterministic nature of shot noise within temporal data \cite{tontini2020}. To evaluate the consistency of generative performance, a second generative model, the convolutional autoencoder (CAE), was considered. The recorded performance of both systems is presented in the experiments and generative modelling sections \cite{su2016,jungerman2022}.


\subsection{Restricted Boltzmann machine (RBM)}
\label{subsec:rbm}

The Restricted Boltzmann Machine (RBM) is an unsupervised generative model inspired by associative memory models, particularly the Hopfield network \cite{hopfield1982}. Unlike Hopfield networks, RBMs employ a bipartite structure comprising visible and hidden layers, enabling efficient learning while maintaining pattern completion capabilities.

RBMs can recall complete patterns from partial or noisy inputs, making them robust for real-world transient image data that often contains noise or incomplete measurements. They generate patterns without requiring exact prior memories by learning probabilistic representations that capture the statistical structure of transient image patterns, enabling the generation of new patterns consistent with the learned distribution. The network comprises stochastic units, making it suitable for learning complex probability distributions and generating diverse outputs. This stochastic nature allows RBMs to model uncertainty, capture data variability, and explore probability spaces effectively, avoiding local minima in the energy landscape.

RBMs are essential for modelling temporal data with non-deterministic distributions, such as Poisson noise in transient images. Transient image data from single-photon detectors exhibits Poisson-distributed photon counts, where variance equals the mean \cite{tontini2020}. The energy-based framework of RBMs can model Poisson distributions, and their stochastic nature accommodates discrete, count-based photon detection events, enabling robust feature extraction even in low-signal conditions.

The network architecture consisted of approximately 250 units in the visible layer and 500 units in the hidden layer. The visible layer represents transient image data directly, with each unit corresponding to a time bin or temporal feature. The 2:1 ratio of hidden to visible units enables overcomplete representations, where multiple hidden units encode different aspects of the same temporal pattern, facilitating more expressive feature learning. In ten training cycles, a reconstruction loss of 0.0012 was achieved for generative tasks. This low error indicates the RBM successfully captured the essential structure of transient image data in its hidden representation, validating that the model learned meaningful features rather than memorizing training examples \cite{ackley1985}.

Physical observation of generated versus real images shows they are hardly distinguishable, providing visual evidence that the model learned the statistical structure of transient image profiles, including material-specific temporal signatures and variability patterns. This result demonstrates the existence of learnable signature patterns in transient image profiles from target materials, indicating that material-specific information is encoded in temporal patterns in a way accessible to machine learning algorithms \cite{su2016,jungerman2022}.


\subsection{Convolutional autoencoder (CAE)}
\label{subsec:cae}

A second autoencoder, the Convolutional Autoencoder (CAE), was employed to further validate the learnability of transient image features. Similar to RBMs, the CAE's generative task involves two complementary operations: encoding transient image sequences into low-dimensional latent vectors, and decoding these encoded features by upsampling into complete sample transient images. The encoding operation compresses high-dimensional transient image data into a compact latent representation that captures essential temporal patterns, while the decoding operation reconstructs the original transient images from this compressed representation, demonstrating that the latent space contains sufficient information to recover material-specific temporal signatures \cite{ackley1985,su2016}.

As the name implies, CAEs employ convolutional kernels to encode input transient images. Unlike fully connected layers in traditional autoencoders, convolutional layers apply learnable filters that detect local temporal patterns and spatial correlations in transient data \cite{lecun1998}. This convolutional approach is particularly effective for transient images because it can capture temporal dependencies, recognise recurring patterns across different time scales, and learn translation-invariant features that are robust to temporal shifts or variations in signal timing. The convolutional architecture enables the model to learn hierarchical representations, where lower layers detect simple temporal features (such as signal peaks or decay rates) and higher layers combine these into complex, material-discriminative patterns \cite{lecun1998,jungerman2022}.

The most popular applications of CAEs are seen in image denoising, where the network learns to separate signal from noise by reconstructing clean images from noisy inputs. However, in this material detection project, CAEs are essential for asserting the claim that transient image data contain learnable features that can not only be encoded in lower dimensions but also decoded to synthesise new sample transient images with minimal reconstruction cost. The ability to successfully encode transient images into compact representations and decode them back to realistic transient profiles demonstrates that the temporal data contain consistent, material-specific patterns that can be captured in low-dimensional latent spaces. This finding validates that transient image features are learnable, structured, and suitable for machine learning applications, rather than being random or unstructured noise \cite{su2016,jungerman2022,becker2023}.

Using the same input data as the RBM model, significantly lower reconstruction loss values ($<3.3877644 \times 10^{-5}$) were recorded for generative operations (Figure~\ref{fig:4_16}).

%\vspace*{\fill}
%\vspace{-1cm}  % Position figure ~1cm from footer (fill pushes to bottom, -1cm positions it)
\begin{figure}[H]
\centering
\includegraphics[width=0.8\textwidth]{figures/4.16.png}
\caption[Visualizations of transient image reconstruction performance from convolutional auto encoder.]{\textbf{Visualizations of transient image reconstruction performance from convolutional auto encoder.}}
\label{fig:4_16}
\end{figure}

This substantially lower reconstruction error compared to the RBM (0.0012) demonstrates that CAEs are particularly effective at capturing and reconstructing transient image patterns. The convolutional architecture's ability to learn local temporal patterns and hierarchical features enables more accurate reconstruction of the complex temporal structures present in transient image data. The extremely low reconstruction loss indicates that the CAE successfully learned a latent representation that preserves nearly all essential information from the original transient images, further validating that transient image data contain learnable, structured features that can be effectively compressed and reconstructed \cite{su2016,tontini2020}.

\section{Future research and development}
\label{sec:future_rd}

Future research will investigate multi-target detection and the associated complexities arising when multiple objects or materials are present simultaneously in the measurement scene. This represents a significant advancement beyond single-target scenarios, introducing spatial, temporal, and material-related challenges that require sophisticated signal processing and classification approaches \cite{su2016,becker2023}.

For scenarios involving identical objects positioned at multiple locations along the axial plane (depth axis), distance-resolved transient images can be retrieved for each instance with relative ease. Temporal separation of signals from objects at different distances enables the system to distinguish between them based on time-of-flight differences \cite{koerner2021}. Each object's transient response can be isolated and analysed independently, as temporal signatures are naturally separated by propagation delays. This scenario, while more complex than single-target detection, remains manageable because objects are identical in material composition, requiring only spatial-temporal separation rather than material discrimination \cite{tontini2020,koerner2021}.

However, when objects are varied along the same plane (at similar distances but with different material compositions), the situation becomes significantly more complex. Variations in reflectivity between different materials can cause their signals to confound with each other, creating occlusions and severe ambiguities that manifest as noise in the detection system \cite{heinz2001,tontini2024}. Materials with high reflectivity may dominate the signal, masking the presence of less reflective materials. Similarly, materials with similar optical properties may produce overlapping temporal signatures that are difficult to separate \cite{su2016,becker2023}. These reflectivity variations can lead to signal interference, where the temporal response from one material interferes with or occludes the response from another, making it challenging to accurately detect and classify all materials present in the scene \cite{heinz2001,schott2007}.

For this reason, possible signal occlusion in single and multi-target ranging will be studied and experimented for practical scenarios. Future research will systematically investigate how signal occlusion occurs, characterise its impact on detection accuracy, and develop methods to mitigate or compensate for occlusion effects \cite{tontini2024,jungerman2022}. This will involve theoretical analysis of signal interference patterns, experimental studies with controlled multi-material scenarios, and the development of deconvolution or separation algorithms that can disentangle overlapping temporal signatures \cite{heinz2001,koerner2021}. Research will explore how occlusion patterns vary with material properties (reflectivity, absorption, scattering), object geometry, spatial arrangement, and distance relationships, enabling the development of occlusion-aware detection algorithms that can handle complex multi-target scenarios \cite{becker2023}.

Another essential future objective is investigating how current classifier models could be modified to detect liquids or fluid objects with higher material inhomogeneity. Current models have been validated primarily on relatively homogeneous materials or liquid materials that have been homogenised (as demonstrated in the milk purity experiments) \cite{nielsen2013,aernouts2015}. However, many real-world materials exhibit significant internal heterogeneity, with varying composition, structure, or properties throughout their volume. Future research will investigate how material inhomogeneity affects temporal signal characteristics, how current feature extraction methods can be adapted to handle heterogeneous materials, and how classification models can be modified to account for or leverage material heterogeneity \cite{su2016,becker2023,elachi2021}. This may involve developing multi-scale feature extraction approaches that capture both local and global material properties, implementing ensemble methods that combine predictions from different material regions, or creating adaptive models that adjust their processing based on detected heterogeneity levels \cite{tafone2023}.

Other complexities such as varying isotropy levels and topologies will be further introduced to sampled targets to study how complex detection models can become. Isotropic materials exhibit uniform properties in all directions, while anisotropic materials have direction-dependent properties that can significantly affect their temporal response characteristics \cite{elachi2021,schott2007}. Future research will systematically vary material isotropy levels, studying how directional dependencies in optical properties manifest in temporal signals and how detection models must be adapted to handle anisotropic materials. Similarly, varying topologies—including surface roughness, internal structure, porosity, and geometric complexity—will be introduced to understand how structural variations affect material detection \cite{su2016,becker2023}. This will involve controlled experiments with materials of known isotropy and topology, characterisation of how these properties influence temporal signatures, and development of models that can account for or exploit these variations for improved material discrimination \cite{tafone2023}.

Further work on single-pixel detection for multi-surface scenarios and long-range distances will be explored. Multi-surface detection presents challenges because a single pixel receives signals from multiple surfaces at different distances, requiring temporal separation and identification of each surface's contribution \cite{jungerman2022,turpin2019}. Future research will investigate how to extract and classify materials from multiple surfaces using single-pixel measurements, potentially leveraging temporal deconvolution, multi-component analysis, or advanced signal processing techniques \cite{koerner2021,tontini2020}. Long-range distance detection introduces additional challenges, including signal attenuation, increased noise, reduced signal-to-noise ratio, and more pronounced atmospheric or environmental effects \cite{koerner2021}. Research will explore how detection performance degrades with distance, what signal processing techniques can extend the effective range, and how models can be adapted to maintain accuracy at longer distances. This will involve systematic distance-dependent studies, development of range-adaptive processing methods, and investigation of fundamental limits on detection range for single-pixel systems \cite{moramartin2021,yang2024}.

Related to this, future work will investigate in greater depth why single-zone transient inputs perform worse than multi-zone stacked inputs when distance and geometry vary. As demonstrated in Chapter~\ref{chap:flat_homogeneous_material_detection}, the single-zone 1D CNN attains high accuracy (approximately 91\%) under fixed-distance conditions, but its accuracy drops to roughly 68--70\% when target distance is varied between 0 and 100\,cm while maintaining orientation variations. Changes in propagation distance alter both the amplitude and temporal shape of the transient response through attenuation, temporal broadening and background contributions governed by the inverse square law \cite{griffiths2017}, and with only a single spatial zone the network cannot reliably disentangle these distance-dependent effects from true material signatures. Different materials at similar distances can therefore exhibit similar temporal profiles, while the same material at different distances can appear dissimilar, confusing the classifier. By contrast, multi-zone stacked inputs provide spatially diverse transients that allow the CNN to learn distance- and angle-invariant representations while preserving material contrast. Understanding these mechanisms in detail will be essential for designing future single-zone architectures that retain their efficiency advantages under fixed-distance conditions while improving robustness when distance and angle are allowed to vary.

Other detectors will also be experimented with to evaluate performance, especially for single-pixel inferencing. Current research has focused on specific detector types (such as SPAD arrays), but different detector technologies offer different advantages and limitations \cite{piron2020}. Future work will systematically compare detector types including different SPAD configurations, photomultiplier tubes (PMTs), avalanche photodiodes (APDs), and emerging detector technologies, evaluating their performance for single-pixel material detection and classification tasks \cite{piron2020,itzler2011}. This comparative evaluation will assess factors such as temporal resolution, sensitivity, dynamic range, noise characteristics, and suitability for single-pixel inferencing applications \cite{koerner2021}. The goal is to identify optimal detector configurations for different material detection scenarios and establish guidelines for detector selection based on application requirements, material properties, and environmental conditions \cite{becker2023,tafone2023}.

\section{Conclusion}
\label{sec:conclusion_chapter4}

This chapter has demonstrated the feasibility of automated remote material detection using low-cost, compact single-photon avalanche diode (SPAD) sensors with non-spectral illumination sources. The experimental framework successfully addressed fundamental challenges in material detection by leveraging temporal signatures encoded in transient histograms, establishing that material-discriminative information can be extracted from time-of-flight signals without requiring spectral decomposition \cite{su2016,becker2023}.

The core contribution of this work lies in demonstrating that non-spectral transient images from flat, homogeneous material surfaces encode signature features that enable reliable material classification across varying acquisition conditions. Through systematic experiments spanning multiple distances (20--400\,cm), orientations, view angles, and perspectives, the multi-zone two-dimensional convolutional neural network (2D CNN) classifier achieved overall accuracy exceeding 99\% when all sample variations were combined \cite{su2016,becker2023}. This performance validates the hypothesis that temporal signal profiles contain material-intrinsic information that remains robust to geometric variations, supporting the perspective invariance of homogeneous material surface composition.

The experimental design addressed critical ambiguities inherent in transient imaging through two complementary preprocessing strategies: multi-pixel transient stacking and depth-integrated single-pixel transients. The multi-zone approach, which processes spatially resolved transient arrays from 16 detector zones, demonstrated superior robustness to distance and view angle variations compared to single-pixel methods. This spatial diversity enables the model to infer distance- and angle-invariant representations while preserving material contrast, effectively resolving cases where single-zone inputs are ambiguous \cite{jungerman2022,koerner2021}. In contrast, the single-zone one-dimensional CNN (1D CNN) achieved high accuracy (approximately 91\%) under fixed-distance conditions but exhibited substantial performance degradation (68--70\%) when distance variation was introduced, highlighting the fundamental trade-off between computational efficiency and robustness to geometric ambiguities \cite{su2016,tafone2023}.

Validation through zero-shot learning on untrained perspective transients yielded approximately 90\% prediction accuracy, confirming that the classifier learns material-intrinsic features rather than view-dependent artefacts \cite{jungerman2022}. This perspective invariance represents a significant advancement over prior work, which typically required tightly controlled acquisition conditions with fixed geometries \cite{becker2023}. The ability to generalise across untrained viewing perspectives demonstrates that the learned temporal signatures capture fundamental material properties that are invariant to sensor--target geometry.

Generative modelling experiments provided independent validation of the learnability and structure of transient image features. Both the restricted Boltzmann machine (RBM) and convolutional autoencoder (CAE) successfully captured material-specific temporal signatures in their latent representations, with the CAE achieving reconstruction loss values below $3.3877644 \times 10^{-5}$ \cite{ackley1985,su2016}. The extremely low reconstruction error indicates that transient image data contain consistent, material-discriminative patterns that can be effectively compressed and reconstructed, validating that the temporal features are learnable, structured, and suitable for machine learning applications rather than being random or unstructured noise \cite{jungerman2022,becker2023}.

The preprocessing pipeline successfully addressed noise sources inherent in SPAD-based time-of-flight measurements, including ambient photon contributions and Poisson shot noise. Custom algorithms for baseline correction and depth extraction enabled reliable feature extraction from raw transient histograms, demonstrating that software-based signal processing can compensate for hardware limitations in low-cost detectors \cite{tontini2024}. The measured signal-to-background ratio of approximately 0.05 reflects challenging low-light conditions, yet the classifier maintained high accuracy through sophisticated denoising and feature extraction methods.

The experimental results establish several key findings. First, non-spectral transient images from homogeneous material surfaces encode sufficient information for reliable material classification, eliminating the computational overhead associated with spectral decomposition while maintaining strong discriminative performance \cite{su2016,becker2023}. Second, multi-zone processing provides essential robustness to distance and view angle variations, making the approach suitable for practical deployments where sensor--target geometry cannot be precisely controlled \cite{fasolino2024}. Third, the perspective invariance demonstrated through zero-shot learning confirms that material-specific temporal signatures are intrinsic properties that persist across viewing conditions \cite{jungerman2022}. Fourth, generative models validate that transient features are learnable and structured, supporting the use of machine learning approaches for material detection tasks.

The work demonstrates that material detection is feasible with compact, low-cost hardware when combined with appropriate signal processing and machine learning algorithms. The lightweight 2D CNN architecture, with approximately 1.06 million parameters, operates efficiently on ultra-low-resolution $16 \times 16$ transient inputs, making it suitable for embedded deployment alongside low-cost dToF sensors such as the VL53L5/VL53L8 families \cite{vl53l5cx2022}. This combination of low-cost hardware and efficient algorithms addresses a critical gap in practical material detection systems, enabling applications in robotics, safety-critical machine perception, and industrial automation where cost and form factor constraints are paramount.

The limitations identified in this work, particularly the sensitivity of single-zone methods to distance variation and the challenges posed by signal occlusion in multi-target scenarios, provide clear directions for future research. The demonstrated trade-off between computational efficiency and robustness highlights the need for hybrid approaches that combine the efficiency advantages of single-pixel methods with the robustness of multi-zone processing \cite{koerner2021,tontini2020}. Future work will explore adaptive algorithms that can dynamically select processing strategies based on detected scene complexity and available computational resources.

In summary, this chapter has established that non-spectral time-of-flight sensing with low-cost SPAD detectors can achieve reliable material detection for flat, homogeneous surfaces across practical ranges and viewing conditions. The combination of sophisticated signal processing, depth-aware preprocessing, and lightweight neural network architectures enables robust material classification that generalises across untrained perspectives and distances. These findings validate the core hypothesis that temporal signatures encode material-discriminative information accessible through machine learning, providing a foundation for extending material detection to more complex scenarios involving inhomogeneous materials, multi-target scenes, and varying material topologies \cite{su2016,becker2023,tafone2023}.

\newpage
\chapter{Spatiotemporal Object Detection}
\label{chap:spatiotemporal_object_detection}

\newpage
\section{Introduction}
\label{sec:spatiotemporal_introduction}

Automatic pattern recognition systems for machine vision have achieved unprecedented levels of sophistication, evolving from intuitive foundation models to large-scale interactive language architectures supported by powerful computational infrastructure. Despite these advances, contemporary computing systems remain substantially inferior to the visual perception capabilities exhibited by advanced biological systems. Guided by the framework (see Chapter 3), this work presents an instrument designed to enhance structural feature extraction in machine vision systems, addressing what is believed to be a fundamental limitation underlying numerous perception quality challenges. Through the integration of a single-photon avalanche diode (SPAD) time-of-flight sensor with an active pixel sensor (APS), this project demonstrates experimental findings that validate the proposed solutions, as illustrated in Figure~\ref{fig:5_1}. Beyond the practical advantages of low cost and compact form factor, the presented instrument yields significant improvements in machine vision quality and the reliability of subsequent decision-making processes.

\begin{figure}[H]
\centering
\includegraphics[width=0.9375\textwidth]{figures/5.1.png}
\caption[The combination of SPAD time-of-flight sensor and active pixel sensor enables improved structural feature extraction for enhanced machine vision quality.]{\textbf{The combination of SPAD time-of-flight sensor and active pixel sensor enables improved structural feature extraction for enhanced machine vision quality.} The combination of single-photon avalanche diode (SPAD) time-of-flight sensor and active pixel sensor (APS) enables improved structural feature extraction for enhanced machine vision quality.}
\label{fig:5_1}
\end{figure}

\quotethree

The experiments demonstrate that, even in the absence of spectral composition analyses—the predominant approach for understanding radiation-matter interactions—and without reliance on high computational power, sophisticated neural network architectures, or multimodal signal sources, the combination of transient and spatially resolved intensity images of the same target enables machines to perceive scenes in a manner that enhances higher-order signal intuition, thereby improving reasoning capabilities compared to contemporary vision systems. While improvements in learning architectures and scaling of processing power remain important considerations for machine learning and computational processing, the results presented herein suggest that these represent less fundamental aspects of constructing machines that require intelligence to function effectively. The most fundamental factor appears to be data quality: the perceived signal must contain correctly sourced spatial and structural features necessary for scene understanding. Reliable machine perception exhibits robustness against adversarial inputs, which fundamentally depends on the correct sourcing and interpretation of target spatial and structural features. However, state-of-the-art vision systems tend to rely exclusively on two-dimensional spatial images for certain structural features that are more accurately inferred from time-resolved signals or transient images. For instance, critical information such as material composition, true target depth, and texture constitute structural details that form the basis for safety-critical decision-making. The limitations of extracting structural features from two-dimensional images are extensively documented in stereoscopy literature \cite{laga2022}, which addresses depth extraction and other structural feature inference from two-dimensional intensity-resolved images. For the first time, this work employs SPAD time-of-flight detectors as the appropriate source for such structural features. One-dimensional time-resolved signals (transient images) are used to complement traditional two-dimensional intensity images as fundamental input sources for feature extraction and inference. The outcomes reveal substantially more reliable and meaningful perception by visual machines.

The main experimental input combines two-dimensional intensity images with one-dimensional transient images. The input quality contains variables for low-level object detection, a challenging task for contemporary machine vision systems. When these inputs are evaluated using state-of-the-art (SOTA) visual perception models, improvements in machine intuition are observed. Machines demonstrate enhanced understanding of visual content with respect to true structural details, which shapes their intuition for making more reliable probabilistic decisions. Further assessments of the solution indicate that the model can address additional challenges including occlusion and other obscured target perception problems in machine vision, scale estimation, and visible light limitations; in low-lighting conditions, an active infrared source remains unaffected. Time-resolved information enables more accurate target localisation in range and achieves this with lower computational complexity. Thus, even without sufficient or accurate spatial details, the transient signal may suffice, particularly for targets with unique structural features. The solution proves to be computationally less expensive, simpler, and significantly more effective. As this project represents the first of its kind in machine perception literature, a custom benchmark that holistically evaluates the presented work across instrumentation, signal source, and pattern recognition quality (model architecture, performance, et cetera) is also presented as an accurate evaluation methodology.

Overall, this work demonstrates, at multiple levels, the indispensable role of SPAD time-of-flight detectors in sourcing structural features such as material composition, true depth, and distance. These features can address long-standing machine vision challenges that spatial imagers producing two-dimensional intensity images have been unable to satisfactorily resolve independently. These solutions indicate a pathway toward designing and constructing truly intelligent machines. Among the problems confronting machine vision systems, the most undermining is their susceptibility to adversarial or deceptive observations, rendering them unreliable under circumstances requiring safety-critical decisions. The core problem stems from the lack of a holistic blueprint or model of perception upon which systems may be designed. Despite these hurdles, machine vision progression has been unstoppable, evolving from hand-crafted rule-based feature extraction algorithms (SIFT, HoG) through automatic pattern recognition models, particularly deep neural networks, that rely on decision theory for inference. Irrespective of rapid growth, it is evident that the field is experiencing greater advancement in learning architectures and processing power while signal quality remains overlooked. The consequences of this imbalance manifest in simplistic, unreliable decisions made by state-of-the-art (SOTA) models, which are often falsely attributed to the fact that these models perform pattern matching. This project, presents the single-photon avalanche diode (SPAD) time-of-flight detector as a complement to the traditional active pixel sensor (APS) in signal processing; specifically, it adopts an instrumentation approach to address the problem of input signal quality for scene understanding tasks. Experiments demonstrate that the combined instruments (time-of-flight sensor and active pixel sensor) can facilitate extraction of correct and higher-quality spatial and structural signal features that can be processed into superior representations of scenes, thereby enhancing machine decision-making and reasoning.

The perfect roadmap to building intelligent machines does not yet exist; nevertheless, researchers agree on the different factors that come into play. Larger vision systems rely on visual foundation models for encoding semantic information, common classes, or labels per point, for example: tree, vegetation, car, buildings, roads, or person. Though such descriptions are not geometric features, they are ambiguous for use as structural information necessary for the level of intuition targeted by these models. For instance, what does the tree feature imply in the scene context? A large tree constructed from wood? A silicone bonsai tree on an office table? A tree drawn on cardboard or merely a toy plastic tree? (Figure~\ref{fig:5_2}) This challenge of distinguishing natural objects from their synthetic representations is central to the experimental design, as detailed in Table~\ref{tab:5_1}. The challenge with modern visual systems is the exclusive reliance on instruments that can only source two-dimensional spatially resolved images. Thus, the vision system is forced to extract structural feature semantics from this two-dimensional intensity image input. This represents a more complex and often inaccurate approach, as evidenced by the works of numerous vision research experts, both past and present. Providing crude structural information is important for better interpretation of signals and correct understanding of target scenes by machines. This structural feature quality problem exists in all visual pattern recognition systems. For decades, experts have sought answers through improvements in learning architectures or increasing power in processing infrastructure; nevertheless, adversarial problems remain unsolved. Researchers who associate the problem with data quality often consider data metrics such as the number of observations, annotation quality, or labelling quality. While these are essential, they remain secondary to the representation quality of the received signal. This is because, if two-dimensional spatially resolved signals without transient information are used to represent a scene, there is a fundamental loss of important features for deriving quantities such as true depth or distance, target composition, and several others. Thus, irrespective of dataset size or annotation quality, the system is bound to make decisions based on limited information. This occurs because the core makeup of the perceived signal remains unchanged at the instrument level, and yet this fundamentally controls the quality of features extracted and the reliability of decisions made. By combining time-resolved signals from SPAD time-of-flight sensors with two-dimensional intensity images from active pixel sensors, the correct signal quality requirement at the instrumentation level of the machine perception framework is met. This consequently sets the system on a path toward enhanced scene interpretation and understanding.

\begin{figure}[H]
\centering
\includegraphics[width=0.75\textwidth]{figures/5.2.png}
\caption[Natural versus synthetic material discrimination challenge.]{\textbf{Natural versus synthetic material discrimination challenge.} Left: bell pepper object of natural material. Right: bell pepper object of synthetic (3D model) material. Both objects perceived from aerial view. Biological vision might find it easy to distinguish between the two objects, but how easily can machines do the same?}
\label{fig:5_2}
\end{figure}


\begin{sidewaystable}[p]
\centering
\adjustbox{max width=\textwidth,center}{
\footnotesize
\begin{tabularx}{\textwidth}{p{3.8cm}p{3.5cm}p{4.2cm}p{4.2cm}}
\hline
\textbf{Natural object} & \textbf{Natural material} & \textbf{Synthetic object} & \textbf{Synthetic material} \\
\textbf{(excludes images,} & \textbf{(excludes images)} & \textbf{(includes images,} & \textbf{(includes images)} \\
\textbf{includes components)} &  & \textbf{includes components)} &  \\
\hline
Bell pepper & Natural red colored skin & Artificial red colored bell pepper. Also paper print or electronic images of object (e.g LED screen) & Plastic, silicone or other 3D moulded material, Also Paper print or electronic screen (e.g LED) \\
Carrot & Natural carrot skin & Silicone moulded carrot. also paper print or electronic images of object (e.g LED screen) & Silicone. Printed or electronic images \\
Ceramic bowl & Bowl in its custom natural form (of ceramic material) & Paper print or electronic images of object (e.g LED screen) & Paper print or electronic screen (e.g LED screen) \\
Teacup & Cup in its custom natural or original form (of ceramic material) & Paper print or electronic images of object (e.g LED screen) & Paper print or electronic screen (e.g LED screen) \\
Plastic bowl & Cup in its custom natural or original purple colored form (of plastic material) & Paper print or electronic images of object (e.g LED screen) & Paper print or electronic screen (e.g LED screen) \\
\hline
\end{tabularx}
}
\caption[Naturally occurring or original objects used in project.]{\textbf{Naturally occurring or original objects used in project that have unique composition and corresponding synthetic objects could be made from different materials.} The term 'natural' is used to represent 'original' form of the tested material. SPAD time of flight detectors can potentially identify natural objects by unique composition even in the absence of spatial features like full geometry, form, etc. Such scenarios occur during scene occlusion or camouflage.}
\label{tab:5_1}
\end{sidewaystable}

State-of-the-art (SOTA) vision research demonstrates impressive work in three-dimensional perception and reconstruction, a task strongly founded on the quality of structural features for depth estimation. Nevertheless, the majority of recognised research in this domain has implemented solutions by depending on spatially resolved signals for all features. Scene flow prediction, three-dimensional image attribute estimation (for example, camera parameters, dense point clouds, multi-view depths), three-dimensional reconstruction with diffusion models, to mention a few, are tasks that continue to advance rapidly along computational performance benchmarks while remaining susceptible to adversarial examples. This challenge underscores the trend of contemporary vision algorithms still suffering from known blind spots, including but not limited to low-level classification (Figure~\ref{fig:5_3}), scale estimation, occlusion, models learning from fewer observations, and others. The presented experiments investigate how sourcing structural and spatial patterns from intensity and time-resolved signals, respectively, may help address these persistent problems. Findings from the conducted experiments demonstrate that a hybrid imager combining SPAD time-of-flight sensors with APS plays a major role in extracting quality patterns from scenes that enhance overall machine perception and reasoning quality. As spatial features are more suited for extracting geometric patterns, time-resolved features are better suited for structural or true depth features; a mix-up of using one source for the other's features may lead to complex processing or even less informative features, as a cross-feature sourcing problem would have been created.

\begin{figure}[H]
\centering
\includegraphics[width=0.975\textwidth]{figures/5.3.png}
\caption[Limitations of 2D spatial image-based object detection.]{\textbf{Limitations of 2D spatial image-based object detection.} Left: illustration of targets correctly detected by visual model without SPAD transient input. Right: same objects with some obscured parts (mammal) getting wrong model predictions, indicating limitations of using only 2D spatial images in training object detection models.}
\label{fig:5_3}
\end{figure}


\subsection{Spatial vs temporal properties}
\label{subsec:spatial_vs_temporal_properties}

Spatial properties enable the extraction of features essential for inferring target shape and form. These properties are typically derived from laterally resolved two-dimensional intensity images, and recent advances in stereoscopic processing demonstrate that spatial images can be transformed into features that enable models to infer target texture, distance, depth~\cite{depthpro2024}, and other structural properties. Nevertheless, existing literature and the results presented in this study reveal that such approaches are less efficient when compared to the performance of models performing the same task using time-resolved signals. The solution presented herein is designed to process target inputs that combine multi-sensor intensity and transient images. Although the implementation employs two neural network classifiers for each sensory input, it should not be mistaken for a novel neural network architecture. Rather, as described in Chapter~3, the solution demonstrates the feasibility of combining both spatial and time-resolved signals. Thus, a non-neural network classifier is technically possible. This clarification is essential to distinguish the solution from categorisations that would classify it solely as a neural network contribution.

To measure shape or other topological patterns, intensity measurements of the scene are collected using an active pixel sensor (APS) for both achromatic and chromatic multispectral (RGB) signals. This follows the described PAWD models for intensity signal measurements. Simultaneously, achromatic temporal signal measurements are acquired from the same target for structural and depth feature extraction. The multi-sensor features are learned using a custom spatiotemporal neural network-based model, which differs from traditional object detection models in both architecture and input format. The model requires paired input of spatially and time-resolved images.

Currently, there exists no universally accepted theory of machine perception. Nevertheless, the research community agrees on a common model of visual perception: for machines to perceive their environment accurately in three dimensions, they must extract both spatial and structural features. Spatial features relate to lines, vectors, angles, geometry, or, at a much higher level, topological properties of the target being perceived. Structural features, on the other hand, consider depth, internal and external composition of the target as it interacts with the source. In this context, two-dimensional spatially resolved images are targeted for spatial features, while transient or time-resolved images are processed for structural features. Corresponding outputs of spatial annotation coordinates, spatial annotation confidence, spatial annotation labels, and transient image labels serve as model inputs during training.

The input format also demonstrated potential to improve detection of perspective, scale, occlusion, visible light challenges, learning from less data, and other known issues in state-of-the-art vision systems. Performance shows substantial improvement in the reliability of machine decisions derived from perceived information, both geometrically and structurally. In space inference or related localisation tasks (often used in object tracking or detection), features of regions of interest are encoded within polygons, lines, or other geometric properties. This makes it possible for labels of different objects with similar geometric features (e.g., identical twins, distant football players wearing the same jersey, or even a human-like statue and a real person both in the same uniform) to be switched unnoticed. Both statue and real person could be classified as a single target, perhaps a human object. Spatial properties of a human-sized statue and a real person may both map to similar polygons or vectors. Thus, classifying both as one human object would be geometrically correct.

Spatial object classification poses no problems at higher-level vision tasks. Nonetheless, for low-level vision (e.g., security surveillance systems that require signature individual identification), each target should be identified uniquely beyond general geometric features. This becomes challenging for vision algorithms trained on spatially resolved images alone. The need for such distinct low-level classification is necessary for intelligent vision. Three-dimensional models of natural objects (e.g., common household items, vegetables, etc.) are used as observation samples. For each localised target, the corresponding transient signal is paired as input, with the expected output being the target object label together with the surface profile label. Predictive outcomes show significantly improved performance for both low- and high-level detection tasks.

At higher levels, most objects are identifiable from physical appearance. During obscured views such as complex perspectives (Figure~\ref{fig:5_4}) or occlusion (Figure~\ref{fig:5_5}), the usual object form may not be perceived, making it harder for machines trained on specific object views (e.g., front perspective) to correctly infer. For synthetics (objects that do not naturally occur, e.g., carved tables), natural objects like trees, humans, or mammals, and corresponding photographs (Figure~\ref{fig:5_6}), there exists a specific material composition property that may be unique to them. Most objects can be synthetically created, which makes physical shapes less unique. Nonetheless, the material composition of most synthetic objects serves as a signature feature.

\begin{figure}[H]
\centering
\includegraphics[width=1.035\textwidth]{figures/5.4.png}
\caption[Complex perspectives in object detection.]{\textbf{Complex perspectives in object detection.} Left: Camouflage of leaf mantis taking on surface patterns of the leaf on which it lays. At higher level with less obvious edges, classical vision models would miss the two different targets and possibly classify them as one. Nonetheless, transient images are able to encode unique composition of the leaf and the insect separately. Right: Camouflage of armadillo taking on the shape of a sphere after rolling up. At all levels, classical vision models are likely to treat both as different targets and possibly classify them independently. Nonetheless, transient images are able to encode unique surface composition of the armadillo and return the same material composition in either state.}
\label{fig:5_4}
\end{figure}

\begin{figure}[H]
\centering
\includegraphics[width=0.5\textwidth]{figures/5.5.png}
\caption[Partial occlusion of targets.]{\textbf{Partial occlusion of targets.} Illustration of partial occlusion of target by different objects. Depending on level of occlusion, it may be difficult to intuitively derive target from exposed spatial features. Nonetheless, SPAD time-of-flight detectors can reconstruct the target based on signal from its surface.}
\label{fig:5_5}
\end{figure}

\begin{figure}[H]
\centering
\includegraphics[width=0.4\textwidth]{figures/5.6.png}
\caption[Discrimination between natural and synthetic materials.]{\textbf{Discrimination between natural and synthetic materials.} Vision models trained only on spatially resolved images would make inferences based on geometric patterns of these three images, thus classifying them all as human objects. However, the presented solution proves that hybrid images including 1D transient signals can infer the geometric object as human and also infer the composition of each object as real human skin, wood, or paper, respectively.}
\label{fig:5_6}
\end{figure}

In a scenario where a human object is perceived through a pinhole (Figure~\ref{fig:5_7}), reflection from this opening could imply the presence of a human, as skin is peculiar to real human objects. Objects of plastic material would give a hint of the corresponding plastic or synthetic object. It may be difficult to infer the exact object without a complete or sufficient surface view. Nevertheless, material composition information may be sufficient under circumstances where composition data is necessary for unique identification or classification of the target. All state-of-the-art vision models are challenged by low-level detection, which, as seen from this work, is enabled by material composition data. The experiments presented serve as proofs of enhancements in machine understanding. By detecting structural composition and true depth with SPAD time-of-flight detectors, machines are better positioned to correctly classify targets based on unique and non-ambiguous features.

\begin{figure}[H]
\centering
\includegraphics[width=0.75\textwidth]{figures/5.7.png}
\caption[Viewing complex objects from different perspectives.]{\textbf{Viewing complex objects from different perspectives.} Illustration of viewing complex objects from different perspectives: full target view versus viewing components of targets through a pinhole. For systems that only process spatially resolved inputs, it may be easier to classify whole objects based on polygons, lines, or edges. For time-resolved systems, the signal can be used to deduce composition of the surface from which it emerged. A combination of these inputs gives hybrid space-time resolved input data that is suitable for the proposed spatiotemporal neural network.}
\label{fig:5_7}
\end{figure}

Questions that define this problem include: how to distinguish between two geometrically similar objects of different composition (Figure~\ref{fig:5_8}); if an object's full shape is significantly obscured, how would machines recognise it? By what means would this be reliably and efficiently accomplished? Perspective problems are common in scenes or targets with complex geometry, such as multiple fused polygons. Consider a simple object like a circle: views from all perspectives (aerial, side, or frontal view) render the same shape of circle. However, a complex object like the human body or other animated body forms are likely to produce significantly different object views from different perspectives. During the experiment, different perspectives and views of objects are presented to spatial object detection models. The performance in state-of-the-art object detection is seen to drop for instances where the presented perspective is ambiguous.

\begin{figure}[H]
\centering
\includegraphics[width=0.4\textwidth]{figures/5.8.png}
\caption[Geometrically similar objects with different material composition.]{\textbf{Geometrically similar objects with different material composition.} Two geometrically similar human objects depicting common causes of challenging low-level detection for machine vision. With similar geometric properties, it is easy for vision models to classify them as one target (human object). But at a low level, these would be two different individuals (e.g., identical twins).}
\label{fig:5_8}
\end{figure}

By introducing temporal features, the machine's understanding is improved beyond benchmarked performances of existing state-of-the-art visual detection systems. Besides low-level and perspective-invariant classification, SPAD time-of-flight detectors also play an important role in other vision tasks. In most vision tasks, scale estimation may not be required. However, in situations where it is necessary, it is also a hard problem for machines; this has been the case for many decades, and several attempts have been made to address the problem in recent years. The performance of some of the best scale estimation algorithms today is still not sufficient, especially when presented with larger scale variations. The fundamental questions that may help solve this problem with SPADs include: how reliably or efficiently would machines distinguish between a huge object and its miniaturised version? How does the ToF-based scale estimation compare to existing filtering-based or stereoscopic scale estimation approaches? SPAD time-of-flight detectors are presented as a potential solution in comparison with existing approaches, in terms of computational complexity and accuracy. ToF scale estimation is considered a more efficient estimation technique for both near and far distances within the detection range. The approach is expected to outperform filtering-based methods in cases of large scale variations (e.g., Baobab versus bonsai tree, see Figure~\ref{fig:5_9}).

\begin{figure}[H]
\centering
\includegraphics[width=0.75\textwidth]{figures/5.9.png}
\caption[Scale estimation challenges.]{\textbf{Scale estimation challenges.} Left: Spatially resolved (SR) images of a bonsai tree (left, 5~cm from detector) versus life-size tree (right, 20~cm from detector); both perceived from different distances. Spatially resolved images barely depict size difference of the two, and such vision systems may find it harder to distinguish between object scales, especially without information about distance or depth. However, time-resolved (TR) signals from SPAD ToF detector indicate that even with similar signal strengths, one item is significantly distant from the other. Thus, one object may be miniaturized in close view to sensor. Right: Spatially resolved (SR) images of a bonsai tree (left, 20~cm from detector) versus life-size natural tree (right, 20~cm from detector); both perceived from the same distance. Spatially resolved images clearly depict size difference of the two, and such vision systems may still find it harder to distinguish between object scales. However, time-resolved (TR) signals from SPAD ToF detector indicate that with variations in signal strengths even at the same distances from sensor, one item is significantly larger than the other. Thus, one object may be miniaturized in close view to sensor.}
\label{fig:5_9}
\end{figure}

For visible light range source systems, the presence of visible light determines the quality of the received and processed signal. Depending on detector sensitivity, too much, too little, or just unwanted lighting can compromise signal quality and consequently affect decisions taken based on them. The impact of visible light quality or availability on visual perception systems is high. For instance, APS may actively illuminate the scene to ensure sufficient visibility. Nonetheless, in unusual events where passive and/or active illumination fails, the resolved image quality is low. The combined imager of SPAD dToF and active pixel sensor is immune to this challenge. Even in the absence of high-quality visible light, SPADs can spatially resolve a scene by intensity (despite coarse lateral resolution) or possibly create a three-dimensional reconstruction with time-resolved signals. Raw transients may also be spatially resolved for pose estimation or other sparsely retrievable information.

The material composition of an undetected material is also still informative for vision tasks. Altogether, a custom vision system that localises and classifies targets using both temporal and spatial patterns is presented. The presented solution answers some questions: could machine vision ever achieve near-human-level perception or even surpass it? What role may time-resolved signals play? Results from the study prove that for machines to make meaningful and reliable decisions, time-resolved signals from SPADs may not be excluded. From low-level object detection, occlusion, scale estimation for cases of massive scale variations, visible light limitations, and others, single-photon avalanche diode detectors have proven from all experiments and study results to be indispensable in reliable and efficient intelligent vision systems. Due to the novel nature of the problem presented, it has been impossible to find an existing benchmark that holistically evaluates this work. Thus, a custom benchmark is used. Unlike other metrics, this benchmark considers instrumentation quality, pattern recognition, and classification as key parts of a fully functional and intelligent machine vision system.


\subsection{Feasibility of SPAD material detection for the tested materials}
\label{subsec:feasibility_spad_material_detection}

The materials selected for evaluating the proposed solution comprise both naturally occurring and original objects, each exhibiting distinct optical properties that generate unique reflectivity profiles suitable for material classification. The feasibility of these materials for the present study stems from their characteristic energy--matter interactions, which produce discriminable temporal signatures in SPAD time-of-flight measurements. The test objects span multiple material categories, including natural specimens such as vegetables that have developed through biological processes, synthetic counterparts represented by three-dimensional models, and image-based representations of both natural and synthetic variants. The most critical test case among these material categories involves image-based representations, whether rendered on printed substrates or displayed on electronic LED screens.

Prior research addressing spatially resolved material detection has been reviewed (see Chapter~2). However, a common limitation observed across these solutions is their failure to evaluate image-based representations of detected objects. This evaluation represents a crucial determinant of solution robustness, as the inability to distinguish between physical objects and their image representations implies that the material detection system may be susceptible to adversarial examples. For instance, if a material detection model cannot differentiate between a vegetable displayed on a screen versus a three-dimensional model or the natural version of the same vegetable, the system may be unable to make appropriate decisions that require distinct interaction strategies for each material type.

In practical applications, robotic systems must recognise material differences to determine appropriate handling and interaction protocols. A glass object may require careful and secure grasping to prevent potential falls and breakage, whereas a wooden spoon may not necessitate the same level of caution during manipulation. Similarly, autonomous vehicles trained to respond to pedestrian presence must distinguish between actual pedestrians and their image representations (e.g., billboard advertisements) to determine when stopping maneuvers are necessary. Most significantly, the ability to distinguish material composition can address sophisticated spoofing attacks and related security vulnerabilities that exploit the inability of vision systems to differentiate between physical objects and their representations.


\subsection{Feature quality challenge of vision}
\label{sec:feature_quality_challenge}

Biologically, the human brain, which neural networks seek to emulate, relies on distributed neuronal activations to process sensory inputs; hierarchical layers generate an exponential number of representational states for a given set of inputs or neuronal populations. While these processes excel at feature extraction, pattern matching, and inference, the quality of decisions rendered depends fundamentally on the fidelity of the received signal and its representational adequacy with respect to the scene. Even for non-automatic pattern recognition systems (e.g., non-ML vision systems that relied on handcrafted features), the quality of features that can be acquired remains constrained by the information content of the received signal. If a detector receives only spatially resolved signals (commonly encoded as RGB images), the feature space is inherently limited to spatial characteristics. Although stereoscopy demonstrates that structural features can be synthesised from such images, empirical evidence from such approaches indicates reduced efficiency compared to direct structural measurement.

Experiments presented in this work demonstrate that by combining a low-cost, compact SPAD time-of-flight sensor with traditional active pixel sensors (APS), more informative signals that better represent and enable understanding of the target are acquired. Without time-resolved signals, vision systems may process structural features from two-dimensional spatially resolved images; an approach that is computationally complex and less effective. Recent works addressing data quality often focus on dataset size (number of observations available for learning), annotation quality (in supervised learning paradigms), augmentation strategies, or automated labeling techniques. The data quality discussed in this project relates to signal informativeness: how well the signal represents the scene in three dimensions. In biological vision, this quality governs the discriminability of different targets and tends to have little dependence on processing speed or learning architecture.

This work addresses fundamental questions concerning how vision systems can be enhanced for safety and reliability. Through the integration of a single-photon avalanche diode (SPAD) time-of-flight (ToF) detector with an active pixel sensor (APS), this project demonstrates experimentally that sourcing signals from both spatial and temporal axes of a scene improves machines' understanding of perceived signals and consequently reduces their susceptibility to adversarial inputs. This enhancement improves reasoning quality for more intelligent decision-making. Specific vision challenges addressed include low-level classification, recognition of targets with complex geometry from perspectives that render them ambiguous, and visible light limitations. Based on experimental observations, the aforementioned challenges may share a common origin: input signal quality and cross-feature sourcing. In vision, targets are represented by spatial and structural properties.

The level of sophistication required to accurately retrieve and process these features determines the complexity of the task. Most vision systems rely solely on spatially resolved images to extract a wide range of topological, structural, depth \cite{depthpro2024}, and even behavioural features. Unlike time-resolved signals (see Chapter~\ref{chap:pawd}), spatial features can be computationally expensive and, under certain circumstances, inefficient when used for depth estimation \cite{depthpro2024} or other structural inferences. Nevertheless, this approach is commonly employed in most vision systems where three-dimensional perception is required \cite{depthpro2024}. As the research community gravitates toward building generally intelligent (AGI) machines, machines' susceptibility to adversarial examples becomes unacceptable. Hence, the focus of this work on less frequently discussed challenges, such as the limitations of relying solely on two-dimensional spatially resolved images for depth and other structural features. The experiments demonstrate that this challenge is fundamental for both foundation and higher-level interactive reasoning models.

Although the described problem is rarely addressed in recent vision literature, it affects all aspects of machine vision, especially autonomous systems. For state-of-the-art models to be trusted with safety-critical tasks, they must be capable of correctly and fully interpreting targets and not be susceptible to adversarial inputs; they should perceive at least at near-human levels of intelligence. A machine's inability to differentiate between a real human target and an identical statue or image (a problem that exists in models trained on spatially resolved signals) can lead to inaccurate decisions. Partial input sources from spatial images may consistently reduce a model's potential to fully or properly understand what they perceive in the physical world. The most common challenge of low-level classification is observed in almost all visual tasks: detection, tracking, or even generative systems. Machines can readily infer higher-level geometric features such as shape.

However, the task becomes more challenging as they encounter different targets with similar or identical geometry. This project presents evidence that time-resolved features can address such low-level classification with improved accuracy, lower computational complexity, and under certain circumstances (see Chapters~\ref{chap:flat_homogeneous_material_detection} and~\ref{chap:material_purity_detection}) with no reliance on spatial features. This is possible as a result of signatures embedded within time-resolved signal profiles (see Chapter~\ref{chap:pawd}). For machines to advance toward near-human-level vision, sole reliance on features from spatial images or even transients alone would lead to partial understanding of the scene. As demonstrated in Chapter~\ref{chap:flat_homogeneous_material_detection}, material classification required only time-resolved signals to extract composition features. Spatial classification tasks also require geometric features only to describe whether they contain a target object or not. Nonetheless, object detection or tracking rely on both structural and spatial features from spatially resolved signals.

Recent monocular depth estimation work \cite{depthpro2024} represents a typical example. Although not observed in the literature, transient data from range (ToF) sensors may be minimally efficient for spatially resolving scene targets. When cross-feature sourcing operations occur, the quality of decisions that can be made with the solution is reduced. For some low-level vision tasks, either structural or spatial features are useful. Material composition classifiers (see Chapter~\ref{chap:flat_homogeneous_material_detection}), for instance, do not require extensive information about the target configuration in space. Spatial object detection models also have minimal use for an object's material composition. Nevertheless, such systems exhibit less intuitive and partial understanding of targets, making them less suitable for high-intelligence vision environments. For machines to make reliable interpretations of their environment, there is a need to eliminate their susceptibility to adversarial inputs and improve their signal perception quality to acquire and process both spatial and time-resolved features.

The most prominent three-dimensional approaches, often employing stereoscopy, extract structural features such as depth and/or texture from spatially resolved images. Other non-stereoscopic methods rely on generic semantic features. Although these approaches are effective and have been responsible for the majority of breakthroughs in three-dimensional machine vision, they contain incomplete feature representations and exhibit several limitations: computational expense, low reliability (susceptibility to adversarial examples), and in some contexts, ill-posed formulations. This renders them generally less efficient for truly intelligent vision. Evidence of this is observed in reported and known limitations of recent works, including three-dimensional reconstruction, scene flow predictions, and others.

\subsection{Spatial vs time resolved image}
\label{subsec:spatially_resolved_image}

A spatially resolved image (see Chapter~\ref{chap:pawd}) represents the point-wise intensity resolution of signal, typically across two-dimensional space. Features extracted from such signals are commonly employed to define target shape and other spatial properties, as they encode geometric information of the scene. Most vision systems that process natural images from active pixel sensors—commonly RGB or infrared images—utilise spatially resolved images. At both lower and higher levels, these features enable localisation or segmentation tasks, which can be employed for decisions such as object detection: a hybrid task combining localisation and classification. Spatial features may be processed to infer other non-spatial attributes of a scene, including structural or depth features such as texture and distance. This process is known to be computationally intensive and complex, yet may yield reduced accuracy. Nonetheless, spatially resolved images serve as the primary input source for most vision systems and object detection models.

A temporal signal, or time-resolved image, represents the photon arrival time-wise distribution of arriving photons from a scene (see Chapter~\ref{chap:pawd}). Although typically employed for depth or distance estimation, the coarse resolution of most time-of-flight sensors can be spatially resolved for spatial tasks such as pose estimation. Additionally encoded within the signal profile are signature patterns that accurately determine the structural composition of material surfaces from which the signal emerged. This information has significantly improved scene understanding algorithms and consequently addressed some of the most long-standing challenges in vision. The presented solution that combines space- and time-resolved signals represents, however, a first in the literature.

\subsection{Objective functions}
\label{sec:objective_functions}

All machine learning systems optimise objective functions to fulfil their assigned tasks. For visual machine learning systems, the objective is the minimisation of a loss function that quantifies the discrepancy between ground truth signals and model predictions. Mathematically, the success of a vision task is measured by the degree to which the loss function is minimised. In the presented spatiotemporal model, two distinct sets of objective functions are optimised: the spatial detection functions, which are identical to those employed by conventional object detection models, and an additional material detection loss function, implemented as a binary cross-entropy loss for optimising the material composition of localised objects.

Traditional object detection models are trained to optimise three objective functions to successfully detect objects: localisation (bounding box) loss, confidence loss, and class label loss. The localisation loss quantifies how well the predicted bounding box coordinates align with the ground truth annotations. This function encompasses the $x$ and $y$ coordinates of the box centre, as well as the width and height dimensions, as expressed in Equations~\ref{eq:localization_coords} and~\ref{eq:localization_dims}. Alternative algorithms may employ corner coordinates as the starting point and expand based on predefined width and height parameters. This function is essential for detecting object location in space and is therefore retained for the spatial detection module of the spatiotemporal model.

\begin{equation}
\label{eq:localization_coords}
\equationfifteen
\end{equation}
\addcontentsline{loe}{equations}{\protect\numberline{\theequation}$\equationfifteen$}

\begin{equation}
\label{eq:localization_dims}
\equationeleven
\end{equation}
\addcontentsline{loe}{equations}{\protect\numberline{\theequation}$\equationeleven$}

The confidence loss, defined in Equation~\ref{eq:confidence_loss}, quantifies the model's certainty regarding the presence of an object within a grid cell within the localisation bounding box. Since this metric relates solely to spatial configuration, it is also retained for the spatial detection module of the model.

\begin{equation}
\label{eq:confidence_loss}
\equationtwelve
\end{equation}
\addcontentsline{loe}{equations}{\protect\numberline{\theequation}$\equationtwelve$}

In current object detection models, the classification loss is optimised to predict the conditional probability of a target within the grid belonging to one of the previously learned object labels, as expressed in Equation~\ref{eq:classification_loss}. In addition, an extra classification loss function, given in Equation~\ref{eq:material_classification_loss}, accounts for the conditional probability of the target's material composition belonging to one of the previously learned material classes.

\begin{equation}
\label{eq:classification_loss}
\equationthirteen
\end{equation}
\addcontentsline{loe}{equations}{\protect\numberline{\theequation}$\equationthirteen$}

\begin{equation}
\label{eq:material_classification_loss}
\equationfourteen
\end{equation}
\addcontentsline{loe}{equations}{\protect\numberline{\theequation}$\equationfourteen$}

In total, four objective functions are optimised for this model, as opposed to the three employed in traditional object detection models.




\subsection{Input format}
\label{sec:input_format}

The novel nature of this application precludes complete reliance on existing file formats. Consequently, a custom data type—implemented as a Python pickle object—was developed to create a tailored data class with attributes and methods that better accommodate the spacetime object representation. The file format employs the \texttt{.sto} extension (spacetime object) to encapsulate the combined transient and spatially resolved data layout. This custom file format enables precise control over the storage and serialization of spatiotemporal data, ensuring that the vision model can correctly retrieve and process relevant information. The \texttt{.sto} file format represents an effective approach for encoding richer signal representations, thereby enhancing machine perception for reliable and safer vision systems. Each file contains both spatially resolved and transient images from the time-of-flight detector for the perceived scene, providing a unified data structure that stores both spatial and temporally resolved signal inputs for visual processing.





\subsection{The spatiotemporal detection benchmark}
\label{sec:spatiotemporal_detection_benchmark}

This work demonstrates that separately extracting spatial and structural features enhances pattern quality for scene understanding. Experimental evidence indicates that combining transient (temporal) signals with spatially resolved images addresses persistent vision challenges more effectively than either modality alone. Existing benchmarks in computer vision primarily evaluate either spatial object detection or material classification in isolation, but not their combined performance in spatiotemporal analysis. This gap is significant because real-world scene understanding requires both identifying what objects are present (spatial detection) and understanding what materials they are composed of (structural and material analysis).

Current benchmarks fail to capture this dual requirement, leading to models that excel in one aspect while underperforming in the other, or that optimise for metrics that do not reflect holistic scene understanding. Without a benchmark that explicitly evaluates both spatial and structural signal quality, it becomes difficult to assess whether models are truly advancing toward comprehensive scene interpretation or merely improving in narrow, isolated tasks. In the absence of a suitable benchmark for evaluating spatiotemporal detection models, this work introduces a custom benchmark that assesses both spatial and structural signal quality. A key limitation in achieving near-human perception levels is the correct interpretation of acquired signals, rather than learning algorithms or processing speed.

Therefore, model evaluation is based on a custom benchmark that explicitly considers both spatial and structural features. The benchmark employs a customisable weighted average of spatial object detection and material composition detection modules. Users can adjust weights based on application requirements: material detection can be weighted higher for material analysis tasks, while object detection can be weighted higher for localisation-focused applications. For general object detection tasks, as demonstrated in the experiments presented here, equal weights (50\% each) have been used to emphasise that both object detection and material composition are equally essential for comprehensive scene understanding. This weighting scheme is subject to the solution purpose and is customisable.


\newpage
\section{Assumptions}
\label{sec:assumptions}

The following assumptions guide the experimental design and analysis of the spatiotemporal object detection system:

\begin{itemize}
    \item SPAD time-of-flight sensors provide structural features (material composition, depth, distance) that address limitations of appearance-based detection systems.
    
    \item Time-resolved signals encode material-specific optical properties, enabling distinction between objects with identical appearance but different material compositions.
    
    \item Spatial and temporal features are complementary: spatial features provide localisation and semantics, while temporal features encode structural properties (material composition, depth, scattering).
    
    \item Multi-material surfaces require Gaussian mixture modelling to account for superposition of material responses, each contributing distinct temporal signatures.
    
    \item Spatial object detection models can integrate with one-dimensional transient data to predict both object presence and material composition simultaneously.
    
    \item Temporal information enables effective learning with limited data, making the approach practical when large labelled datasets are unavailable.
    
    \item The modular architecture supports integration of diverse vision architectures with the material detection module.
\end{itemize}

\newpage
\section{Experiment overview}
\label{sec:experiment_overview}

This experiment investigates low-level object detection using a combined imaging system comprising an active pixel sensor (APS) and a single-photon avalanche diode (SPAD) time-of-flight detector. Object detection and tracking represent higher-level machine perception tasks that combine object localisation and classification to achieve comprehensive scene understanding. While significant advances have been made in this domain, contemporary systems continue to exhibit limitations in reliable and accurate predictions at lower levels of abstraction. Notable challenges include identity mismatches for geometrically similar objects, where surveillance systems performing object tracking for human targets may erroneously switch identity labels between targets exhibiting similar two-dimensional geometric patterns.

Such identity confusion can arise in several scenarios: identical twins (Figure~\ref{fig:5_8}), distant individuals wearing identical uniforms, or even between a target and an identical statue. In instances of high target density or crossover movements between multiple targets, identity labels may be switched without the system's awareness. Although such outcomes are undesirable from an application perspective, the model's behaviour remains technically functional within its operational context. Most visual detection algorithms are trained to predict targets based on learned patterns derived from spatial annotations. By combining two-dimensional spatially resolved images with corresponding transient signals, true structural information is acquired, rendering the representation more informative and thereby mitigating these limitations.

The presented object detection architecture processes two-dimensional images paired with time-resolved signals, simultaneously learning both spatial and structural features through dedicated processing pathways. In high-performing detection or tracking models, it is common for systems to predict objects based on semantic categories such as human, car, or tree. These semantic labels may prove insufficient under conditions requiring human-level intuitive understanding of the target. Consequently, reliable structural predictions from the presented model complement these semantic labels with intuitive details including target composition, distance from detector, target foreground composition, target background composition, pose, and other structural attributes. With such information, machines can deduce richer meanings from perceived objects, translating to higher-quality interpretation and more reliable decision-making. For instance, a prediction of both the object category (e.g., human) and material composition (e.g., bronze) enables deductions such as ``human-shaped bronze statue.'' As the model predicts object material and form (e.g., shape), it becomes easier for machines to infer lower-level attributes. Especially in scenes where differentiation of synthetic objects matters, the composition information helps disambiguate the exact type of object present. An autonomous vehicle configured to stop for humans may not be spoofed by synthetic human objects. Another application domain is medical in-vivo surgery, where robots tasked with operating on specific tissue require material composition information to reliably and correctly operate on the target; mistakes in this context could be critical. Correct structural information can guide machine vision systems in the safe and reliable execution of tasks.

Both humans and electronic imagers rely on visible light to illuminate and perceive objects within the environment. Under circumstances of limited visible light, the quality of perceived objects may suffer, consequently leading to poor signal interpretation. To ensure availability of illumination sources, most systems rely on active illumination. This experiment demonstrates the advantages of infrared SPAD illumination sources and how they complement visible light sources for APS. The combined instruments ensure that even under circumstances of insufficient visible illumination, the infrared source can recover a more reliable signal for further visual processing. Another aspect of the experiment is to showcase how SPAD time-of-flight detectors can assist in object detection for obscured views or perspectives that conceal the usual geometry of the object from recognition. To test the study claims, the project defines two object categories: synthetic and natural.

Natural objects, as implied by the name, suggest non-man-made objects that occur naturally in the environment. Natural objects exhibit material composition features that can uniquely describe them even when the full form is absent. For instance, if one observes human skin through a pinhole without seeing the face or size, they may not be certain whether it is synthetic or organic human skin. However, by perceiving and analysing the signal with a SPAD detector, this distinction can be readily made. Examples of natural objects include real animals, real humans, real trees, and others, each of which can be synthetically duplicated. Synthetic objects are non-naturally occurring and are often made from natural object materials for specific purposes. For all objects, material composition uniqueness and relevance are subject to how they are designed to function. Even for two objects of the same material but different geometry (e.g., a live tree and a wooden chair), the composition of the tree as seen in the living tree significantly differs from the composition of the tree as seen in the chair, which is not alive.

Under constrained partial views (Figure~\ref{fig:5_5}), spatial occlusion, or significant geometric transformations, the structural composition retrieved from the time-resolved signal becomes essential for deducing the signature of the object, even under obscured view circumstances. Spatial occlusion represents a common challenge in two-dimensional intensity imaging. This is because full or partial occlusion of targets may distort the perceived signal, leading to ambiguous interpretation and less reliable decision-making. In objects with less complex or consistent geometry throughout the scene, it may be straightforward to employ methods such as depth map segmentation, multi-view scanning, or generative machine learning architectures that can reconstruct the expected full shape from the partially seen form. However, these methods fail for complex geometry transformations or massive obscurities of the full form (e.g., pinhole views). This experiment leverages structural features acquired by SPAD sensors to retrieve unique composition signatures of targets as a means to complement the observed partial spatial image. This task is crucial for visual systems where crude identification of targets is necessary and semantics are insufficient. Autonomous driving systems and surgical robots represent examples of such applications.


\subsection{Instruments and setup}
\label{subsec:instruments_setup_chapter5}

Spatially resolved images were acquired from seven everyday objects, encompassing a diverse range of thirty-four distinct material compositions. Comprehensive details regarding these objects and their corresponding material classifications are provided in Table~\ref{tab:5_2} and Table~\ref{tab:5_3}.

\begin{sidewaystable}[p]
\centering
\adjustbox{max width=\textwidth,center}{
\footnotesize
\begin{tabularx}{\textwidth}{p{3.2cm}p{2.8cm}p{2.4cm}p{2.4cm}p{2.8cm}p{2.8cm}}
\hline
\textbf{Objects} & \textbf{Consolidation} & \textbf{Natural} & \textbf{3d model} & \textbf{LED Image (natural)} & \textbf{LED Image (3d model)} \\
\hline
Bell\_pepper\_red & Bell pepper & $\sim$2000 (x2) & $\sim$2000 (x2) & $\sim$2000 (x2) & $\sim$2000 (x2) \\
Bell\_pepper\_green &  & $\sim$2000 (x2) & $\sim$2000 (x2) & $\sim$2000 (x2) & $\sim$2000 (x2) \\
Bell\_pepper\_yellow &  & $\sim$2000 (x2) & $\sim$2000 (x2) & $\sim$2000 (x2) & $\sim$2000 (x2) \\
Teacup & Teacup & $\sim$2000 (x2) & NA & $\sim$2000 (x2) & NA \\
Eggplant & Eggplant & $\sim$2000 (x2) & $\sim$2000 (x2) & $\sim$2000 (x2) & $\sim$2000 (x2) \\
Bowl\_whiteceramic & Bowl & $\sim$2000 (x2) & NA & $\sim$2000 (x2) & NA \\
Bowl\_purpleplastic &  & $\sim$2000 (x2) & NA & $\sim$2000 (x2) & NA \\
Carrot & Carrot & $\sim$2000 (x2) & $\sim$2000 (x2) & $\sim$2000 (x2) & $\sim$2000 (x2) \\
Tomato & Tomato & $\sim$2000 (x2) & $\sim$2000 (x2) & $\sim$2000 (x2) & $\sim$2000 (x2) \\
Potato & Potato & $\sim$2000 (x2) & $\sim$2000 (x2) & $\sim$2000 (x2) & $\sim$2000 (x2) \\
\hline
\end{tabularx}
}
\caption[Table of objects used in training spatiotemporal object detection model.]{\textbf{Table of objects used in training spatiotemporal object detection model.} Objects were either of their natural form or 3d model form. X2 indicates the samples were collected for both transient and spatially resolved images.}
\label{tab:5_2}
\end{sidewaystable}

\begin{sidewaystable}[p]
\centering
\footnotesize
\begin{tabular}{p{7.5cm}p{7.5cm}p{7.5cm}}
\hline
\textbf{34 classes} & \textbf{18 classes} & \textbf{12 classes} \\
\hline
bowlLEDimagepurpleplastic & LEDscreen & LEDscreen \\
bowlLEDimagewhiteceramic & LEDscreen & LEDscreen \\
carrotLEDimage3dmodel & LEDscreen & LEDscreen \\
carrotLEDimagenatural & LEDscreen & LEDscreen \\
eggplantLEDimage3dmodel & LEDscreen & LEDscreen \\
eggplantLEDimagenatural & LEDscreen & LEDscreen \\
greenpepperLEDimage3dmodel & LEDscreen & LEDscreen \\
greenpepperLEDimagenatural & LEDscreen & LEDscreen \\
potatoLEDimage3dmodel & LEDscreen & LEDscreen \\
potatoLEDimagenatural & LEDscreen & LEDscreen \\
redpepperLEDimage3dmodel & LEDscreen & LEDscreen \\
redpepperLEDimagenatural & LEDscreen & LEDscreen \\
tomatoLEDimage3dmodel & LEDscreen & LEDscreen \\
tomatoLEDimagenatural & LEDscreen & LEDscreen \\
teacupLEDimageceramic & LEDscreen & LEDscreen \\
yellowpepperLEDimage3dmodel & LEDscreen & LEDscreen \\
yellowpepperLEDimagenatural & LEDscreen & LEDscreen \\
eggplant3dmodel & eggplant3dmodel & 3D model \\
greenpepper3dmodel & greenpepper3dmodel & 3D model \\
potato3dmodel & potato3dmodel & 3D model \\
redpepper3dmodel & redpepper3dmodel & 3D model \\
carrot3dmodel & carrot3dmodel & 3D model \\
tomato3dmodel & tomato3dmodel & 3D model \\
yellowpepper3dmodel & yellowpepper3dmodel & 3D model \\
yellowpeppernatural & yellowpeppernatural & yellowpeppernatural \\
redpeppernatural & redpeppernatural & redpeppernatural \\
teacupceramic & teacupceramic & teacupceramic \\
tomatonatural & tomatonatural & tomatonatural \\
bowlpurpleplastic & bowlpurpleplastic & bowlpurpleplastic \\
bowlwhiteceramic & bowlwhiteceramic & bowlwhiteceramic \\
carrotnatural & carrotnatural & carrotnatural \\
eggplantnatural & eggplantnatural & eggplantnatural \\
greenpeppernatural & greenpeppernatural & greenpeppernatural \\
potatonatural & potatonatural & potatonatural \\
\hline
\end{tabular}
\caption[Class consolidation strategies for material classification.]{\textbf{Class consolidation strategies for material classification.} The 34-class configuration represents the full material classification set, while 18-class and 12-class configurations consolidate related materials to improve classification performance.}
\label{tab:5_3}
\end{sidewaystable}

The experimental design incorporated systematic variation across multiple detection variables to assess robustness: random orientation, view angles, view perspectives, lighting conditions, and sensor--target distance. Four distinct orientations were employed, denoted as \texttt{pu}, \texttt{pd}, \texttt{pr}, and \texttt{pl} (portrait upright, downward, rightward, and leftward, respectively). View angles were systematically varied across \texttt{0dtsn}, \texttt{+15dtsn}, and \texttt{-15dtsn} (degrees relative to target surface normal). Material observations were conducted from three orthogonal axes: x-axis, y-axis, and z-axis. Bright illumination conditions were provided through both passive and active sources, including natural daylight. An equivalent number of observations were acquired under dark conditions to evaluate performance under low-light scenarios. Ranging experiments were conducted across multiple distances spanning 20\,cm through 60\,cm. No temporal frame averaging was applied due to the substantial sample size acquired. Corresponding RGB images were captured synchronously for each range cycle to enable cross-modal analysis.

A single experimental configuration was established using a SPAD detector for multi-pixel classification, with extracted features processed by a simple neural network classifier. The enabling instrumentation comprises a dual imager integrating a single-photon avalanche diode (SPAD) detector with an active pixel sensor (APS) employing RGB encoding. The combined imager enables multi-sensor resolution of both temporal and spatial signals, facilitating simultaneous acquisition of time-resolved and spatially resolved features. It is hypothesized that scene representation through such combined spatial and temporal features significantly enhances decision algorithm performance, manifesting as improvements in both processing speed and inferential quality. While instrumental capability fundamentally influences decision quality, the material and behavioural complexity of target scenes, as well as the characteristics of signal measurements, also play critical roles in determining overall system performance.

Time-resolved features acquired along axial axes (see Chapter~\ref{chap:pawd}) are frequently useful for inferring structural or depth-related characteristics of targets. Analogous to spatially resolved images, transients from multiple diodes can be spatially resolved to infer geometric features such as target pose estimation or nonlinear object detection. The resulting outcomes demonstrate that the instrument enables reliable classification and localisation of targets through combined spatial and temporal attributes. A custom SPAD sensor (VL53L8, MZ-Plus) operating in 4×4 zone mode with coarse lateral resolution and non-spectral illumination source is employed for detection. The setup features an active illumination source: a ~950\,nm laser with 15\,ns pulse width. The SPAD detector exhibits a 25° field of view, while the illumination source field of view is 35°. In addition to the SPAD detector, a Sony RGB camera is co-mounted with greater than 80\% field of view intersection with the SPAD detector on a Raspberry Pi platform. The Sony RGB camera exhibits a 35° field of view. An external Wi-Fi adapter mounted on a four-direction remote-controlled robot enables remote access to the sensor serial interface. Both SPAD and RGB sensors operate at 30\,fps (2\,ms exposure time) with a signal-to-background ratio (SBR) of approximately 0.05.

The data collection process (Figures~\ref{fig:5_12}, \ref{fig:5_13}, \ref{fig:5_14}, and \ref{fig:5_15}) was designed to source ground truth images of both spatial and structural features from the combined imager. A manual ranging experiment was established in a realistic natural scene setting comprising foreground and background targets. The transient image acquisition component of the setup consists of a SPAD module (VL53L8, MZ-Plus) with non-spectral illumination source. The module comprises a SiPM array with integrated TCSPC/TDC for generating timestamped photon counts. The system outputs 144-bin photon arrival time histograms, with return signals actively sourced from targets under several controlled conditions: scale, orientation, lighting, distance, and perspective. The spatially resolved image acquisition component of the setup consists of a Raspberry Pi camera module: an APS employing RGB multispectral encoding. Time-resolved data is read in hexadecimal format via the module's serial interface, sampling at approximately 19\,MHz through an I2C interface to the Raspberry Pi host. Readout data is processed through the same preprocessing pipeline (see Chapter~\ref{chap:flat_homogeneous_material_detection}) to generate clean feature representations of target surfaces.

\begin{figure}[H]
\centering
\includegraphics[width=0.75\textwidth]{figures/5.12.png}
\caption[Data collection setup from sample target.]{\textbf{Data collection setup from sample target.} Data collection setup from sample target (natural bell pepper) from ~30\,cm range distance. Left image indicates signal accumulation over 16 zones.}
\label{fig:5_12}
\end{figure}

\begin{figure}[H]
\centering
\includegraphics[width=0.75\textwidth]{figures/5.13.png}
\caption[Data collection setup of sample objects.]{\textbf{Data collection setup of sample objects.} Left: Data collection setup of a sample object (bell pepper) suspended from an invisible thread. Imager views from below while capturing both spatially resolved 2D images (RGB camera for raspberry pi module) and transient image (4×4 zones SPAD sensor). Right: data collection setup of a sample object (teacup) suspended from an invisible thread. Imager views tilted from sideways while capturing both spatially resolved 2D images (RGB camera for raspberry pi module) and transient image (4×4 zones SPAD sensor).}
\label{fig:5_13}
\end{figure}

\begin{figure}[H]
\centering
\includegraphics[width=0.75\textwidth]{figures/5.14.png}
\caption[Combined SPAD-APS data collection setup.]{\textbf{Combined SPAD-APS data collection setup.} Targets suspended from an invisible string from ceiling within detectors' field of view. Setup was within a natural less controlled environment with different lighting and ambient sources (daylight, tungsten bulb and low or no light). For each datapoint, the transient and corresponding spatially resolved image is sourced simultaneously.}
\label{fig:5_14}
\end{figure}

\begin{figure}[H]
\centering
\includegraphics[width=0.75\textwidth]{figures/5.15.png}
\caption[Intensity and transient imaging of perceived images of objects from an LED screen.]{\textbf{Intensity and transient imaging of perceived images of objects from an LED screen.} Intensity and transient imaging of perceived images of objects from an LED screen. Similar imaging techniques would be imaging the object from other persisted media: printed photograph, cardboard, or even wood painting. The active pixel sensor captures spatial features of the object. Hence the consistent performance to predict the same object on screen as though it was physically present. The SPAD records transient histograms which are used to predict the material composition of the surface from which the object is viewed. This unlocks low-level detection of knowing whether the object is natural material or an image version of it.}
\label{fig:5_15}
\end{figure}


\subsubsection{Spatiotemporal model process flow}
\label{subsec:spatiotemporal_model_process_flow}

To facilitate comprehensive demonstration and user interaction, a web application interface has been developed for all experimental configurations. Three state-of-the-art vision models were trained on the project dataset to elucidate the limitations inherent in conventional spatial-only detection approaches. These models were subsequently integrated into the spatial detection modules of the dual detection architecture. The material detection module employs a custom lightweight convolutional neural network trained specifically on transient image data. The complete spatiotemporal model is accessible for online testing (contact authors for web link). The following description outlines the process flow of the dual detection model (Figure~\ref{fig:5_22}).

\begin{figure}[H]
\centering
\includegraphics[width=0.75\textwidth]{figures/5.22.png}
\caption[Process diagram of dual detection web application.]{\textbf{Process diagram of dual detection web application.}}
\label{fig:5_22}
\end{figure}

The workflow commences with file upload in STO format, enabling immediate preview of the two-dimensional intensity image in the left panel. Optional preprocessing operations (cropping, rotation, flipping, or reset) may be applied to the preview, with these transformations affecting exclusively the spatial detection input stream. Users subsequently select both the spatial detection module and material detection module configurations. Upon activation of the scene detection function, the system executes parallel inference operations: the spatial head processes the edited two-dimensional intensity image through the selected model architecture (YOLOv3, YOLOv8, or DINOv3), while the material head processes a 16×16 transient image (STO index 0) through the selected 12- or 18-class classifier. The interface dynamically displays only the relevant dropdown menus corresponding to the chosen head configuration; the material head loads weights from preconfigured 12- or 18-class model paths. Both head and weight selections undergo validation prior to enabling the detection operation.

The backend executes both inference pipelines concurrently, returning detection results, processing times, and a composite result image (with bounding boxes where applicable). The visual display presents the final input image (reflecting the exact edited version) alongside inference outputs from both models. For spatial object detection architectures such as YOLOv3 and YOLOv8, the resulting image display includes drawn bounding boxes. For self-supervised image classifiers such as DINOv3, no bounding boxes are returned or rendered on the resulting image. The results section displays the top-three predictions for both detection modules. The system returns ``Found no objects'' if the detection set is empty or if all confidence values fall below the 5\% threshold, which is user-configurable.

Performance and architecture sections are dynamically updated with the following metrics: ``spatial detection architecture'', ``spatial detection inference time'', ``material detection architecture'', and ``material detection inference time''. A whiteboard area becomes visible following inference completion, incorporating an overlay caption that synthesises the detection results into customized natural language descriptions. When a new image is selected, the system resets the final image display, performance metrics, caption text, and hides any animated GIF content until the next inference cycle.


\subsection{Modality and fusion approach}
\label{subsec:modality_fusion_approach}

The dual detection model embodies both a multimodal and multi-sensor architectural framework. As the nomenclature implies, the multi-sensor component incorporates both SPAD and APS image data as input streams. The multimodal structure integrates two-dimensional RGB intensity imagery with one-dimensional transient signals. This configuration bears conceptual similarity to other multimodal input paradigms such as RGB-D and D-AI systems, among others. For existing multimodal input processing methodologies, the design space encompasses early and late fusion strategies \cite{gadzicki2020}. Early fusion enables merged inputs for unified processing, an approach particularly well-suited for compatible input sources whose low-level features can be combined without substantial information loss. However, the present work does not benefit from such flexibility due to the inherent complexity of the multi-sensor information employed.

The fundamental challenge arises from the disparate nature of the input modalities: RGB imagery represents a higher-level two-dimensional intensity representation of the scene, whereas transient signals encode low-level one-dimensional reflectivity information. The preprocessing requirements for extracting low-level features from the two-dimensional image substantially exceed those necessary for feature extraction from transient data. Consequently, the only viable approach for accurate input processing within the dual detection model is late fusion: the separate processing of RGB and one-dimensional transient images by independent neural networks prior to combining their results to render a decision. While late fusion represents the sole feasible option for the project's data formats, it confers several distinct advantages. First, it ensures superior error handling capabilities. In scenarios where one model receives corrupted input or fails to generate a prediction, the complementary model will likely succeed. A second benefit lies in its facilitation of enhanced machine intelligence through cross-modal inference validation. For instance, if a natural bell pepper is misclassified by the spatial model as a bowl, while the material classifier accurately detects natural bell pepper material, such inference discrepancies can be employed to audit the vision system: if the object is indeed a bowl, it cannot simultaneously exhibit natural pepper texture characteristics. This cross-validation mechanism can trigger further diagnostic procedures.

While multimodal and multi-sensor inputs are widely considered as performance enhancement strategies in neural network architectures \cite{gadzicki2020}, this project adopts a distinct perspective on the specific multimodal--multi-sensor data employed. The time-resolved signal is regarded as a rich source of structural features. Thus, the quality of the input does not merely stem from the fact that it emerges from different sensors or different data formats. Instead, the quality is enhanced specifically by the time-resolved features, which encode both target distance and material reflectivity information. When combined with any two-dimensional spatially resolved intensity image (whether RGB, infrared, or other modalities), this fusion yields a rich data representation for the corresponding scene. As a result, the performance of the presented solution is attributed to the spatiotemporal input characteristics rather than the multimodal nature per se. This distinction is evident from several other multimodal systems (see Chapter~\ref{chap:literature_review}) that have attempted to perform similar material detection tasks but have not achieved the level of success demonstrated by the solutions presented in this work.

\newpage
\section{The spatiotemporal model}
\label{subsec:signal_classifier}

The term \emph{spatiotemporal} describes the dual nature of the solution, characterising its capacity to process both spatially resolved and time-resolved signals for decision-making. The architectural framework employs a modular design that facilitates the integration of state-of-the-art vision models (including object detection and regression architectures) alongside material detection classifiers, as demonstrated in Chapters~\ref{chap:flat_homogeneous_material_detection} and~\ref{chap:material_purity_detection}. This modularity enables flexible deployment of complementary processing pipelines while maintaining computational efficiency and preserving the distinct advantages of each modality.

\subsubsection{Spatial detection module}
\label{subsec:spatial_detection_module}

The spatial detection module implements a decoupled localisation and classification architecture that enables the integration of multiple state-of-the-art object detection models with diverse architectural designs into the spatiotemporal framework. Three high-performing vision models were integrated for experimental validation: YOLOv3, YOLOv8, and DINOv3. This modular architectural design facilitates the seamless integration of alternative vision models alongside the material detection module, thereby providing flexibility in model selection and deployment. The web application interface demonstrates the implementation of all described models and their integration within the spatiotemporal solution.


\paragraph{YOLOv3 architecture}
\label{par:yolov3_architecture}

The first spatial detection network evaluated within the spatiotemporal framework was YOLOv3 (Figure~\ref{fig:5_16}), an open-source single-shot detector architecture employing a Darknet-53 backbone with residual block connections. The network extracts multi-scale features through a hierarchical feature pyramid, coupled with detection modules that predict bounding box coordinates and class probabilities at three distinct scales. Each detection module operates on predefined anchor boxes per grid cell, enabling efficient localisation and classification across varying object sizes. The objective function comprises three constituent loss terms: localisation loss for bounding box regression, objectness confidence quantifying the presence of an object within the predicted bounding box, and classification probability representing the likelihood that the detected object belongs to one of the learned class categories.

\begin{figure}[H]
\centering
\includegraphics[width=0.75\textwidth]{figures/5.16.png}
\caption[Web interface of spatial detection module using YOLOv3 weights.]{\textbf{Web interface of spatial detection module using YOLOv3 weights.} Web interface of spatial detection module using YOLOv3 weights. Same implementation is possible with other SOTA spatial detection models like DINOv3, YOLOv8 or even handcrafted vision models of the past.}
\label{fig:5_16}
\end{figure}

A pre-trained YOLOv3 model comprising approximately 26 million parameters was fine-tuned on a custom dataset of eight object classes: \texttt{bell\_pepper}, \texttt{bowl}, \texttt{carrot}, \texttt{eggplant}, \texttt{poor\_vis}, \texttt{potato}, \texttt{teacup}, and \texttt{tomato}. The \texttt{poor\_vis} label was introduced to represent images acquired under low-light conditions, where object detection models exhibit reduced predictive capability. This class serves to demonstrate that the material detection module can successfully identify materials from corresponding transient images captured under identical low-light conditions, thereby highlighting the complementary nature of temporal and spatial modalities. The model was trained on 640 images with 160 validation samples, achieving validation accuracy in the range of 94\%--99\%. Both training and inference procedures were configured to execute on CPU hardware. For a 224$\times$224 pixel input, inference time was approximately 99\,ms; while substantially slower than GPU-accelerated implementations, this performance was sufficient for the development and validation tests conducted in this work.


\paragraph{YOLOv8 architecture}
\label{par:yolov8_architecture}

The second spatial detection network integrated into the spatiotemporal framework was YOLOv8 (Figure~\ref{fig:5_17}), a more recent, proprietary, and stable iteration of the YOLOv3 architecture. The network employs a Cross Stage Partial (CSP) style block architecture that partitions feature maps, processes them through multiple lightweight convolutional pathways with partial residual connections, and subsequently concatenates the processed features. This design strategy, in contrast to heavier residual stacks, serves to enhance gradient flow and feature reuse while maintaining a minimal parameter count, thereby constraining the number of learned weights (e.g., convolution filters, kernels, and biases) that determine model size and memory footprint, which directly impacts overfitting propensity, and minimising floating-point operations (FLOPs) required for a single forward pass.

\begin{figure}[H]
\centering
\includegraphics[width=0.75\textwidth]{figures/5.17.png}
\caption[Web interface of spatial detection module using YOLOv8 weights.]{\textbf{Web interface of spatial detection module using YOLOv8 weights.} Web interface of spatial detection module using YOLOv8 weights. Same implementation is possible with other SOTA spatial detection models like DINOv3, YOLOv3 or even handcrafted vision models of the past.}
\label{fig:5_17}
\end{figure}

YOLOv8 employs an anchor-free detection module, with its objective function optimising three constituent loss terms: distributional box regression optimised by Distribution Focal Loss (DFL), objectness confidence optimised by Intersection over Union (IoU), and classification logits for objects within bounding boxes optimised by binary cross-entropy (BCE). To achieve enhanced detection performance, a pre-trained YOLOv8 model was custom-trained on the same dataset comprising eight object classes (\texttt{bell\_pepper}, \texttt{bowl}, \texttt{carrot}, \texttt{eggplant}, \texttt{poor\_vis}, \texttt{potato}, \texttt{teacup}, and \texttt{tomato}), yielding a mean Average Precision (mAP) of approximately 96.3\%. The model was trained on 640 images with 160 validation samples, identical to the YOLOv3 configuration. When executed on the same CPU hardware, inference speed was approximately 1.5$\times$ faster than YOLOv3, achieving an inference time of approximately 68\,ms for a 224$\times$224 pixel input.


\paragraph{DINOv3 architecture}
\label{par:dinov3_architecture}

To further demonstrate that the proposed solution exhibits no inherent limitations with respect to learning architectures, whether supervised or unsupervised, a self-supervised vision transformer framework (DINOv3, Figure~\ref{fig:5_18}) was integrated into the spatial detection module. DINOv3 employs a teacher-student knowledge distillation paradigm wherein two networks process different views of the same image input. The student network receives a patched or partially masked image, while the teacher network processes the complete target object image. The student network is subsequently trained to predict the full image representation, enabling the learning of robust visual features through self-supervision without requiring explicit class labels during pre-training.

\begin{figure}[H]
\centering
\includegraphics[width=0.75\textwidth]{figures/5.18.png}
\caption[Web interface of spatial detection module using DINOv3 weights.]{\textbf{Web interface of spatial detection module using DINOv3 weights.} Web interface of spatial detection module using DINOv3 weights. Same implementation is possible with other SOTA spatial detection models like YOLOv8, YOLOv3 or even handcrafted vision models of the past.}
\label{fig:5_18}
\end{figure}

Similar to YOLOv3 and YOLOv8, a pre-trained DINOv3 model comprising approximately 86 million parameters was custom-trained on eight object classes. The objective function employed cross-entropy loss, optimised using the AdamW optimiser with a learning rate of $1.001 \times 10^{-3}$ and a weight decay of $1 \times 10^{-4}$. Input samples were normalised 224$\times$224 pixel two-dimensional intensity RGB images. The output layers utilised softmax activation to transform logits into probability distributions, from which class predictions were inferred. The model achieved impressive performance, recording approximately 97\% accuracy (see Figure~\ref{fig:5_19}).

\begin{figure}[H]
\centering
\includegraphics[width=0.75\textwidth]{figures/5.19.png}
\caption[Confusion matrix of DINOv3 self supervised visual detection model.]{\textbf{Confusion matrix of DINOv3 self supervised visual detection model.} Confusion matrix of DINOv3 self supervised visual detection model. Custom trained and plugged into dual model as spatial detection module.}
\label{fig:5_19}
\end{figure}

The models described above represent a subset of spatial vision systems that may be integrated into the spatial detection module of this dual-model architecture. Beyond machine learning models, alternative non-machine-learning-based vision systems that learn spatial patterns can be equally integrated into this framework. Nevertheless, only machine learning models were employed in the experimental demonstrations, as they represent the current state-of-the-art approach to machine vision applications.


\subsubsection{Material detection module}
\label{subsec:material_detection_module}

The material detection module constitutes a custom-built classifier designed to integrate seamlessly within the detection modulespace. While alternative material classification approaches (including physics-based models or alternative machine learning architectures) may be employed, the present implementation employs a specialised convolutional neural network architecture optimised for material discrimination from transient image inputs. The classifier accepts processed 16×16 pixel images, converted to tensor format and normalised prior to feature extraction.

Three distinct model variants were trained according to different class consolidation strategies, yielding 34-class, 18-class, and 12-class configurations (see Table~\ref{tab:5_3}). Performance evaluation revealed that validation accuracy approached 99\% for consolidated class sets, whereas the full 34-class configuration achieved approximately 60\% accuracy (Figures~\ref{fig:5_20} and \ref{fig:5_21}). This substantial performance differential highlights the trade-off between classification granularity and discriminative capability inherent in material classification tasks.

\begin{figure}[H]
\centering
\includegraphics[width=0.75\textwidth]{figures/5.20.png}
\caption[Confusion matrix showing test results for 18 consolidated material classes.]{\textbf{Confusion matrix showing test results for 18 consolidated material classes.} Confusion matrix showing test results for 18 consolidated material classes; ~98\%.}
\label{fig:5_20}
\end{figure}

\begin{figure}[H]
\centering
\includegraphics[width=0.75\textwidth]{figures/5.21.png}
\caption[Confusion matrix showing test results for 12 consolidated material classes.]{\textbf{Confusion matrix showing test results for 12 consolidated material classes.} Confusion matrix showing test results for 12 consolidated material classes; ~98\%.}
\label{fig:5_21}
\end{figure}

The material classifier employs a decoupled convolutional neural network architecture optimised for 16×16 pixel image inputs. The network architecture comprises a three-block convolutional structure incorporating padding, batch normalization, ReLU activation functions, and 2×2 max pooling operations. Feature channels progressively increase through the network layers: 1, 3, 32, 64, and finally 128 channels, enabling hierarchical feature extraction from low-level edge detection to high-level material-specific patterns. Rather than employing traditional flattening operations, the architecture utilises global average pooling to reduce spatial dimensions to unity, thereby mitigating overfitting tendencies and substantially reducing the number of trainable parameters while preserving discriminative feature representations.

The classification head consists of three fully connected layers with ReLU activations, interspersed with dropout layers (dropout rates of 0.5 and 0.3, respectively) for regularisation. This architectural design balances feature extraction capability with computational efficiency, leveraging batch normalization to stabilize training dynamics and global average pooling to generate robust, translation-invariant feature representations.

Initial training was conducted on all 34 material classes derived from seven distinct objects (Table~\ref{tab:5_2}), yielding validation accuracy of approximately 57\%. However, when the architecture was optimised for reduced class sets (specifically 12-class or 18-class material detection configurations), the model achieved validation accuracy approaching 99\%, demonstrating its effectiveness for distinguishing between different material types across the seven object categories (Table~\ref{tab:5_2}). Inference performance was exceptional among all evaluated models, achieving inference times in the range of 1--5\,ms on CPU hardware (Figure~\ref{fig:5_24}). While GPU inference performance was not explicitly evaluated, execution on GPU hardware is expected to yield further performance improvements due to parallel processing capabilities inherent in convolutional operations.

\begin{figure}[H]
\centering
\includegraphics[width=0.75\textwidth]{figures/5.24.png}
\caption[Web interface of material detection module using custom classifier weights.]{\textbf{Web interface of material detection module using custom classifier weights.} Web interface of material detection module using custom classifier weights. Same implementation is possible with any ML, pattern recognition algorithm, or even physics-based models.}
\label{fig:5_24}
\end{figure}

\subsubsection{Guiding principles and models}
\label{subsec:ch5_guiding_principles_models}

In alignment with the experimental configuration detailed in Section~\ref{subsec:instruments_setup} of Chapter~\ref{chap:flat_homogeneous_material_detection}, this investigation employs a 940\,nm class 1 laser source with achromatic processing. However, in contrast to the homogeneous planar materials examined in Chapter~\ref{chap:flat_homogeneous_material_detection}, the common everyday objects investigated herein exhibit inherent inhomogeneity and non-planar surface geometries. Consequently, the temporal signal is modelled using a per-bin Gaussian approximation of non-spectral signals from multi-material surfaces, as formalized by Equation~\eqref{eq:gaussian_mixture}. This Gaussian mixture formulation provides per-bin approximations that account for the superposition of multiple material responses, where each constituent material contributes a distinct temporal signature characterised by its signal strength $S_m$, standard deviation $\sigma_m$, and peak arrival time $t_{0,m}$.


\subsection{Results}
\label{sec:results}

Three spatial object detection models and a custom material classification model were evaluated to assess their individual performance and effectiveness when combined for dual detection tasks. The evaluation framework demonstrates that spatial object detection models can seamlessly integrate with one-dimensional transient image data to make meaningful predictions about both object presence and material composition~\cite{depthpro2024,gyongy2021}.

\subsubsection{Individual Model Performance}
\label{subsec:individual_model_performance}

YOLOv8 achieved the highest object detection accuracy at approximately 98\%. As the latest iteration of the YOLO architecture~\cite{redmon2016,redmon2018}, YOLOv8 incorporates improvements in feature extraction, anchor-free detection, and multi-scale processing that contribute to its superior performance. The model's high accuracy demonstrates its effectiveness in detecting and localising objects in spatial images, providing reliable spatial information for scene understanding.

DINOv3, an unsupervised self-supervised learning model~\cite{simeoni2025}, achieved approximately 95\% accuracy. Despite being unsupervised, DINOv3's strong performance demonstrates the effectiveness of self-supervised learning for object detection tasks and highlights the potential for reducing annotation requirements in material detection applications.

YOLOv3 achieved approximately 94\% validation accuracy (see Section~\ref{par:yolov3_architecture})~\cite{redmon2018}. While an earlier version of the YOLO architecture, YOLOv3 still achieves competitive performance, demonstrating the robustness of the YOLO detection framework across different iterations. The model's performance validates that established object detection architectures can be effectively adapted for material detection applications.

The models were trained on fewer input images than typical object detection systems~\cite{edozie2025}, as the intended purpose was to demonstrate their ability to combine seamlessly with one-dimensional transient images to make meaningful predictions. Traditional object detection models are typically trained on large datasets with thousands or millions of images, but this evaluation focused on demonstrating the feasibility of integration rather than achieving maximum individual performance. The reduced training data approach validates that the models can learn effective representations with limited data when combined with complementary temporal information from transient images, making the approach more practical for applications where large labelled datasets may not be available.

\subsubsection{Combined Dual Detection Performance}
\label{subsec:combined_dual_detection_performance}

Table~\ref{tab:5_4} summarises the performance per spatial model and the corresponding performance in combination with the material classifier. The combined performance is calculated as a simple average, as the benchmark assumes equal relevance for both material and object classification correctness~\cite{gadzicki2020}. This equal weighting reflects the dual nature of the detection task, where both identifying what objects are present (spatial detection) and understanding what materials they are composed of (material classification) are equally important for comprehensive scene understanding.

YOLOv3 validation performance was approximately 94\%. When combined with the material classifier achieving approximately 99\% validation accuracy, the estimated dual detection performance is an average of the two (see Table~\ref{tab:5_4}), resulting in approximately 94.5\% combined performance. This combined metric reflects the system's ability to simultaneously detect objects and classify their materials, providing a comprehensive assessment of scene understanding capability.

Similar dual detection performance is estimated for YOLOv8 and DINOv3. YOLOv8 combined performance and speed are approximately 98.5\% and $\sim$8\,ms, respectively. The high combined performance (averaging YOLOv8's $\sim$98\% object detection with the material classifier's $\sim$99\% accuracy) demonstrates that the integration maintains high accuracy while providing dual detection capabilities. The $\sim$8\,ms processing time indicates that the combined system can operate in near-real-time, making it suitable for practical applications requiring rapid scene analysis~\cite{gyongy2021}.

The speed of the material classifier was highest, ranging from approximately 1\,ms to 5\,ms for all validation samples. This exceptional speed performance is attributed to the efficient processing of one-dimensional transient image data, which requires less computational overhead compared to two-dimensional spatial image processing. The material classifier's speed advantage, combined with its high accuracy ($\sim$99\%), makes it particularly well-suited for real-time applications where both speed and accuracy are critical. The speed range (1--5\,ms) reflects variability based on transient image complexity and signal characteristics, but remains consistently fast across all validation samples.

\subsection{Spatiotemporal model performance justification}
\label{subsec:spatiotemporal_model_performance_justification}

The high validation accuracies achieved across multiple independent model architectures (YOLOv3: 94\%, YOLOv8: 98\%, DINOv3: 95\%) provide evidence that the spatiotemporal approach learns generalisable representations rather than overfitting to training data. Importantly, the high validation performance of the spatial models (YOLO and DINO) is consistent with publicly recorded performance by the model developers~\cite{redmon2018,simeoni2025}, and thus these results are not unexpected. The consistency of high performance across architecturally distinct models indicates that the observed accuracy stems from the complementary nature of spatial and temporal features rather than model-specific memorization.

The material classifier's exceptional performance (99\% validation accuracy) is particularly justified given the well-characterised resultant signal profile described by Equation~\eqref{eq:gaussian_mixture}. This Gaussian mixture formulation accounts for the superposition of multiple material responses, where each constituent material contributes a distinct temporal signature characterised by its signal strength $S_m$, standard deviation $\sigma_m$, and peak arrival time $t_{0,m}$. The material classifier's ability to achieve high accuracy while operating at extremely fast inference speeds (1--5\,ms) demonstrates that the model has learned efficient, generalisable feature representations from these well-modelled transient signal characteristics, rather than complex, overfitted decision boundaries.

The material classifier's performance is also consistent with previous relevant work in material classification using time-of-flight sensors. The Princeton study by Su et al.~\cite{su2016} demonstrated that machine learning models can achieve high accuracy (80.9\% validation accuracy using RBF SVM) for material classification using raw time-of-flight measurements, establishing the feasibility of learning material-discriminative temporal patterns. Similarly, the MDPI Sensors paper by Becker and Koerner~\cite{becker2023} reported classification performance of 96\% accuracy for plastic material classification using optical parameter features from direct time-of-flight depth sensors. The consistency of high performance (approximately 99\% in this work) across these independent studies, employing different machine learning approaches and material classification tasks, provides strong evidence that the observed accuracy stems from the inherent material-discriminative information encoded in transient time-of-flight signals, validating the material classifier's performance as representative of the general capability of machine learning methods for time-of-flight-based material classification.

\newpage
\section{Future research and development}
\label{sec:future_research_development}

This study employs the SPAD sensor as a cost-effective and compact instrument to enhance vision systems' capability in detecting material composition of targets, thereby improving state-of-the-art (SOTA) object detection systems. However, several additional research directions and applications remain to be explored with this technology. In parallel, the benchmarking framework used to evaluate spatiotemporal models will be further developed to incorporate more informative and efficient metrics beyond simple averaging, enabling finer-grained assessment of performance across tasks, object categories, and operating conditions.

\subsection{Multi-object spatiotemporal detection}
\label{subsec:multi_object_spatiotemporal_detection}

The spatiotemporal vision system will be evaluated for multiple object detection scenarios. Current state-of-the-art vision systems can detect multiple objects spatially using laterally resolved images. As described in Chapter~\ref{chap:flat_homogeneous_material_detection}, two input processing methods were presented: the multi-zone input, used in this work, is better suited for detecting single targets in range, while single-zone inputs can identify different targets within the SPAD field of view. The primary challenge lies in the sensor's coarse lateral resolution. Future research could explore higher lateral resolution SPAD arrays to test multi-spatiotemporal object detection using single-zone inputs. A larger dataset for this application will be published.

\subsection{Camouflage detection}
\label{subsec:camouflage_detection}

More practical experiments on camouflage detection represent a promising research direction. Camouflage involves blending objects into their environment such that they become less perceivable by observers. In camouflage detection scenarios, there is high complexity in detecting different geometric or compositional transformations. Common camouflage techniques include armadillo folding into a round object, leaf mantis resembling exact leaves on which they lie, chameleon skin patterns changing seasonally, and military personnel applying mud to their skin in muddy environments. While most biological organisms use this as a survival mechanism, it poses significant challenges for reliable machine vision, especially systems trained to classify objects based on spatial and semantic information.

For objects that blend into their environment or change color, semantic features such as colour patterns or edges may become less obvious. Consequently, most vision models fail to make correct inferences. For instance, the reflection profile from armadillo skin before and after folding may vary: from the ball surface it could return circular reflection, while from the unrolled side it could exhibit diffused reflection. Camouflage in natural objects is inherently difficult and often elusive even for human observers. However, when machines are trained to identify objects from semantics and spatial information, their interpretation of perceived signals becomes unreliable under complex camouflages such as leaf mantis, chameleon, or other major geometric transformations seen in armadillo. Generative model solutions may struggle to reconstruct the original shape from the new shape, which has little or no resemblance to the original after camouflage. Thus, structural features become essential for reliable detection.

Stereoscopic methods such as depth or texture extraction are complex, expensive, and often less reliable. The use of SPAD time-of-flight sensors to acquire signature structural features represents an important research area. Misclassifications by vision models that rely on semantic and/or spatial information (which provide barely valid structural data) are common in most state-of-the-art systems. For stereoscopic structural feature extraction, previous works~\cite{hart1973,depthpro2024} account for limitations~\cite{depthpro2024} and complexity overhead, which makes them less attractive for designing reliable autonomous vision systems. With SPAD time-of-flight detectors, these tests may be conducted to demonstrate how camouflage and other visual illusions in natural objects may be addressed for autonomous systems that require high levels of accuracy in signal perception. From hidden soldiers in mud to chameleons changing colors by season, SPAD time-of-flight detectors could make it impossible for even the most clever illusions to evade detection.

\subsection{Scale estimation}
\label{subsec:scale_estimation}

Scale estimation is highly relevant in machine vision and represents another area for intensive future research and development. Without accurate judgment of target size and relative distance to the sensor, scale estimation problems may arise. A detector's inability to differentiate miniaturized versions of objects from actual ones may lead to poor decisions. For instance, a close-up view of a bonsai tree may appear similar to a distant view of a baobab tree. Thus, typical object detection systems without accurate depth estimation may assign the same inference to both objects. Correct target scale estimation is often needed in systems that simultaneously move and/or interact with different objects in a scene, and is important for enhancing how machines manipulate, navigate, or make interaction decisions. Especially for applications in three-dimensional reconstruction, targets must not be represented as smaller or larger than they are relative to other objects. Otherwise, the model becomes inaccurate and can lead to navigation errors or accidents that threaten the safety of agents using it.

Despite significant progress in recent years, the problem remains difficult due to factors such as partial occlusion, deformation, motion blur, fast motion, illumination variation, background clutter, and others. Most existing approaches also report low performance when encountering large scale variations in complex image sequences~\cite{depthpro2024}. From stereo triangulation to sensor fusion approaches with machine learning, limited progress has been achieved. In future experiments, a simple, low-budget, and compact solution may be explored. Unlike stereoscopic methods, the SPAD sensor measures true target distance (in centimeters or meters) as objects move in real time, which can improve the accuracy of estimated scale and support fast-changing environments as seen in autonomous vehicles. The small size of the sensor and infrared source also makes it portable and easy to integrate with other systems, even without visible light sources.

To compute accurate scale (width and height) of targets, the active pixel sensor (APS) and SPAD would be used to simultaneously capture images of targets. The transient would be used to estimate true distance, which is mapped to a sequence of centre-of-mass (CoM) values for each pixel. The target scale could be measured as a ratio of the distance to a computed index from CoM value(s) that are above a defined threshold. The threshold signifies pixel relevance; that is, whether a pixel has received sufficient signal to be considered as having detected the target. Without relying on any metrics from the two-dimensional intensity-resolved image, this method will be one of the least complex and yet reliable solutions in current literature.

Additionally, a common problem of data privacy could be addressed: two-dimensional images may be too invasive, but one-dimensional transient images of the same target may not be. The problem of working in domains where image privacy issues exist may be tackled with one-dimensional signals containing structural information or spatially resolving them to contain pose information.

\subsection{Spatial and temporal occlusion}
\label{subsec:spatial_and_temporal_occlusion}

Finally, more specialised experiments concerning the perception of two or more co-occurring signals (both spatial and temporal occlusion) at a common point could also be addressed. This represents a challenging problem for machine vision systems. Occlusion problems are well studied in spatially resolved image processing. However, occlusion may occur along either the spatial or temporal axis and could be full or partial depending on targets' relative configuration to occluding object(s). Future studies would tackle spatial occlusion in both temporal and two-dimensional space-resolved images using SPAD detectors. Such experiments could answer questions on: how accurately machines perceive targets of interest in space during full or partial scene occlusion; and what roles could time-of-flight SPAD detectors play in addressing these challenges?

\newpage
\section{Conclusion}
\label{sec:conclusion_chapter5}

This chapter has demonstrated that a compact, low-budget SPAD sensor is relevant in machine vision through practical demonstration. The key contribution of this work is the proof that reliance on time-resolved signals rather than spatial images for structural feature processing is essential. This has been proven via comparison of spatial material detection systems in current literature to the spatiotemporal material detection solution presented.

The evaluation of three state-of-the-art spatial object detection models (YOLOv3, YOLOv8, and DINOv3) integrated with a custom material classification model demonstrates that spatiotemporal fusion significantly enhances scene understanding capabilities. The spatial models achieved validation accuracies of 94\%--98\%, consistent with publicly recorded performance by their developers, while the material classifier achieved exceptional performance (99\% accuracy) justified by the well-characterised Gaussian mixture signal profile described by Equation~\eqref{eq:gaussian_mixture}. The combined dual detection system achieved overall performance of 94.5\%--98.5\%, with YOLOv8 achieving the highest combined performance (98.5\%) at near-real-time processing speeds ($\sim$8\,ms), demonstrating the practical viability of the approach for real-world applications.

One quality of the human visual system that machines struggle to achieve is precise judgment of depth, a capability attributed to three-dimensional vision. Stereoscopic depth sensing is known to require more complex setups and higher processing overhead in algorithm performance. This complexity tends to hinder the quality of vision systems. The impact of these challenges is often seen in the quality of decisions made, which consequently pose security and system reliability problems, especially in safety-critical applications such as autonomous driving systems. The presented spatiotemporal approach addresses these limitations by providing true depth information directly from time-resolved signals, reducing computational complexity while improving accuracy.

State-of-the-art vision models have excellent learning architectures and computational performance. However, the quality of understanding and interpretation of perceived signals, even by the smartest models of today, is far from that of seemingly less intelligent forms of life. Most discussions around data quality have been fixated on the size of observations, class distributions, and several others. None of which truly touches on the integrity of what information is actually in the scene in comparison to what information is actually recorded by the instrument or perceived by the machine. For instance, an interactive reasoning model with spatial image input might be able to contextually infer from an image of a cat with its paws pointing to the sky that it is laying on its back, but cannot tell the difference between a toy cat and a real living cat. As a basic example, if such systems are used to power an autonomous robot to deliver treats to living cats, it is likely to deliver treats to toy cats within the space. This work demonstrates that time-resolved signals encode material-specific optical properties that enable machines to distinguish between real objects and their representations, addressing this fundamental limitation of appearance-based systems~\cite{depthpro2024,gyongy2021}.

Empirically, most natural data have non-linear patterns. Machine learning systems operate under the assumption that by learning to extract features from large datasets, they can encode or learn these hidden patterns and thus make accurate predictions. This work is in agreement with such assumptions. Nonetheless, it argues that for high-quality three-dimensional perception, structural features may be sourced from time-resolved signals. It is also claimed that structural feature sourcing from two-dimensional space-resolved signals is a key factor in the low-quality predictions that affect the fluidity of higher-level decision making. The experiments described demonstrate that time-resolved signals are better suited for structural feature processing. They tend to perform intuitive tasks often related to depth, target composition, texture, true distance, and others. In material or material purity detection (see Chapter~\ref{chap:material_purity_detection}), only structural patterns may be necessary. However, full scene understanding tasks, as seen in vision models for object detection, require complete spatial and structural features to learn patterns at different levels of representation. Without such extent of learning, performing higher tasks like reasoning could be rife with inaccuracies and less reliable decisions. For this reason, time-resolved signals are considered a much more reliable source for structural features.

Progress in vision models has often focused on spatially resolved signal inputs. As a result, the quality of reasoning remains partial. Also, as research continues to advance, the race to build expensive computational infrastructure to meet the growing data and processing demands is on. But the rapid growth surrounding this could be making some unresolved problems of vision remain overlooked. This project presents a combined imaging sensor of an active pixel sensor and a SPAD time-of-flight detector to address such under-publicized problems of structural feature sourcing in both reasoning and foundation models. Coupled with experimental evidence presented, this work shows that as part of the journey towards building generally intelligent vision systems, certain fundamentals must be met prior to higher-level tasks like foundation intuitive, interactive reasoning, or even world models. While the described problem is barely tackled by most state-of-the-art vision research, it may not be a new problem in the field; it is seen to have close association with previously described inverse problems of optics~\cite{hart1973,koerner2021} or ill-posed depth estimation~\cite{depthpro2024} in other contexts.

The key contribution of this work lies in demonstrating that spatial object detection models can seamlessly integrate with one-dimensional transient image data to make meaningful predictions about both object presence and material composition. This integration addresses fundamental limitations of appearance-based detection systems by incorporating material-specific optical properties encoded in transient signals, enabling machines to distinguish between objects with identical appearance but different material compositions, a capability that has remained elusive in conventional vision systems~\cite{depthpro2024,gyongy2021,gadzicki2020}. The successful fusion of spatial and temporal modalities validates the complementary nature of these information sources, where spatial features provide localisation and semantic understanding while temporal features encode structural properties such as material composition, true depth, and scattering characteristics.

The experimental results demonstrate that the spatiotemporal approach learns generalisable representations rather than overfitting to training data, as evidenced by consistent high performance across architecturally distinct models ranging from supervised object detectors (YOLOv3, YOLOv8) to self-supervised feature extractors (DINOv3). Importantly, the models achieved high validation accuracy despite being trained on fewer input images than typical object detection systems, validating that the approach can learn effective representations with limited data when combined with complementary temporal information from transient images. This makes the approach particularly practical for applications where large labelled datasets may not be available. The material classifier's exceptional performance, operating at extremely fast inference speeds (1--5\,ms), demonstrates that efficient processing of one-dimensional transient image data requires less computational overhead compared to two-dimensional spatial image processing, making it well-suited for real-time applications where both speed and accuracy are critical~\cite{gyongy2021}. The modular architecture design enables seamless integration of diverse vision architectures alongside the material detection module, demonstrating the versatility of the spatiotemporal approach across different learning paradigms.

Before machines are tasked with navigating the world, predicting the future, or running on horsepower processing speeds, their ability to make such structural distinctions like true depth, material composition, and other lower-level perceptual abilities (as seen in humans) need to be accounted for. This problem has not only left the machine vision systems of today susceptible, but also caused evaluation of traditional pattern recognition to be based on monotonic benchmarking standards. This work may serve as a wake-up call for a new paradigm in building foolproof or more intelligent vision systems.

This work establishes SPAD time-of-flight detectors as an appropriate source for structural features such as material composition, true depth, and distance, features that can address long-existing machine vision challenges that spatial imagers outputting two-dimensional intensity images have not been able to address solely and satisfactorily. The demonstrated improvements in machine vision quality and decision-making reliability, combined with the low-cost and compactness of the presented instrument, indicate a pathway toward designing and building more intelligent vision systems that can make safety-critical decisions with greater confidence and accuracy.

\newpage
\chapter{Material purity detection}
\label{chap:material_purity_detection}


\newpage
\section{Introduction}
\label{sec:chapter6_introduction}

Material purity detection is fundamental across all material applications. Purity represents the extent to which the integrity of base components and structure within a material substance is preserved \cite{klocek1991}. In all utilisation of matter, from scientific research through electronic device manufacturing \cite{goodman1986} to food production \cite{walstra1999}, purity assessment is crucial for ensuring credibility, quality, usability, and performance. Materials may be considered >99.9\% pure if composed almost entirely of a single element or compound, such as mercury (Hg), pure alcohol (ethanol), or bromine (Br). Conversely, hybrid materials measure purity according to specific criteria, such as the proportions of constituent elements necessary for normal or expected functioning. For instance, perfume is a mixture of fixatives, fragrant essential oils, and solvents that combine in specific proportions to produce a unique scent; the slightest alteration to expected proportions can significantly change the resulting scent, rendering it impure. Purity of mixtures, especially in fluid materials, is subjective and dependent upon the specific composition necessary for function. This project takes an optical approach to purity assessment of homogenised fresh milk in a consumable state using a low-budget SPAD detector and non-spectral illumination source.

As a multi-element, multi-compound mixture, milk in its crude form is particularly complex, implying possible distinct transients of the same material under changing temperatures or structural transformation over time. Nonetheless, this project focuses on a single profile of interest as ground truth and detects deviations from this as impure; that is, patterns that are consistent in all consumable milk states. The majority of consumable milk purity measurement systems are based on spectral analyses of electromagnetic wave interaction, often in the near-infrared (NIR), ultraviolet, and visible ranges. Optical detection approaches are preferred due to their less invasive nature compared to chemical-based detection systems. Review of milk purity research across decades shows that, similar to all other materials, milk detection faces three key challenges: material complexity and instrumentation (see Chapter~4). In the context of material complexity, it is difficult to fully infer exact purity detection variables from source signal scattering in the target using a single calibration model \cite{aernouts2015,melfsen2012}. The root cause of this challenge is directly linked with fluid isotropy and homogeneity. While material purity is less correlated with structural composition, it tends to be easily measured in homogeneous samples (see Chapter~4) due to structural consistency. Foundations of this claim are rooted in some of the earliest milk purity studies \cite{chandra2024}, where milk purity instruments had homogenizing components. As noted by Barbano et al.:

\noindent\quotesix

Though milk homogenisation is a food process aimed at improving consumption experience, the resulting consistent structure doubles as a key property enabling efficient profiling of optically sourced transient images. Milk samples for purity experiments are sourced from pasteurized homogenised fresh milk products commonly found in consumer stores. In this regard, the experiments are not significantly impacted by the described material complexity. As Aernouts et al. observed:

\noindent\quoteseven

The immediate or eventual stakes of consumable material impurity are high. Yet, instrumentation for most milk purity assessment systems is often expensive, bulky, or complicated to implement, if not unreliable. With the emergence of compact and commercialized (see Chapter~4) low-cost single-photon avalanche diode (SPAD) detectors, this project presents a low-budget optical detection solution.

Practically, consumers may physically taste, smell, or feel (by touch) milk to infer purity. However, this approach of direct biological sensing can be unreliable in events of subtle or barely differentiable significant contamination. Also, under circumstances where the resulting mixture is hazardous, coming into direct contact with impure substances, even with instruments, may pose safety risks. Optical remote sensing (ideally through unopened packaging) is non-invasive and enables safe, reliable detection of all contamination forms.

The choice of signal measurement approach is as essential as detection hardware. Often in NIR spectroscopy, signal intensity is expressed as a function of different variables, including how intensely certain signal wavelengths are absorbed or scattered by different components of the fluid. These quantitative processes tend to be computationally expensive, not to mention the overhead of multiple calibration models necessary for analysing signal interaction with milk composition. Purity assessments that resolve visible light to intensities (spectrophotometric, e.g., spatially resolved multispectral (RGB) approaches) have even more critical drawbacks: ambient lighting conditions, image quality, and inability to detect invisible transformations \cite{chandra2024,joolaei2025} are among them. High-quality images are often required for feature extraction; sometimes, microscopes are used, which could be both expensive and impractical for the everyday milk consumer. In the absence of visible light (dark or low-light conditions), target signal cannot be accurately resolved. Spectrophotometric systems typically look for color-based changes of purified milk samples in the corresponding contaminated sample to make an inference. In events where the milk-contaminant interaction leaves no physically observable changes, the inferential quality of the spectrophotometer is challenged. In effect, such color-image-based purity detection methods are less reliable.

A promising solution, however, is with non-spectral time-of-flight detectors (SPADs) (see Chapters~3 and 4). Similar work with time-of-flight detectors \cite{nielsen2013} occurred for the first time within the past decade where spectroscopic analyses were performed. This project relies on non-spectral time-resolved transient signals for characterisation of microcomponents' interaction with contaminants of either matter state. With reduced material-of-interest (MOI) complexity, different studies are conducted to investigate and understand the efficiency of the SPAD detector for purity assessment, from homogenisation impact through milk manufacturer, fat content, and product generalisability of the presented solution. Machine learning methods like principal component analysis are used.

Material purity describes the degree to which materials are free from external matter that compromises their original properties for expected functioning. In both synthetic and natural materials, purity is a significant parameter for defining true composition. The essence of accurately assessed material purity directly impacts safety, performance, and credence. Similar to other substances of matter, most fluids encountered in daily life are mixtures whose purity assessment is subject to the properties and/or criteria or standards defined for usability. Purity in a hybrid mixture is subject to intrinsic structural state for which the material is useful. This work embarks on low-cost, non-spectral, time-resolved purity assessment of homogenised milk sourced from common consumer groceries. The quest to quantitatively, reliably, and economically assess milk purity spans decades \cite{biggs1967,barbano1989}. Most producers of consumer milk engage in some level of processing to make it suitable for human consumption. Often such processes (e.g., pasteurization, homogenisation, skimming, etc.) tend to shift original properties from crude colloidal composition toward a more homogeneous state. While supposed effects \cite{michalski2007,walstra1999} of homogenisation have been studied, the resulting impact of the process (see Chapter~4) \cite{barbano1989} for effective characterisation of milk particles \cite{aernouts2015,melfsen2012} cannot be overlooked. The low complexity of the presented solution in this work is strongly founded on this basis: the structural consistency achieved from the milk homogenisation process.

The flux form of milk makes it incapable of stability. This necessitates the use of a solid container for target positioning during ground-truth transient image sourcing. Choice of containers depends on signal penetrability. Transparent containers have full transmittance (see Chapter~3), which is ideal for high-penetrability remote sensing. However, most milk comes packaged in translucent containers. Thus, both materials are considered during data collection. Retrieved signal is labelled as emerging from material of interest (MOI): a hybrid signal of both the milk and light-penetrable solid container.

Many purity detection methods (especially with spectrometric methods) have been developed over the years \cite{joolaei2025}. Most are invasive chemical approaches or unique contaminant-specific solutions targeted at individual impurities, making them less generalisable. To achieve higher generalisability of the presented method, MOI is regarded as a homogenised liquid mixture interacting with external materials of either matter state. This enables much lower-level definitions of resulting impure mixtures described by intrinsic material properties.

In-depth examination of fluid impurities reveals two categorical sources of contaminants: material and non-material agents (see Figure~\ref{fig:6_1}). Either contamination agent can be actively or passively introduced. Non-material agents are mostly passive and may stem from product storage conditions, manufacturing processes, temperature variations in product environment, or sometimes just by mere passage of time (e.g., product expiration period). Non-material active contamination events may not differ significantly from the passive alternatives; however, their occurrence is certainly planned or intended. Deliberate fermentation of fresh milk samples in this work is an example. In the experiments, active non-material contamination is achieved by subjecting MOI to poor temperature and allowing the product expiration date to pass.

\begin{figure}[htbp]
    \centering
    \includegraphics[width=0.9\textwidth]{figures/6.1.png}
    \caption[Flow diagram of material impurity sources and categories and examples.]{\textbf{Flow diagram of material impurity sources and categories and examples.}}
    \label{fig:6_1}
\end{figure}

Material agents are impurities that take the form of matter. They can be actively or passively introduced. The majority of milk impurity cases involve active agents, for instance, dilution, preservation, and other milk adulteration cases. Nonetheless, material agents can also be passively introduced into milk (see Figure~\ref{fig:6_1}). Optical detection is an efficient remote sensing technique (see Chapter~4). There are many instances where visibly contaminated materials are easy to spot by the human eye (or consumer). However, due to the undeniable fallibility of biological senses, detection efficiency with human vision could be lower or error-prone under higher-speed, large-scale processing circumstances (e.g., milk manufacturing quality checks). Optical impurity changes manifest as alterations in fluid volume, structural consistency, or even color. In some circumstances, contamination is barely obvious (see Figure~\ref{fig:6_2}). Such invisible contamination of milk poses a much higher safety risk, even under low-speed processing scenarios or domestic settings. Low doses of hydrogen peroxide (H$_2$O$_2$) solution added to milk \cite{gregory1961} make for a typical example, with the resulting mixture barely showing any changes in consistency, volume, or color. Yet, the impurity is significant and may reduce nutritional value \cite{gregory1961}, which could possibly pose severe problems to consumers.

\begin{figure}[htbp]
    \centering
    \includegraphics[width=0.7\textwidth]{figures/6.2.png}
    \caption[(experiment 3) left: control sample; center: impure mixture containing hydrogen peroxide contaminant; right: impure mixture with lemon contaminant]{\textbf{(experiment 3) left: control sample; center: impure mixture containing hydrogen peroxide contaminant; right: impure mixture with lemon contaminant}}
    \label{fig:6_2}
\end{figure}

Current purity assessment solutions tend to be expensive, invasive, and/or complex to implement. Leading technology for most material detection systems is the single-photon avalanche diode (SPAD) detector (see Chapter~4). As a result, a low-budget, non-spectral remote detection system is the chosen instrument for this task. The choice of remote optical detection instrument considers reliable purity detection even for subtle or barely obvious changes. To properly investigate and prove the efficiency of the chosen device, a series of experiments are conducted. The envisioned thoroughness of this project guided how data is sourced, processed, and analysed. During data collection, the detector is set up in close range to the target at approximately 20~cm distance (see Figure~\ref{fig:6_3}). In comparison to long-distance, multi-angle ranging (see Chapter~4), there are minimal ambiguities to handle. Nonetheless, expected nonlinearities in transients are learnable by machine learning algorithms. The preprocessing pipeline from Chapter~4 is adopted to compose transient images.

\begin{figure}[htbp]
    \centering
    \includegraphics[width=0.7\textwidth]{figures/6.3.png}
    \caption[Simple data collection setup featuring target MOI, SPAD sensor and wireless access point for remotely accessing detector and data.]{\textbf{Simple data collection setup featuring target MOI, SPAD sensor and wireless access point for remotely accessing detector and data.}}
    \label{fig:6_3}
\end{figure}

In the first study, the detector is used to source transients to investigate the extent of milk inhomogeneity on signal scattering. Principal component analysis (PCA) was conducted for two milk samples from the same manufacturer (brand0) and fat content (14\%): one sample is homogenised while the other is not. Transients are passed through an autoencoder and upsampled to larger vectors. The resulting encoded transient is fit to a principal component analysis model. Principal components are extracted and visualised in a 3D scatter plot (see Figure~\ref{fig:6_4}). It is observed that purity patterns for fresh homogenised milk do not strongly overlap with those of inhomogenised milk. Thus, it may not easily generalise to fresh unprocessed milk. On these grounds, only homogenised pure samples have been used in all remaining purity assessment experiments for consistency.

\begin{figure}[htbp]
    \centering
    \includegraphics[width=0.7\textwidth]{figures/6.4.png}
    \caption[(experiment 01) 2D Visualization of fresh unprocessed milk (14\% fat) vs fresh homogenised milk (14\% fat).]{\textbf{(experiment 01) 2D Visualization of fresh unprocessed milk (14\% fat) vs fresh homogenised milk (14\% fat).} Spatial distribution of features indicate linear separability}
    \label{fig:6_4}
\end{figure}

A second experiment is designed to understand the extent to which all materials and MOI samples (pure and impure) are separable. Transient images are sourced from two control (pure milk) samples, and seventeen impure (material and non-material contaminated milk) samples are used. A multi-label CNN classifier is trained to successfully learn features for a total of nineteen classes. The third experiment is a general purity assessment for all impurity matter categories. A binary classifier is trained to distinguish pure control milk samples from all other impure samples (material, non-material, active, and passive). Using a portion of data samples from experiment three, experiment four investigates contamination of whole milk with solid and liquid impurities that leave invisible or barely obvious contamination effects (e.g., hydrogen peroxide or sugar). Experiments two to four were based on a specific homogenised fresh milk product (brand0). To answer the manufacturer generalisability of the binary model for the same product, the fifth experiment is conducted. The aim was to investigate whether SPAD-sourced transient signals encode purity metrics of milk that are similar across different manufacturers of the same milk product.

As milk comes in different fat contents, it was also essential to investigate in the sixth experiment whether purity metrics from transients vary across different fat concentrations. Four fat concentrations were studied: whole milk (14\%), whole milk (3.7\%), skimmed milk (0.3\%), and semi-skimmed milk (1.8\%). Transient images from these samples were collected with SPAD and processed as inputs for an autoencoder. Extracted features were then fitted to a PCA model. The first few (2--3) principal components were visualised for all five samples (see Figures~\ref{fig:6_5} and~\ref{fig:6_6}). In both the second and third dimensions, the group variables strongly overlap. This indicates non-linearly separable features or purity patterns among the homogenised milk samples of different fat contents. In comparison with experiment one, the features for the inhomogeneous samples were linearly separable from the homogenised sample. To further understand nonlinearities, deep neural networks are considered. The control sample was set to 3.7\% fat.

\begin{figure}[htbp]
    \centering
    \includegraphics[width=0.7\textwidth]{figures/6.5.png}
    \caption[(experiment 6) 3D plot of PCA of four homogenised whole milk products(14\%, 3.9\%, 3.7\%, 1.8\% and 0.3\% fat content).]{\textbf{(experiment 6) 3D plot of PCA of four homogenised whole milk products(14\%, 3.9\%, 3.7\%, 1.8\% and 0.3\% fat content).} Overlap of variable per class indicate high similarity in purity patterns. Explained Variance Ratio:[0.18696326, 0.15745316]}
    \label{fig:6_5}
\end{figure}

\begin{figure}[htbp]
    \centering
    \includegraphics[width=0.7\textwidth]{figures/6.6.png}
    \caption[(experiment 6) 2D plot of PCA of four homogenised whole milk products(14\%, 3.9\%, 3.7\%, 1.8\% and 0.3\% fat content).]{\textbf{(experiment 6) 2D plot of PCA of four homogenised whole milk products(14\%, 3.9\%, 3.7\%, 1.8\% and 0.3\% fat content).} Full overlap of variable per class indicate high similarity in purity patterns. Explained Variance Ratio:[0.18696326, 0.15745316]}
    \label{fig:6_6}
\end{figure}

The final question asked was how purity metrics of one milk product (e.g., fresh liquid milk) would generalise to other consumable or non-consumable milk products. Experiments 7a and 7b were conducted to investigate the structural composition of different milk products. In 7a, two consumable milk products other than fresh homogenised milk are considered: condensed and evaporated milk. The essence of this experiment was to confirm the claim that milk products can be represented as different mixtures or materials and thus may have different purity evaluation metrics. PCA analyses of features from the control sample and other milk products (condensed and evaporated milk) show stronger overlap between condensed milk, evaporated milk, and control samples (see Figure~\ref{fig:6_7}), but the association between evaporated whole milk and condensed whole milk seems stronger. As this implies the presence of non-linear patterns, individual validation tests are conducted per class using the pre-trained classifier. The impurity class from experiment three is maintained.

\begin{figure}[htbp]
    \centering
    \includegraphics[width=0.7\textwidth]{figures/6.7.png}
    \caption[(experiment 7a) PCA of three different whole milk (>3.7\% fat) products in 3D plot; condensed milk, evaporated milk and fresh whole milk.]{\textbf{(experiment 7a) PCA of three different whole milk (>3.7\% fat) products in 3D plot; condensed milk, evaporated milk and fresh whole milk.} Overlap indicative of similarity in feature variables but low linear separability.}
    \label{fig:6_7}
\end{figure}

Assessing material purity is essential in every utilisation of matter. From electronic applications \cite{barbano1989} to consumable food production, material purity is the foundation of several metrics including quality, performance, and safety. This project leverages a low-budget SPAD detector to conduct purity assessments of homogenised milk through logical and accurate characterisation of liquid particles. Both material and non-material contamination sources are evaluated. Typically, material formed from multiple elements is hybrid; its purity subject to expected proportions of the constituent materials that make it function as expected. For the specific case of milk in this project, purity is defined by milk composition as acquired within manufacturer-defined validity period and storage requirements. Low-budget and optical remote characterisation of milk composition with electromagnetic sources has existed for decades. In spite of this, the problem is barely solved. Review of trends \cite{joolaei2025} in most methods shows some relevant challenges which have inspired novel ways of re-examining the milk purity detection problem in this work.

\subsection{The milk material}
\label{subsec:the_milk_material}

Milk is fundamentally a hybrid mixture of water, fats, and proteins in proportions and forms that render it consumable. Numerous quantitative methods exist for assessing milk purity, each designed with consideration of milk as a standardized product. In this work, the material of interest (MOI) is considered a hybrid material a homogenised mixture in a solid container rather than as a general product with standard composition. The essence of such custom definitions can be seen from both past \cite{chandra2024} and recent (see Chapter~4) studies.

First, milk as a product is generally described as a colloidal mixture with proteins and dispersed fat globules that settle on the surface as it sits still. While this definition holds for raw unprocessed liquid milk, it barely fits the profile of general consumable liquid milk or other milk products, particularly the milk used in this work. The majority of consumable liquid milk (whole or skimmed) is processed \cite{walstra1999} into pasteurized and/or homogenised mixtures (see Figure~\ref{fig:6_8}). Thus, fat globules in homogenised milk are significantly reduced (approximately 1~$\mu$m) and uniformly dispersed. Particles do not suspend or settle over time (see Figure~\ref{fig:6_8}), making it a valid homogeneous fluid rather than a colloid. As noted by Barbano et al.:

\vspace{1cm}
\noindent\begin{minipage}{\textwidth}
\setlength{\parindent}{0pt}%
\setlength{\leftskip}{0pt}%
\setlength{\rightskip}{0pt}%
\setlength{\emergencystretch}{3em}%
\setstretch{2.0}%
\footnotesize%
\ttfamily%
\color{textgrey}%
\justifying%
\sloppy%
\noindent\quotesix
\end{minipage}

Without this clarity, it is easy to misinterpret experimental results, especially in relation to illumination scattering, in ways that could be very difficult to detect or rectify.

\begin{figure}[htbp]
    \centering
    \includegraphics[width=0.7\textwidth]{figures/6.8.png}
    \caption[(experiment 01) Left: fresh homogenised whole milk (14\% fat). Right: fresh heterogenous (not homogenised, colloidal) whole milk (14\% fat)]{\textbf{(experiment 01) Left: fresh homogenised whole milk (14\% fat). Right: fresh heterogenous (not homogenised, colloidal) whole milk (14\% fat)}}
    \label{fig:6_8}
\end{figure}

In the experiments, milk is defined as the material of interest (MOI) and interpreted in terms of composed matter. This approach to viewing MOI helps investigate purity assessment from a generalizable contaminant context. Another observation from milk purity literature shows that most research work makes generic definitions of product features with less consideration of physical product variations. For instance, condensed milk (see Figure~\ref{fig:6_9}) and evaporated milk (see Figure~\ref{fig:6_10}) are both liquid milk products of different physical composition, though chemically they may be similar, and thus should be treated as different materials in this context.

\begin{figure}[htbp]
    \centering
    \includegraphics[width=0.176\textwidth]{figures/6.9.png}
    \caption[(experiment 7a) consumable condensed milk (6\% fat content), homogenised]{\textbf{(experiment 7a) consumable condensed milk (6\% fat content), homogenised}}
    \label{fig:6_9}
\end{figure}

\begin{figure}[htbp]
    \centering
    \includegraphics[width=0.7\textwidth]{figures/6.10.png}
    \caption[(experiment 07a) Evaporated milk sample in homogenised state]{\textbf{(experiment 07a) Evaporated milk sample in homogenised state}}
    \label{fig:6_10}
\end{figure}

These experiments are conducted to ensure standard purity definitions that suit the generalisation of purity from one manufacturer of the same product to another become feasible. Liquid milk and yoghurt are different products from the same source (milk), yet depending on the processes applied to each, their resulting properties such as concentration, refractive indices, and attenuation \cite{gravert1968} factors during electromagnetic interactions differ. As Aernouts et al. observed:

\vspace{1cm}
\noindent\begin{minipage}{\textwidth}
\setlength{\parindent}{0pt}%
\setlength{\leftskip}{0pt}%
\setlength{\rightskip}{0pt}%
\setlength{\emergencystretch}{3em}%
\setstretch{2.0}%
\footnotesize%
\ttfamily%
\color{textgrey}%
\justifying%
\sloppy%
\noindent\quoteseven
\end{minipage}

Optical purity metrics that would pass for homogenised milk could possibly fail for yoghurt. Another benefit of viewing milk through the matter lens is the realisation that as a fluid, milk particles may be treated as a continuum that deforms when acted upon by force.

The reliable way to source learnable signals from these particles in a stable state is to leverage a solid signal-penetrable container that can be treated as an intrinsic part of the MOI during ground truth sourcing (see Section~\ref{subsec:role_of_solid_container}). Numerous previous research efforts have attempted the problem of milk purity. However, over decades, the most relevant solutions lean towards partial impurity-specific tests per material, which could be limiting given that there are infinite numbers of possible impurities that could contaminate milk. Material-based milk definitions helped establish finite categories of matter--matter interactions (see Table~\ref{tab:6_1}) for impurity generalisation.

\begin{sidewaystable}[p]
    \centering
    \footnotesize
    \begin{tabularx}{\textwidth}{p{1.8cm}p{1.5cm}p{2.5cm}Xp{1.8cm}X}
    \toprule
    \textbf{Material Impurity} & \textbf{Color Change} & \textbf{Consistency Change} & \textbf{Contaminant} & \textbf{Fluid Container Type} & \textbf{ID} \\
    \midrule
    \multirow{6}{*}{\textbf{SOLID}} & \multirow{3}{*}{Yes} & Homogenous & Instant coffee & & hm\_as\_tp\_sch\_ic \\
    \cmidrule{3-6}
    & & Inhomogeneous colloidal & Bread crumbs (skin) & & hm\_as\_tp\_scic\_bc \\
    \cmidrule{3-6}
    & & Inhomogeneous suspension & Food wrappers & & hm\_as\_tp\_scis\_fw \\
    \cmidrule{2-6}
    & \multirow{3}{*}{No} & Homogenous & Grounded sugar & & hm\_as\_tp\_suh\_gs \\
    \cmidrule{3-6}
    & & Inhomogeneous colloidal & Soft white bread (core) & & hm\_as\_tp\_suic\_swb \\
    \cmidrule{3-6}
    & & Inhomogeneous suspension & Transparent soft plastic bag & & hm\_as\_tp\_suis\_tp \\
    \midrule
    \multirow{6}{*}{\textbf{LIQUID}} & \multirow{3}{*}{Yes} & Homogenous & TCP antiseptic & \multirow{6}{*}{Transparent} & hm\_al\_tp\_sch\_tcp \\
    \cmidrule{3-4}\cmidrule{6-6}
    & & Inhomogeneous colloidal & Fresh lemon juice & & hm\_al\_tp\_scic\_lj \\
    \cmidrule{3-4}\cmidrule{6-6}
    & & Inhomogeneous suspension & Olive oil & & hm\_al\_tp\_scis\_oo \\
    \cmidrule{2-4}\cmidrule{6-6}
    & \multirow{3}{*}{No} & Homogenous & Hydrogen peroxide & & hm\_al\_tp\_suh\_hp \\
    \cmidrule{3-4}\cmidrule{6-6}
    & & Inhomogeneous colloidal & Viscous petroleum jelly & & hm\_al\_tp\_suic\_pj \\
    \cmidrule{3-4}\cmidrule{6-6}
    & & Inhomogeneous suspension & Clear liquid perfume & & hm\_al\_tp\_suis\_cp \\
    \midrule
    \multirow{3}{*}{\textbf{GAS \& LIQUID}} & \multirow{3}{*}{NA} & Inhomogeneous colloidal & Air (foamy milk) & & hm\_ag\_tp\_scic\_air \\
    \cmidrule{3-6}
    & & Inhomogeneous colloidal & Spray paint & & hm\_ag\_tp\_scic\_spaint \\
    \cmidrule{3-6}
    & & Inhomogeneous suspension & Spray perfume & & hm\_ag\_tp\_scis\_sperf \\
    \bottomrule
    \end{tabularx}
    \caption{Categorization of material contaminants by matter state and resulting mixture attributes including color and homogeneity. Materials of samples used in the experiments are listed under contaminants.}
    \label{tab:6_1}
\end{sidewaystable}

\subsection{Relevant milk processes}
\label{subsec:relevant_milk_processes}

Fresh unprocessed milk is considered whole with at least 3.25\% fat content. In its native state, milk exists as a colloid with less dense fat globules that rise to the surface when left undisturbed. Native milk fat globules range from 0.1~$\mu$m to 20~$\mu$m in diameter, with an average size of 3--5~$\mu$m \cite{walstra1999}. Unlike fresh milk, consumable milk from retail sources \cite{walstra1999} may undergo skimming after pasteurization to reduce fat percentage. While different processes serve distinct purposes for consumption, milk homogenisation serves a dual role: it improves product consistency and enables reliable automated pattern recognition for purity assessment.

Both skimming and homogenisation alter the intrinsic milk texture to produce a smoother, more consistent structure. Skimming removes a significant portion of dense suspended fat particles, improving uniformity within the fluid without reducing fat globule size. In both skimmed and unprocessed whole milk, homogenisation \cite{walstra1999,michalski2007} can be performed to reduce fat globules to approximately 1~$\mu$m by pressurizing and dispersing macromolecular elements uniformly throughout the fluid.

The nutritional impact of homogenisation processes has been studied \cite{michalski2007}. Nonetheless, the homogenised state of milk is essential \cite{aernouts2015,melfsen2012} for signal processing applications (see Chapter~4) \cite{chandra2024}. In homogeneous mixtures, there is structural and compositional consistency per unit volume. When milk settles in a colloidal state, particles at certain regions of the fluid differ significantly in size from those at other regions; the concentration of fat globules at one position also varies in density and signal penetrability compared to surrounding water molecules. Unprocessed milk with unstable colloid may be considered heterogeneous.

Optically, such variations result in highly inconsistent transient matter states at different times (Figure~\ref{fig:6_11}), making them barely usable as ground truths. Coupled with heterogeneity, colloids exhibit highly volatile states such that particle distribution or configuration changes under the slightest fluid motion. While this mercurial property is learnable, it represents a major complexity overhead for automatic pattern recognition algorithms. In this regard, homogenised milk provides an ideal fluid medium for encoding purity patterns.

\begin{figure}[htbp]
    \centering
    \includegraphics[width=0.7\textwidth]{figures/6.11.png}
    \caption[Different colloid mixture states showing unsuitability for transient GT image sourcing.]{\textbf{Different colloid mixture states showing unsuitability for transient GT image sourcing.} Left: stable colloid mixture with uniformly dispersed particles through water. Less suitable state for transient GT image sourcing. Center: unstable colloid mixture with dispersed particles through water. Less suitable state for transient GT image sourcing. Right: suspensive colloid state with both suspended and dispersed particles through water. Less suitable state for transient GT image sourcing.}
    \label{fig:6_11}
\end{figure}

\subsection{Role of the solid container}
\label{subsec:role_of_solid_container}

Milk consists of a continuum of particles with distinct optical and physical properties \cite{walstra1999}. As a liquid, milk exhibits loosely bonded particles that enable free flow, necessitating containment for practical handling and measurement. To ensure adequate light penetrability for optical sensing, translucent and transparent containers with high signal transmittance are employed \cite{klocek1991}.

During signal detection, the container and milk form an occluded scene where the fluid cannot be detected independently of its container. Consequently, the transient signal represents the combined optical response of milk as perceived through the container material. When a transparent container is illuminated, full transmittance occurs, with subsequent scattering through the uniformly distributed particles within the milk \cite{aernouts2015,melfsen2012}.

Most fresh milk products are packaged in translucent materials, where illumination sources encounter scattering contributions from both the container walls and the milk itself \cite{walstra1999}. The returned signal may exhibit relatively lower intensities compared to transparent containers, yet the signal profile alterations remain less significant, preserving the essential material-dependent signatures required for purity assessment \cite{nielsen2013}. Ground truth sourcing from both container types enables generalisation of purity assessment across common packaging standards.

\subsection{Material impurities}
\label{subsec:material_impurities}

A common presumption in milk purity solution design is the finite categorization of contaminants. Impurities are typically classified as chemical or biological agents for which spectroscopic or chemical tests are designed \cite{joolaei2025}. While these methods are well-established, they tend to challenge the generalisability of solutions, as evidenced by the ever-growing number of milk purity detection systems \cite{joolaei2025,chandra2024}. Prior to experiment design, a preliminary investigation was conducted to address this limitation by understanding the broader scope of impurities that can potentially contaminate milk. This study established that fluids (or materials generally) may be contaminated by material or non-material agents, either of which could occur actively or passively (Figure~\ref{fig:6_1}).

For non-material agents, most impurities are passively introduced, though active agents are possible. Passive non-material contamination is defined by unintentional alterations to the structure and/or components of the material of interest (MOI). This could occur from activities such as manufacturing processes, storage conditions, temperature variations, or the mere passage of time (e.g., expired milk) \cite{lott2023,burgwald1947}. Despite the prevalence of passive agents, certain active non-material contamination events are deliberate. Often useful in product manufacturing or transformation, milk may be intentionally fermented to create other variations of the product. Transformations such as exposing milk to temperature levels that dry it into powdery form or make cheese could pass for passive non-material contamination within the context of liquid milk as the ground truth MOI.

With material agents, most impurities are actively instigated, though passive material contamination exists. Active material contamination (which accounts for the majority of milk impurity cases) refers to impurities that are purposefully introduced \cite{joolaei2025}. Some common active contamination processes include additives, dilution, or adulteration, as seen in most cases of milk contamination \cite{joolaei2025,chandra2024}. Doping, for instance, involves the intentional addition of foreign materials to another substance in order to enhance certain properties. In food preservation, certain substances may be added to enhance shelf life. Another common event is dilution: the addition of water to a fluid in order to increase volume and/or lower its concentration \cite{joolaei2025}.

There are also passive material agents that contaminate materials. In manufacturing settings, certain food ingredients may accidentally contain unexpected traces of other foreign ingredients due to manufacturing environment or process errors \cite{goodman1986}. Non-material agents are mostly passive. Natural expiration of a product, unintentional poor storage, and/or temperature conditions can significantly alter fluid properties \cite{lott2023,burgwald1947}. Both tangible (material) and intangible (non-material) contaminants produce resulting mixtures that take on either state of matter: solid or fluid (gas, liquid, plasma). However, the extent of material transformation is measured from optical attributes such as scattering by particle distribution and material consistency of the incident source.


\subsection{ Signal profiles of the milkcontainer}
\label{subsec:signal_profiles_of_the_milkcontainer}

The milk container constitutes an independent material system with distinct optical properties \cite{klocek1991}. Characterising the signal profiles of the empty container provides a baseline for isolating the absolute scattering contribution of the milk content \cite{aernouts2015,melfsen2012}. However, since measurements were conducted under ambient atmospheric conditions rather than vacuum, an empty container measurement captures the combined optical response of both the container material and the intervening air molecules as detected through the container walls. Consequently, the ground truth composition is defined by the composite light scattering signature of milk as perceived through the container system as an integrated optical path \cite{nielsen2013,walstra1999}.

\subsection{Classifier models}
\label{subsec:classifier_models_chapter_6}

Neural networks excel as pattern recognition algorithms, particularly for scenarios involving unavoidable noise and inherent ambiguities in acquired transient images \cite{lecun1998,su2016} (see Chapter~\ref{chap:flat_homogeneous_material_detection}). Solid contaminants produced impure mixtures that manifested either visibly through color or texture changes, or remained invisible in certain cases. For visible contamination, the resulting impurity manifestation was classified into three fluid texture categories: homogeneous, heterogeneous (colloid), or heterogeneous (suspension). In cases where contamination was invisible or less obvious, the same fluid composition classes were assigned. This categorization scheme was applied to solid, liquid, and gaseous impurities.

An initial multi-label classifier was employed to investigate relationships between all materials involved in the study prior to binary classification \cite{bishop2006}. The impurity samples comprised 17 classes (Table~\ref{tab:6_1}) from both material and non-material contaminants, sourced from both transparent and translucent containers. The intent of the multi-label model was to investigate which features were most similar and to what extent the algorithm could distinguish between classes. Control and non-material impure mixtures from transparent and translucent containers exhibited overlapping feature distributions. Most features (Figure~\ref{fig:6_12}) could be predicted except for variables with highly similar or overlapping characteristics.

With only one valid state of purity defined, binary CNN classifiers were designed to model all possible impure states from six classes \cite{lecun1998}. The first classifier distinguished pure control samples from all other samples, enabling transient images to be interpreted as either pure consumable milk or contaminated. The persisted model was employed for several validation tests. First, the classifier sought to distinguish elusive contamination or impurities that do not manifest as color or consistency changes. Impure classes fitting such elusive manifestations were identified as uncolored soluble solid and uncolored soluble liquid contaminants (e.g., hydrogen peroxide \cite{gregory1961}) from active material contamination classes. Both contaminants showed no color or consistency change after mixing with pure samples, resulting in invisible impurity states. Control samples were maintained from the former retrained binary model.

Additional validation tests evaluated the generalisability of purity features across other pure samples of varying fat content, as well as pure samples of the same fat content sourced from different manufacturers. Finally, different milk products (milk products other than fresh homogenised milk) were validated. Two consumable milk products (one with viscous cream texture and another with liquid cream texture) and one non-consumable milk product (white viscous texture) were validated using the pretrained binary classifier.

% NOTE: Citation flag - The original text cited [6] for hydrogen peroxide, but reference [6] in ch6_refmaps.rtf corresponds to michalski2007 (milk homogenisation), while reference [5] corresponds to gregory1961 (hydrogen peroxide effect on milk). The citation has been corrected to gregory1961

\begin{figure}[htbp]
    \centering
    \includegraphics[width=0.8\textwidth]{figures/6.12.png}
    \caption[(experiment 2) Confusion matrix for converging multi label classifier for all materials considered for the study.]{\textbf{(experiment 2) Confusion matrix for converging multi label classifier for all materials considered for the study.} Both pure and impure materials.}
    \label{fig:6_12}
\end{figure}

\newpage
\section{Assumptions}
\label{sec:assumptions_chapter6}

The following hypotheses served as the foundation for formulating the research questions and guided the design of the conducted experiments:

\begin{enumerate}
    \item Homogenized milk, due to its uniform particle distribution, may be more readily characterised using electromagnetic sources compared to inhomogeneous milk.
    
    \item Visible contamination of homogenised milk that affects composition can be detected through analysis of time-resolved transient signals.
    
    \item Invisible contamination of homogenised milk that affects composition can be detected through analysis of time-resolved transient signals.
    
    \item Achromatic signal processing may offer greater computational efficiency than spectroscopic methods for milk purity assessment, as inferred from experimental processing costs (spectroscopy was not directly assessed).
    
    \item Optical milk purity detection can be performed economically and reliably using time-of-flight (ToF) detectors.
\end{enumerate}

\newpage
\section{Experiments' overview}
\label{sec:experiments_overview_chapter6}

Following the problem definition established in the preceding sections, experimental preparation begins with close-range signal acquisition from the material of interest (MOI) at approximately 20~cm distance. This range selection is motivated by two primary considerations: (1) practical product-to-user proximity estimates for real-world deployment scenarios, and (2) minimisation of measurement ambiguities that pose significant challenges in long-range detection (see Chapter~3) \cite{jungerman2022,tontini2020}. At this close range, distance--albedo ambiguities and view-angle dependencies are substantially reduced compared to multi-meter ranging configurations, enabling more reliable signal profile extraction \cite{becker2023,su2016}.

Three pattern recognition approaches are evaluated across the experimental suite: principal component analysis (PCA) for dimensionality reduction and feature visualisation \cite{wold1987}, binary neural network classifiers for purity assessment, and multi-label neural network classifiers for multi-class material discrimination \cite{lecun1998,bishop2006}. Each experiment is validated across multiple signal processing pipelines to ensure consistency and reliability of results, with preprocessing strategies adapted from Chapter~4 to handle ambient noise estimation, baseline correction, and transient histogram normalization \cite{tontini2024}.


\subsection{Materials}
\label{subsec:materials_chapter6}

This section details the non-control sample materials employed in the experimental framework. Two primary contaminant categories were prepared: material-based and non-material contaminants, both of which were actively introduced for experimental purposes \cite{joolaei2025,chandra2024}. Each category encompasses subcategories of material and non-material impurities, enabling systematic evaluation of detection capabilities across diverse contamination scenarios.

\subsubsection{Material contaminant samples}
\label{subsec:material_contaminant_samples_chapter6}

Material impurity classification was established through systematic investigation of contamination mechanisms affecting fluid samples. Contaminants were categorised by their physical state (solid, liquid, or gas) and further classified according to the optical properties of the resulting mixture, specifically color and consistency. This dual classification scheme enabled unique impurity class assignments for each contaminant state, with six subclasses derived for active material contamination.

Solid impurities were classified into six categories based on their interaction with the control sample. Colored solid contaminants produced three distinct mixture consistencies: homogeneous, heterogeneous colloidal, and heterogeneous suspension. Similarly, uncolored solid contaminants also yielded three mixture consistencies: homogeneous, heterogeneous colloidal, and heterogeneous suspension. 

The first solid contaminant (hm\_as\_tp\_sch\_ic) was a minimal quantity of colored solid impurity that, upon uniform mixing and settling, produced an instant color change with consistent texture, resulting in a homogeneous colored mixture. The second contaminant (hm\_as\_tp\_suh\_gs) was an uncolored solid impurity that, after uniform mixing and settling, exhibited no observable color or consistency changes, classified as a solid contaminant with uncolored consistent texture. The third contaminant (hm\_as\_tp\_suic\_swb) was an uncolored insoluble solid that produced an inhomogeneous colloidal mixture without visible color changes. When replaced with a colored insoluble solid (hm\_as\_tp\_scic\_bc), a similar inhomogeneous colloidal mixture was achieved, but with visible color changes in the resulting impure mixture. The fifth and sixth solid contaminants (hm\_as\_tp\_scis\_fw, hm\_as\_tp\_suis\_tp) differed in color but were both insoluble solid materials that remained suspended in the control sample after mixing.

Liquid impurities were classified using the same matter--matter interaction metrics, resulting in six distinct mixture types: hm\_al\_tp\_sch\_tcp, hm\_al\_tp\_scic\_lj, hm\_al\_tp\_scis\_oo, hm\_al\_tp\_suh\_hp, hm\_al\_tp\_suic\_pj, and hm\_al\_tp\_suis\_cp. Gas impurities exhibited fundamentally different behaviour due to their unusual nature as milk contaminants \cite{joolaei2025,chandra2024}. Consequently, gas impurity classes were limited to three categories, all producing heterogeneous mixtures without regard to color: hm\_ag\_tp\_scic\_air, hm\_ag\_tp\_scic\_spaint, and hm\_ag\_tp\_scis\_sperf. Comprehensive details of all active material impurities are provided in Table~\ref{tab:6_1}.

\subsubsection{Non-material contaminant samples}
\label{subsec:non-material_contaminant_samples_chapter6}

Non-material contamination arises from storage conditions and temporal degradation rather than the introduction of foreign matter. Active non-material impurity samples were prepared through two distinct mechanisms. The first category (hm\_UC\_tp) involved storing milk under suboptimal temperature conditions, resulting in physical transformations such as clotting and suspension that alter the fluid texture (Figure~\ref{fig:6_14}) \cite{lott2023,burgwald1947}. These temperature-induced changes compromise the structural integrity of the milk, rendering it impure despite the absence of external contaminants. The second category (hm\_UC\_tl) comprised fresh samples stored at recommended temperatures (approximately 5--9~°C) but maintained beyond their best-before or expiry dates \cite{lott2023,burgwald1947}. Unlike the temperature-compromised samples, these temporally degraded samples maintain visual similarity to control samples, appearing pure despite potential quality deterioration that may not be immediately observable.

\begin{figure}[htbp]
    \centering
    \includegraphics[width=0.21\textwidth]{figures/6.14.png}
    \caption[(experiment 3) active non-material contaminated milk sample by poor temperature resulting to a colloid]{\textbf{(experiment 3) active non-material contaminated milk sample by poor temperature resulting to a colloid}}
    \label{fig:6_14}
\end{figure}

\subsubsection{Control samples}
\label{subsec:control_samples_chapter6}

In addition to material and non-material contaminants, standard control material was sourced from fresh homogenised whole milk with 3.7\% fat content (brand 02). Test samples of 100~mL were prepared in transparent plastic bottles. Following validation of the standard binary model for purity assessment, several validation tests were conducted using untrained samples to evaluate model generalisability.

Validation experiments investigated two primary aspects of model performance. First, generalisability across different manufacturers was assessed by sourcing control samples from three additional consumable milk manufacturers beyond the standard control. Second, performance across varying fat contents was evaluated using control samples with fat contents of 0.3\%, 1.8\%, and 14\%, in addition to the standard 3.7\% fat control \cite{walstra1999,aernouts2015,melfsen2012}.


\subsection{Non-consumable milk products}
\label{subsec:non-consumable_milk_products_chapter6}

At the outset of this work, it was hypothesized that defining milk as a product is inherently ambiguous. This claim is supported by product-based impurity assessments that are impurity-specific and difficult to generalise. Different milk products may undergo significant transformation that alters their intrinsic properties as milk. Consequently, the extent to which similar but different milk products differ in composition was investigated. Condensed (see Figure~\ref{fig:6_9}) and evaporated (see Figure~\ref{fig:6_10}) milk, along with non-consumable milk products, were sourced as alternative controls in validation experiments \cite{walstra1999,nielsen2013}.

Approximately eight different experiments were conducted to investigate the study claims. The initial experiment established proof of fluid homogeneity impact on modelling as well as the learnability of different feature groups in the selected control materials with SPAD. The outcome of this study was instrumental in deciding that all samples be taken from homogenised fluid sources. Additionally, the combination of certain control samples was validated from outcomes of feature learnability. Subsequent experiments explored general purity metrics for visible and invisible contaminations. Further tests investigated purity of material across different fat contents, manufacturers, and non-consumable milk products \cite{aernouts2015,melfsen2012,joolaei2025}.

\subsection{Milk composition integrity}
\label{subsec:milk_composition_integrity_chapter6}

The complex composition of milk renders it highly sensitive to temperature deviations, with even slight departures from recommended storage conditions (<4~°C) potentially inducing significant state transformations that compromise both safety and compositional integrity \cite{walstra1999}. Such transformations include reduced enzyme activity and accelerated bacterial growth, which threaten the validity of ground truth data. To ensure credible ground truth sourcing, all milk samples were maintained at recommended refrigeration temperatures (2--3~°C) and removed from storage only at the commencement of data collection.

Each SPAD illumination cycle and signal return recording requires a maximum of 1~second (typically <1~s), with 2000 observations collected per sample. This yields a total data collection duration of at most 2000~seconds per sample. Data collection was conducted at ambient room temperature (~23--25~°C). Significant compositional changes in consumable milk at these temperatures typically commence after approximately two hours \cite{lott2023,burgwald1947}. The maximum collection duration of 2000~seconds (0.556~hours, or ~33~minutes) falls well within this safety window, ensuring that no material changes capable of compromising milk safety or compositional integrity occur during ground truth acquisition. All ground truth samples were sourced within the manufacturer's stated use-by or expiry dates.

\subsection{Instruments and setup}
\label{subsec:instruments_setup_chapter6}

Several optical methods exist for inferring fluid purity, each with distinct limitations. Human sensory assessment is slow, complex, and potentially unsafe. While chemical purity tests are effective, their consumptive nature makes them invasive compared to optical alternatives. Near-infrared (NIR) spectroscopy, the most prominent optical technique, tends to be complex and costly to implement \cite{aernouts2015,melfsen2012,biggs1967,barbano1989}. Spectrophotometric methods, such as spatially resolved visible-range RGB image classification, face additional limitations: high image resolution requirements and inability to detect elusive impurities where contaminated mixtures exhibit no observable color or consistency changes \cite{chandra2024,joolaei2025}.

Physically, contaminated mixtures may manifest changes in color, consistency, or volume levels. However, due to the fallible nature of human vision, very subtle contaminations can be missed. A typical example is hydrogen peroxide contamination \cite{gregory1961}, which human vision or visible light detection can hardly detect. From contaminant analyses (see Section~\ref{subsec:material_impurities}), it is evident that not all impure mixtures exhibit color or structural manifestations in visible light. This poses a severe challenge to systems that rely solely on visual changes (spatially resolved RGB or spectrophotometric solutions) to predict impurity.

Single-photon avalanche diodes (SPADs) (the same instrument used in Chapter~4) with an invisible illumination source are presented to address these limitations. With pulsed illumination and very high gain, SPAD sensors deliver extreme sensitivity levels that ensure reliable perception of even the most subtle changes in composition, whether visible to human eyes or not. Time-resolved signals encode reflectivity and albedo that translate to unique energy--matter interactions between source and target. Unlike NIR spectroscopy, a single calibration model (see Chapter~4) is sufficient to represent purity variables from non-spectral source scattering \cite{nielsen2013}. Temporal transient images also provide important information about the micro-structural properties of the material-of-interest (MOI) particles in a way that is complex to achieve by traditional RGB imaging \cite{su2016}. The instrument offers a non-invasive, reliable, low-budget, and easy-to-implement means to distinguish pure consumable milk state from every contamination source, whether material or non-material.

\subsection{Data collection}
\label{subsec:data_collection_chapter6}

Milk samples were collected on random days from multiple manufacturers across different fat contents. For each product, the ground truth state of safe consumption served as the qualifier for experimental sampling. Key details regarding the maintenance of ground truth composition integrity for all samples are described in Section~\ref{subsec:milk_composition_integrity_chapter6}. Eight experiments were conducted to investigate the extent to which SPAD sensors can detect impurities in milk products.

Data were prepared from pure and impure milk samples across different product types. Impure samples contained both material and non-material impurities. Practically, impurity detection thresholds are measured in parts per million (ppm) as the standard unit for tracing contaminants. The detection range spans from 100~ppm (approximately 100~mg per 1~kg of milk, or 0.01~g per 100~g of milk, equivalent to 0.01\% impurity) to 500~ppm (approximately 500~mg per 1~kg of milk, or 0.05~g per 100~g of milk, equivalent to 0.05\% impurity). This range is based on the assumption that milk samples containing $\geq$0.05\% impurity are considered contaminated \cite{goodman1986}. The stated thresholds are useful for extreme sensitivity purity tests. However, to ensure strong variations in ground truth signatures of contamination, wider concentration ranges were employed.

For material contaminants, approximately 75\% of 100~mL transparent bottles and 85\% of 1~L original packaging bottles were filled with homogenised whole milk samples in ground truth safe-to-consume states. Impurities were added in increments from 500~ppm to 10,000~ppm (0.05\%--1\%) for sample data collection. For each of three matter states, contaminants were subcategorised by heterogeneity and color profile of the resulting contamination. In total, six standard contaminant classes were sampled per matter state.

Non-material impurities were actively sourced from intentional storage of milk under poor temperature conditions and until long after the best-before date \cite{lott2023,burgwald1947}. In addition to contaminated mixtures, pure samples were obtained from different milk manufacturers and fat concentrations. Four commercial milk manufacturers were selected for pure samples. Two of the pure sample sets had three fat categories: whole milk, semi-skimmed, and skimmed. The other two sets included one unhomogenised sample and another whole milk sample \cite{walstra1999,michalski2007}.

The transient image acquisition setup consists of a SPAD sensor with an integrated non-spectral source positioned approximately 20~cm from the target material (see Figure~\ref{fig:6_3}). Setup lighting was minimal. Approximately two thousand transient images were acquired per class for preprocessing.

\subsection{Data preprocessing}
\label{subsec:data_preprocessing_chapter6}

Despite close-range signal detection (see Figure~\ref{fig:6_3}), Poisson and ambient noise contributions remain unavoidable in SPAD-based transient measurements \cite{piron2020}. As established in Chapter~\ref{chap:pawd}, transient images can exhibit ambiguities arising from view angle variations, material-of-interest (MOI) container permeability, and surface curvature effects \cite{jungerman2022,tontini2024}. Data preprocessing addresses these limitations through algorithmic noise estimation and subtraction: ambient noise is estimated from initial histogram bins where target signals do not accumulate, while Poisson noise is mitigated through mean value estimation and subtraction across all time bins \cite{tontini2020}. These preprocessing steps substantially reduce noise contributions and ambiguity-related artefacts. However, subtle nonlinearities inherent in the temporal signal distributions are expected to persist, necessitating nonlinear learning algorithms for effective feature extraction and classification \cite{lecun1998,su2016}.

\subsection{Learning algorithm}
\label{subsec:learning_algorithm_chapter6}

As established in Section~\ref{subsec:data_preprocessing_chapter6}, the inherent complexity of noise and ambiguities within transient images naturally lends itself to non-linear modelling techniques such as deep neural networks. While preprocessing algorithms effectively mitigate ambient and Poisson noise contributions through baseline subtraction and mean estimation \cite{tontini2020}, residual non-linearities persist in the temporal signal distributions due to the quantum nature of photon detection and the complex interactions between material optical properties and measurement geometry \cite{jungerman2022,tontini2024}. These non-linear relationships arise from the convolution of multiple signal components—laser pulse profiles, detector response functions, and material-dependent scattering characteristics—which cannot be fully captured by linear dimensionality reduction or feature extraction methods \cite{lecun1998,su2016}. Deep neural networks provide the necessary non-linear mapping capabilities to extract material-discriminative features from these complex temporal signatures, enabling robust classification despite the presence of measurement ambiguities and residual noise artefacts \cite{ackley1985,bishop2006}.

\subsection{The PCA technique}
\label{subsec:the_pca_technique}

Transient images are first processed through an autoencoder, which upsamples them into larger-dimensional feature vectors \cite{ackley1985,lecun1998,su2016}. The resulting encoded transient vectors are then fitted to a Principal Component Analysis (PCA) model \cite{wold1987,chandra2024}. PCA is an unsupervised dimensionality reduction technique that identifies orthogonal directions of maximum variance in the data \cite{wold1987,bishop2006}. The algorithm computes the covariance matrix of the encoded transient vectors, performs eigendecomposition to extract principal components (eigenvectors), and orders them by their corresponding eigenvalues, which represent the amount of variance explained by each component \cite{wold1987}. The first principal component captures the direction of maximum variance in the encoded transient space, the second captures the next highest variance orthogonal to the first, and so on \cite{bishop2006}.

The top two to three principal components are extracted and used for visualisation, as they typically capture the majority of the variance in the transient signal space while remaining interpretable in a low-dimensional representation \cite{wold1987}. These principal components represent the dominant temporal patterns that distinguish different materials or objects in the transient domain. By projecting the encoded transient vectors onto these principal components, each sample is represented as a point in a reduced-dimensional space, where the coordinates correspond to the projections (scores) along each principal component axis \cite{bishop2006}.

These projections are visualised in a 3D scatter plot (Figure~\ref{fig:6_4}), where each axis represents one of the top principal components, and each point corresponds to a transient image sample. The spatial distribution of points in this 3D space reveals clustering patterns that indicate material similarity, with samples from the same material type typically forming distinct clusters \cite{chandra2024}. This visualisation enables qualitative assessment of material separability in the transient domain and provides insight into how well the temporal characteristics of different materials can be distinguished using the extracted principal components \cite{wold1987,bishop2006}.

\newpage
\section{Multi-zone milk classifier (2D CNN)}
\label{sec:multi_zone_milk_classifier_chapter6}

As described in Section~\ref{sec:multizone_signal_classifier}, the same multizone input classifier architecture is implemented for milk purity assessment. The proposed network is a feedforward 2D CNN designed for material classification of planar surfaces. The input tensor has shape $(N_b, 3, 16, 16)$, where $N_b$ is the batch size and the three channels correspond to three stacked transient-derived modalities per zone \cite{su2016,lecun1998}. For the contaminants' classifier, the model outputs logits for all purity and impurity classes (Experiment 2, Figure~\ref{fig:6_12}) \cite{bishop2006}. The remaining models are binary classifiers (Experiments 3--7) designed to assess purity across different dimensions \cite{bishop2006}. To ensure methodological validity, complementary analytical techniques such as principal component analysis (PCA) have been employed to validate relationships between different variables and establish confidence in the results \cite{wold1987,chandra2024}.

\subsubsection{FoV of camera relative to milk container}
\label{subsec:fov_of_camera_relative_to_milk_container}

The field of view (FoV) configuration of the SPAD sensor relative to the milk container is critical for ensuring accurate ground truth data acquisition while minimising contamination from background targets. As established in Chapter~\ref{chap:flat_homogeneous_material_detection} (Section~\ref{subsec:instruments_setup}), the SPAD illumination source exhibits a field of view of approximately 30° (see Figure~\ref{fig:4_12}), which defines the angular coverage of the active illumination beam. During data collection, the SPAD sensor was positioned at randomized distances ranging from 20~cm to 30~cm from the milk container surface, ensuring variation in acquisition geometry while maintaining consistent signal quality.

At the minimum distance of 20~cm, a 30° field of view subtends a circular coverage area with a diameter of approximately 10.7~cm ($2 \times 20 \times \tan(15°) \approx 10.7$~cm), while at the maximum distance of 30~cm, the coverage diameter extends to approximately 16.1~cm ($2 \times 30 \times \tan(15°) \approx 16.1$~cm). This angular coverage ensures that the SPAD sensor captures a sufficient portion of the milk container surface—typically covering the central region of standard 100~mL transparent bottles (diameter $\approx$ 5--6~cm) and 1~L packaging bottles (width $\approx$ 7--9~cm)—while maintaining adequate signal intensity for reliable transient histogram acquisition. The field of view is sufficiently wide to encompass the primary region of interest where milk material properties are most representative, accounting for the container's curvature and ensuring that multiple zones of the 4×4 SPAD array receive signal returns from the milk surface.

Conversely, the 30° field of view is deliberately narrow enough to exclude background targets that could contaminate the ground truth measurements. In typical laboratory or controlled environment settings, background objects such as walls, tables, or other containers are positioned at distances substantially greater than the 20--30~cm sensor--target separation. At these extended distances, even objects within the angular field of view contribute minimal signal due to the inverse square law governing light propagation \cite{griffiths2017}, and their temporal signatures are temporally separated from the milk container returns in the transient histogram due to the direct time-of-flight measurement principle. The narrow angular coverage thus ensures that the dominant signal contributions originate from the milk container surface, while background artefacts are effectively suppressed through both geometric isolation and temporal discrimination. This configuration is essential for maintaining ground truth integrity, as any contamination from background targets would introduce spurious temporal features that could compromise the material classification accuracy of the purity assessment models.


\section{Experiment 1 Homogenous vs inhomogeneous milk}
\label{sec:experiment_1_chapter6}

This study investigated the effect of material homogeneity on purity assessment using SPAD-sourced transient images. Homogeneity plays a critical role in material analysis, particularly in relation to signal interference, as established in Chapters~\ref{chap:pawd} and~\ref{chap:flat_homogeneous_material_detection}. Beyond interference effects, inhomogeneous milk exhibits colloidal state volatility, where the material texture at time $t_0$ may differ substantially from that at time $t_1$ following slight movement or positional changes. This temporal variability is not observed in homogenised milk, which maintains consistent particle distribution across time and position \cite{walstra1999,michalski2007,aernouts2015}.

To understand the impact of homogeneity on purity assessment, a fundamental question is proposed: what do pattern recognition algorithms have in common with milk production and consumption? The answer is consistency: both pattern recognition systems and milk producers require consistent material properties. For consumers, consistency eliminates the need for manual agitation and ensures uniform distribution of milk constituents. For learning algorithms and milk analysis systems \cite{chandra2024}, consistency is essential for accurate signal measurements that are not compromised by volatile scattering from fat globules or temporal variations in particle distribution.

At the beginning of this project, features were extracted from both homogenised (14\% fat) and non-homogenised (14\% fat) milk samples from the same manufacturer (brand 0), as shown in Figure~\ref{fig:6_8}. Extracted features were upsampled using a neural network encoder, and the resulting encoded vectors were fitted to a PCA model. Visualization of the first three principal components (Figure~\ref{fig:6_15}) revealed minimal overlap between homogenised and non-homogenised milk features, despite both samples originating from the same producer. In contrast, Figure~\ref{fig:6_16} demonstrates greater feature overlap among samples of the same type, indicating more generalizable purity metrics within homogeneous samples. Based on these findings, only homogenised samples were considered in subsequent purity assessment experiments.

It is important to clarify that material inhomogeneity is not inherently a limitation; this is demonstrated in Chapter~\ref{chap:spatiotemporal_object_detection}, where multimaterial targets have been successfully modelled. However, the behaviour of inhomogeneous solids differs fundamentally from that of inhomogeneous liquids. Inhomogeneous solids maintain the same material composition configuration across different positions and over time. In contrast, inhomogeneous liquid compositions undergo continuous reconfiguration whenever any force acts upon the material, whether through container shaking or positional displacement. This fundamental difference explains why modelling inhomogeneous milk has been challenging since the early days of milk purity testing \cite{biggs1967,barbano1989,walstra1999}.

\begin{figure}[htbp]
    \centering
    \includegraphics[width=0.8\textwidth]{figures/6.15.png}
    \caption[Fresh homogenised vs inhomogenised milk (same manufacturer, 14\% fat each) showing weak overlap of features from either classes.]{\textbf{Fresh homogenised vs inhomogenised milk (same manufacturer, 14\% fat each) showing weak overlap of features from either classes.} Explained variance ratio of three principal components;[0.21196705, 0.15826207, 0.11945989]}
    \label{fig:6_15}
\end{figure}

\begin{figure}[htbp]
    \centering
    \includegraphics[width=0.8\textwidth]{figures/6.16.png}
    \caption[(Experiment 5) PCA of two homogenised whole milk products samples (3.7\% fat each) from different manufacturers indicating strong overlap of group variables.]{\textbf{(Experiment 5) PCA of two homogenised whole milk products samples (3.7\% fat each) from different manufacturers indicating strong overlap of group variables.} Explained variance ratio of three principal components; [0.18387789, 0.1645287, 0.10170276]}
    \label{fig:6_16}
\end{figure}

\section{Experiment 2  Feature learnability of pure and impure  samples}
\label{sec:experiment_2_chapter6}

While deep neural networks are powerful pattern recognition tools, they can only learn features that actually exist in the data \cite{lecun1998,bishop2006}. The accuracy and reliability of recognising existing patterns are fundamentally dependent on the consistency and distinctiveness of these features. Without accurate feature distinction, the reliability of any purity detection system built upon these features will be compromised.

In this experiment, SPAD detectors were used to collect transient images from all material samples, both pure and impure, with the objective of accurately classifying them to assess the reliability of the purity detection model. Seventeen 100~mL samples (brand 2) were reserved for active contamination. Two impurity classes were non-material contaminants, while the remainder comprised material agents from solid, liquid, and gaseous matter. The resulting colour and texture of each contaminated mixture were documented.

The experiment employed a multi-label classifier using all material-contaminated samples, non-material-contaminated samples, and control samples as distinct classes. The classifier achieved approximately 80\% convergence (Figure~\ref{fig:6_12}). Gaseous impurity mixtures containing liquid particles (hm\_ag\_tp\_scic\_air, hm\_ag\_tp\_scic\_spaint, hm\_ag\_tp\_scis\_sperf) exhibited feature interchangeability, which was attributed to the similarity of their resulting mixture compositions.

An additional interesting overlap occurred between transparent and translucent control samples (hm\_UC\_tl, hm\_UC\_tp, hm\_C\_tl, hm\_UC\_tp). Both sample types utilise containers with higher signal penetrability, suggesting that scattering effects from the milk container may have minimal or less significant impact on the actual profile of the milk content \cite{klocek1991,walstra1999}. This understanding formed the basis for successfully combining both transparent and translucent containers of milk samples in the binary purity detection model.

Other weak overlaps were observed during intermediate training stages. However, these overlaps disappeared once the model reached full convergence (Figure~\ref{fig:6_17}). To ensure that the high model accuracy was not due to overfitting, variations among materials were further verified through principal component analysis (PCA) \cite{wold1987,chandra2024}. The results indicate learnability of all features, even among material-of-interest (MOI) samples with strong feature similarities, confirming the discriminative capability of the learned representations.

\begin{figure}[htbp]
    \centering
    \includegraphics[width=0.8\textwidth]{figures/6.17.png}
    \caption[(experiment 2) Confusion matrix of fully converged multi label classifier for all materials considered for the study.]{\textbf{(experiment 2) Confusion matrix of fully converged multi label classifier for all materials considered for the study.}}
    \label{fig:6_17}
\end{figure}

\section{Experiment 3 Purity detection of homogenised milk with SPADs}
\label{sec:experiment_3_chapter6}

Milk purity assessment has been an enduring concern in dairy science and quality control, driven by the need to detect both material contaminants and degradation of otherwise consumable products \cite{walstra1999,goodman1986}. Among available technologies, optical techniques (including visible and near-infrared (NIR) spectroscopy and imaging-based methods) are widely adopted because they offer rapid, non-destructive measurements of bulk optical properties \cite{aernouts2015,melfsen2012,joolaei2025}. However, conventional NIR spectrometers are often costly and operationally complex, while image-based approaches remain sensitive to illumination conditions, camera characteristics, and subtle changes in appearance that may not correlate reliably with contamination state \cite{sun2009,chandra2024}.

This experiment investigates whether a high-gain time-of-flight detector based on a Single-Photon Avalanche Diode (SPAD) can provide a more robust basis for purity assessment by exploiting transient photon transport rather than solely relying on spectral or RGB intensity cues \cite{nielsen2013,gyongy2021}. Material impurities were drawn from the standardized contaminant taxonomy described in Table~\ref{tab:6_1}, spanning solids, liquids, and gases at controlled concentrations. Non-material impurities were introduced via two mechanisms: (i) exposure to poor storage temperatures, producing visible clotting and suspension effects (Figure~\ref{fig:6_14}), and (ii) prolonged storage beyond the best-before date under recommended refrigeration, leading to temporal degradation that can remain visually subtle (Figure~\ref{fig:6_18}) \cite{lott2023,burgwald1947}. Control samples comprised homogenised whole milk with 3.7\% fat in both transparent and translucent containers, ensuring consistent bulk optical properties while preserving the container permeability conditions established in earlier experiments.

In this context, ``pure'' milk is operationally defined as the consumable ground-truth state delivered by the manufacturer, while any deviation induced by material or non-material impurities is labelled ``impure''. The purity label is therefore binary: each sample is assigned a single true state (pure or impure), and all alternatives are treated as false. A binary convolutional neural network (CNN) classifier, using the multizone 2D architecture described in Section~\ref{sec:multi_zone_milk_classifier_chapter6}, was trained to recognise optical transients corresponding to the consumable (pure) state as the positive class, with all contaminated or degraded states forming the negative class \cite{lecun1998,bishop2006}. Training and evaluation were carried out on homogenised milk only, leveraging the homogeneity advantages established in Experiment~1 for stable temporal scattering signatures \cite{walstra1999,michalski2007,aernouts2015}.

The classifier achieved approximately 98\% training accuracy after convergence, indicating that the transient features extracted from SPAD measurements were highly separable for the pure versus impure classes (Figure~\ref{fig:6_19}). Validation on previously unseen homogenised milk samples yielded an accuracy of about 93.33\%, demonstrating that the learned decision boundary generalises beyond the training set while maintaining a clear separation between consumable and non-consumable states (Figure~\ref{fig:6_20}). To probe model behaviour under out-of-distribution conditions, invalid transient data from solid surfaces were supplied to the classifier. In these tests, the network assigned low purity scores, with confidence exceeding 80\% that the inputs did not correspond to valid milk transients (Figure~\ref{fig:6_21}), reinforcing its ability to reject non-milk inputs and to act as a robust front-end purity detector.

The achieved performance (approximately 98\% training accuracy and 93.33\% validation accuracy) is consistent with previous relevant work in material classification using time-of-flight measurements. The Princeton study by Su et al.~\cite{su2016} demonstrated that machine learning models can achieve high accuracy (80.9\% validation accuracy using RBF SVM) for material classification using raw time-of-flight measurements, validating the feasibility of learning material-discriminative features from temporal signal characteristics. Similarly, the MDPI Sensors paper by Becker and Koerner~\cite{becker2023} reported classification performance of 96\% accuracy for plastic material classification using optical parameter features measured with direct time-of-flight depth sensors. The consistency of high performance across these independent studies, employing different machine learning approaches and material classification tasks, provides strong evidence that the observed accuracy in this work stems from the inherent material-discriminative information encoded in transient time-of-flight signals, validating the purity classifier's performance as representative of the general capability of machine learning methods for time-of-flight-based material classification tasks.

\begin{figure}[htbp]
    \centering
    \includegraphics[width=0.7\textwidth]{figures/6.18.png}
    \caption[(experiment 03) active non-material contaminated milk sample by expiration resulting to an unstable colloid mixture]{\textbf{(experiment 03) active non-material contaminated milk sample by expiration resulting to an unstable colloid mixture}}
    \label{fig:6_18}
\end{figure}

\begin{figure}[htbp]
    \centering
    \includegraphics[width=0.7\textwidth]{figures/6.19.png}
    \caption[(experiment 3) confusion matrix for classification of material and non material impurities (~98\%)]{\textbf{(experiment 3) confusion matrix for classification of material and non material impurities (~98\%)}}
    \label{fig:6_19}
\end{figure}

\begin{figure}[htbp]
    \centering
    \includegraphics[width=0.7\textwidth]{figures/6.20.png}
    \caption[validation test on pre-trained classifier with untrained samples.]{\textbf{validation test on pre-trained classifier with untrained samples.} (~93.33\%)}
    \label{fig:6_20}
\end{figure}

\begin{figure}[htbp]
    \centering
    \includegraphics[width=0.7\textwidth]{figures/6.21.png}
    \caption[(experiment 3) confusion matrix for validation of experiment 3 model with invalid data; samples were transient from solid surfaces.]{\textbf{(experiment 3) confusion matrix for validation of experiment 3 model with invalid data; samples were transient from solid surfaces.} ~18\% show model is over 80\% confident that input data is invalid.}
    \label{fig:6_21}
\end{figure}

\section{Experiment 4 Purity detection of invisible contamination of homogenised milk with SPADs}
\label{sec:experiment_4_chapter6}

Building on the binary purity classifier developed in Experiments~2--3, this experiment probes one of the most critical limitations of conventional spectrophotometric and imaging-based purity assessment: the inability to reliably detect contamination that produces no discernible change in visible colour or texture.\footnote{Here, ``invisible'' contamination refers to impurities that leave the macroscopic appearance of the bulk sample visually indistinguishable from the pure reference state.} In contrast to chemical assays or spectrometric methods that exploit wavelength-dependent absorption or scattering signatures \cite{aernouts2015,joolaei2025}, spatially resolved RGB imaging depends almost entirely on changes in apparent surface or bulk appearance. As a result, its performance degrades sharply when (i) image resolution and camera characteristics are suboptimal, (ii) ambient illumination is poorly controlled, or (iii) contaminants alter the microstructure or composition of the fluid without producing a visible macroscopic cue.

To isolate this challenge, two representative contaminants were selected from the standardized taxonomy: a solid-phase impurity (hm\_as\_tp\_suh\_gs) and a liquid-phase impurity (hm\_al\_tp\_suh\_hp). For each $\sim$100~ml aliquot of pure homogenised milk, approximately 5~g~($\pm$2~g) of the chosen impurity was added, after which the mixture was allowed to mix thoroughly and then settle. Throughout this process, the resulting impure mixtures remained visually indistinguishable from the corresponding pure controls when inspected under standard ambient lighting, thereby emulating an ``invisible'' contamination scenario that is challenging for purely photometric or colour-based methods.

Transient SPAD measurements were then acquired from both the invisibly contaminated mixtures and their matched pure controls, using the same multizone acquisition protocol and binary labeling strategy as in Experiment~2. Each contaminated sample was assigned to the impure class, while its corresponding uncontaminated homogenised milk sample served as a pure control. The previously trained SPAD-based purity classifier was reused without architectural modification, and the invisibly contaminated measurements were treated as an additional out-of-distribution regime for evaluation rather than for retraining.

On this validation set, the classifier achieved an accuracy of approximately 96\% in discriminating pure from invisibly impure homogenised milk, despite the absence of any obvious visual cue in the contaminated mixtures. To better understand the feature-space structure underlying this performance, transient feature vectors from pure and invisibly impure samples were projected into a principal component analysis (PCA) subspace \cite{wold1987}. Visualization of the first three principal components revealed a noticeable but incomplete separation between classes, with a region of overlap between pure and contaminated samples (Figure~\ref{fig:6_13}), indicating that the class boundary is moderately non-linear in the raw feature space.

\begin{figure}[htbp]
    \centering
    \includegraphics[width=0.7\textwidth]{figures/6.13.png}
    \caption[(experiment 4) 3D PCA of features from pure sample and invisibly impure milk samples (hydrogen peroxide impurity and sugar impurity).]{\textbf{(experiment 4) 3D PCA of features from pure sample and invisibly impure milk samples (hydrogen peroxide impurity and sugar impurity).} Higher dimensions reveal more separation of hydrogen peroxide contamination in comparison with sugar.}
    \label{fig:6_13}
\end{figure}

This overlap is consistent with the expectation that subtle microscopic changes in composition may produce only modest perturbations in the time-resolved scattering signatures, such that linear projections alone are insufficient for complete class separation. The strong performance of the binary SPAD-based classifier in this setting therefore underscores the value of deep neural networks in learning non-linear decision boundaries in high-dimensional transient feature spaces, enabling the detection of contamination that is effectively invisible to conventional RGB imaging. Alternative non-photometric purity assessment systems, such as NIR spectroscopy, can in principle detect similarly elusive adulteration patterns via their spectral fingerprints \cite{aernouts2015,joolaei2025}; however, these instruments are typically expensive, less economical for low-cost deployments, and demand a higher level of specialist expertise for calibration and operation. Within this context, the proposed SPAD-based approach offers a practical and scalable route to detecting invisible contamination in homogenised milk, while maintaining a compact form factor and low per-sample operational overhead.

\section{Experiment 5 Manufacturer generalisability of purity detection model}
\label{sec:experiment_5_chapter6}

Milk samples used in Experiments~2--4 were all sourced from a single manufacturer (brand~0). In practical settings, however, consumable milk is produced by multiple companies, each supplying nominally similar products (for example, homogenised whole milk with 3.7\% fat) but with potentially subtle process-dependent differences in microstructure and bulk optical properties \cite{walstra1999,aernouts2015}. This raises a natural question: does the binary purity detection model developed in Experiment~2 rely on manufacturer-specific signatures, or does it learn more generalizable purity features that transfer across brands producing the same product type?

To investigate this, a fifth experiment examined the generalisability of the encoded purity features across manufacturers. Principal component analysis (PCA) was first applied to transient feature vectors from two brands of homogenised milk with identical labelled fat content and similar fluid texture (brand~1, 3.7\% fat; brand~2, 3.7\% fat). The resulting projection showed substantial overlap between the feature distributions of both brands (Figure~\ref{fig:6_16}), indicating that the encoded purity metrics form a shared, non-linearly separable manifold rather than separating cleanly by manufacturer. This behaviour is consistent with the expectation that, for homogenised milks of comparable composition, the dominant time-resolved scattering response is largely invariant to manufacturer-specific processing within the operating range of the SPAD-based system \cite{nielsen2013}.

While this overlap suggests that the feature space is shared across brands, it does not by itself quantify how well the purity detection model trained on one brand will perform when evaluated on another. To address this, the original binary classifier from Experiment~2 was re-validated using a modified test set in which pure samples from brand~2 were replaced by pure samples from a different manufacturer (brand~1), while the impurity mixtures were kept unchanged. Under this cross-manufacturer configuration, the classifier achieved an accuracy of approximately 99\% (Figure~\ref{fig:6_22}), closely matching its performance on the original validation set and confirming that the learned decision boundary remains effective when the source of pure milk is changed.

These results jointly indicate that the proposed SPAD-based purity detection model exhibits strong manufacturer generalisability for homogenised milk with matched fat content. In practical terms, this means that the system can be deployed across a range of commercial milk brands without requiring brand-specific retraining or extensive recalibration, provided that the target products share similar compositional specifications.

\begin{figure}[htbp]
    \centering
    \includegraphics[width=0.8\textwidth]{figures/6.22.png}
    \caption[(experiment 5) validation test on pre-trained binary classifier with untrained homogenised milk samples of different manufacturers as pure milk.]{\textbf{(experiment 5) validation test on pre-trained binary classifier with untrained homogenised milk samples of different manufacturers as pure milk.} (~99\%)}
    \label{fig:6_22}
\end{figure}

\section{Experiment 6 Fat content generalisability of purity detection model}
\label{sec:experiment_6_chapter6}

Milk fat content significantly influences signal scattering in energy-matter interactions \cite{aernouts2015,walstra1999}. In non-homogenised whole milk, regions with concentrated or larger fat globules exhibit more intense scattering of the source signal. However, homogenised milk maintains consistent fat globule configuration (~1~$\mu$m) throughout the fluid \cite{walstra1999,michalski2007}, though the total fat percentage per unit volume may vary across samples, even within the same brand \cite{aernouts2015,melfsen2012}.

To investigate the relationship between purity variables across fat content variations, samples were varied across five milk fat levels: 14\%, 3.9\%, 3.7\%, 1.8\%, and 0.3\% (Table~\ref{tab:6_2}). Features were extracted and fitted to a principal component analysis (PCA) model \cite{wold1987,chandra2024}. Results indicated no linear separation among the samples' features, suggesting that fat content alone does not create distinct feature clusters in the purity detection space.

The purity detection model was subsequently tested for generalisability across homogenised milk samples of different fat contents. In the first validation test, control samples were replaced with homogenised milk containing 14\% fat. The model achieved a prediction accuracy of approximately 99\% (Figure~\ref{fig:6_23}, left). In a second test, the control was swapped with a sample of 1.8\% fat, recording 97\% model performance (Figure~\ref{fig:6_23}, centre). Finally, samples with low fat composition (0.3\%) were tested, yielding an accuracy of approximately 91\% (Figure~\ref{fig:6_23}, right).

The observed decline in accuracy with decreasing fat content is consistent with the Beer--Lambert law, which describes the linear relationship between the concentration of an absorbing or scattering species and the intensity of transmitted or returned illumination \cite{holt2016,elachi2021}. In this context, variations in milk fat concentration directly influence the scattering and absorption intensity of the returned source illumination, affecting the transient signatures captured by the SPAD detector \cite{nielsen2013,aernouts2015}. The reduced fat content in low-fat samples (0.3\%) results in diminished scattering interactions, leading to weaker transient signals that may approach the detection limits of the system, thereby reducing classification accuracy relative to higher-fat samples.

\begin{table}[htbp]
    \centering
    \begin{tabular}{lc}
    \toprule
    \textbf{Fat content (\%)} & \textbf{Validation accuracy (\%)} \\
    \midrule
    0.3 & $\sim$91 \\
    1.8 & $\sim$97 \\
    3.7 & $\sim$98 \\
    14.0 & $\sim$99.9 \\
    \bottomrule
    \end{tabular}
    \caption{Validation accuracy of the binary purity detection model across representative fat content levels in Experiment~6, corresponding to the results summarised in Figure~\ref{fig:6_23}.}
    \label{tab:6_2}
\end{table}

\begin{figure}[htbp]
    \centering
    \includegraphics[width=0.9\textwidth]{figures/6.23.png}
    \caption[Left: (experiment 6) validation test on pre-trained binary classifier with untrained homogenised fresh milk sample (14\% fat) as pure milk. (~99\%). Center: (experiment 6) validation test on pre-trained binary classifier with untrained homogenised fresh milk sample (1.8\% fat) as pure milk. (~97\%). Right: (experiment 6) validation test on pre-trained binary classifier with untrained homogenised fresh milk sample (0.3\% fat) as pure milk. (~91\%)]{\textbf{Left: (experiment 6) validation test on pre-trained binary classifier with untrained homogenised fresh milk sample (14\% fat) as pure milk. (~99\%). Center: (experiment 6) validation test on pre-trained binary classifier with untrained homogenised fresh milk sample (1.8\% fat) as pure milk. (~97\%). Right: (experiment 6) validation test on pre-trained binary classifier with untrained homogenised fresh milk sample (0.3\% fat) as pure milk. (~91\%)}}
    \label{fig:6_23}
\end{figure}

\section{Experiment 7a Product generalisability of purity detection model (consumable milk)}
\label{sec:experiment_7a_chapter6}

In the final stage of this work, the generalisability of the binary purity detection model (trained exclusively on fresh homogenised milk) is evaluated on two additional consumable products: condensed milk and evaporated milk (see Figures~\ref{fig:6_9} and~\ref{fig:6_10}). Principal component analysis of the extracted purity features from the control (fresh milk), condensed milk, and evaporated milk shows substantial overlap among the three classes (Figure~\ref{fig:6_7}), consistent with their shared underlying milk composition despite differences in processing and physical structure~\cite{aernouts2015,nielsen2013}. This overlap suggests limited linear separability in feature space and indicates the presence of non-linear decision boundaries across products.

To probe product-level generalization, a zero-shot evaluation is performed by applying the pre-trained binary purity model directly to evaporated and condensed milk without retraining. For evaporated milk, the model attains a prediction accuracy of approximately 99\%, as illustrated in Figure~\ref{fig:6_24a}. In contrast, the same model achieves only about 65\% accuracy on condensed milk (Figure~\ref{fig:6_24b}). These results are notable given that evaporated milk exhibits a relatively low viscosity and flow behaviour that are closer to the fresh control sample than to condensed milk, whereas the visual colour of condensed milk more closely resembles that of the control than that of evaporated milk.

Taken together, these observations suggest that the detector is more sensitive to texture- and structure-related cues (e.g., viscosity and surface morphology) than to colour differences when operating under the present experimental conditions~\cite{nielsen2013}. The high performance on evaporated milk is therefore consistent with its physical similarity to the control sample in terms of flow and surface appearance, while the markedly lower performance on condensed milk aligns with its distinct viscous texture and altered optical response. Consequently, the reduced accuracy for condensed milk is an expected outcome of both the underlying material differences and the model's learned decision strategy.

\begin{figure}[htbp]
    \centering
    \includegraphics[width=0.8\textwidth]{figures/6.24a.png}
    \caption[(experiment 7a) Validation test on pre-trained binary classifier with untrained evaporated milk sample as pure milk]{\textbf{(experiment 7a) Validation test on pre-trained binary classifier with untrained evaporated milk sample as pure milk}}
    \label{fig:6_24a}
\end{figure}

\begin{figure}[htbp]
    \centering
    \includegraphics[width=0.8\textwidth]{figures/6.24b.png}
    \caption[(experiment 7a) validation test on pre-trained binary classifier with untrained condensed milk sample as pure milk.]{\textbf{(experiment 7a) validation test on pre-trained binary classifier with untrained condensed milk sample as pure milk.} (~65\%)}
    \label{fig:6_24b}
\end{figure}

\section{Experiment 7b Product generalisability of purity detection model (non-consumable milk)}
\label{sec:experiment_7b_chapter6}

The broad definition of milk products necessitates thoroughness in designing purity assessment systems. Similar to Experiment~7a, this final test evaluates how the purity model trained on liquid milk performs when applied to a non-consumable liquid cosmetic milk product. The diverse range of products marketed as ``milk'' (spanning consumable dairy items to cosmetic formulations) underscores the importance of understanding model generalisability across product categories with fundamentally different compositions and intended uses \cite{walstra1999,aernouts2015}.

Principal component analysis of extracted features from both control (fresh homogenised milk) and cosmetic milk samples reveals clear linear separation in the projected feature space (Figure~\ref{fig:6_25}), as expected given the substantial compositional differences between consumable dairy milk and cosmetic formulations. This separation contrasts with the substantial overlap observed in Experiment~7a between different consumable milk products, reflecting the greater divergence in material properties between dairy and cosmetic applications.

To quantify the model's performance on this non-consumable product, a zero-shot validation test was conducted using the pre-trained binary purity detection model. The classifier achieved an accuracy of approximately 65\% (Figure~\ref{fig:6_26}), which is notably lower than the performance on consumable milk variants. This reduced accuracy is consistent with expectations, as cosmetic milk products differ substantially from dairy milk in their base composition, optical properties, and scattering characteristics \cite{nielsen2013,aernouts2015}. The model, having been trained exclusively on dairy milk transients, encounters a distribution shift when applied to cosmetic formulations, leading to the observed performance degradation. This outcome highlights the importance of product-specific training or domain adaptation strategies when extending purity detection systems to fundamentally different material categories.

\begin{figure}[htbp]
    \centering
    \includegraphics[width=0.8\textwidth]{figures/6.25.png}
    \caption[(experiment 7b) PCA of two homogenised milk products; cosmetic milk (non-consumable) and whole milk (consumable).]{\textbf{(experiment 7b) PCA of two homogenised milk products; cosmetic milk (non-consumable) and whole milk (consumable).}}
    \label{fig:6_25}
\end{figure}

\begin{figure}[htbp]
    \centering
    \includegraphics[width=0.8\textwidth]{figures/6.26.png}
    \caption[(experiment 7b) validation test on pre-trained binary classifier with untrained cosmetic milk sample as pure milk.]{\textbf{(experiment 7b) validation test on pre-trained binary classifier with untrained cosmetic milk sample as pure milk.} (~67\%)}
    \label{fig:6_26}
\end{figure}

\begin{figure}[htbp]
    \centering
    \includegraphics[width=0.8\textwidth]{figures/6.27.png}
    \caption[2D plot of PCA of two whole milk samples; pure (homogenous) and expired (colloid).]{\textbf{2D plot of PCA of two whole milk samples; pure (homogenous) and expired (colloid).} Partial overlap of variables per class indicate little similarity in purity patterns. Explained Variance Ratio:[0.18096039, 0.16276842]}
    \label{fig:6_27}
\end{figure}


\subsection{Guiding principles and models}
\label{subsec:guiding_principles_models_chapter6}

The material of interest (MOI), comprising homogenised milk and its solid container, exhibits translucent optical properties \cite{walstra1999,klocek1991}. Consequently, the same guiding principles established in Section~\ref{subsec:guiding_principles_models} of Chapter~\ref{chap:flat_homogeneous_material_detection} apply to this experimental configuration. Given the homogeneous composition of the target materials, the temporal signal adheres to the per-bin Gaussian approximation model described by Equation~\eqref{eq:pure_material_gaussian} \cite{tontini2020,koerner2021,su2016}. This single-peak Gaussian formulation accurately provides per-bin representations of photon arrival times for uniform material surfaces, where the temporal distribution of reflected photons follows a normal distribution centred at the time-of-flight $t_0$ with standard deviation $\sigma$ \cite{becker2023,nielsen2013}.


\subsection{Multizone milk experiments' results}
\label{subsec:multizone_experiment_results_chapter6}

Across the eight multizone milk experiments, the time-of-flight SPAD instrument consistently demonstrated efficient performance for purity assessment of dairy products. Because samples were sourced from multiple manufacturers, fat contents, and acquisition days, the learned purity representation generalises to similar commercial products acquired under different conditions. To guard against misleading performance due to overfitting, every experiment was cross-validated using principal component analysis (PCA), which produced cluster structures and separation boundaries that were consistent with the decisions of the trained machine learning models \cite{wold1987}.

Experiment~1 confirmed that reliable feature extraction requires homogenised milk. In non-homogenised samples, transient responses changed with small perturbations in the liquid state, so ground-truth labels assigned at one instant could not be guaranteed to remain valid after the container was moved. This temporal instability of colloidal particle configurations makes it difficult for a single calibration model to assess the purity of inhomogeneous milk, aligning with earlier reports on the challenges of optical milk analysis \cite{aernouts2015,melfsen2012,walstra1999}. Consequently, subsequent purity experiments were restricted to homogenised milk, where particle distributions remain stable over time and position.

Experiment~2 showed that all impurity classes (including solid, liquid, gaseous, and non-material contaminants) are separable (albeit non-linearly) in the transient domain. This was further corroborated by the successful detection of invisible contamination from hydrogen peroxide using a binary classifier, demonstrating that even elusive impurities yield learnable feature shifts in the time-resolved signal space. Without this established separability, proceeding with Experiments~3 and~4 would not have been justified, as there would be no guarantee that features from contaminated and control samples were distinguishable by the learning models.

Experiment~5 demonstrated that the learned purity features generalise across manufacturers, indicating that users are not constrained to a specific brand for effective deployment of the solution. Experiment~6 yielded more nuanced behaviour: fat-content variation produced largely inseparable clusters in PCA space, suggesting that the model primarily captures energy--matter interactions governed by the underlying milk matrix rather than fat concentration. As a result, low- and high-fat variants exhibit similar transient responses despite differences in absolute fat levels. Finally, Experiment~7 showed that the trained model does not generalise to structurally dissimilar consumable milk products or to non-consumable cosmetic milk of comparable viscosity or colour, underscoring the need for product-specific calibration when extending the method beyond fresh homogenised drinking milk.

\subsection{Multizone milk models' performance justification}
\label{subsec:multizone_model_performance_justification_chapter6}
%
The accuracy figures reported for the multizone milk classifiers arise from a combination of experimental design choices and independent validation analyses that collectively reduce the risk of overfitting. First, all models were trained exclusively on homogenised milk, as established in Experiment~1, ensuring that the underlying material state remained stable over time and position. This constraint eliminates label drift due to colloidal reconfiguration and yields consistent ground-truth transients, which is critical for learning reliable decision boundaries \cite{walstra1999,michalski2007}.

Second, Experiments~2 and~3 demonstrated that purity-relevant features are intrinsically learnable and separable in the transient domain rather than being artefacts of a particular network configuration. The multi-label classifier in Experiment~2 successfully distinguished all 17 impurity and control classes, and principal component analysis (PCA) revealed consistent clustering among these classes, confirming that class separability is present at the feature level independent of the neural network architecture \cite{wold1987,bishop2006}. The subsequent binary purity classifier in Experiment~3 achieved high training and validation accuracies while operating on these same feature distributions, indicating that it is exploiting genuine structure in the data rather than memorising noise \cite{lecun1998,bishop2006}.

Third, generalisation was explicitly probed through out-of-sample validation scenarios. Experiment~5 evaluated performance on control samples from manufacturers unseen during training, while Experiment~6 assessed behaviour across different fat contents drawn from multiple brands. In both cases, the binary classifier maintained strong accuracy, and PCA projections showed that pure samples with different packaging and fat levels still occupied a compact manifold in feature space, consistent with a shared notion of ``purity'' rather than brand-specific signatures \cite{aernouts2015,nielsen2013}. Additional zero-shot tests on non-consumable cosmetic milk and structurally transformed consumable products (Experiment~7) produced systematically degraded performance and clear PCA separation from the training distribution, as expected for true domain shifts, further supporting the view that the model is not overfitted to incidental properties but instead captures a robust representation of fresh homogenised drinking milk.

Finally, multiple safeguards were employed during training and evaluation to mitigate classical overfitting modes. Data were collected across different days, containers, and acquisition sessions, increasing intra-class variability and reducing the likelihood that models could rely on spurious correlations such as acquisition order or micro-alignment of the setup. Model selection was based on held-out validation sets rather than training loss alone, and convergence behaviour was monitored to avoid over-parameterised solutions that fit noise. Combined with cross-technique consistency between PCA and CNN-based classifiers, these measures justify the reported accuracies as credible estimates of real-world performance rather than optimistic artefacts of overfitting.

\newpage
\section{Future research and development}
\label{sec:future_research_development_chapter6}
\noindent
Future research will extend the mixture-model framework introduced in Chapter~3 to purity assessment in heterogeneous fluids, including raw milk. The experiments in this thesis showed that the prototype instrument can quantify purity across manufacturers, fat contents and acquisition days, but also revealed limitations that motivate further work. Building on Experiment~1, which established that homogenisation is essential for stable feature extraction, future studies will investigate advanced homogenisation strategies and their impact on measurement repeatability. In inhomogeneous milk, the challenge arises not only from spatially varying fat globules and particles, but also from the liquid nature of the sample, which causes particle configurations to evolve over time. This temporal instability complicates the definition of consistent ground-truth labels, as shaking or moving the container alters the internal microstructure. Future work will therefore explore dynamic calibration procedures and online or adaptive models that track and compensate for time-varying particle configurations in real time, addressing a fundamental difficulty that has also been noted in optical characterisation of untreated milk and in-line milk analysis \cite{aernouts2015,melfsen2012,walstra1999}.

Experiment~2 demonstrated that material and contaminant classes are non-linearly separable in the learned feature space, and that invisible hydrogen peroxide contamination can be detected reliably using a binary classifier. This foundational result opens the door to more complex contamination scenarios, such as multi-contaminant mixtures or lower concentration regimes. Future studies will broaden the evaluation to a wider range of milk products and to temperature conditions under which products remain consumable (for example, refrigerated storage close to freezing or brief heating). Temperature affects particle dynamics, fat distribution and optical properties, which may both expose new purity-relevant features and introduce additional sources of variability that models must accommodate \cite{lott2023,burgwald1947}.

The generalisation behaviour observed in Experiment~5 suggests that users are not restricted to specific brands, as the binary purity classifier maintained performance on manufacturers unseen during training. Future work will extend this analysis to a broader set of products, including organic milk, lactose-free variants and milks with a wider range of fat percentages, to delineate the practical operating domain of the system. Experiment~6 further indicated that fat content was largely inseparable in principal-component space, suggesting that the model responds primarily to energy--matter interactions with particles rather than to bulk concentration alone. Subsequent research will test whether this apparent particle-level sensitivity can be exploited for more sophisticated purity assessments, such as detecting structural changes or subtle molecular-level contaminants beyond simple concentration shifts.

Experiment~7 revealed that the current solution does not automatically transfer to consumable products with markedly different structure and/or colour, nor to non-consumable liquids with similar viscosity, colour or structure. This limitation motivates a line of work on transfer learning and domain adaptation to translate the learned representations to other milk types and, more broadly, to other fluid systems. Promising directions include domain-adversarial or feature-alignment methods, multi-task learning frameworks that jointly model multiple products, and mixture-model approaches that explicitly encode product clusters while preserving purity-detection accuracy within each domain.

Across all experiments, principal component analysis was used to cross-check the behaviour of the learned classifiers and to guard against overfitting. Future research will incorporate more rigorous validation protocols, including repeated and stratified cross-validation across different temporal periods, external validation on fully independent datasets, and benchmarking against established analytical techniques from dairy science and spectroscopy. These steps will help further quantify the robustness and generalisability of the proposed purity-assessment pipeline.

Finally, future work will investigate how this measurement and classification framework can be integrated into practical workflows, including in-line quality control in dairy processing facilities, portable or consumer-facing test devices and regulatory monitoring infrastructures. Realising these applications will require addressing constraints on acquisition time, computational latency and cost, as well as designing interfaces that support non-expert users while still exposing sufficient diagnostic information for expert oversight.
\newpage
\section{Conclusion}
\label{sec:conclusion_chapter6}
\noindent
The stakes of consumable-material purity are high, particularly for staple products such as milk where contamination can compromise safety, nutritional value and trust in supply chains \cite{gregory1961,goodman1986}. This chapter has presented and validated an economical, low-complexity and non-invasive model for milk purity assessment that addresses these stakes by overcoming key limitations of traditional analytical methods and laboratory-based assays \cite{joolaei2025}. By exploiting transient time-of-flight measurements rather than visual inspection or bulk composition alone, the system delivers efficient and reliable detection of purity deviations in settings ranging from domestic kitchens to industrial quality-control lines.

The model was rigorously tested for its ability to detect invisible contamination, one of the most pressing challenges in milk safety. Experiment~2 demonstrated that non-linearly separable impurity classes, including invisible hydrogen peroxide contamination, remain distinguishable in the learned feature space using binary classification. This confirms that contaminants which do not alter the apparent colour or structure of milk still induce systematic, learnable shifts in the transient response, aligning with established evidence that even low levels of hydrogen peroxide can materially affect milk quality and safety \cite{gregory1961}. Together with the broader literature on spectroscopic and optical techniques for detecting adulteration, these results position the proposed model as a practical route to routine detection of otherwise elusive contaminants \cite{aernouts2015,joolaei2025}.

Generalisation was evaluated across multiple, practically relevant axes. Experiment~5 showed that the learned purity representation transfers across manufacturers, indicating that the system does not rely on brand-specific signatures and can therefore be deployed in diverse supply chains. Experiment~6 revealed that the model generalises across fat compositions: principal component analysis (PCA) suggested that the classifier responds primarily to energy--matter interactions at the particle level rather than to concentration-dependent bulk features, making performance robust to compositional variations in fat content. The experimental design, which deliberately sourced products from multiple manufacturers, with different labelled contents and, critically, across different acquisition days, mitigated temporal and batch effects that commonly confound in-line milk analysis \cite{melfsen2012,lott2023}. As a result, the findings can be confidently extended to similar products obtained on different occasions, addressing concerns about day-to-day and batch-to-batch variability in real-world use.

The model's operating domain and limitations were also explicitly characterised. Across the main study, the solution was tested on different drinking-milk products (see discussion of results), demonstrating versatility while setting clear boundaries. Experiment~7 showed that the trained classifier does not automatically generalise to consumable products with markedly different structure and/or colour, or to non-consumable cosmetic milk with similar viscosity, colour or structure. This delineates a practical applicability region centred on fresh, homogenised drinking milk and highlights the need for product-specific calibration or domain adaptation when extending the approach to structurally dissimilar matrices. Experiment~1 further established that homogenisation is essential for stable feature extraction, as it suppresses temporal and spatial uncertainties inherent in inhomogeneous milk \cite{walstra1999,michalski2007}. By enforcing homogenisation, the study ensured that observed differences in transient responses reflected genuine material and contamination effects rather than uncontrolled microstructural rearrangements over time.

To guard against deceptive results or overfitting, the experimental campaign combined machine-learning evaluation with independent statistical validation. PCA was used throughout to verify that separability and clustering patterns observed in the learned feature space were also present in a model-agnostic representation of the data \cite{wold1987}. Consistent findings between PCA and the convolutional neural network classifiers confirm that the reported accuracies are supported by underlying structure in the measurements rather than by model-specific artefacts. This cross-technique agreement strengthens confidence that the system captures a robust and physically plausible notion of ``purity'' in the transient domain.

Overall, the eight comprehensive experiments, supported by independent validation, establish the proposed solution as a rigorous and reliable framework for milk purity assessment. The model detects invisible contamination, generalises across manufacturers and fat compositions, and maintains consistency across temporal variations in acquisition, positioning it as a viable candidate for both domestic and commercial deployment. Future work on mixture models, temperature-dependent behaviour and domain adaptation to other milk types and heterogeneous fluids (as outlined in the preceding section) will extend this foundation and help translate the approach into broader classes of consumable materials \cite{burgwald1947,lott2023}. Nonetheless, the current implementation already addresses a critical need for economical, non-invasive purity assessment in milk and demonstrates how time-resolved optical sensing, coupled with machine learning, can provide actionable guarantees about the integrity of everyday consumables.

\newpage
\chapter{Summary of Contributions}
\label{chap:Summary of contributions}

\newpage
\section{Introduction}
\label{sec:chapter7_introduction}
This thesis investigates how compact, low-cost SPAD time-of-flight (ToF) sensors can be used to enhance machine vision. Building on the literature reviewed in Chapter~2, it extends prior work by proposing and experimentally validating approaches that target a diverse set of machine-vision tasks while remaining cost-effective and accessible. The core contribution is a set of practical solutions that address key limitations of current vision systems without relying on expensive or bulky hardware.

To support these developments, a sensing and signal-processing framework is introduced that formalises core concepts from engineering and computer-vision literature. The framework provides a theoretical bridge between optical sensing principles and machine-learning models by defining standardised procedures for processing time-resolved signals, extracting structural features, and integrating temporal information with spatial image data for scene understanding. In doing so, it supports reproducible research and offers a foundation for future work on spatiotemporal vision systems.

Solutions One (Chapter~4) and Three (Chapter~6) demonstrate that time-resolved signals alone can reveal internal and surface properties of targets through characteristic temporal signatures. Solution One shows that transient image analysis can extract material-specific temporal patterns that enable accurate material classification without spectral information, capturing optical properties, scattering behaviour, and internal structure not readily accessible from spatial intensity images. Solution Three extends this capability to material-purity assessment in liquids, showing that time-resolved signals can distinguish pure from contaminated materials, detect otherwise invisible contaminants such as hydrogen peroxide, and generalise across manufacturers and fat compositions.

Solution Two (Chapter~5) addresses higher-level computer-vision tasks such as scene understanding and the development of intelligent visual machines. By combining spatial object detection with temporal material classification, it enables joint reasoning about which objects are present and what materials they are composed of, yielding richer semantic representations than traditional object-detection pipelines. This combined view supports more informative decision-making in autonomous systems.

The results further show that certain structural properties of objects can be inferred more effectively from time-resolved than from intensity-resolved images, as tested in the RGB case. Time-resolved measurements encode depth, material composition, and structural characteristics through temporal patterns that complement spatial intensity information. At the same time, state-of-the-art models that rely solely on 2D RGB data confirm that intensity-resolved images remain the primary source for spatial features such as object location, shape, size, and spatial relationships. Together, these findings suggest that optimal perception systems should combine both modalities: spatial images for spatial structure and temporal signals for material and structural cues.

Across the three solutions, the work shows that low-budget SPAD ToF detectors, even without spectral analysis, can substantially improve machine perception. Temporal information offers a cost-effective alternative to spectroscopic systems while achieving comparable or superior performance in material-detection tasks. The thesis argues that robust treatment of such sensing fundamentals is a prerequisite for building larger, more complex models and infrastructures that aspire to human-level visual reasoning.

The proposed approaches address several specific problems, including low-level object classification via material detection and improved scene understanding from combined SPAD and RGB data. The representation includes target depth (for example, centre-of-mass estimates based on histogram-bin position), true distance in metric units, and explicit foreground–background information in the prediction output, enabling a three-dimensional and material-aware view of the scene. On this basis, the work points towards perception systems capable of supporting safer and more reliable applications such as household robots, autonomous driving systems, and large-scale reasoning models.

Finally, the food-safety solution introduces a compact, low-budget SPAD-based method for assessing milk purity. Unlike many optical systems that require bulky or expensive setups, and non-optical methods that are often invasive or inconvenient, the proposed system is non-invasive, compact, and accessible to non-experts. By enabling end users to detect contamination, assess purity, and generalise across products, it offers a simple and affordable means for consumer-level food-safety monitoring and has the potential to strengthen everyday food-safety practices and consumer trust.

\section{Novelty of contributions}
\label{sec:novelty_of_contributions_chapter7}

This project introduces several novel contributions that advance the state of the art in machine vision and material detection. Building on the overview in Chapter~1, it proposes new theoretical frameworks, demonstrates previously unexplored sensing and inference capabilities, and establishes foundations for future research in spatiotemporal vision systems.

Chapter~3 formalises core machine-vision and signal-processing principles into four models within the Particle Wave Detection (PAWD) framework, unifying spatial-spectral, spatial non-spectral, temporal-spectral, and temporal non-spectral regimes. This framework provides a common theoretical foundation that bridges optical sensing, signal processing, and machine learning, defining standardised approaches for acquiring and processing time-resolved signals, extracting structural features, and integrating temporal information with spatial images for scene understanding (Figure~\ref{fig:1_1}). The four models jointly address feature extraction from time-resolved signals, material characterisation through temporal patterns, the integration of spatial and temporal cues, and evaluation metrics for combined detection systems, including mixture-based formulations for multi-material signals (e.g., Equation~\ref{eq:gaussian_mixture}).

Within this framework, the thesis introduces a theory of how machines can perceive single objects of multiple material composition as unified identities. When a target contains several materials (for example, a ceramic cup with a metal handle or a wooden table with metal legs), conventional pipelines often struggle to preserve object identity while simultaneously resolving constituent materials. The proposed theory models the temporal signal from such multi-material regions as a coherent mixture, where the object identity emerges from the joint temporal response of its components rather than being obscured by material diversity (Equation~\ref{eq:gaussian_mixture}). This theory underpins the multimaterial reasoning in Chapter~5, providing the basis for simultaneous object detection and material classification from hybrid SPAD--RGB measurements (Figure~\ref{fig:5_1}).

For the first time in computer-vision and engineering literature, Chapter~5 demonstrates joint detection of object identity and material composition using only 2D RGB images and co-registered 1D time-resolved signals from the target surface. This dual-detection capability goes beyond existing approaches that typically address object detection and material classification in isolation, showing that fundamentally different signal types can be integrated into a single, end-to-end trainable system (Figures~\ref{fig:5_1} and~\ref{fig:5_2}). The results further show that the system maintains material-discrimination performance across arbitrary target positions, orientations, and distances, including challenging views and occlusions (Figure~\ref{fig:5_4}), thereby demonstrating that material understanding can be made robust to typical deployment conditions.

The novelty lies not only in the multimodal combination but in demonstrating that RGB images and time-resolved histograms contribute complementary, synergistic information. RGB images provide spatial context, object boundaries, and appearance features, while time-resolved signals supply material-specific temporal signatures, depth cues, and structural characteristics (Figure~\ref{fig:5_3}). Their successful fusion shows that spatial and temporal features can be jointly exploited to achieve superior scene understanding compared with either modality alone, particularly for multi-material and non-planar targets governed by mixture responses such as those modelled in Equation~\ref{eq:gaussian_mixture}.

Among related literature, the presented solution uniquely enhances the quality and trustworthiness of material detection by enabling a clear distinction between original physical materials and their 3D models or images. Time-resolved signals encode material-specific temporal responses that arise from true optical properties, scattering behaviour, and internal structure, which cannot be replicated by photographs or rendered models (Figure~\ref{fig:5_6}). This capability addresses a critical limitation of appearance-only vision systems, which often cannot distinguish real physical objects from visually similar representations, and instead grounds decisions in physically constrained signals.

The ability to distinguish real materials from representations constitutes a breakthrough for safety-critical security systems and other robotic or autonomous visual platforms. In such settings, machines must reliably differentiate genuine objects from photographs, projected images, decoys, or visually convincing replicas. By jointly detecting object identity and material composition, and by rejecting visually plausible but physically inconsistent stimuli, the proposed system provides a more reliable perceptual layer for applications such as autonomous driving, medical robotics, and security monitoring.

Beyond improving material-detection quality, the thesis argues that accurate material-composition detection is a necessary step toward more intelligent vision systems. Material understanding is treated as a core component of scene understanding rather than an auxiliary attribute: detecting what objects are present is coupled with inferring their physical properties, likely behaviours, and affordances. This enables higher-level reasoning, such as predicting how objects will interact, selecting appropriate manipulation strategies, and making contextually appropriate decisions based on both geometry and composition.

By establishing material detection as a foundational capability and demonstrating its feasibility with compact, low-cost SPAD time-of-flight sensors and RGB imagers, this work provides a practical pathway toward materially informed, depth-aware visual intelligence. The presented frameworks, mixture models, and hybrid sensing architectures collectively show that incorporating temporal material cues alongside spatial appearance moves machine vision closer to human-like perception and more reliable decision-making in complex, real-world environments.

\section{generalisability of contributions}
\label{sec:generalisability_of_contributions_chapter7}
\noindent The experiments presented in this thesis were deliberately conducted in relatively unconstrained settings, yet each solution is linked to a specific problem class and experimental configuration. Assessing how far these solutions generalise beyond their original domains is therefore essential for understanding their scope of applicability, the main limitations that currently bound their use, and the most promising directions for extending them to broader and more challenging scenarios.

In Chapter~4, generalisation was studied primarily on flat, planar translucent and opaque materials under a range of viewing distances and orientations. The multizone SPAD classifier maintained high accuracy across variations in standoff distance and view angle (e.g.\ Figures~\ref{fig:4_8}--\ref{fig:4_11},~\ref{fig:4_14},~\ref{fig:4_15} and~\ref{fig:4_16}), indicating robustness to moderate changes in sensor--target geometry and supporting deployment in settings where precise positioning cannot be guaranteed. At the same time, all experiments were restricted to single planar targets in controlled indoor environments. Real-world scenes often contain multiple, non-planar surfaces, self-occlusions, and complex backgrounds, together with outdoor factors such as uncontrolled illumination, atmospheric effects, temperature variation, and environmental noise. These factors are likely to introduce additional ambiguities (for example, multipath interference or signal mixing across surfaces) that are not captured by the current planar-target study. A natural avenue for future work is therefore to evaluate the transient-based material classifier in complex indoor scenes with multiple objects and in outdoor environments where natural lighting and environmental variability are present.

\paragraph{More material types to be tested in future.}
In Chapter~5, temporal material features were integrated with spatial object detection to improve discrimination between real objects and their representations (Figures~\ref{fig:5_1}--\ref{fig:5_3}). The demonstrated system focused on single objects observed one at a time, often of hybrid material composition but isolated from other targets in the scene (Figures~\ref{fig:5_4}--\ref{fig:5_6}). Under these conditions, the joint SPAD--RGB detector generalises well to individual objects whose transient signals can be analysed without interference from neighbouring items, and it reliably separates genuine physical targets from photographs, renderings, or decoys. Extending this framework to multi-object scenes remains an open challenge. When several objects are present simultaneously, additional failure modes arise, including mutual occlusion, spatial overlap of temporal signals, confusion between materials with similar signatures, and increased computational load from jointly processing multiple regions of interest. Addressing these issues will be important for deploying the hybrid detector in realistic applications such as autonomous navigation, robotic manipulation, and surveillance, where multiple interacting objects are the norm rather than the exception.

Chapter~6 introduced a SPAD-based milk purity assessment system evaluated on homogenised consumable milk from several local manufacturers (Figures~\ref{fig:6_1} and~\ref{fig:6_12}). Within the experimental design, the learned decision boundaries generalise across brands and batches that satisfy the fresh consumable milk standard, while reliably rejecting a broad range of material and non-material contaminants (Figures~\ref{fig:6_19}--\ref{fig:6_21}). In this sense, the ``pure'' class is tied to a specific physical state of the fluid: any deviation from this state, including heating, freezing, or storage-induced degradation, is treated as an impurity because it alters the underlying transient response and optical properties (Figures~\ref{fig:6_13},~\ref{fig:6_14}, and~\ref{fig:6_18}). The system's sensitivity to such transformations is essential for robust purity assessment, but it also implies that generalisation is conditioned on a well-defined notion of purity and on the range of states represented during training.

Several extensions follow naturally from this work. First, testing purity across a wider span of temperatures would quantify how thermal variation affects transient signatures and classification performance, which is crucial for practical deployments in supply chains where milk is stored and consumed at different temperatures. Second, most experiments relied on observing the fluid through transparent or translucent containers, which simplifies deployment for typical packaged products. Although this container-based sensing demonstrates that non-invasive purity assessment is feasible in realistic packaging, it leaves open the question of direct observation in cases where containers are opaque or bulk storage systems preclude side access. Direct measurement of the milk volume itself may require alternative geometries or illumination strategies, and could reveal additional material properties that are partially masked when viewing through a container.

Finally, extending the purity framework beyond milk requires defining suitable ground-truth ``pure'' states for new materials. Unlike consumable milk, where established standards and industrial practice provide clear reference definitions, other materials may not have a single canonical pure state, or their purity may depend on application-specific functional properties. For each new material, it will therefore be necessary to specify reference samples, purity criteria (for example, based on composition, optical behaviour, or functional performance), and calibration procedures that account for material-specific variability. These steps may, in turn, motivate adaptations of the detection architecture or training protocol to accommodate different transient regimes, container types, or structural properties.

Overall, the generalisability assessment indicates that the proposed solutions perform strongly within their tested regimes and already tolerate substantial variability in geometry, distance, orientation, and manufacturing source. At the same time, they highlight clear opportunities for extension to more complex scenes, richer material sets, and broader operational conditions. The frameworks developed in this thesis therefore provide a solid foundation that can be systematically expanded through future work, while preserving the core advantages of compact instrumentation, low cost, and non-invasive measurement.

\section{Future research and development}
\label{sec:future_research_development_chapter7}

Future research may proceed along three broad, complementary directions. First, at the algorithmic level, there is a clear need to extend the current frameworks to more complex scenes, including multi-target configurations, heterogeneous and anisotropic materials, and long-range or outdoor deployments. A particular focus will be on understanding why single-zone transient inputs underperform relative to multizone stacked inputs when distance and geometry vary, and on developing distance- and angle-aware single-zone architectures that retain the efficiency advantages of single-pixel sensing while closing the robustness gap observed between 0 and 100\,cm. In parallel, the benchmarking protocols introduced in Chapters~\ref{chap:flat_homogeneous_material_detection} and~\ref{chap:spatiotemporal_object_detection} will be refined with more informative metrics beyond simple averages, enabling finer-grained evaluation across tasks, object categories, material classes, and operating conditions.

Second, application-specific developments will extend the demonstrated solutions to new domains. For spatiotemporal object detection, future work will include multi-object scenes, camouflage and illusion-rich environments, improved scale estimation, and joint spatial--temporal occlusion handling. For purity assessment, the mixture-model framework will be generalised from homogenised milk to more heterogeneous fluids, broader product families and temperature regimes, supported by richer validation protocols and domain-adaptation techniques. Across all applications, deployment-oriented research will target compact embedded implementations on single-board computers or microcontrollers, so that the proposed instruments can be translated from laboratory prototypes into reliable, field-ready systems.

Finally, some of the unique concepts and architectures developed in this thesis (including hybrid SPAD--RGB sensing strategies, transient-based purity taxonomies and multimaterial reasoning frameworks) may warrant intellectual property protection and industrial collaboration. Patenting appropriate components and seeking targeted sponsorship would help secure the resources required to mature these ideas into robust products, accelerate their adoption in safety-critical and industrial settings, and sustain a longer-term research programme on materially informed, depth-aware machine perception.

\section{Conclusion}
\label{sec:conclusion_chapter7}

This thesis has shown that compact, low-cost SPAD time-of-flight sensors, when coupled with appropriate signal processing and learning frameworks, can provide a unifying basis for materially informed, depth-aware machine perception. Building on the conceptual foundations of the Particle Wave Detection framework and the literature review in Chapter~\ref{sec:conclusion_chapter2}, the experimental chapters demonstrated that non-spectral transient signals from flat surfaces enable reliable material classification, that hybrid SPAD--RGB sensing supports joint object and material detection in complex scenes, and that the same sensing principles can be extended to economically assess milk purity, including invisible contamination, across manufacturers and fat compositions. Together, these results establish temporal sensing as a practical alternative or complement to purely spatial and spectral approaches, bridging low-level structural perception (depth, material composition, purity) with higher-level scene understanding in a way that is compatible with embedded, resource-constrained platforms. The work therefore provides both theoretical and empirical evidence that addressing foundational sensing questions is a necessary step toward building more robust, trustworthy and generally intelligent vision systems.

\newpage
\bibliographystyle{IEEEtranN}
\bibliography{references}

\newpage
\appendix

\newpage
\chapter{Extra images}
\label{sec:extra_images}

\begin{figure}[htbp]
    \centering
    \includegraphics[width=0.43125\textwidth]{figures/5.10.png}
    \caption[Target regions of similar 2D geometric patterns.]{\textbf{Target regions of similar 2D geometric patterns.}}
    \label{fig:5_10}
\end{figure}

\begin{figure}[htbp]
    \centering
    \includegraphics[width=0.525\textwidth]{figures/5.11.png}
    \caption[F1 score and confidence levels for custom trained models.]{\textbf{F1 score and confidence levels for custom trained models.} (Top): F1 score and level of confidence with which the various class labels make predictions. This curve gives a fair idea of the quality of custom trained models that plug into the spatial detection module of the dual model. (Bottom): High performance state-of-the-art vision model (YOLOv8) struggling to distinguish between two different materials of similar geometric shape: carrots and parsnips. Model classifies everything as carrot due to limitation of sole spatial feature processing. On the contrary, transient signal encodes unique optical features that can distinguish between the different material composition. Hence the presented instrument that combines space and time resolved features.}
    \label{fig:5_11}
\end{figure}

\begin{figure}[htbp]
    \centering
    \includegraphics[width=0.75\textwidth]{figures/5.23.png}
    \caption[Confusion matrix of YOLOv3 spatial detection training.]{\textbf{Confusion matrix of YOLOv3 spatial detection training.} Confusion matrix of YOLOv3 spatial detection training (~94\% accuracy).}
    \label{fig:5_23}
\end{figure}

\begin{figure}[htbp]
    \centering
    \includegraphics[width=0.75\textwidth]{figures/5.25.png}
    \caption[Side-by-side comparison of the same object from different materials.]{\textbf{Side-by-side comparison of the same object from different materials.} Left: purple bowl object perceived from LED screen surface. Right: purple bowl object perceived from plastic material surface. These two objects have same prediction by spatial vision model but different predictions by material classifier.}
    \label{fig:5_25}
\end{figure}

\begin{figure}[htbp]
    \centering
    \includegraphics[width=0.75\textwidth]{figures/5.26.png}
    \caption[Project web application demo section interface.]{\textbf{Project web application demo section interface.}}
    \label{fig:5_26}
\end{figure}

\begin{figure}[htbp]
    \centering
    \includegraphics[width=0.75\textwidth]{figures/5.27.png}
    \caption[Spatiotemporal detection of a natural potato from physical space.]{\textbf{Spatiotemporal detection of a natural potato from physical space.}}
    \label{fig:5_27}
\end{figure}

\begin{figure}[htbp]
    \centering
    \includegraphics[width=0.75\textwidth]{figures/5.28.png}
    \caption[Spatiotemporal detection of 3D model potato from LED screen.]{\textbf{Spatiotemporal detection of 3D model potato from LED screen.}}
    \label{fig:5_28}
\end{figure}

\begin{figure}[htbp]
    \centering
    \includegraphics[width=0.75\textwidth]{figures/5.29.png}
    \caption[Spatial training samples of everyday objects.]{\textbf{Spatial training samples of everyday objects.}}
    \label{fig:5_29}
\end{figure}


\newpage
\chapter{Extra Tables}
\label{sec:extra_tables}

\begin{table}[H]
\centering
\footnotesize
\begin{tabularx}{\textwidth}{lX X X}
\hline
\textbf{Model} & \textbf{Object Detection Accuracy (\%)} & \textbf{Material Classification Accuracy (\%)} & \textbf{Combined Performance} \\
\hline
YOLOv3 & 94.0 & 99.0 & 94.5 \\
YOLOv8 & 98.0 & 99.0 & 98.5 \\
DINOv3 & 95.0 & 99.0 & 97.0 \\
Material Classifier (standalone) & -- & 99.0 & -- \\
\hline
\end{tabularx}
\caption[Combined dual detection performance metrics.]{\textbf{Combined dual detection performance metrics.} Performance summary for spatial object detection models individually and in combination with the material classifier. Combined performance is calculated as the average of object detection and material classification accuracies, reflecting equal weighting of both tasks.}
\label{tab:5_4}
\end{table}

\begin{table}[H]
\centering
\footnotesize
\begin{tabular}{lccc}
\hline
 & \textbf{Train (\%)} & \textbf{Validation (\%)} & \textbf{ZSL invalid data (\%)} \\
\hline
Experiment 3 & $\sim$98.9 & $\sim$93.3 & $\sim$18 \\
Experiment 4 & -- & $\sim$95 & -- \\
\hline
\end{tabular}
\caption[Training and validation results of purity detection classifier.]{\textbf{(experiment 3, experiment 4) training and validation results of purity detection classifier for all impurities vs invisible impurities.}}
\label{tab:6_3}
\end{table}

\begin{table}[H]
\centering
\footnotesize
\begin{tabular}{lcc}
\hline
\textbf{Brand} & \textbf{Fat content} & \textbf{Consistency} \\
\hline
Brand0a & 14\% & Homogenized \\
Brand0b & 14\% & Inhomogenised \\
Brand1 & 3.7\%, 1.8\%, 0.3\% & Homogenized \\
Brand2 & 3.7\%, 1.8\%, 0.3\% & Homogenized \\
Brand3 & 3.9\% & Homogenized \\
\hline
\end{tabular}
\caption[Metadata of milk samples used in experiment.]{\textbf{Metadata of milk samples used in experiment; brands, fat content and consistency.}}
\label{tab:6_4}
\end{table}

\begin{table}[H]
\centering
\footnotesize
\begin{tabularx}{\textwidth}{X X X X}
\hline
\textbf{NONMATERIAL IMPURITY} & \textbf{CONSISTENCY CHANGE} & \textbf{CONTAMINANT} & \textbf{FLUID CONTAINER TYPE} \\
\hline
Active nonmaterial impurity agent & Inhomogeneous colloidal & Poor temperature conditions & Transparent \\
 & Inhomogeneous suspension & Transluscent &  \\
Active nonmaterial impurity agent & Inhomogeneous colloidal & Product expiration under suitable temperature & Transparent \\
\hline
\end{tabularx}
\caption[Categorization of non-material contaminants.]{\textbf{Categorization of non-material contaminants by product environment conditions and expiration with resulting mixture attributes.}}
\label{tab:6_5}
\end{table}

\begin{table}[H]
\centering
\footnotesize
\begin{tabular}{lccccc}
\hline
\textbf{BRAND} & \textbf{MANUFACTURER} & \textbf{FAT CONTENT} & \textbf{CONSISTENCY} & \textbf{CONTAINER} & \textbf{ID} \\
\hline
Brand0a & Graham's family dairy & 14\% & Homogenized & Transluscent & hm\_C\_tl \\
Brand0b & Graham's family dairy & 14\% & Inhomogenised & Transluscent & im\_C\_tl \\
Brand1 & Sainsbury & 3.7\%, 1.8\%, 0.3\% & Homogenized & Transluscent & hm\_C\_tl \\
Brand2 & Tesco & 3.7\%, 1.8\%, 0.3\% & Homogenized & Transluscent & hm\_C\_tl \\
Brand2 & Tesco & 3.7\% & Homogenized & Transparent & hm\_C\_tp \\
Brand3 & Lancashire diaries & 3.9\% & Homogenized & Transluscent & hm\_C\_tl \\
\hline
\end{tabular}
\caption[Milk sources per brand, consistency, fat content and container.]{\textbf{Milk sources per brand, consistency, fat content and container.}}
\label{tab:6_6}
\end{table}





\end{document}


